% Options for packages loaded elsewhere
\PassOptionsToPackage{unicode}{hyperref}
\PassOptionsToPackage{hyphens}{url}
\documentclass[
]{article}
\usepackage{xcolor}
\usepackage{amsmath,amssymb}
\setcounter{secnumdepth}{-\maxdimen} % remove section numbering
\usepackage{iftex}
\ifPDFTeX
  \usepackage[T1]{fontenc}
  \usepackage[utf8]{inputenc}
  \usepackage{textcomp} % provide euro and other symbols
\else % if luatex or xetex
  \usepackage{unicode-math} % this also loads fontspec
  \defaultfontfeatures{Scale=MatchLowercase}
  \defaultfontfeatures[\rmfamily]{Ligatures=TeX,Scale=1}
  % unicode-math maps \setminus to U+29F5 (REVERSE SOLIDUS OPERATOR), which
  % latinmodern-math.otf lacks (tectonic/lualatex emit "Missing character
  % U+29F5"). Remap to U+2216 (SET MINUS) -- the glyph the font provides and
  % that matches the pdfTeX/amssymb rendering. Output-preserving; silences the
  % warning. Applies to all \setminus occurrences globally.
  \AtBeginDocument{\renewcommand{\setminus}{\mathbin{\char"2216}}}
\fi
\usepackage{lmodern}
\ifPDFTeX\else
  % xetex/luatex font selection
\fi
% Use upquote if available, for straight quotes in verbatim environments
\IfFileExists{upquote.sty}{\usepackage{upquote}}{}
\IfFileExists{microtype.sty}{% use microtype if available
  \usepackage[]{microtype}
  \UseMicrotypeSet[protrusion]{basicmath} % disable protrusion for tt fonts
}{}
\makeatletter
\@ifundefined{KOMAClassName}{% if non-KOMA class
  \IfFileExists{parskip.sty}{%
    \usepackage{parskip}
  }{% else
    \setlength{\parindent}{0pt}
    \setlength{\parskip}{6pt plus 2pt minus 1pt}}
}{% if KOMA class
  \KOMAoptions{parskip=half}}
\makeatother
\setlength{\emergencystretch}{3em} % prevent overfull lines
\providecommand{\tightlist}{%
  \setlength{\itemsep}{0pt}\setlength{\parskip}{0pt}}
\usepackage{bookmark}
\IfFileExists{xurl.sty}{\usepackage{xurl}}{} % add URL line breaks if available
\urlstyle{same}
\hypersetup{
  hidelinks,
  pdfcreator={LaTeX via pandoc}}

\author{}
\date{}
\author{}
\date{}
\begin{document}

\textbf{Ratushnyi Neural Repair Coding, Part II: Random Access,
Finite-Block Surprisal, and the Deviation Hierarchy}

Pavlo Ratushnyi

\emph{Independent Research --- Draft v0.8, 2026}

\textbf{Abstract}

This is Part II of the Ratushnyi Neural Repair Coding (RNR) series. Where
Part I fixes the framework, the five coder types, and their achievability /
converse matrix, this Part develops the finite-block behaviour of RNR
archives. We treat three linked questions. First, \emph{random access}: we
show (Theorems 7.2--7.5) that per-position sequential RNR coding under a
context-bounded predictor attains the source entropy rate while supporting
byte-granular seek-and-read, and we quantify the sync-point tax that
byte-granular access levies on a stationary source. Second, the
\emph{exact finite-block entropy} of a hidden-state predictor (the
Section 7.15 arc): we locate precisely where computing the coded length of
a block is polynomial, where it becomes hard as a coefficient-query
problem, and how a causal or fixed-slice encoder sidesteps the wall; we
also give neural realisations (Theorems 7.27$'$ and 7.30) in which a
genuine hidden-state predictor escapes the model-free succinct-index
regime. Third, the \emph{surprisal / deviation hierarchy} (Theorems
7.31--7.33, 7.35--7.36): the overflow error exponent of a lossless archive
(a large-deviations law whose Markov endpoint is a maximum cycle mean), the
universality dividend of a mixture predictor, the moderate-deviation fade
of the random-access penalty, an extreme-value law for the worst-case
query cost, and a functional central limit theorem governing the archive's
buffers. Random-access deployment as a block device closes the Part
(Section 10.7--10.8).

\textbf{Keywords:} lossless compression, neural side information, random
access, byte-granular seek, finite-block dispersion, varentropy,
large deviations, moderate deviations, extreme-value theory, functional CLT,
hidden Markov entropy, high-order empirical entropy.

\emph{Numbering follows the unified RNR series; this Part contains Sections
7.2--7.33, 7.35--7.36, and 10.7--10.8 (retaining their original numbers).
Results proved in a companion Part are cited as [RNR-I] (Part I:
framework and core theory), [RNR-II] (this Part; self-references keep their
bare numbers), and [RNR-III] (Part III: operational rate--distortion
dispersion). Cross-references to Sections and results outside the ranges
carried here resolve in the companion Parts.}

\emph{A note on numbering collisions.} The unified series preserves the
original identifiers, so two distinct objects legitimately carry the label
``7.2'': the subsection \textbf{7.2 Runtime parameter selection} (a
mode-selection discussion, in [RNR-I]) and \textbf{Theorem 7.2} (the
entropy-rate achievability result opening this Part). They are different
objects; context disambiguates.

\textbf{Scope of this Part.} Section 7 opens with the composition and
mode-selection results of [RNR-I, Section 7] (Theorems 7.1/7.1b/7.1c and
the runtime-parameter discussion 7.2); here we develop the finite-block
behaviour from Theorem 7.2 onward. The reader is assumed familiar with the
RNR repair principle and the five coder types; the prerequisites below
restate, statement-only, exactly the companion results this Part invokes.

\textbf{P. Prerequisites (restated from Part I)}

This Part is self-contained given the following results, all proved in
[RNR-I]. Statements are reproduced verbatim where practical; proofs are in
the companion Part. We also fix the evaluation frame (Definition 2.1) and
the coding tools (Sections 2.2--2.7) that every claim below is scoped by.

\textbf{P.1 Evaluation frame and coding tools (from [RNR-I, Section 2])}

\textbf{Definition 2.1 (Conditionally advantageous on a data class).}
For a source $X \sim D$ and a baseline compressor $B$, an RNR coder
$\mathsf{Enc}$ is \emph{conditionally advantageous} on $D$ if
\[
\mathbb{E}[L(\mathsf{Enc}(X))] < \mathbb{E}[L(B(X))], \quad X \sim D.
\tag{2.1}
\]
Throughout, $C$ denotes \emph{public context} (the input to the predictor
$M$, so $Y = M(C)$), $D$ the source distribution, and
$\mathsf{Enc}/\mathsf{Dec}$ the encoder/decoder pair.

We use, from [RNR-I, Sections 2.2--2.7], the following standard tools
(statements only; see the companion Part for discussion):
\emph{side-information coding} --- with decoder-reproducible side
information $Y = M(C)$ available to both parties, the relevant
achievability is ordinary conditional source coding at rate
$H(X \mid Y)$ (Cover \& Thomas 2006, Theorem 5.4.1), not Slepian--Wolf
binning; \emph{entropy coding} --- an arithmetic / ANS coder driven by an
explicit model $Q$ with $Q>0$ on the support codes a length-$n$ sequence in
$\lceil -\log_2 Q(x)\rceil + \varepsilon_n$ bits with
$\varepsilon_n / n \to 0$; \emph{enumerative coding} --- an element of a
decoder-reproducible finite class $\mathcal{C}$ costs $\log_2|\mathcal{C}|$
bits up to the $\varepsilon_n$ redundancy, with the two combinatorial
formulas $\log_2\binom{N}{k}$ (subset given $k$) and
$\log_2 M(N;k_0,\dots,k_{15}) = \log_2 \frac{N!}{k_0!\cdots k_{15}!}$
(16-class partition given the count vector), each preceded by an
$O(\log N)$ header for its data-dependent parameters; \emph{MDL} (Rissanen
1978), used only in its two-part admission form
$L(M) + L(X\mid M) < L_{\mathrm{base}}(X)$ as a design criterion, not an
optimality theorem; and \emph{bits-back / BB-ANS} (Hinton--van Camp 1993;
Frey 1997; Townsend--Bird--Barber 2019), whose expected net length is the
negative ELBO, equal to the source entropy only when the model marginal
matches the data and the variational posterior is exact.

\textbf{P.2 The RNR framework and the five types (from [RNR-I, Section 3])}

\emph{The repair principle.} A neural generator alone cannot losslessly
compress, because generator and source disagree on some bits with
non-negligible probability and a single disagreeing bit is fatal. RNR
treats the model output $Y = M(C)$ not as a reconstruction but as
\emph{side information} and stores a separate \emph{repair} description
that, combined with $Y$, reconstructs $X$ exactly. A coder is parameterised
by the predictor $M$, the representation target (1-D stream, 2-D positional
field, or higher tensor), and the repair encoding (error mask, enumerative
rank, selector, per-position probability, or mode flag).

\textbf{Definition 3.1 (RNR Type-I).} An RNR Type-I coder predicts a
one-dimensional damaged stream $Y = M_I(C)$ and stores a repair stream
$R_I$ together with a verification value $v = H_v(X)$. The decoder forms
$\hat{X} = \mathrm{DecRepair}(Y, R_I)$ and accepts only if
$H_v(\hat{X}) = v$; for a valid archive $\hat{X} = X$. The canonical
instance is Code12, mapping each source byte to a 12-bit codeword and
entropy-coding the predicted error mask by Hamming-weight class.

\textbf{Definition 3.2 (RNR Type-II).} An RNR Type-II coder predicts a
positional marginal field $Q_v(i) \approx \Pr(i \in M_v \mid C)$,
$v \in \{0,\dots,15\}$, $\sum_v Q_v(i) = 1$, and stores a partition
residual $R_\Pi$ recovering the actual partition
$\Pi(X) = (M_0,\dots,M_{15})$ into disjoint position sets with
$\sum_v |M_v| = N$; the partition determines $X$.

\textbf{Definition 3.3 (RNR Type-III-A).} An RNR Type-III-A coder predicts
an entropy field $U_m(t)$ over tiles $t$ and modes $m$, selects
$\hat{m}(t) = \min\{m : U_m(t) = \min_{m'} U_{m'}(t)\}$ with a fixed
canonical tie-break, and stores per-tile residual streams plus override
flags where the encoder departs from the predicted argmin. The $U_m(t)$
must be computed in the same deterministic representation on encoder and
decoder for reproducibility (Theorem 10.6).

\textbf{Definition 3.4 (RNR Type-III-B).} An RNR Type-III-B coder predicts
a normalised byte tensor $Q_{h,l}(i) \approx \Pr(h_i=h, l_i=l \mid C)$
over high/low nibbles, $\sum_{h,l} Q_{h,l}(i) = 1$, and a
decoder-reproducible mode choice selects one of three coding strategies
(the list (i)/(ii)/(iii) restated in P.4). The archive stores enough
residual to reconstruct $X$.

\textbf{Definition 3.5 (RNR Type-III-C).} An RNR Type-III-C coder is
parameterised by a learned token dictionary
$D = \{\tau_0,\dots,\tau_{K-1}\}$, $K = a\cdot b$, hierarchically indexed by
family and variant, predicts $Q(k,i) \approx \Pr(\tau_k\text{ begins at }i
\mid C)$, and stores a tokenization residual $R_C$ recovering the exact
token-interval cover.

\emph{The repair-length decomposition.} Every RNR coder has
\[
L_{\text{RNR}}(X) = L(M) + L(R \mid Y) + L(V) + L(\Theta),
\tag{3.1}
\]
with $L(M)$ the model description length (paid once, amortised over the
$m$ blocks encoded by the same model, effective per-block cost $L(M)/m$),
$L(R\mid Y)$ the repair-stream length, $L(V)$ verification overhead, and
$L(\Theta)$ metadata overhead. The conditional-advantage criterion is
$\mathbb{E}_D[L(M) + L(R\mid Y) + L(V) + L(\Theta)] < \mathbb{E}_D[B(X)]$.
Throughout we assume the amortised regime.

\textbf{P.3 Type-I converse baseline, universality, verification}

\textbf{Theorem 4.2 (Type-I converse and ideal achievability)} (proved in
[RNR-I]). \emph{Let $X \sim D$ over length-$N$ byte blocks and
$Y = M_I(C)$ be decoder-reproducible side information; fix an injective
Code12 map $E_{12}$ so that $E = E_{12}(X)\oplus Y$ is in bijection with $X$
given $Y$. Then (C, converse) any uniquely decodable lossless code whose
decoder uses only the codeword and $Y$ has expected length at least
$H(X\mid Y)$; and (A, ideal achievability) an ideal class-and-payload
coder driven by the true conditional law of $E$ achieves
$H(X\mid Y) + \varepsilon_N$ plus the $O(1)$ headers
$L(M_I)+L(V)+L(\Theta)$.} A concrete coder with an approximate class model
or a uniform within-class payload pays the corresponding cross-entropy,
matching the bound only when the models match the true conditional law.

\textbf{Theorem 8.1 (No lossless universal compression)} (proved in
[RNR-I]). \emph{No injective map from $\{0,1\}^n$ to binary strings of
length less than $n$ assigns every length-$n$ input a strictly shorter
codeword.} (Pigeonhole: $2^n$ inputs, $2^n - 1$ shorter strings.) RNR's
claim is therefore conditional and distributional, never a universal
dominance; Corollary 7.31a below invokes this baseline.

\textbf{Theorem 9.1 (No cross-subblock context propagation, with explicit
framing)} (proved in [RNR-I]). \emph{Assume the archive header stores, per
subblock $i$, an explicit offset into the repair stream (framing), so
subblock $i{+}1$'s slice is parsed independently of the bytes recovered in
subblock $i$; suppose the decoder uses the verified-only context-update
rule --- extend the predictor context with $\hat{X}_i$ only when
$H_v(\hat{X}_i) = v_i$, else abort or mark unverified --- and $H_v$ is a
cryptographic hash (keyed with a per-archive random key, or a random-oracle
abstraction) with collision probability at most $\delta$. Then the
probability that an incorrect $\hat{X}_i \ne X_i$ is used to update the
predictor context or repair-stream parsing for subblock $i{+}1$ is at most
$\delta$.} For truncated SHA-256 to $t$ bits, $\delta \le 2^{-t}$; a 64-bit
truncation costs 8 bytes per subblock. This theorem is used at Theorems
7.16 and 7.24 and in Section 10.8.

\textbf{P.4 Cold-context lemmas (from [RNR-I, Section 5]; load-bearing for
Theorems 7.2/7.3/7.29 and Section 10.7)}

These five lemmas characterise the entropy-rate achievability and the
random-access tax of per-position sequential RNR coding. They are the
backbone of Theorems 7.2--7.5, of the sub-block converse invoked by Remark
7.29a, and of the Section 10.7 sync-overhead accounting. Definitions 10.2
($W$-bounded predictor) and 10.3 (sync point) referenced below are restated
in the body of Section 10.7.

\textbf{Lemma 5.1a (Per-position sequential recoding achieves the
$W$-Markov rate plus sync overhead, with byte-granular random access)}
(proved in [RNR-I]). \emph{Let $X \sim D_N$ over a $\sigma$-symbol
alphabet and $M$ a $W$-bounded predictor outputting the order-$W$ Markov
conditional $P(x_i \mid x_{i-W},\dots,x_{i-1},C)$. Without sync points,
per-position sequential arithmetic coding achieves expected total length}
\[
H_W(D_N) := \sum_{i=1}^N H(X_i \mid X_{\max(1,i-W)},\dots,X_{i-1}, C)
\le N \log_2 \sigma,
\]
\emph{plus $\varepsilon_N \cdot N$ redundancy. With sync points at spacing
$K$ (Definition 10.3) it additionally pays
$\Delta_{\mathrm{sync,total}} \le m \cdot S \cdot W \cdot \log_2(1/\eta)$
bits (the per-sync cold-context excess of Section 10.7, $S$ symbols/byte,
$\eta$ the probability floor, $m$ the sync count). In general
$H_W(D_N) \ge H(D_N)$ with equality iff $D_N$ is order-$W$ Markov given
$C$; combined with sync points the coding admits byte-granular random
access via Theorem 10.3 at $O(\log L(R) + (K+k)\cdot\mathrm{cost}(M))$ per
query.}

\textbf{Lemma 5.1b (Logarithmic $W$ suffices under exponential
conditional-MI decay)} (proved in [RNR-I]). \emph{Let $X$ be strictly
stationary on a finite alphabet with
$I(X_n; X_{1:n-W-1} \mid X_{n-W:n-1}) \le C\rho^W$ for constants $C>0$,
$\rho\in(0,1)$. Then $0 \le H_W(D_N) - H(D_N) \le N C \rho^W$, and choosing
$W_N := \max(1, \lceil \log_{1/\rho}(NC)\rceil)$ gives
$H_{W_N}(D_N) - H(D_N) \le 1$; the Lemma 5.1a scheme then achieves total
expected length $H(D_N) + O(1) + \Delta_{\mathrm{sync,total}}(W_N) +
\varepsilon_N N$.} The overhead is sub-linear only when sync-point spacing
$K = \omega(\log N)$. The hypothesis holds strictly for finite-order Markov
sources and is plausible for geometrically ergodic HMMs.

\textbf{Lemma 5.1c (Tight cold-context excess for order-$W'$ Markov,
$W' \le W$)} (proved in [RNR-I]). \emph{For strictly stationary order-$W'$
Markov $X$ and a $W$-bounded predictor with $W \ge W'$ outputting the true
conditional, the expected cold-context excess per sync point equals
exactly}
\[
\delta_{W'} := \sum_{i=1}^{W'}[H(X_i \mid X_1,\dots,X_{i-1}) - h(X)]
= H(X^{W'}) - W' h(X) \ge 0,
\]
\emph{a source-specific constant independent of $W$ (for $W \ge W'$) and of
$K$ (for $K \ge W'$); for order-1 Markov, $\delta_1 = H_{\mathrm{marg}} -
h(X) = E$. With $m$ sync points and $N = mK$, the total sync excess over
$H(D_N)$ is exactly $L_{\mathrm{synced}} - H(D_N) = (m-1)\delta_{W'}$.}

\textbf{Lemma 5.1d (Cold-context excess for exponentially mixing
stationary sources)} (proved in [RNR-I]). \emph{For strictly stationary
$X$ satisfying the decay (5.8) with $W$-bounded conditional rate
$h_W := H(X_{W+1}\mid X_1,\dots,X_W)$, the expected sub-block length on
$K \ge W$ satisfies $\mathbb{E}[L_{\mathrm{sub}}(K)] = K h_W +
\delta_W^{(h_W)} \le K h_W + \delta_\infty$, where
$\delta_W^{(h_W)} := \sum_{i=1}^W [H(X_i\mid X_{<i}) - h_W] \ge 0$ and
$\delta_\infty := \sum_{n\ge 1}[H(X_n\mid X_{<n}) - h(X)]$ is the
Crutchfield--Feldman excess entropy.} With $W = W_N$ the total excess over
$H(D_N)$ is $(m-1)\delta_\infty + O(1)$. Lemma 5.1c is the order-$W'$-Markov
special case $\delta_\infty = \delta_{W'}$.

\textbf{Lemma 5.1e (Matching information-theoretic lower bound for
sub-block byte-granular random access)} (proved in [RNR-I]). \emph{For any
byte-granular scheme that (i) encodes the archive as $m$ independently
decodable sub-blocks of size $K$ ($N = mK$), (ii) codes each sub-block's
payload as a uniquely decodable codeword for its realised data alone, and
(iii) uses only source-independent shared metadata, one has
$L_{\mathrm{synced}} \ge m\cdot H(D_K) = m\cdot H(X_1,\dots,X_K)$; for
order-$W'$ Markov $X$ with $K \ge W'$ this is matched by Lemma 5.1c at
$L_{\mathrm{synced}} = m(K h(X) + \delta_{W'})$, so the total excess
$(m-1)\delta_{W'}$ is information-theoretically tight.} Condition (iii) is
essential: a source-dependent global header can evade the naive bound.

\emph{Corollary 5.1f (Open-Problem-1 resolution within the sub-block
paradigm)} (proved in [RNR-I]) combines Lemmas 5.1a--e: for stationary
exp-mixing sources, per-position sequential coding with $W = O(\log N)$ and
$K \ge W$ achieves total length $H(D_N) + (m-1)\delta_W + o(N)$, and this
$(m-1)\delta_W$ excess is the information-theoretically optimal cost of
byte-granular random access under the sub-block paradigm.

\textbf{P.5 Type-III-B coding strategies (from [RNR-I, Section 6.3];
needed by Section 10.7)}

Type-III-B extends Type-II to byte alphabets by modelling the high and low
nibbles jointly, $Q_{h,l}(i) \approx \Pr(h_i=h, l_i=l \mid C)$. Three
coding strategies are available: \textbf{(i)} \emph{direct partition
coding} treats the alphabet as 256 flat classes; \textbf{(ii)}
\emph{hierarchical partition coding} first encodes the high-nibble
partition $M^H_h = \{i : h_i = h\}$, then the low-nibble conditional
partition $M^{L\mid H=h}_l = \{i \in M^H_h : l_i = l\}$ per $h$ (preferred
when $Q_{h,l}(i) = Q^H_h(i)\cdot Q^{L\mid H}_l(i,h)$); and \textbf{(iii)}
\emph{per-byte sequential coding} encodes, at each position $i$ in source
order, the high nibble under $Q^H(\cdot\mid i,C)$ then the low nibble under
$Q^{L\mid H}(\cdot\mid i, h_i, C)$ via per-position arithmetic coding.
Strategies (i)--(ii) achieve the multi-information rate but use whole-block
combinatorial-rank encodings; strategy (iii) is the per-position sequential
stream that admits Theorem 10.3's byte-granular random access. Section 10.7
refers to these three strategies by number.

\emph{Companion note.} Section 10.7 also draws on the Section 6.9 cluster
(Lemma 6.9, Remark 6.9a, Corollary 6.9b) and Theorem 5.4, all proved in
[RNR-I]; the Section 7.15 arc invokes the Section 6.9 high-order-entropy
cluster of [RNR-I]. Where these are named in the body below, the reader
should consult [RNR-I].

\textbf{Random-access and deviation theory: the Theorem 7.2 arc}

\textbf{Theorem 7.2 (RNR achieves the Shannon entropy rate with
byte-granular random access on stationary mixing sources).}
\emph{Let $X = (X_n)_{n \geq 1}$ be a strictly stationary process on
a finite alphabet satisfying the exponential conditional-MI decay
(5.8) of Lemma 5.1b. Assume the following encoding conventions:
(C1) the predictor $M$ is treated as a pre-shared model with}
$L(M) = O(1)$ \emph{amortized into a fixed header (or, equivalently,
counted separately in $L(M)$ overhead rather than in the repair
stream); (C2) the predictor outputs an arithmetic-coder-compatible
probability distribution with a uniform floor $\eta > 0$ such that
the per-sub-block-stream cross-entropy excess over the true conditional
is at most $\varepsilon_{\mathrm{ac}}^{(j)} = O(\log(1/\eta))$ for each
of the $m$ independently-coded sub-block streams (the standard practical
assumption under high-quality predictors, e.g.\ Cover \& Thomas 2006
Theorem 13.4 on arithmetic coding with finite-precision floors), so the
total $\varepsilon_{\mathrm{ac}}^{\mathrm{total}} := \sum_j \varepsilon_{\mathrm{ac}}^{(j)}
= O((N/K) \cdot \log(1/\eta))$; (C3) auxiliary metadata $L(\Theta)$ is $\mathrm{polylog}(N)$.
Then for any sub-block spacing $K$ in the range
$W_N \leq K = \mathrm{polylog}(N)$ (with
$W_N = O(\log N)$ chosen per Lemma 5.1b), the RNR framework
(per-position sequential coding under a $W_N$-bounded predictor
outputting true-conditional probabilities clipped to the $\eta$-floor,
sync points at uniform $K$-spacing, and Theorem 10.3 byte-granular
access) achieves expected total repair-stream length}
\[
L^{\mathrm{rep}}_{\mathrm{RNR}}(X^N) \;=\; N \cdot h(X) + o(N) + \varepsilon_{\mathrm{ac}}^{\mathrm{total}},
\]
\emph{with per-query random-access cost $O(\log L(R) + (K + k) \cdot
\mathrm{cost}(M))$, polylog in $N$ when $K, k, \mathrm{cost}(M),
\log L(R)$ are all polylog in $N$. In particular, the per-position
repair-stream rate $L^{\mathrm{rep}}_{\mathrm{RNR}}/N \to h(X)$ as
$N \to \infty$, attaining the Shannon entropy rate (under the
encoding conventions C1-C3).}

\emph{\textbf{Proof.}} The RNR repair stream decomposes into
$m = \lceil N/K \rceil$ independent sub-block codings, each encoded
by the $W_N$-bounded predictor with arithmetic coding. Per sub-block
of length $K$, the expected bit cost is $H_{W_N}(D_K) +
\varepsilon_{\mathrm{ac}}^{(j)}$ with
$\varepsilon_{\mathrm{ac}}^{(j)} = O(\log(1/\eta))$ per stream
(convention C2). Summing,
\[
L^{\mathrm{rep}}_{\mathrm{RNR}}(X^N)
\;=\;
m \cdot H_{W_N}(D_K) + \varepsilon_{\mathrm{ac}}^{\mathrm{total}},
\quad
\varepsilon_{\mathrm{ac}}^{\mathrm{total}} = O\bigl((N/K) \cdot \log(1/\eta)\bigr).
\]

(i) [Per-sub-block bound via Lemma 5.1d] Since $K \geq W_N$, Lemma 5.1d
under exp-mixing (5.8) gives
$H_{W_N}(D_K) = K \cdot h_W + \delta_W^{(h_W)}$
with $\delta_W^{(h_W)} \leq \delta_\infty < \infty$ and $h_W := h_{W_N}$
the $W_N$-bounded conditional rate. Equivalently,
\[
H_{W_N}(D_K) \;\leq\; K \cdot h(X) + K \cdot (h_W - h) + \delta_\infty.
\]

(ii) [Rate convergence via Lemma 5.1b] Under exp-mixing (5.8) with
constants $C, \rho$ and the choice $W_N = O(\log N)$ from Lemma 5.1b,
$h_W - h \leq C \rho^{W_N} = O(1/N)$ (Lemma 5.1d
remark). Hence the per-sub-block W-bounded rate excess summed
over $m$ sub-blocks satisfies $m \cdot K \cdot (h_W - h) =
N \cdot (h_W - h) = O(1) = o(N)$.

(iii) [Total bit cost] Summing the per-sub-block bound of (i):
\[
m \cdot H_{W_N}(D_K) \;\leq\; N \cdot h(X) + N \cdot (h_W - h) + m \cdot \delta_\infty.
\]
With $K \geq W_N = \Omega(\log N)$ and $\delta_\infty < \infty$,
$m \cdot \delta_\infty = (N/K) \cdot \delta_\infty = O(N/\log N) = o(N)$.
With $\log(1/\eta) = O(1)$ (convention C2),
$\varepsilon_{\mathrm{ac}}^{\mathrm{total}} = O(N/K) = O(N/\log N) = o(N)$.
Combined with (ii):
\[
L^{\mathrm{rep}}_{\mathrm{RNR}}(X^N) \;\leq\; N \cdot h(X) + o(N),
\]
and dividing by $N$: $L^{\mathrm{rep}}_{\mathrm{RNR}}(X^N) / N \to h(X)$.

(iv) [Per-query RA cost] By Theorem 10.3, the byte-granular
random-access cost is $O(\log L(R) + (K + k) \cdot \mathrm{cost}(M))$
per query, polylog in $N$ when $K, k, \mathrm{cost}(M), \log L(R) =
\mathrm{polylog}(N)$.
\hfill$\blacksquare$

\emph{Scope and significance.} Theorem 7.2 is the headline asymptotic
claim of the RNR framework: \emph{any} stationary mixing source --- the
canonical natural-data setting (Lemma 5.6c covers order-$W$ Markov
chains, Lemma 5.6d covers HMMs, Lemma 5.6b covers latent-class
mixtures, all under exp-mixing) --- is compressed to its Shannon
entropy rate $h(X)$ asymptotically by RNR, \emph{with byte-granular
random access} at polylog per-query cost. This unifies the four
TC characterizations of §5.6 with the random-access framework of
§5.1 and Theorem 10.3, and demonstrates that RNR's framework
operates at the information-theoretic limit on natural-data sources
while preserving the random-access property that distinguishes it
from sequential-only compressors (NNCP, PixelCNN++, etc.).
\emph{Open extensions:} non-stationary sources, non-summable mixing,
and sources outside the exp-CMI-decay class fall outside this
theorem's scope; for such sources weaker (or no) Shannon-rate
attainability is currently provable. Corollary 7.2a (below) extends
to polynomial-mixing sources with a sub-linear (rather than polylog)
RA cost trade-off; Lemma 5.6f extends to AMS non-stationary sources
in a different sense (TC linearity, not Shannon-rate attainability).

\textbf{Corollary 7.2a (Polynomial-mixing extension of Theorem 7.2:
mixing-vs-RA-cost trade-off).} \emph{Let $X = (X_n)_{n \geq 1}$ be
a strictly stationary process on a finite alphabet satisfying the
\emph{polynomial} conditional-MI decay (replacing exp-mixing (5.8))}
\[
I\bigl(X_n; X_{1:n-W-1} \mid X_{n-W:n-1}\bigr) \;\leq\; C \cdot W^{-c}
\]
\emph{for constants $C > 0$ and $c > 1$ (the summability threshold
for Crutchfield--Feldman excess $\delta_\infty < \infty$). Then the
$W_N$ choice analogous to Lemma 5.1b becomes
$W_N := \lceil (C \cdot N)^{1/c} \rceil$, ensuring
$H_{W_N}(D_N) - H(D_N) \leq 1$, and Theorem 7.2 holds with the
modified sub-block range}
\[
W_N \;=\; \Theta(N^{1/c}) \;\leq\; K \;\leq\; \mathrm{poly}(N),
\]
\emph{giving expected total repair-stream length
$L^{\mathrm{rep}}_{\mathrm{RNR}}(X^N) = N \cdot h(X) + o(N) +
\varepsilon_{\mathrm{ac}}^{\mathrm{total}}$ as in Theorem 7.2, but
with per-query random-access cost
$O(\log L(R) + (K+k) \cdot \mathrm{cost}(M)) = \Theta(N^{1/c})$
when $K = \Theta(N^{1/c})$ at the lower edge of the sub-block range,
which is \emph{sub-linear} in $N$ but \emph{not} polylog.}

\emph{Trade-off characterisation.} As the mixing exponent $c$
decreases from $\infty$ (exp-mixing) toward $1^+$ (barely summable):
RA cost grows as $K^* \sim N^{1/c}$:
\begin{itemize}
\item $c \to \infty$ (exp-mixing): $K^* \to W_N = O(\log N)$, RA cost
polylog. Theorem 7.2 regime.
\item $c = 2$ (quadratic mixing): $K^* = \Theta(\sqrt{N})$, RA cost
$\Theta(\sqrt{N})$.
\item $c = 3$: $K^* = \Theta(N^{1/3})$, RA cost $\Theta(N^{1/3})$.
\item $c \to 1^+$ (barely summable): $K^* \to N^{1-\varepsilon}$,
RA cost approaches whole-block (sequential coding).
\end{itemize}
\emph{Sources with $c \leq 1$ (non-summable conditional-MI excess):
$\delta_\infty$ diverges, and Theorem 7.2's $m \cdot \delta_\infty = o(N)$
analysis fails. Such sources require either weaker Shannon-rate
attainability claims (e.g., asymptotic-rate only, not bounded excess)
or different proof techniques.}

\emph{\textbf{Proof.}} The polynomial-mixing hypothesis replaces the
exp-mixing constants $(C, \rho)$ with $(C, c)$. The Lemma 5.1b proof
becomes:
$H_W(D_N) - H(D_N) \leq N \cdot C \cdot W^{-c}$ (replacing $C \cdot \rho^W$).
Choosing $W_N := \lceil (C N)^{1/c} \rceil$ gives the bound
$\leq C \cdot N / (CN) = 1$, matching Lemma 5.1b's $\leq 1$ conclusion.

The Theorem 7.2 proof carries through verbatim with this $W_N$ choice,
replacing $\log N$ scaling of $W_N$ by $N^{1/c}$. Total repair-stream
length analysis is unchanged: $N \cdot h(X) + m \cdot \delta_\infty +
O(1) + \varepsilon_{\mathrm{ac}}^{\mathrm{total}}$ with
$m \cdot \delta_\infty = (N/K) \cdot \delta_\infty = o(N)$ for
$K = \omega(1)$ (any growing $K$). Crucially the constraint
$K \geq W_N$ from Lemma 5.1d's hypothesis now gives
$K \geq N^{1/c}$ rather than $K \geq \log N$, raising the lower edge
of the K range. Per-query cost from Theorem 10.3 is
$O(\log L(R) + (K+k) \cdot \mathrm{cost}(M))$, which at $K = N^{1/c}$
is $\Theta(N^{1/c})$.
\hfill$\blacksquare$

\emph{Scope.} Corollary 7.2a closes the ``non-summable mixing''
open-extension hook of Theorem 7.2 in the partial sense: it shows
that polynomial mixing (sub-class of non-exponential, sub-class of
non-summable) preserves Shannon-rate attainability at the cost of
sub-linear (not polylog) RA queries. Non-summable mixing ($c \leq 1$)
remains genuinely open --- the Crutchfield--Feldman excess diverges
and the $m \cdot \delta_\infty$ accounting fails.

\textbf{Theorem 7.3 (Information-theoretic tightness of Theorem 7.2:
matching lower bound).} \emph{Under the hypotheses of Theorem 7.2
(stationary $X$ with exp-mixing (5.8), sub-block spacing
$W_N \leq K = \mathrm{polylog}(N)$, conventions C1--C3) and the
additional assumption that any shared metadata is source-independent
(Lemma 5.1e condition (iii)), any byte-granular random-access coding
scheme that emits $m = \lceil N/K \rceil$ independently-decodable
sub-blocks (Lemma 5.1e conditions (i)--(ii)) has expected
repair-stream length}
\[
L^{\mathrm{rep}}(X^N) \;\geq\; m \cdot H(D_K) \;=\; N \cdot h(X) + m \cdot \delta_K,
\]
\emph{where $\delta_K := \sum_{n=1}^K [H(X_n \mid X_{<n}) - h(X)] \geq 0$
is the cumulative conditional-entropy excess at length $K$
(Lemma 5.6c). Combined with the Theorem 7.2 upper bound
$L^{\mathrm{rep}}_{\mathrm{RNR}}(X^N) \leq N \cdot h(X) +
m \cdot \delta_\infty + O(1) + \varepsilon_{\mathrm{ac}}^{\mathrm{total}}$
and the convergence $\delta_\infty - \delta_K = O(\rho^K)$ under
exp-mixing, the matching upper and lower bounds coincide modulo
the inherent finite-precision arithmetic-coding slack:}
\[
L^{\mathrm{rep}}_{\mathrm{RNR}}(X^N) - L^{\mathrm{rep}}_{\min}(X^N)
\;=\; O(1) + \varepsilon_{\mathrm{ac}}^{\mathrm{total}} + o(1),
\]
\emph{where $L^{\mathrm{rep}}_{\min}$ denotes the lower bound above.
Hence RNR is information-theoretically optimal among byte-granular
random-access schemes on exp-mixing stationary sources, up to the
arithmetic-coding finite-precision overhead.}

\emph{\textbf{Proof.}} \emph{Lower bound.} Conditions (i)--(iii) of
Lemma 5.1e are satisfied by hypothesis; the lemma gives
$L^{\mathrm{rep}}(X^N) \geq m \cdot H(D_K)$. By Lemma 5.6c (stationary
$X$, conditional entropy chain rule), $H(D_K) = K \cdot h(X) + \delta_K$
with $\delta_K \geq 0$. Multiplying by $m = \lceil N/K \rceil$ and
using $mK = N$ (or $N + O(K)$ for the ceiling case, absorbed in $O(1)$),
$m \cdot H(D_K) = N \cdot h(X) + m \cdot \delta_K$.

\emph{Convergence rate.} Under exp-mixing (5.8), by Crutchfield--Feldman
summability $\delta_K \to \delta_\infty < \infty$ at exponential rate:
$\delta_\infty - \delta_K = \sum_{n > K} [H(X_n \mid X_{<n}) - h(X)]
\leq C \rho^K / (1 - \rho)$ for the same constants $C, \rho \in (0, 1)$
of (5.8) (the tail of the Crutchfield--Feldman series inherits the
exp-mixing rate). At $K \geq W_N$ with Lemma 5.1b's choice
$W_N = \lceil \log_{1/\rho}(N C) \rceil$, $\rho^K \leq \rho^{W_N} \leq
1/(N C)$, hence
\[
m \cdot (\delta_\infty - \delta_K)
\;\leq\; (N/K) \cdot \tfrac{1}{N (1 - \rho)}
\;=\; \tfrac{1}{K (1 - \rho)}
\;=\; O(1/\log N) \;=\; o(1).
\]

\emph{Match.} Theorem 7.2 gives the upper bound
$L^{\mathrm{rep}}_{\mathrm{RNR}} \leq N \cdot h(X) + m \cdot \delta_\infty +
O(1) + \varepsilon_{\mathrm{ac}}^{\mathrm{total}}$. The lower bound is
$L^{\mathrm{rep}}_{\min} = N \cdot h(X) + m \cdot \delta_K$. Their
difference is $m \cdot (\delta_\infty - \delta_K) + O(1) +
\varepsilon_{\mathrm{ac}}^{\mathrm{total}} = o(1) + O(1) +
\varepsilon_{\mathrm{ac}}^{\mathrm{total}}$, as claimed.
\hfill$\blacksquare$

\emph{Scope.} Theorem 7.3 is the matching lower bound for Theorem 7.2:
the upper-lower bound pair coincides at the $m \cdot \delta_\infty$
excess level, not merely at the leading $N \cdot h(X)$ order. The
gap is the inherent arithmetic-coding finite-precision slack
$\varepsilon_{\mathrm{ac}}^{\mathrm{total}}$, which can be made
$O(\log(1/\eta))$-per-sub-block via the standard floor convention (C2).
This generalizes Corollary 5.1f's tight matching from order-$W'$
Markov sources to all exp-mixing stationary sources within the
Theorem 7.2 framework, and demonstrates that the
$(m-1) \cdot \delta_\infty$ ``sync overhead'' identified in Lemma 5.1d
is information-theoretically irreducible --- it cannot be removed by
any byte-granular RA scheme satisfying Lemma 5.1e's independence
conditions. \emph{Open question:} if source-dependent shared metadata
is allowed (relaxing Lemma 5.1e (iii)), can the sync overhead be
reduced below $(m-1) \cdot \delta_\infty$?

\emph{Attempted closure (Theorem 7.3a), retracted.} A natural Shannon-converse
attempt would lower-bound $L^{\mathrm{rep}}$ by $m \cdot H(D_K) -
L(\Theta)$ via the sub-additivity claim $\sum_j I(X_{\text{sub-block } j};
\Theta) \leq I(D_N; \Theta) \leq L(\Theta)$. \emph{This sub-additivity
of mutual information FAILS for general dependent sub-blocks}: e.g.,
$D_N = (X_1, X_2)$ independent bits and $\Theta = X_1 \oplus X_2$
gives $I(X_1; \Theta) = I(X_2; \Theta) = 0$ but $I(D_N; \Theta) =
H(\Theta) = 1$ --- the SUM is strictly LESS than the joint MI, not
$\leq$. Using the alternative trivial bound $\sum_j I(D_K^j; \Theta)
\leq m \cdot L(\Theta) = \Theta(N)$ for $K = \mathrm{polylog}(N)$
makes the resulting lower bound vacuous (negative). The
Shannon-converse route does not close the question; the answer
remains \emph{genuinely open} whether source-dependent metadata
can reduce sync overhead from $m \cdot \delta_\infty$ to a smaller
quantity. We document the failed attempt here to prevent recurrence
of the same approach by future readers, and leave the question for
§13.2.

\textbf{Theorem 7.4 (Robustness of Theorem 7.2 to non-ideal predictor:
explicit KL-divergence accounting).} \emph{Let $X$ satisfy the
hypotheses of Theorem 7.2 (stationary, exp-mixing (5.8),
$W_N \leq K = \mathrm{polylog}(N)$, conventions C1--C3). Let $M'$ be
any (possibly learned) $W_N$-bounded predictor with expected per-position
Kullback--Leibler divergence}
\[
\varepsilon_{\mathrm{KL}}
\;:=\;
\mathbb{E}_{X \sim D_N}\Bigl[\,\tfrac{1}{N} \sum_{i=1}^N
  D_{\mathrm{KL}}\!\bigl(P_{\mathrm{true}}(\cdot \mid X_{i-W_N : i-1})
  \,\big\|\, P_{M'}(\cdot \mid X_{i-W_N : i-1})\bigr)
\,\Bigr] \;\geq\; 0
\]
\emph{from the true $W_N$-bounded conditional, averaged across positions
(by stationarity, equivalent to a per-position divergence; we additionally
require $\varepsilon_{\mathrm{KL}} < \infty$, ensured in practice by
$M'$ inheriting convention C2's uniform $\eta$-floor so that
$P_{M'}(x \mid \cdot) \geq \eta > 0$ for every symbol). Then RNR
encoding with $M'$ instead of the ideal predictor achieves expected
total repair-stream length}
\[
L^{\mathrm{rep}}_{\mathrm{RNR}}(X^N; M')
\;\leq\;
N \cdot h(X) + N \cdot \varepsilon_{\mathrm{KL}} + m \cdot \delta_\infty + O(1) + \varepsilon_{\mathrm{ac}}^{\mathrm{total}}.
\]
\emph{Consequently:}

\emph{(a) Asymptotically ideal predictor ($\varepsilon_{\mathrm{KL}} =
o(1)$ as $N \to \infty$): $L^{\mathrm{rep}}_{\mathrm{RNR}} / N \to h(X)$;
Shannon rate is attained.}

\emph{(b) Constant-gap predictor ($\varepsilon_{\mathrm{KL}} = c > 0$):
$L^{\mathrm{rep}}_{\mathrm{RNR}} / N \to h(X) + c$; the rate exceeds
Shannon by exactly the predictor's cross-entropy gap.}

\emph{\textbf{Proof.}} For any predictor $M'$ and any position $i$,
the expected arithmetic-coding bit cost using $M'$'s distribution is
\[
\mathbb{E}[-\log_2 P_{M'}(X_i \mid X_{i-W_N : i-1})]
= H(X_i \mid X_{i-W_N : i-1})
  + D_{\mathrm{KL}}(P_{\mathrm{true}} \,\|\, P_{M'}),
\]
the exact cross-entropy decomposition (Cover \& Thomas 2006 §2.6 on
relative entropy / §5.4 on suboptimal codes). Summing
over all $N$ positions within their respective sub-blocks:
\[
\mathbb{E}[L^{\mathrm{rep}}_{\mathrm{RNR}}(X^N; M')]
= m \cdot H_{W_N}(D_K) + N \cdot \varepsilon_{\mathrm{KL}} + \varepsilon_{\mathrm{ac}}^{\mathrm{total}}.
\]
By the Theorem 7.2 analysis,
$m \cdot H_{W_N}(D_K) \leq N \cdot h(X) + m \cdot \delta_\infty + O(1)$
(combining Lemma 5.1d's per-sub-block bound with Lemma 5.1b's
$h_{W_N} - h(X) = O(1/N)$). Adding the $N \cdot \varepsilon_{\mathrm{KL}}$
term gives the claimed bound. The two corollaries follow by dividing
by $N$ and taking $N \to \infty$: $m \cdot \delta_\infty / N = O(1/\log N) =
o(1)$ for $K \geq W_N = \Omega(\log N)$, $O(1)/N \to 0$, and
$\varepsilon_{\mathrm{ac}}^{\mathrm{total}} / N = o(1)$ under bounded
$\log(1/\eta)$.
\hfill$\blacksquare$

\emph{Scope and significance.} Theorem 7.4 bridges Theorem 7.2's
idealization (true conditional probabilities) to the practical setting
of learned predictors. For high-quality language models on natural text
the empirical per-position cross-entropy $H + \varepsilon_{\mathrm{KL}}$
is roughly $1$ bit/byte (compare Del{\'e}tang et al.\ 2024
\emph{Language Modeling Is Compression}, ICLR, arXiv:2309.10668, which
reports neural-network LM compression ratios on enwik9 in the
sub-1-bpb regime) versus Shannon's lower-bound estimate
$h(X) \approx 0.6$--$1.3$ bits/byte for English (Shannon 1951
\emph{Bell System Technical Journal} 30(1):50--64, ``Prediction and
Entropy of Printed English''; see Cover \& King 1978
\emph{IEEE Trans.\ Inf.\ Theory} 24(4):413 for the tightened gambling
estimator), giving $\varepsilon_{\mathrm{KL}}$ on the order of tenths
of a bit per byte in practice. As predictor quality improves with model
scale and training data, $\varepsilon_{\mathrm{KL}} \to 0$ (per the
universal-approximation property of neural networks
--- Cybenko 1989 \emph{Math.\ Control Signals Syst.}\ 2(4):303;
Hornik et al.\ 1989 \emph{Neural Networks} 2(5):359 --- combined with
consistency of MLE and the empirical scaling laws of Kaplan et al.\
2020 \emph{Scaling Laws for Neural Language Models}, arXiv:2001.08361,
and Hoffmann et al.\ 2022 \emph{Training Compute-Optimal Large Language
Models}, arXiv:2203.15556), bringing RNR's empirical rate arbitrarily
close to Shannon's lower bound $h(X)$. \emph{Practical interpretation:}
the gap $h(M') - h(X) = \varepsilon_{\mathrm{KL}}$ is the ``predictor
suboptimality cost'' borne by every position; reducing it by training
better LMs directly reduces $L^{\mathrm{rep}}_{\mathrm{RNR}}/N$ by the
same amount, position-for-position.

\textbf{Theorem 7.5 (Optimal sub-block spacing: storage--access
trade-off).} \emph{Under the Theorem 7.2 hypotheses, with the
polylog$(N)$ upper bound on $K$ relaxed to $K \geq W_N$ (the bit-cost
analysis of T7.2 holds without this upper bound; per-query cost
becomes $O(\log L(R) + K)$ rather than polylog). Let $Q \geq 1$ be
the expected number of byte-granular random-access queries against
the archive over its lifetime, $\alpha > 0$ the cost-weight of query
operations relative to storage bits, and let}
$\gamma := \delta_\infty + \log(1/\eta)$
\emph{denote the per-sub-block overhead constant (Crutchfield--Feldman
excess plus arithmetic-coding floor). The total expected operational
cost}
\[
C(K) \;:=\; \mathbb{E}[L^{\mathrm{rep}}_{\mathrm{RNR}}(K)]
        + \alpha \cdot Q \cdot \mathrm{cost}_{\mathrm{query}}(K)
       \;=\; N \cdot h(X) + \frac{\gamma \cdot N}{K} + \alpha \cdot Q \cdot K + O(1) + N \cdot \varepsilon_{\mathrm{KL}}
\]
\emph{is minimized at}
\[
K^* \;=\; \sqrt{\frac{\gamma \cdot N}{\alpha \cdot Q}},
\]
\emph{yielding optimal cost}
\[
C(K^*) \;=\; N \cdot h(X) + 2 \sqrt{\gamma \cdot N \cdot \alpha \cdot Q}
            + N \cdot \varepsilon_{\mathrm{KL}} + O(1).
\]
\emph{For balanced storage--access weight ($\alpha \cdot Q = \Theta(1)$),
the unconstrained optimum is $K^* = \Theta(\sqrt{N})$; the storage
overhead $C(K^*) - N \cdot h(X) - N \cdot \varepsilon_{\mathrm{KL}} =
\Theta(\sqrt{N})$ and the per-query cost is also $\Theta(\sqrt{N})$.
This is the AM--GM-optimal $\sqrt{N}$ balance between two reciprocal
cost terms, structurally analogous to the sampling-rate parameter in
compressed suffix arrays (Sadakane 2003 \emph{J.\ Algorithms} 48(2):294;
Grossi--Vitter 2005 \emph{SIAM J.\ Comput.}\ 35(2):378), where a
sampling rate $s$ gives $O(N/s)$ storage and $O(s \log N)$ query, yielding
$s = \Theta(\sqrt{N})$ at balance --- here the trade-off is expressed
in the RNR compression setting with $K$ playing the role of the
sampling parameter.}

\emph{Constraint on $K^*$.} The formula $K^* = \sqrt{\gamma N / (\alpha Q)}$
is the \emph{unconstrained} optimum and is valid only when $K^* \geq W_N$.
For extremely access-saturated workloads $\alpha Q > \gamma N / W_N^2$,
the formula would give $K^* < W_N$, at which point the lower constraint
$K \geq W_N$ from Theorem 7.2 binds and the constrained optimum is
$K = W_N$, with cost $C(W_N) = N \cdot h(X) + \gamma N / W_N +
\alpha Q \cdot W_N + N \varepsilon_{\mathrm{KL}} + O(1)$. The boundary
threshold $\alpha Q = \gamma N / W_N^2 = \Theta(N / \log^2 N)$ separates
the interior $\sqrt{N}$ regime from the constraint-bound regime.

\emph{\textbf{Proof.}} The bit-cost component of $C(K)$ follows from
the proof of Theorem 7.2 with $m = N/K$ sub-blocks: $L^{\mathrm{rep}}_{\mathrm{RNR}}
\leq N h(X) + m \delta_\infty + N \varepsilon_{\mathrm{KL}} + O(1) +
\varepsilon_{\mathrm{ac}}^{\mathrm{total}} = N h(X) + \gamma N/K +
N \varepsilon_{\mathrm{KL}} + O(1)$, absorbing $\varepsilon_{\mathrm{ac}}^{\mathrm{total}}
= O(N/K \cdot \log(1/\eta))$ into $\gamma N/K$. The per-query cost
$O(\log L(R) + K)$ from Theorem 10.3 contributes $\alpha Q (\log L(R) + K)$.
The $\alpha Q \log L(R)$ term is independent of $K$ (a $K$-independent
constant in the minimization), and the $\alpha Q \cdot K$ term is what
trades off against $\gamma N / K$.

Setting $\partial C / \partial K = 0$:
$-\gamma N / K^2 + \alpha Q = 0$, so $K^* = \sqrt{\gamma N / (\alpha Q)}$.
Second-order: $\partial^2 C / \partial K^2 = 2 \gamma N / K^3 > 0$ at
$K^*$, confirming a true minimum. Substituting:
$\gamma N / K^* + \alpha Q \cdot K^* = \sqrt{\gamma N \alpha Q} +
\sqrt{\gamma N \alpha Q} = 2 \sqrt{\gamma N \alpha Q}$, hence
$C(K^*) = N h(X) + 2 \sqrt{\gamma N \alpha Q} + N \varepsilon_{\mathrm{KL}} +
\alpha Q \log L(R) + O(1)$. The $\alpha Q \log L(R)$ is a $K$-independent
overhead from indexing into the repair stream; for $L(R) = \mathrm{poly}(N)$
it grows as $\alpha Q \log N$, which is lower-order than the
$\sqrt{N \alpha Q}$ trade-off term whenever $\alpha Q = o(N / (\log N)^2)$,
i.e., for all but extremely query-saturated regimes.
\hfill$\blacksquare$

\emph{Scope and significance.} Theorem 7.5 is the operational-cost
companion to Theorems 7.2--7.4. It establishes that RNR's natural
parameter $K$ is governed by a $\sqrt{N}$ time--space trade-off, with
the constant of proportionality determined by the source's
Crutchfield--Feldman excess $\delta_\infty$ and the operational
$\alpha \cdot Q$ ratio. Three regimes:

\emph{(a) Storage-dominated ($\alpha \cdot Q \ll 1$, archival writes
once, rarely read):} $K^* \gg \sqrt{N}$, sync overhead $\gamma N /
K^* \to 0$, near-Shannon-rate storage; tolerates slow RA.

\emph{(b) Access-dominated ($\alpha \cdot Q \gg 1$, hot data, many
queries):} $K^* \ll \sqrt{N}$, fast RA at cost of larger sync overhead;
storage rate $h(X) + \gamma / K^*$ exceeds Shannon by a constant.

\emph{(c) Balanced ($\alpha \cdot Q = \Theta(1)$):} $K^* = \Theta(\sqrt{N})$,
both storage overhead and per-query cost are $\Theta(\sqrt{N})$ ---
the classical sub-linear trade-off.

The polylog$(N)$ upper bound on $K$ in Theorem 7.2 corresponds to
$K^* = \mathrm{polylog}(N)$, which from the formula above requires
$\alpha \cdot Q = \Theta(N / \mathrm{polylog}^2(N))$; this is the
``query-heavy'' regime where the user expects $\Omega(N / \mathrm{polylog})$
queries to a roughly $N$-byte archive. The $\sqrt{N}$ regime
$K^* = \Theta(\sqrt{N})$ is the balanced setting most appropriate for
typical archive use.

\textbf{Theorem 7.6 (Universal online RNR without pre-trained predictor,
for source classes of polylog parametric dimension).}
\emph{Let $X$ satisfy the hypotheses of Theorem 7.2 (stationary
exp-mixing on finite alphabet, $W_N \leq K = \mathrm{polylog}(N)$), and
suppose $X$ is realisable by some predictor in a model class
$\mathcal{M}_d$ of parametric dimension $d = d(N) = \mathrm{polylog}(N)$
(e.g., order-$W'$ Markov for some $W' = O(\log \log N)$ giving $d =
|\Sigma|^{W'} = \mathrm{polylog}(N)$, or a low-rank HMM with bounded
hidden-state count, or any compactly-parameterized sequence model).
Then RNR with an \emph{online universal predictor} over $\mathcal{M}_d$
--- e.g., the Bayesian-mixture / Krichevsky--Trofimov estimator
(Krichevsky \& Trofimov 1981 \emph{IEEE Trans.\ Inf.\ Theory}
27(2):199--207 ``The performance of universal encoding''; Rissanen 1996
\emph{IEEE Trans.\ Inf.\ Theory} 42(1):40--47 ``Fisher information
and stochastic complexity'' for MDL-style universal-mixture regret;
Cover \& Thomas 2006 Ch.\ 13 ``Universal Source Coding'') ---
incurs regret at most}
\[
\mathrm{Regret}_{\mathcal{M}_d}(N) \;\leq\; d \cdot \log(N) + O(1),
\]
\emph{giving expected per-position cross-entropy excess
$\varepsilon_{\mathrm{KL}}(N) \leq \mathrm{Regret}/N =
O(d \log(N) / N)$. Combined with Theorem 7.4's KL accounting:}
\[
\mathbb{E}\bigl[L^{\mathrm{rep}}_{\mathrm{RNR}}(X^N;\,M'_{\text{universal}})\bigr]
\;\leq\;
N \cdot h(X) + O\bigl(d(N) \cdot \log N\bigr) + m \cdot \delta_\infty + O(1) + \varepsilon_{\mathrm{ac}}^{\mathrm{total}}.
\]
\emph{For $d(N) = \mathrm{polylog}(N)$: $d \cdot \log N = \mathrm{polylog}(N) =
o(N)$, so $L^{\mathrm{rep}}_{\mathrm{RNR}} = N \cdot h(X) + o(N)$,
attaining Shannon rate \emph{without any pre-trained predictor}
(no $L(M)$ header beyond the $O(1)$ mixture-encoder code).}

\emph{\textbf{Proof.}} The Krichevsky--Trofimov regret bound for
$d$-parameter exponential-family / order-$k$-Markov estimators is
classical (Krichevsky--Trofimov 1981 for Markov; Rissanen 1996
\emph{IEEE Trans.\ Inf.\ Theory} 42(1):40--47 ``Fisher information
and stochastic complexity'' for general MDL): with the Jeffreys'
prior (or a comparable smooth prior on $\mathcal{M}_d$), the online
Bayesian mixture predictor has expected cumulative log-loss within
$(d/2) \cdot \log(N) + O(1)$ of the offline optimum in $\mathcal{M}_d$.
For source $X \in \mathcal{M}_d$, the offline optimum equals
$N \cdot h(X)$ (the source's entropy rate, since the true predictor
is in the model class). Hence the per-position cross-entropy excess
$\varepsilon_{\mathrm{KL}} \leq d \cdot \log(N) / N$ (up to constants).

The universal-predictor estimator is itself encoded in $O(1)$ bits
(the mixture-estimator's algorithm description), so its $L(M)$
contribution is $O(1)$ amortised; this distinguishes it from a
pre-trained large-parameter predictor (which would have $L(M)
\sim d$ unless amortised over an exponential number of archives).

Substituting $\varepsilon_{\mathrm{KL}} = O(d \log N / N)$ into
Theorem 7.4:
$L^{\mathrm{rep}}_{\mathrm{RNR}} \leq N \cdot h(X) +
N \cdot \varepsilon_{\mathrm{KL}} + m \cdot \delta_\infty + O(1) +
\varepsilon_{\mathrm{ac}}^{\mathrm{total}}
\leq N \cdot h(X) + O(d \log N) + m \cdot \delta_\infty + O(1) +
\varepsilon_{\mathrm{ac}}^{\mathrm{total}}$.

For $d = \mathrm{polylog}(N)$: $d \log N = \mathrm{polylog}(N)$ is
$o(N)$. With Theorem 7.2's analysis $m \cdot \delta_\infty = o(N)$
and Theorem 7.2's $\varepsilon_{\mathrm{ac}}^{\mathrm{total}} = o(N)$:
the bound becomes $L^{\mathrm{rep}}_{\mathrm{RNR}} = N \cdot h(X) +
o(N)$. Shannon rate attained.
\hfill$\blacksquare$

\emph{Scope and limitations.} Theorem 7.6 closes the
``universal online RNR'' question for source classes of polylog
parametric dimension. The hypothesis $X \in \mathcal{M}_d$ with
$d = \mathrm{polylog}(N)$ is the key restriction: it requires the
source to be \emph{approximable} by a polylog-parameter model
(e.g., low-order Markov, sparse HMM, mixture of $\mathrm{polylog}$
simple components). Generic exp-mixing sources with no compact
parametric representation (e.g., requiring $|\Sigma|^{O(\log N)} =
\mathrm{poly}(N)$ Markov parameters) fall outside this scope --- for
such sources the universal-predictor regret is $\Omega(\log N)$
per position, dominating the Shannon term and preventing
rate-optimality without pre-training.

The result complements Theorem 7.4 (pre-trained predictor with
empirical $\varepsilon_{\mathrm{KL}}$) by giving an explicit regret
bound for the online-learning regime: pre-trained models pay $L(M)$
amortised but achieve potentially small $\varepsilon_{\mathrm{KL}}$;
online universal models pay no $L(M)$ but incur regret-based
$\varepsilon_{\mathrm{KL}} = O(d \log N / N)$. The trade-off is
between predictor storage and adaptation cost.

\textbf{Theorem 7.6$'$ (Tight minimax regret for universal online RNR:
the leading constant is $d/2$, not $d$, and it is unimprovable).}
\emph{Under the hypotheses of Theorem 7.6, with $\mathcal{M}_d$ a
$d$-dimensional smooth parametric family (an exponential family or
order-$W'$ Markov class with $d = |\Sigma|^{W'}(|\Sigma|-1)$
free parameters, satisfying the Clarke--Barron regularity conditions:
finite and nonsingular Fisher information $I(\theta)$ on a compact
parameter set, twice-differentiable log-likelihood), the online
universal predictor of Theorem 7.6 may be taken to be the
Bayesian mixture under Jeffreys' prior --- equivalently the
Krichevsky--Trofimov estimator on each Markov context --- and then
the regret of Theorem 7.6 sharpens to}
\[
\mathrm{Regret}_{\mathcal{M}_d}(N)
\;=\;
\tfrac{d}{2}\,\log_2 N \;+\; C_{\mathcal{M}_d} \;+\; o(\log N),
\qquad
C_{\mathcal{M}_d} = \log_2\!\!\int_\Theta\!\!\sqrt{\det I(\theta)}\,d\theta
- \tfrac{d}{2}\log_2(2\pi e) + o(1),
\]
\emph{with $C_{\mathcal{M}_d} = O(1)$ in $N$. Moreover this constant is
\emph{tight}: for any online RNR predictor whatsoever (any sequential
probability assignment $q$ over $\Sigma^N$, hence any prefix-code
RNR encoder), and for Lebesgue-almost every source parameter
$\theta \in \Theta$,}
\[
\mathrm{Regret}_{q}(N;\theta) \;\geq\; \tfrac{d}{2}\,(1-\varepsilon)\,\log_2 N
\quad\text{for every }\varepsilon>0\text{ and all sufficiently large }N.
\]
\emph{Consequently the online-RNR per-position cross-entropy excess is
$\varepsilon_{\mathrm{KL}}(N) = \tfrac{d}{2}\log_2 N / N \,(1+o(1))$,
exactly half the bound implied by Theorem 7.6's
$d\log N$, and no online scheme can do asymptotically better on a
typical source. The RNR archive size satisfies}
\[
\mathbb{E}\bigl[L^{\mathrm{rep}}_{\mathrm{RNR}}(X^N;M'_{\mathrm{universal}})\bigr]
= N h(X) + \tfrac{d}{2}\log_2 N + O(1)
+ m\,\delta_\infty + \varepsilon_{\mathrm{ac}}^{\mathrm{total}},
\]
\emph{the Shtarkov/Rissanen minimax pointwise-redundancy form.}

\emph{\textbf{Proof.}}
\emph{(Upper bound.)} Fix the Markov order $W'$ (the i.i.d.\ case is
$W'=0$). Within each of the $|\Sigma|^{W'}$ contexts the source emits
i.i.d.\ symbols over $\Sigma$, an exponential family of dimension
$|\Sigma|-1$; the joint family $\mathcal{M}_d$ has
$d = |\Sigma|^{W'}(|\Sigma|-1)$ parameters. Use the Bayesian mixture
$q(x^N) = \int_\Theta p_\theta(x^N)\,w_J(\theta)\,d\theta$ with the
Jeffreys prior $w_J(\theta) \propto \sqrt{\det I(\theta)}$; on a
product-of-Dirichlet parametrisation this is exactly the
Krichevsky--Trofimov rule
$q(a\mid\text{context }c,\text{history}) = (n_{c,a}+\tfrac12)/(n_c + |\Sigma|/2)$
applied per context (Krichevsky \& Trofimov 1981, \emph{IEEE
Trans.\ Inf.\ Theory} 27(2):199--207). Clarke \& Barron (1990,
\emph{IEEE Trans.\ Inf.\ Theory} 36(3):453--471, ``Information-theoretic
asymptotics of Bayes methods''), under the stated regularity
conditions, give the exact mixture redundancy
\[
\mathbb{E}_\theta\!\bigl[\log_2 p_\theta(X^N)/q(X^N)\bigr]
= \tfrac{d}{2}\log_2\tfrac{N}{2\pi e}
+ \log_2\!\!\int_\Theta\!\!\sqrt{\det I}\,d\theta
- \log_2 w_J(\theta) + o(1),
\]
and with the Jeffreys choice $w_J = \sqrt{\det I}/\!\int\!\sqrt{\det I}$
the $\theta$-dependent terms cancel, leaving the uniform constant
$C_{\mathcal{M}_d}$ above. Since the regret against the offline ML
optimum equals this redundancy up to $o(1)$ (the ML codelength and
$\log_2 p_\theta$ differ by $O(1)$ in expectation by the same
Laplace expansion), $\mathrm{Regret}_{\mathcal{M}_d}(N) =
\tfrac{d}{2}\log_2 N + C_{\mathcal{M}_d} + o(\log N)$. The
worst-case (pointwise) constant is the Shtarkov/NML value
$\tfrac{d}{2}\log_2\tfrac{N}{2\pi} + \log_2\!\int\!\sqrt{\det I}$
(Rissanen 1996, \emph{IEEE Trans.\ Inf.\ Theory} 42(1):40--47,
``Fisher information and stochastic complexity''; Shtarkov 1987),
which differs from the expected constant only by the $O(1)$ term
$\tfrac{d}{2}\log_2 e$ and likewise has leading coefficient $d/2$.
Substituting $\varepsilon_{\mathrm{KL}} = \tfrac{d}{2}\log_2 N / N
(1+o(1))$ into Theorem 7.4's KL accounting yields the stated archive
size. The mixture predictor is still described in $O(1)$ bits (the
KT/Jeffreys algorithm), so the $L(M)$ contribution remains $O(1)$,
preserving Theorem 7.6's no-pre-training conclusion.

\emph{(Lower bound / matching converse.)} The leading coefficient
$d/2$ cannot be reduced. By Rissanen's converse (Rissanen 1984,
\emph{IEEE Trans.\ Inf.\ Theory} 30(4):629--636, ``Universal coding,
information, prediction, and estimation''; sharpened by Merhav \&
Feder 1995, \emph{IEEE Trans.\ Inf.\ Theory} 41(3):714--722, ``A
strong version of the redundancy--capacity theorem of universal
coding''), for \emph{any} sequence of distributions $q$ on $\Sigma^N$
and every $\varepsilon>0$, the set of parameters
$\theta\in\Theta$ for which
$\mathbb{E}_\theta[\log_2 p_\theta(X^N)/q(X^N)] <
\tfrac{d}{2}(1-\varepsilon)\log_2 N$ infinitely often has Lebesgue
measure zero. Every online RNR predictor is precisely such a $q$
(the chain rule turns a sequential probability assignment into a
joint law on $\Sigma^N$, and conversely any prefix-code RNR encoder
induces a $q$ by Kraft's inequality), so the redundancy bound applies
verbatim. Finally the regret (excess over the offline ML optimum)
dominates this redundancy in expectation, since
$\max_{\theta'} p_{\theta'}(X^N) \geq p_\theta(X^N)$ gives
$\mathbb{E}_\theta[\log_2 \max_{\theta'} p_{\theta'}(X^N)/q(X^N)]
\geq \mathbb{E}_\theta[\log_2 p_\theta(X^N)/q(X^N)]$; hence
for almost every source the online-RNR regret is at least
$\tfrac{d}{2}(1-\varepsilon)\log_2 N$ eventually. Combined with the
upper bound, $d/2$ is the exact leading constant. The
redundancy--capacity identity (Merhav--Feder) further shows this
$d/2$ equals the normalised channel capacity
$\lim_N C_N/\log_2 N$ of the family $\{p_\theta\}$, the
information-theoretic price of \emph{not} knowing $\theta$.
\hfill$\blacksquare$

\emph{Scope and significance.} Theorem 7.6$'$ does three things beyond
Theorem 7.6. (i) It halves the regret constant from $d$ to $d/2$ by
committing to the Jeffreys/KT mixture rather than a generic
$O(1)$-per-bit universal estimator; the factor of two is not
cosmetic --- for an order-$W'$ Markov class with $d = |\Sigma|^{W'}
(|\Sigma|-1)$ parameters it removes $\tfrac{d}{2}\log_2 N$ bits from
the archive, which for $|\Sigma|=256$, $W'=2$, $N=10^9$ is on the
order of $10^7$ bits. (ii) It supplies a \emph{matching lower bound},
upgrading the result from an achievability statement to a
characterisation: $\varepsilon_{\mathrm{KL}} = \tfrac{d}{2}\log_2 N/N$
is optimal for online RNR on a typical source, so no cleverer
predictor (Vovk aggregation, mirror descent, follow-the-regularised-
leader) can beat the KT mixture by more than $o(\log N)$. (iii) It
identifies the constant as the Shtarkov/NML minimax value, locating
online RNR exactly within the universal-coding hierarchy: the
$d/2\log N$ term is the same parametric ``model cost'' that appears
in MDL stochastic complexity (§2.5), now paid online and amortised
into the repair stream rather than as a header. The verification
script \texttt{thm\_7\_6prime\_shtarkov\_tight\_regret.py} runs the KT
mixture on i.i.d.\ and order-$1$ Markov sources and confirms the
measured leading constant is $d/2$ (e.g.\ each doubling of $N$ adds
exactly $d/2$ bits of redundancy, not $d$). The limitation inherited
from Theorem 7.6 is unchanged: the result needs $d=\mathrm{polylog}(N)$
for the $\tfrac{d}{2}\log N$ term to stay $o(N)$, and the converse is a
``for almost every $\theta$'' statement (atypical, e.g.\ near-degenerate,
sources may have smaller regret, but cannot be identified in advance
without prior knowledge of $\theta$, which is exactly what the online
setting forbids).

\textbf{Theorem 7.7 (Lossy RNR via Shannon rate-distortion: byte-granular
random access with controlled per-position distortion).}
\emph{Let $X = (X_n)_{n \geq 1}$ satisfy the hypotheses of Theorem 7.2
(stationary, exp-mixing (5.8), finite alphabet $\Sigma$, sub-block
spacing $W_N \leq K = \mathrm{polylog}(N)$, conventions C1--C3),
and let $d_H : \Sigma \times \Sigma \to \{0, 1\}$ denote Hamming
distortion. For any per-position distortion budget $D \in [0,
1 - 1/|\Sigma|]$ (the Hamming-distortion range), define the
Shannon rate-distortion function of $X$:}
\[
R(D) \;:=\; \lim_{n \to \infty} \tfrac{1}{n} \min_{P_{\hat{X}^n \mid X^n} :\, \mathbb{E}[d_H(X^n, \hat{X}^n)/n] \leq D}
I(X^n; \hat{X}^n),
\]
\emph{which exists and is continuous in $D$ for stationary ergodic
$X$ (Cover \& Thomas 2006 §10.5 ``Computation of the rate-distortion
function''; Berger 1971 \emph{Rate Distortion Theory}, Prentice-Hall).
RNR with lossy per-sub-block rate-distortion coding (replacing
per-position lossless arithmetic coding of Theorem 7.2 with the
Shannon-optimal rate-distortion encoder at rate $R(D)$ per position)
achieves expected repair-stream length}
\[
\mathbb{E}\bigl[L^{\mathrm{rep}}_{\mathrm{RNR}, D\text{-lossy}}(X^N)\bigr]
\;\leq\;
N \cdot R(D) + m \cdot \delta_\infty^D + O(1) + \varepsilon_{\mathrm{ac}}^{\mathrm{total}},
\]
\emph{with per-position reconstruction satisfying $\mathbb{E}[d_H(X^N,
\hat{X}^N)/N] \leq D$. The quantity $\delta_\infty^D$ is the
$D$-distortion analog of the Crutchfield--Feldman excess
$\delta_\infty$, finite under exp-mixing (5.8) (Cover \& Thomas 2006
Theorem 10.4.1 for the per-sub-block rate-distortion).}

\emph{Limit cases:} \emph{(a) $D \to 0$: $R(D) \to h(X)$ and
$\delta_\infty^D \to \delta_\infty$, recovering Theorem 7.2's lossless
bound; (b) $D \to 1 - 1/|\Sigma|$: $R(D) \to 0$ (constant prediction
$\hat{X}_n = $ mode$_X$ minimises Hamming distortion at zero rate),
$L^{\mathrm{rep}}_{\mathrm{RNR}, D\text{-lossy}} \to O(1)$
(archive compresses to a constant.}

\emph{\textbf{Proof.}} The per-sub-block bound is the classical Shannon
rate-distortion theorem applied to a $K$-length block of the
stationary source (Cover \& Thomas 2006 Theorem 10.4.1):
for any per-sub-block distortion budget $D$, there exists an encoder
that achieves expected codeword length
$K \cdot R(D) + \delta_K^D$, where $\delta_K^D$ is the
finite-block-length excess (analog of Lemma 5.1d's $\delta_W^{(h_W)}$
in the rate-distortion regime). Under exp-mixing (5.8), the
conditional rate-distortion excess summable
$\sum_n [R(X_n \mid X_{<n}; D) - R(D)] = \delta_\infty^D < \infty$
(rate-distortion analog of Crutchfield--Feldman for finite-rate
encoding; Berger 1971 §6 on stationary rate-distortion).
Summing over $m = N/K$ sub-blocks: $m \cdot K \cdot R(D) + m \cdot
\delta_K^D \to N \cdot R(D) + m \cdot \delta_\infty^D$ as $K \to
\infty$. The arithmetic-coding finite-precision slack
$\varepsilon_{\mathrm{ac}}^{\mathrm{total}}$ is unchanged from
Theorem 7.2.
\hfill$\blacksquare$

\emph{Scope and significance.} Theorem 7.7 extends RNR from the
lossless regime (Theorem 7.2) to the lossy regime via the Shannon
rate-distortion function. Three concrete consequences:
\emph{(a) Smooth interpolation:} the bound is continuous in $D$
from $0$ (lossless = T7.2) to $1 - 1/|\Sigma|$ (trivial constant
encoding). Practical lossy RNR can choose $D$ to trade reconstruction
fidelity for storage cost.
\emph{(b) Empirical compression beats lossless Shannon:} for natural
sources where small Hamming distortion is acceptable (e.g., natural
images at JPEG quality, text with occasional typo, audio at MP3
quality), the lossy rate $R(D) < h(X)$ gives strict improvement.
\emph{(c) Per-query reconstruction precision:} byte-granular RA still
returns sub-block-decoded data; the distortion guarantee is
per-position average, so individual queries may have $d_H = 0$ or $1$
on each byte, with average $\leq D$.

\emph{Limitations.} The classical rate-distortion machinery requires
the source's $R(D)$ function to be COMPUTABLE; closed-form $R(D)$
exists for memoryless sources (binary Bernoulli: Blahut algorithm
Cover \& Thomas 2006 Example 10.3.1) and for some structured
sources (Gaussian: Cover \& Thomas 2006 Theorem 10.3.2), but for
arbitrary exp-mixing $X$ the $R(D)$ function may need numerical
approximation. Hamming distortion is the canonical choice for finite
alphabets; other distortion measures (e.g., $\ell_2$ for continuous
embeddings) require separate analysis.

\textbf{Theorem 7.8 (Slepian--Wolf distributed RNR: joint compression
of correlated sources with byte-granular random access on each).}
\emph{Let $X = (X_n)_{n \geq 1}$ and $Y = (Y_n)_{n \geq 1}$ be jointly
stationary on finite alphabets $\Sigma_X, \Sigma_Y$, with per-position
joint law $p_{XY}(x, y)$ and per-position mutual information $I :=
I(X_1; Y_1) > 0$. Assume the joint process $(X, Y)$ satisfies the
exp-mixing condition (5.8) applied to the product alphabet
$\Sigma_X \times \Sigma_Y$. Then for any rate pair $(R_X, R_Y)$ in the
Slepian--Wolf rate region}
\[
R_X \geq h(X \mid Y), \quad R_Y \geq h(Y \mid X), \quad
R_X + R_Y \geq h(X, Y),
\]
\emph{there exist distributed RNR encoders $\mathcal{E}_X, \mathcal{E}_Y$
(each operating on \emph{only} its own source) producing repair
streams of expected lengths}
\[
\mathbb{E}\bigl[L^{\mathrm{rep}}_X(X^N)\bigr] \leq N \cdot R_X + m \cdot \delta_\infty^{X} + O(1) + \varepsilon_{\mathrm{ac}},
\;\;
\mathbb{E}\bigl[L^{\mathrm{rep}}_Y(Y^N)\bigr] \leq N \cdot R_Y + m \cdot \delta_\infty^{Y} + O(1) + \varepsilon_{\mathrm{ac}},
\]
\emph{such that a joint decoder, given both archives, losslessly
reconstructs $(X^N, Y^N)$. The sum repair-stream length}
\[
\mathbb{E}[L^{\mathrm{rep}}_X + L^{\mathrm{rep}}_Y]
\;\leq\; N \cdot h(X, Y) + 2m \cdot \delta_\infty^{\mathrm{joint}} + O(1) + 2 \varepsilon_{\mathrm{ac}}
\]
\emph{saves $N \cdot I = N \cdot I(X_1; Y_1)$ bits over independent
RNR encoding of each source ($N \cdot h(X) + N \cdot h(Y) + \cdots$).
Byte-granular random access is preserved on each archive via
Theorem 10.3.}

\emph{\textbf{Proof.}} The Slepian--Wolf random-coding theorem
(Slepian \& Wolf 1973 \emph{IEEE Trans.\ Inf.\ Theory} 19(4):471--480
``Noiseless coding of correlated information sources'') establishes
that the stated rate region is achievable via random binning: each
encoder partitions $\Sigma_X^N$ (resp.\ $\Sigma_Y^N$) into $2^{N R_X}$
(resp.\ $2^{N R_Y}$) bins uniformly at random and sends the
bin-index of the realised source. The joint decoder, given both
bin-indices, finds the unique pair $(X^N, Y^N)$ jointly typical and
consistent with both indices.

Within the RNR framework: each encoder $\mathcal{E}_X, \mathcal{E}_Y$
operates per sub-block of size $K$. Per-sub-block bin-index encoding
plus the Theorem 7.2 per-sub-block sync overhead gives the stated
$N \cdot R_X + m \cdot \delta_\infty^X$ expected length for $X$
(analogous for $Y$). Joint decoder uses both bin-indices per
corresponding pair of sub-blocks to recover $(X_j, Y_j)$ jointly.

For \emph{constructive} achievement (deterministic codes instead of
random binning), structured codes such as LDPC-Wyner-Ziv
(Liveris--Xiong--Georghiades 2002 \emph{IEEE Comm.\ Letters}
6(10):440) or polar-Wyner-Ziv (Korada--Urbanke 2010 \emph{IEEE Trans.\
Inf.\ Theory} 56(4):1571) provide explicit binning rules with the
required typicality decoding. These integrate cleanly with the
RNR sub-block structure since each sub-block can use an independent
LDPC/polar code.

Sum-rate savings: independent RNR encoding gives $N \cdot h(X) + N
\cdot h(Y)$ baseline. Slepian--Wolf achieves $N \cdot h(X, Y) =
N \cdot [h(X) + h(Y) - I(X_1; Y_1)]$. Savings $= N \cdot I(X_1; Y_1)$
bits, scaling linearly with archive size whenever $I > 0$.
\hfill$\blacksquare$

\emph{Scope and limitations.} Theorem 7.8 generalises RNR from
single-source archives to multi-source archives where the encoders
do \emph{not} need to communicate during encoding (distributed
setting). Three remarks:

\emph{(a) Achievability vs constructivity.} Slepian--Wolf 1973 is an
existence result via random binning; explicit constructive codes
(LDPC, polar) achieve the bound up to vanishing gap. For RNR
implementation, structured codes are preferred.

\emph{(b) Byte-granular RA per-source.} The decoder, given a query
for byte $b$ in $X^N$, reads $X$'s sub-block containing $b$ AND the
corresponding sub-block of $Y$ to recover $X$ jointly. Per-query
cost: $O(\log L(R_X) + \log L(R_Y) + (K + k) \cdot \mathrm{cost}(M))$,
polylog in $N$ if both archives are polylog-RA. Independent decoding
(without $Y$) is INFEASIBLE in the Slepian--Wolf-optimal regime
$R_X = h(X | Y) < h(X)$: $X$ alone has been compressed below its
own entropy, requiring $Y$ for disambiguation.

\emph{(c) Generalisation to $k$ sources.} Slepian--Wolf extends to
$k$-source distributed coding with rate region $\sum_{i \in S} R_i
\geq H(X_S \mid X_{\bar{S}})$ for every $S \subseteq [k]$
(Csisz{\'a}r \& K{\"o}rner 1981/2011 \emph{Information Theory: Coding
Theorems for Discrete Memoryless Systems}, Cambridge Univ.\ Press,
Theorem 15.4.1; Cover 1975 \emph{IEEE Trans.\ Inf.\ Theory}
21(2):226--228 establishes the 2-source ergodic generalisation, with
the $k$-source case following by similar coding-theoretic
techniques). RNR with $k$ correlated archives achieves the
corresponding $k$-source sum-rate $H(X_{[k]})$.

\textbf{Theorem 7.9 (Successive-refinement RNR: anytime variable-quality
access via layered rate-distortion codes).}
\emph{Let $X = (X_n)_{n \geq 1}$ satisfy the hypotheses of Theorem 7.7
(stationary exp-mixing on finite alphabet $\Sigma$, sub-block range
$W_N \leq K = \mathrm{polylog}(N)$, conventions C1--C3), under Hamming
distortion. Suppose $X$ satisfies the \emph{Equitz--Cover
successive-refinement condition} (Equitz \& Cover 1991
\emph{IEEE Trans.\ Inf.\ Theory} 37(2):269--275
``Successive refinement of information''): for every nested distortion
sequence $D_0 > D_1 > \cdots > D_{k-1} \geq 0$, there exists a Markov
chain $X \rightarrow \hat{X}_{k-1} \rightarrow \hat{X}_{k-2} \rightarrow
\cdots \rightarrow \hat{X}_0$ where $\hat{X}_i$ achieves the rate-distortion
optimum $R(D_i)$ at distortion level $D_i$, and each $\hat{X}_{i-1}$
is a (possibly random) function of $\hat{X}_i$. Then there exist
multi-layer RNR codes with $k$ refinement layers $c_0, c_1, \ldots,
c_{k-1}$ such that}
\begin{enumerate}
\item \emph{(Layer rates) Reading layers $0$ through $i$ yields
reconstruction $\hat{X}^N_i$ with expected per-position distortion
$\mathbb{E}[d_H(X^N, \hat{X}^N_i)/N] \leq D_i$ and cumulative encoded
length}
\[
\sum_{j = 0}^{i} L^{\mathrm{rep}}_{\mathrm{layer}_j}(X^N)
\;\leq\;
N \cdot R(D_i) + (i+1) m \cdot \delta_\infty^{D_i} + O(1) + (i+1) \varepsilon_{\mathrm{ac}}.
\]
\item \emph{(No layering penalty) Total archive length (all $k$ layers
concatenated) equals the single-shot $R(D_{k-1})$-rate code length up
to lower-order terms:}
\[
\sum_{j = 0}^{k-1} L^{\mathrm{rep}}_{\mathrm{layer}_j}(X^N)
\;\leq\;
N \cdot R(D_{k-1}) + k m \cdot \delta_\infty^{D_{k-1}} + O(1) + k \varepsilon_{\mathrm{ac}}.
\]
\item \emph{(Anytime byte-granular RA) Decoder queries can specify
desired distortion $D_i \in \{D_0, \ldots, D_{k-1}\}$ per byte and read
only the layer slices $0..i$ for the sub-block containing the byte.
Per-query cost: $O\bigl((i+1) (\log L(R) + (K+k) \mathrm{cost}(M))\bigr)$,
polylog in $N$ for $i$ bounded and $K, k, \mathrm{cost}(M),
\log L(R)$ polylog.}
\end{enumerate}

\emph{\textbf{Proof.}} The Equitz--Cover 1991 successive-refinement
theorem establishes that under the stated Markov-chain condition,
there exist nested codebooks $\mathcal{C}_0 \subset \mathcal{C}_1
\subset \cdots \subset \mathcal{C}_{k-1}$ achieving the
rate-distortion optimum at each level: $|\mathcal{C}_i| = 2^{N R(D_i)}$.
Each codeword in $\mathcal{C}_{i+1}$ refines a unique codeword in
$\mathcal{C}_i$ (i.e., the refinement layer encodes the residual
between $\hat{X}_i$ and $\hat{X}_{i+1}$ in $R(D_{i+1}) - R(D_i)$ bits).
This is the ``embedded'' or ``nested'' codebook structure
(Rimoldi 1994 \emph{IEEE Trans.\ Inf.\ Theory} 40(1):253--259
``Successive refinement of information: characterization of the
achievable rates'' for the converse and tight characterisation).

Within the RNR framework: each sub-block of size $K$ uses the layered
codebook structure. Layer $i$'s codeword for sub-block $j$ encodes
the residual at rate $R(D_i) - R(D_{i-1})$ bits per position (with
$R(D_{-1}) := 0$ for layer 0). The sync-overhead $\delta_\infty^{D_i}$
at each sub-block is the rate-distortion analog of T7.7's
sync overhead. Summing over $m = N/K$ sub-blocks:
$L^{\mathrm{rep}}_{\mathrm{layer}_i} = N \cdot (R(D_i) - R(D_{i-1})) +
m \cdot \delta_\infty^{D_i} + O(1)$.

Cumulative layers $0..i$:
$\sum_{j=0}^{i} L^{\mathrm{rep}}_{\mathrm{layer}_j} = N \cdot R(D_i) +
\sum_{j=0}^{i} m \cdot \delta_\infty^{D_j} + O(1) \leq N \cdot R(D_i) +
(i+1) m \cdot \delta_\infty^{D_i} + O(1)$ (since
$\delta_\infty^{D_j} \leq \delta_\infty^{D_0}$ for $D_j \geq D_i$ when
$j \leq i$ --- coarser distortion has smaller excess).

Per-query: layers $0..i$ each contribute $O(\log L(R) + (K+k)
\mathrm{cost}(M))$ to the byte-decoding cost.
\hfill$\blacksquare$

\emph{Scope and significance.} Theorem 7.9 enables \emph{anytime}
variable-quality archives: a single archive stores multiple distortion
levels, and the decoder chooses the level per byte query without
re-encoding. Three concrete consequences:

\emph{(a) Examples of refinable sources}: Binary Bernoulli source
under Hamming distortion is successively refinable (Equitz--Cover 1991
Theorem 1); Gaussian source under squared-error is also refinable
(Equitz--Cover 1991 Theorem 2). Memoryless sources with discrete
alphabets and Hamming distortion satisfy the condition under mild
regularity.

\emph{(b) Non-refinable sources}: For sources \emph{not} satisfying
the Equitz--Cover condition, layered RNR incurs an additional rate
penalty per layer: cumulative layer rates exceed $R(D_i)$ by a
source-dependent gap (Rimoldi 1994 cited above gives the complete
characterisation of refinable vs.\ non-refinable cases; for the side-
information variant see Steinberg \& Merhav 2004
\emph{IEEE Trans.\ Inf.\ Theory} 50(8):1636--1654 ``On successive
refinement for the Wyner-Ziv problem''). The penalty is bounded but
non-zero.

\emph{(c) Practical applications}: video streaming with adaptive
quality (mobile networks: low-quality on weak signal, high-quality
on strong); progressive image rendering (JPEG-2000-style); scientific
data archives with quick-look at low fidelity + drill-down at high
fidelity. All built on RNR's byte-granular RA, with layered codes
enabling per-query distortion selection.

\textbf{Theorem 7.10 (Wyner-Ziv RNR: lossy compression of $X$ with
correlated side-information $Y$ available only at the decoder).}
\emph{Let $X = (X_n)_{n \geq 1}$ and $Y = (Y_n)_{n \geq 1}$ be jointly
stationary on finite alphabets, with $(X, Y)$ satisfying the joint
exp-mixing condition (5.8) on the product alphabet. Assume only $X$
needs reconstruction with Hamming distortion budget
$D \in [0, 1 - 1/|\Sigma_X|]$; $Y$ is available at the decoder as
side-information but NOT at the encoder. Define the Wyner-Ziv
rate-distortion function}
\[
R_{\mathrm{WZ}}(D) \;:=\; \min_{\substack{U : U - X - Y \text{ Markov chain} \\ \hat{X} = f(U, Y) \\ \mathbb{E}[d_H(X, \hat{X})] \leq D}} I(X; U) - I(Y; U),
\]
\emph{the Wyner-Ziv 1976 rate-distortion function with side-information
at decoder (Wyner \& Ziv 1976 \emph{IEEE Trans.\ Inf.\ Theory}
22(1):1--10 ``The rate-distortion function for source coding with
side information at the decoder''). Then there exists an RNR encoder
$\mathcal{E}_X^{\mathrm{WZ}}$ (operating on $X^N$ alone, without
access to $Y^N$) such that the joint decoder, given both archives,
reconstructs $\hat{X}^N$ with $\mathbb{E}[d_H(X^N, \hat{X}^N)/N]
\leq D$ and}
\[
\mathbb{E}\bigl[L^{\mathrm{rep},\mathrm{WZ}}_X(X^N)\bigr]
\;\leq\;
N \cdot R_{\mathrm{WZ}}(D) + m \cdot \delta_\infty^{\mathrm{WZ}, D} + O(1) + \varepsilon_{\mathrm{ac}}.
\]

\emph{Closed form for binary doubly symmetric source.} For
$X \sim \mathrm{Bern}(1/2)$, $Y = X \oplus Z$ with $Z \sim \mathrm{Bern}(p)$
independent of $X$, and Hamming distortion:
\[
R_{\mathrm{WZ}}(D) \;=\; \begin{cases} h(p) - h(D) & \text{if } D \in [0, p], \\ 0 & \text{if } D \geq p, \end{cases}
\]
\emph{where $h(\cdot)$ is binary entropy (Wyner--Ziv 1976 §3.5
``Doubly symmetric binary source'').}

\emph{Limits and relations to T7.7/T7.8.}
\emph{(a) $D \to 0$: $R_{\mathrm{WZ}}(0) = h(X \mid Y)$, recovering
the Slepian--Wolf one-sided rate (Theorem 7.8 with only $X$
reconstructed).}
\emph{(b) $D \to D_{\max}$ (where $D_{\max} = p$ for binary doubly
symmetric): $R_{\mathrm{WZ}} \to 0$. Decoder outputs $\hat{X} = Y$
directly, distortion equals $\mathbb{E}[d_H(X, Y)] = p$. No encoding
needed.}
\emph{(c) Comparison with no-side-info T7.7: $R_{\mathrm{WZ}}(D) \leq
R(D)$ always; for correlated $(X, Y)$ the gap $R(D) - R_{\mathrm{WZ}}(D)$
scales with $I(X; Y)$.}

\emph{\textbf{Proof.}} Wyner--Ziv 1976 Theorem 2 establishes the
$R_{\mathrm{WZ}}(D)$ rate-distortion function as the achievable rate
with side-info at decoder. The achievability uses random binning of
$\mathcal{U}$-typical sequences (where $\mathcal{U}$ is the auxiliary
alphabet) combined with the typicality decoder selecting the unique
$U^N$ jointly typical with the decoder's $Y^N$.

Within RNR: sub-block-level Wyner-Ziv coding applies per sub-block of
size $K$. Each sub-block's $X_K$ is mapped to a $U_K$ codeword
(quantization), then $U_K$ is binned at rate $I(X; U) - I(Y; U)$ bits
per position. Joint decoder uses sub-block of $Y$ for that index range
to recover $U_K$, then computes $\hat{X}_K = f(U_K, Y_K)$.

Sync overhead $\delta_\infty^{\mathrm{WZ}, D}$ is the Crutchfield--Feldman
analog under the joint mixing, finite for exp-mixing $(X, Y)$.

For constructive codes: DISCUS (Pradhan--Ramchandran 1999/2003
\emph{IEEE Trans.\ Inf.\ Theory} 49(3):626--643 ``Distributed source
coding using syndromes (DISCUS): design and construction'') gives
explicit syndrome-based binning that integrates with sub-block
structure.
Modern constructions use LDPC or polar codes for the binning
(Liveris--Xiong--Georghiades 2002 cited above).
\hfill$\blacksquare$

\emph{Scope and significance.} Theorem 7.10 completes the multi-source
RNR quadrant:

\begin{center}
\begin{tabular}{l|c|c}
              & Lossless                          & Lossy (distortion $D$)            \\ \hline
Single-source & Theorem 7.2 ($N h(X)$)            & Theorem 7.7 ($N R(D)$)            \\
Multi-source  & Theorem 7.8 ($N h(X, Y)$)         & Theorem 7.10 ($N R_{\mathrm{WZ}}(D)$) \\
\end{tabular}
\end{center}

\emph{Numerical example.} For $X \sim \mathrm{Bern}(1/2)$, $Y = X
\oplus \mathrm{Bern}(0.1)$, Hamming $D = 0.05$:
$R_{\mathrm{WZ}}(0.05) = h(0.1) - h(0.05) = 0.469 - 0.286 = 0.183$ bits/symbol.
Compared to $R(0.05) = 1 - h(0.05) = 0.714$ (no side-info, T7.7):
\emph{74\% savings} from side-info.
Compared to lossless Slepian-Wolf $h(X \mid Y) = h(0.1) = 0.469$:
\emph{61\% savings} from lossy.

\emph{Applications}: distributed video coding where reference frame
$Y$ available at decoder (e.g., previous frame, server-side);
sensor networks with predictable side-info (e.g., GPS plus inertial);
deduplication with similarity-aware updates.

\textbf{Corollary 7.11 (Lossy RNR plus error-correcting code equals
lossless RNR: Shannon source--channel separation in the RNR framework).}
\emph{For any stationary exp-mixing source $X$ under Hamming
distortion budget $D \in [0, D_{\max}]$ (where $D_{\max} = 1 -
\max_x \pi(x)$ is the no-information distortion --- the smallest
distortion achievable by a constant predictor), the rate-distortion-plus-
error-correction decomposition}
\[
L^{\mathrm{rep}}_{\mathrm{lossy}}(X^N; D) + L^{\mathrm{ECC}}(N, D)
\;\geq\;
N \cdot h(X) + 2m \cdot \delta_\infty + O(1),
\]
\emph{where $L^{\mathrm{rep}}_{\mathrm{lossy}}$ is the Theorem 7.7
lossy-RNR repair stream and $L^{\mathrm{ECC}}(N, D) = N \cdot h(D) +
O(\sqrt{N})$ is the Shannon-optimal binary-symmetric-channel ECC
correcting $D$-fraction Hamming errors. For \emph{memoryless} $X$
(e.g., $X \sim \mathrm{Bern}(p)$) the inequality is TIGHT:}
\[
L^{\mathrm{rep}}_{\mathrm{lossy}}(X^N; D) + L^{\mathrm{ECC}}(N, D)
\;=\;
N \cdot [R(D) + h(D)] + O(1)
\;=\;
N \cdot h(X) + O(1)
\;=\;
L^{\mathrm{rep}}_{\mathrm{lossless}}(X^N) + O(1)
\]
\emph{(Theorem 7.2 directly). For sources with memory, the inequality
may be strict: $R(D) + h(D) > h(X)$ for some $D$, reflecting that
Shannon's source--channel separation can be suboptimal under
memory (Cover \& Thomas 2006 §11.7 ``Source--channel separation
theorem'' and counterexample).}

\emph{\textbf{Proof.}} The lossy-RNR rate is $N \cdot R(D) + m
\delta_\infty^D + O(1)$ (Theorem 7.7). The ECC for $D$-fraction
Hamming errors on a length-$N$ string requires at least $N \cdot h(D)$
bits asymptotically (Shannon's binary-symmetric channel-capacity
theorem $C = 1 - h(D)$ implies the ECC overhead $1 - C = h(D)$ per
bit; Cover \& Thomas 2006 §7.1). Summing:
\[
L^{\mathrm{lossy}} + L^{\mathrm{ECC}} \geq N \cdot R(D) + N \cdot h(D) + 2m \delta_\infty + O(1).
\]
For memoryless $X \sim \mathrm{Bern}(p)$ at Hamming distortion in
the valid range $D \in [0, \min(p, 1-p)]$: $R(D) = h(p) - h(D)$
(Shannon 1959; Cover \& Thomas 2006 Theorem 10.3.1 ``Bernoulli
rate-distortion''). Sum: $R(D) + h(D) = h(p) = h(X)$, matching
$L^{\mathrm{lossless}}$ from Theorem 7.2. For $D > \min(p, 1-p)$
the rate-distortion function $R(D) = 0$ (constant predictor
$\hat{X} = \mathrm{argmax}_x \pi(x)$ achieves distortion $\min(p,1-p)$
at zero rate); the identity $R(D) + h(D) = h(X)$ no longer holds in
this regime --- it transitions to $R(D) + h(D) = h(D)$, which can be
less than $h(X)$.

For memory sources, Shannon's separation theorem (Cover \& Thomas 2006
§11.7) gives $R(D) + h(D) \geq h(X)$ with equality conditional on
memorylessness. Strict inequality is possible for certain joint
source--channel settings (Vembu, Verd{\'u} \& Steinberg 1995
\emph{IEEE Trans.\ Inf.\ Theory} 41(1):44--54 ``The source--channel
separation theorem revisited''). For the canonical Hamming-distortion
binary-Markov case, separation remains optimal up to vanishing terms;
strict inequality is a feature of more exotic distortion/channel
combinations (e.g., joint Gaussian source over Gaussian channel under
mean-squared-error, where joint coding beats separation).
\hfill$\blacksquare$

\emph{Practical implication.} Lossy RNR (Theorem 7.7) is genuinely
useful only when:
\emph{(a)} the application tolerates distortion $D > 0$ in the
reconstruction (e.g., perceptual quality in video/audio), or
\emph{(b)} the ECC infrastructure is unavailable or costlier than
direct lossless coding.
For applications requiring exact reconstruction, Theorem 7.2 (direct
lossless RNR) is at least as efficient as Theorem 7.7 + ECC, and
strictly better for memory sources where separation is suboptimal.
This grounds the framework: \emph{lossless RNR is the asymptotic
information-theoretic optimum for byte-granular random-access exact
compression}, with lossy and side-info variants (Theorems 7.7--7.10)
exploring controlled relaxations of the exactness constraint.

\textbf{Theorem 7.12 (Functional RNR: compressing $f(X)$ via the
Orlitsky-Roche characteristic-graph entropy when the decoder needs
only a function of the source).}
\emph{Let $X = (X_n)_{n \geq 1}$ and $Y = (Y_n)_{n \geq 1}$ be jointly
stationary on finite alphabets $\Sigma_X, \Sigma_Y$ with $(X, Y)$
satisfying joint exp-mixing (5.8), and let
$f : \Sigma_X \times \Sigma_Y \to \Omega$ be a deterministic
function. Suppose the decoder, given the encoded $X^N$ and the
side-information $Y^N$, needs only to compute $f(X_n, Y_n)$ at each
position $n$ (not recover $X^N$ itself). Define the
\emph{characteristic graph} $G_f$ of $f$ given $Y$: vertices are
$\Sigma_X$, with an edge $(x, x')$ if and only if there exists
$y \in \Sigma_Y$ with $f(x, y) \neq f(x', y)$ and
$p_{XY}(x, y), p_{XY}(x', y) > 0$ (Witsenhausen 1976
\emph{IEEE Trans.\ Inf.\ Theory} 22(5):592--593 ``The zero-error side
information problem and chromatic numbers''). Define the
\emph{conditional graph entropy}}
\[
H_{G_f}(X \mid Y) \;:=\; \min_{W : W - X - Y,\, W \in \mathrm{stable\text{-}sets}(G_f)} I(X; W \mid Y),
\]
\emph{the Orlitsky--Roche 2001 rate (Orlitsky \& Roche 2001
\emph{IEEE Trans.\ Inf.\ Theory} 47(3):903--917 ``Coding for
computing'') minimised over auxiliary random variable $W$ taking
values in independent (stable) sets of $G_f$, with the Markov chain
$W - X - Y$. Then there exists an RNR encoder $\mathcal{E}_X^f$ such
that}
\[
\mathbb{E}\bigl[L^{\mathrm{rep},f}_X(X^N)\bigr]
\;\leq\;
N \cdot H_{G_f}(X \mid Y) + m \cdot \delta_\infty^f + O(1) + \varepsilon_{\mathrm{ac}},
\]
\emph{with byte-granular random access to the encoded $X^N$, such
that the joint decoder reconstructs $f(X_n, Y_n)$ at every position
$n$ losslessly.}

\emph{Comparison with previous theorems.}
\emph{(a) $f = \mathrm{identity}$ (decoder needs $X$ itself):
$H_{G_f}(X \mid Y) = H(X \mid Y)$, recovering Slepian--Wolf
one-sided rate (Theorem 7.8).}
\emph{(b) Non-trivial $f$:} $H_{G_f}(X \mid Y) \leq H(X \mid Y)$,
with strict inequality possible when $f$ identifies stable-set
equivalence classes in $\Sigma_X$.

\emph{Example: parity function on $\Sigma_X = \{0, 1, 2, 3\}$.}
Let $X$ be uniform on $\{0,1,2,3\}$, $Y$ arbitrary (or null
side-info), and $f(x, y) = x \bmod 2$ (decoder needs only the parity
of $X$). The characteristic graph $G_f$ has edges only between
\emph{different-parity} vertices: $\{(0,1), (0,3), (2,1), (2,3)\}$;
no edges within even $\{0, 2\}$ or odd $\{1, 3\}$. Stable sets:
$\{0, 2\}$ and $\{1, 3\}$ (the parity classes). Auxiliary $W :=
X \bmod 2$ achieves the Orlitsky-Roche bound:
$H_{G_f}(X \mid Y) = H(W) = h(0.5) = 1$ bit/symbol, versus
$H(X \mid Y) = H(X) = 2$ bits/symbol (no side-info case). \emph{Saving:
$1$ bit/symbol = $50\%$ of the Slepian-Wolf rate.}

\emph{\textbf{Proof.}} Orlitsky-Roche 2001 Theorem 1 establishes that
$H_{G_f}(X \mid Y)$ is achievable via random binning of stable-set
codewords of $G_f$: each encoder partition maps $X^N$ to a
typical-stable-set-sequence index, with rate matching the
characteristic-graph entropy. The joint decoder, given $Y^N$ and the
bin-index, recovers a stable-set sequence in $G_f$ consistent with
$Y^N$; since stable sets correspond to $f$-equivalence classes, this
yields $f(X_n, Y_n)$ at every position.

Within RNR: each sub-block of size $K$ uses Orlitsky-Roche encoding
on the sub-block's $(X_K, Y_K)$ pair. Sync overhead
$\delta_\infty^f$ is the analog of T7.8's joint sync overhead,
finite under exp-mixing (5.8). Sum over $m = N/K$ sub-blocks gives
the stated bound.

For \emph{constructive} achievement: Doshi, Shah, M{\'e}dard \&
Effros 2010 \emph{IEEE Trans.\ Inf.\ Theory} 56(8):3901--3917
``Functional compression through graph coloring'' (extended journal
version of the 2007 ISIT proceedings; Doshi, Shah \& M{\'e}dard, Proc.\
IEEE ISIT 2007, pp.\ 1501--1505 ``Source coding with distortion
through graph coloring'') provides explicit graph-colouring-based
encoders; modern coding constructions (LDPC, polar) integrate with
sub-block structure.
\hfill$\blacksquare$

\emph{Scope and significance.} Theorem 7.12 extends the multi-source
RNR family to scenarios where the decoder only needs a FUNCTION of
the source. Practical applications include:

\emph{(a) Aggregate queries on encrypted/compressed databases:}
the decoder computes sums, maxima, parities of encrypted records
without reconstructing them individually.

\emph{(b) Compressed neural-network inference:} the decoder evaluates
$f(\text{model weights}) = \text{predicted label}$ without reconstructing
the full weight tensor.

\emph{(c) Privacy-preserving compression:} the decoder learns
$f(X) = $ aggregated statistic, but cannot reconstruct individual
records $X_n$, achieving differential-privacy-style protection at
$O(1)$-bit-per-position cost.

\emph{(d) Federated computation:} encoders distributed across devices
encode their local $X_n$ at rate $H_{G_f}(X | Y)$; central server
computes $f(X, Y)$ aggregates without seeing $X$.

\emph{Limitations.} $H_{G_f}$ in general requires solving an
information-theoretic optimisation (minimisation over auxiliary $W$).
Closed-form evaluation is possible for specific $f$ structures
(parity, threshold functions, structured graph automorphisms);
generic $f$ requires numerical optimisation. Doshi--Shah--M{\'e}dard--Effros
2010 gives polynomial-time achievable schemes within $O(\log |\Sigma_X|)$
of optimal.

\textbf{Theorem 7.13 (Error-resilient RNR: joint source--channel
coding for byte-granular RA on noisy storage).}
\emph{Let $X = (X_n)_{n \geq 1}$ satisfy the Theorem 7.2 hypotheses
(stationary exp-mixing, $W_N \leq K = \mathrm{polylog}(N)$,
conventions C1--C3), and suppose the encoded archive is stored on a
binary symmetric channel (BSC) with bit-flip rate
$\varepsilon \in (0, 1/2)$ per stored bit (memoryless). Then for any
$\delta > 0$, there exists an encoder-ECC composition achieving
archive size}
\[
L^{\mathrm{archive}}_{\mathrm{RNR}, \varepsilon}(X^N)
\;\leq\;
\frac{L^{\mathrm{rep}}_{\mathrm{RNR}}(X^N) + L(M) + L(\Theta)}{1 - h(\varepsilon)} + O(\sqrt{N})
\]
\emph{such that every byte-granular RA query is decodable to the
correct source byte with probability at least $1 - 2^{-\Omega(K)}$
(decoding failure exponential in sub-block size). The factor
$1/(1 - h(\varepsilon))$ is the Shannon BSC channel-capacity overhead:
the channel can carry $C = 1 - h(\varepsilon)$ bits of information per
channel use. The penalty term $O(\sqrt{N})$ holds when
$K = \Omega(\sqrt{N})$ (balanced T7.5 regime); for polylog $K$ the
penalty is $O(N/\sqrt{K}) = O(N/\sqrt{\mathrm{polylog}(N)})$,
sub-linear in $N$ but larger than $\sqrt{N}$ (Polyanskiy--Poor--Verd{\'u}
2010 \emph{IEEE Trans.\ Inf.\ Theory} 56(5):2307--2359 ``Channel coding
rate in the finite blocklength regime'').}

\emph{Limits and trade-offs.}
\emph{(a) $\varepsilon \to 0$ (clean storage):} $1/(1 - h(\varepsilon))
\to 1$. Archive size $\to L^{\mathrm{rep}}_{\mathrm{RNR}}$, recovering
Theorem 7.2 without ECC overhead.
\emph{(b) $\varepsilon \to 1/2$ (random storage):} $C \to 0$. Archive
size $\to \infty$. No information can be reliably stored on a fully
random channel.
\emph{(c) Constant $\varepsilon = c$:} archive expansion factor
$1/(1 - h(c))$ is a CONSTANT. RNR's polylog per-query cost is
preserved (each query reads $O((K+k) \cdot \mathrm{cost}(M)/(1-h(\varepsilon)))$
bits, still polylog).

\emph{\textbf{Proof.}} Compose Theorem 7.2 with the Shannon
channel-coding theorem (Cover \& Thomas 2006 §7 ``Channel Capacity'').

\emph{Source-encoding stage.} Apply Theorem 7.2 to produce repair
stream of length $L^{\mathrm{rep}}_{\mathrm{RNR}}(X^N) = N \cdot h(X) +
m \cdot \delta_\infty + O(1) + \varepsilon_{\mathrm{ac}}$ bits.

\emph{Channel-coding stage.} Apply an ECC at rate $R_{\mathrm{ECC}} =
1 - h(\varepsilon) - \xi$ for any small $\xi > 0$ (achievable by random
linear codes; Cover \& Thomas 2006 Theorem 7.7.1). The ECC-encoded
archive has length $L^{\mathrm{rep}}/R_{\mathrm{ECC}} = L^{\mathrm{rep}}/(1
- h(\varepsilon) - \xi)$. Per-block decoding error probability is
$\leq 2^{-\Omega(K \xi^2)}$ (random-coding bound; Cover \& Thomas
2006 Theorem 7.7.1 proof).

For \emph{per-query} reliable decoding: the ECC is applied at the
\emph{sub-block} level. Each sub-block's encoded payload (length
$\sim K \cdot h(X) + \delta_\infty$ bits) is ECC-encoded
independently to length $\sim K \cdot h(X) / (1 - h(\varepsilon))$
bits. Per-query decoding: read the ECC-encoded sub-block,
ECC-decode, then source-decode via the RNR sub-block decoder.

The finite-blocklength penalty accumulates over $m = N/K$ sub-blocks,
each with $O(\sqrt{K \log K})$ penalty (Polyanskiy--Poor--Verd{\'u}
channel-dispersion bound): total
$m \cdot \sqrt{K \log K} = (N/K) \cdot \sqrt{K \log K} = N \sqrt{\log
K / K}$. For $K = \mathrm{polylog}(N)$: penalty scales as $N /
\sqrt{\mathrm{polylog}(N)}$, sub-linear in $N$ but \emph{larger}
than $\sqrt{N}$. To achieve $O(\sqrt{N})$ total penalty, require
$K = \Omega(\sqrt{N})$ (matching Theorem 7.5's balanced regime).
For Theorem 7.13's polylog-$K$ regime, the more honest bound is
$O(N / \sqrt{K})$.

For constructive ECC: LDPC, polar, or BCH codes at sub-block size
$K = \mathrm{polylog}(N)$ achieve the rate $1 - h(\varepsilon) - O(1/K)$
with polynomial encoding/decoding complexity (Arıkan 2009 \emph{IEEE
Trans.\ Inf.\ Theory} 55(7):3051 ``Channel polarization'' for polar
codes; Richardson \& Urbanke 2008 \emph{Modern Coding Theory},
Cambridge Univ.\ Press, for LDPC).
\hfill$\blacksquare$

\emph{Scope and significance.} Theorem 7.13 establishes that RNR's
byte-granular RA property extends to noisy storage with a Shannon-
optimal expansion factor $1/(1 - h(\varepsilon))$. Three regimes:

\emph{(a) DRAM / SSD storage} ($\varepsilon \sim 10^{-15}$): negligible
expansion ($1 + 10^{-15} \cdot \log 2$), essentially Theorem 7.2.

\emph{(b) Magnetic tape / cold storage} ($\varepsilon \sim 10^{-6}$):
$h(\varepsilon) \approx 2 \cdot 10^{-5}$, expansion factor
$\sim 1 + 2 \cdot 10^{-5}$, $0.002\%$ overhead.

\emph{(c) Quantum-decoherence-prone or radiation-hardened storage}
($\varepsilon \sim 10^{-2}$): expansion factor $\sim 1.09$ (9\%
overhead), still very tractable.

\emph{(d) Heavy-noise channels} ($\varepsilon \sim 0.1$): expansion
factor $\sim 1/(1 - 0.469) = 1.88$, nearly doubling archive size.
Near $\varepsilon = 0.5$ the framework breaks down.

\emph{Relation to Corollary 7.11.} Both invoke Shannon source-channel
separation, but with opposite directionality:
Corollary 7.11 \emph{decomposes} lossless into lossy + ECC for a
fictitious channel (the ``distortion budget'' as effective noise);
Theorem 7.13 \emph{composes} lossless source coding with ECC for
an actual physical channel (storage bit-flips). Together they show
that the framework's $h(X) + h(\varepsilon)$ Shannon bound is the
asymptotic capacity for both interpretations.

\textbf{Theorem 7.14 (Concentration of $L^{\mathrm{rep}}_{\mathrm{RNR}}$
around its expectation: high-probability $O(\sqrt{N})$ deviation).}
\emph{Let $X = (X_n)_{n \geq 1}$ satisfy the Theorem 7.2 hypotheses
(stationary, exp-mixing (5.8), $W_N \leq K = \mathrm{polylog}(N)$,
conventions C1--C3 including the $\eta$-floor predictor distribution).
Define the per-position log-loss}
$\ell_n := -\log_2 P_M(X_n \mid X_{n-W_N : n-1})$,
\emph{the bits arithmetic-coded for the $n$-th source symbol under the
$W_N$-bounded predictor. Then the total repair-stream length
$L^{\mathrm{rep}}_{\mathrm{RNR}}(X^N) = \sum_n \ell_n + (\text{sub-block
sync overhead}) + \varepsilon_{\mathrm{ac}}^{\mathrm{total}}$
concentrates around its expectation: for every $t > 0$,}
\[
\Pr\Bigl(\,\bigl|\,L^{\mathrm{rep}}_{\mathrm{RNR}}(X^N) - \mathbb{E}[L^{\mathrm{rep}}_{\mathrm{RNR}}(X^N)]\,\bigr| > t\,\Bigr)
\;\leq\;
2 \exp\!\Bigl(-\frac{t^2}{2 N \cdot \bigl(W_N + 1 + \tfrac{\rho}{1-\rho}\bigr)^2 \cdot \log_2^2(1/\eta)}\Bigr).
\]
\emph{The $(W_N+1)$ factor reflects the direct dependency of $\ell_k$
($k \in [n, n+W_N]$) on $X_n$ via the $W_N$-bounded predictor context,
and the additional $\rho/(1-\rho)$ term tracks the source-mediated
mixing-tail contribution: under exp-mixing (5.8), revealing $X_n$
alters the conditional distribution of $X_{n+W_N+1}, X_{n+W_N+2}, \ldots$
at geometric rate $\rho^j$, summing to $\rho/(1-\rho)$. See the
Remark below for the Doob-vs-compensated-martingale distinction
and the constant-$\rho$ asymptotic tightness.
Equivalently, for $t = c \cdot \sqrt{N} \cdot (W_N + 1 + \rho/(1-\rho))
\cdot \log_2(1/\eta)$ the deviation probability is $\leq 2 \exp(-c^2/2)$.
Taking $c = \sqrt{2\log(2/\delta)}$ gives a $1 - \delta$ confidence
bound: with probability at least $1 - \delta$,}
\[
\bigl|\,L^{\mathrm{rep}}_{\mathrm{RNR}}(X^N) - N \cdot h(X) - m \cdot \delta_\infty - O(1)\,\bigr|
\;\leq\;
\sqrt{2 N \log(2/\delta)} \cdot \bigl(W_N + 1 + \tfrac{\rho}{1-\rho}\bigr) \cdot \log_2(1/\eta).
\]

\emph{\textbf{Proof.}} The per-position log-loss
$\ell_n = -\log_2 P_M(X_n \mid \text{context}_n)$ is bounded above by
$\log_2(1/\eta)$ (the $\eta$-floor convention C2 guarantees
$P_M \geq \eta$, so $\ell_n \leq \log_2(1/\eta)$) and below by $0$
(probabilities are at most $1$). Hence $\ell_n \in [0, \log_2(1/\eta)]$.

The sum $\sum_n \ell_n$ is a function of the stationary process
$(X_n)$. Under exp-mixing (5.8), this sum is approximately a sum of
weakly-dependent bounded random variables. Specifically, define
$S_N := \sum_n \ell_n$ and the Doob martingale
$M_n := \mathbb{E}[S_N \mid X_1, \ldots, X_n] - \mathbb{E}[S_N]$.
The martingale differences $D_n := M_n - M_{n-1}$ capture the change
in conditional expectation induced by revealing $X_n$, with two
contributions:
(i) \emph{predictor-context}: $X_n$ enters the context window of $\ell_k$
for $k \in [n, n + W_N]$, so the direct change to $\mathbb{E}[\ell_k \mid
\mathcal{F}_n] - \mathbb{E}[\ell_k \mid \mathcal{F}_{n-1}]$ at each such
$k$ is bounded by $\log_2(1/\eta)$, summing to
$(W_N+1) \cdot \log_2(1/\eta)$.
(ii) \emph{source-mixing tail}: for $k > n + W_N$ the predictor's
context window no longer contains $X_n$, but under exp-mixing (5.8)
the conditional distribution $\Pr(X_k \in \cdot \mid \mathcal{F}_n)$ vs.\
$\Pr(X_k \in \cdot \mid \mathcal{F}_{n-1})$ differs in total variation
by at most $C \rho^{k-n}$ for some constant $C$ absorbed into the
exp-mixing constant (W.l.o.g.\ $C = 1$ after re-scaling $\rho$). The
contribution to $D_n$ from each such $k$ is bounded by
$\log_2(1/\eta) \cdot \rho^{k-n}$, summing over $k > n + W_N$ to at most
$\log_2(1/\eta) \cdot \rho/(1-\rho)$.
Combining (i) and (ii):
$|D_n| \leq (W_N + 1 + \rho/(1-\rho)) \cdot \log_2(1/\eta)$.

By Azuma--Hoeffding's inequality (Azuma 1967 \emph{Tohoku Math.\ J.}\
19(3):357; Hoeffding 1963 \emph{J.\ Amer.\ Stat.\ Assoc.}\
58(301):13):
\[
\Pr(|M_N| > t) \;\leq\; 2 \exp\!\Bigl(-\frac{t^2}{2 \sum_n D_n^2}\Bigr)
\;\leq\;
2 \exp\!\Bigl(-\frac{t^2}{2 N \cdot \bigl(W_N + 1 + \tfrac{\rho}{1-\rho}\bigr)^2 \cdot \log_2^2(1/\eta)}\Bigr),
\]
giving the stated concentration. The exp-mixing assumption (5.8)
is essential for (ii) above; without exp-mixing, the
source-mixing-tail term would be ill-defined and the bound would
require either an iid-input bounded-differences argument
(McDiarmid 1989 \emph{Surveys in Combinatorics} Cambridge
148:148--188 ``On the method of bounded differences'') or
specialized concentration for mixing processes (Rio 2000
\emph{Th\'eorie asymptotique des processus al\'eatoires faiblement
d\'ependants}; Samson 2000 \emph{Ann.\ Probab.} 28(1):416).

The sub-block sync overhead $m \cdot \delta_\infty$ and the
finite-precision slack $\varepsilon_{\mathrm{ac}}^{\mathrm{total}}$
are deterministic terms not affecting the concentration of the
random part $\sum_n \ell_n$.
\hfill$\blacksquare$

\emph{Scope and significance.} Theorem 7.14 strengthens Theorem 7.2's
expectation bound to a high-probability concentration bound. Three
practical consequences:

\emph{(a) Predictable archive sizing.} For an archive of $N$ source
bytes, the encoded archive size is $N \cdot h(X) + m \cdot \delta_\infty +
O(\sqrt{N \log(1/\delta)})$ with probability $\geq 1 - \delta$. For
$N = 1\,\mathrm{GB}$, $\eta = 2^{-16}$, $\delta = 10^{-6}$, predictor
context $W_N = O(\log N)$, exp-mixing rate $\rho$ bounded away from
$1$: the rigorous Azuma deviation is
$\sqrt{2 \cdot 10^9 \cdot \log(2/\delta)} \cdot (W_N + 1 + \rho/(1-\rho))
\cdot \log_2(1/\eta)$, which for $W_N = 30, \rho = 0.5$ gives
$\sqrt{5.6 \cdot 10^{10}} \cdot 32 \cdot 16 \approx 1.2 \cdot 10^8$ bits
$\approx 15\,\mathrm{MB}$; the CLT-based estimate via empirical
$\sigma_{\ell_1}^2$ is typically much smaller (per T7.15's worked example,
$\sim 43\,\mathrm{kB}$ at the same $\delta$ for English-text-like
sources). Storage planners typically use the CLT estimate plus a
safety margin; the Azuma bound serves as a worst-case proof of
boundedness.

\emph{(b) Adaptive worst-case sizing.} Storage systems can pre-allocate
$N \cdot h(X) + \mathrm{quantile}_\delta$ space and accept rare overruns
with frequency $\delta$. The deviation is symmetric, so both upper and
lower bounds apply.

\emph{(c) Streaming-encoder buffer sizing.} An online encoder
(Theorem 7.2's per-sub-block scheme) needs to buffer at most
$\mathbb{E}[L^{\mathrm{rep}}_{\mathrm{RNR}}/K] \cdot K + O(\sqrt{K})$ bits
per sub-block with high probability, sizing the encoder's working
memory predictably.

\emph{Relation to CLT.} For finite-variance per-position log-losses,
the Central Limit Theorem additionally gives that
$(L^{\mathrm{rep}} - \mathbb{E}[L^{\mathrm{rep}}])/\sqrt{N \cdot \sigma^2}
\xrightarrow{d} \mathcal{N}(0, 1)$ where $\sigma^2 = \mathrm{Var}(\ell_1)$
under stationarity (martingale CLT, McLeish 1974 \emph{Ann.\ Probab.}\
2(4):620--628). Theorem 7.14 is the non-asymptotic Azuma-Hoeffding
version; the CLT gives sharper asymptotic constants but matches the
$\sqrt{N}$ scaling.

\emph{Remark on the $(W_N+1)$ factor and source-mixing tail.}
The $(W_N+1)$ factor in the bound above tracks the
\emph{direct predictor-context} influence of $X_n$ on the
subsequent log-losses $\ell_n, \ldots, \ell_{n+W_N}$ (positions whose
predictor context-window contains $X_n$). A second, source-mediated
contribution must also be accounted for: under the exp-mixing
hypothesis (5.8), conditioning on $X_n$ also alters the conditional
distribution of $X_{n+W_N+1}, X_{n+W_N+2}, \ldots$ at geometric rate
$\rho^k$, contributing an additional mixing-tail factor
$\sum_{k \geq 1} \rho^k = \rho/(1-\rho)$ in the bound on the
Doob-martingale increment. The full bound is therefore
\[
|D_n| \;\leq\; \bigl(W_N + 1 + \tfrac{\rho}{1-\rho}\bigr)
\cdot \log_2(1/\eta),
\]
\emph{and the concentration becomes}
\[
\bigl|\,L^{\mathrm{rep}}_{\mathrm{RNR}}(X^N) - \mathbb{E}[L^{\mathrm{rep}}_{\mathrm{RNR}}(X^N)]\,\bigr|
\;\leq\;
\sqrt{2 N \log(2/\delta)} \cdot \bigl(W_N + 1 + \tfrac{\rho}{1-\rho}\bigr)
\cdot \log_2(1/\eta)
\]
\emph{with probability at least $1 - \delta$. For exp-mixing with
$\rho$ bounded away from $1$ (the standard assumption (5.8) with
constant $\rho$), $\rho/(1-\rho) = O(1)$ and is dominated by
$W_N = O(\log N)$, so the $(W_N+1)$ leading-order bound stated in
the theorem is asymptotically tight. For polynomially-slowly-mixing
processes (where $1/(1-\rho)$ becomes super-constant in $N$), the
mixing-tail factor would dominate; the present statement therefore
implicitly assumes exp-mixing with constant decay rate. A compensated
martingale $\tilde{M}_N = \sum_{k=1}^N (\ell_k - \mathbb{E}[\ell_k
\mid \mathcal{F}_{k-1}])$ has Azuma increments bounded by
$\log_2(1/\eta)$ without either factor, but concentrates a different
quantity (a centered sum, not $L^{\mathrm{rep}} -
\mathbb{E}[L^{\mathrm{rep}}]$), and so does not bound the deviation
of the encoder output directly.}

\textbf{Theorem 7.15 (End-to-end performance bound for neural-LM-based
RNR: synthesis of T7.2, T7.4, T7.5, T7.14 with practical numerical
prediction).}
\emph{Let $X = (X_n)_{n \geq 1}$ be a stationary exp-mixing source on
finite byte alphabet $|\Sigma| = 256$, and let $M'$ be a pre-trained
neural language model (e.g., transformer, RNN, state-space model)
producing the per-position conditional distribution
$P_{M'}(X_n \mid X_{n - W_N : n - 1})$ with empirical per-position
cross-entropy $H_{M'}$ bits on $X$ (so $H_{M'} = h(X) +
\varepsilon_{\mathrm{KL}}$ with $\varepsilon_{\mathrm{KL}} \geq 0$).
Assume bit-exact integer-arithmetic computation (Theorem 10.6
conformance), $\eta$-floor probability quantization at $p$-bit
precision (so $\eta = 2^{-p}$), and sub-block spacing $K$ in the
balanced regime $K = \Theta(\sqrt{N})$ (Theorem 7.5 with $\alpha
\cdot Q = \Theta(1)$). Then for any $\delta > 0$, with probability
at least $1 - \delta$, the encoded archive size satisfies}
\[
\boxed{
\;
L^{\mathrm{rep},M'}_{\mathrm{RNR}}(X^N)
\;=\;
\underbrace{N \cdot H_{M'}}_{\text{predictor cross-entropy}}
\;+\;
\underbrace{m \cdot \delta_\infty}_{\substack{\text{sync overhead} \\ O(\sqrt{N} \cdot \delta_\infty)}}
\;+\;
\underbrace{m \cdot p \cdot O(1)}_{\substack{\text{AC slack} \\ O(\sqrt{N} \cdot p)}}
\;\pm\;
\underbrace{\sqrt{2 N \log(2/\delta)} \cdot \bigl(W_N + 1 + \tfrac{\rho}{1-\rho}\bigr) \cdot p}_{\substack{\text{concentration} \\ O(\sqrt{N} \log(1/\delta) \cdot W_N p)}}
\;+\;
\underbrace{O(1)}_{\text{header}}
\;
}
\]
\emph{bits, with byte-granular random-access query cost
$O(\sqrt{N} \cdot \mathrm{cost}(M'))$ per query (where
$\mathrm{cost}(M') = O(W_N \cdot d)$ is the neural-network inference
time for $W_N = O(\log N)$ context and $d$ parameters).}

\emph{Numerical prediction for enwik9-scale archive
($N = 10^9$ bytes, English Wikipedia text)}: with a
Chinchilla-70B-class transformer compressing enwik9 to raw rate $8.3\%$
of original size, i.e.\ $H_{M'} \approx 0.664$ bits/byte (Del{\'e}tang
et al.\ 2024 ICLR arXiv:2309.10668, Table 1), $\eta = 2^{-32}$
($p = 32$), $K = \sqrt{N} \approx 31623$, $\delta_\infty \approx 1$ bit
(English text empirical), $\delta = 10^{-6}$:
\begin{itemize}
\item Predictor cross-entropy: $N \cdot H_{M'} = 10^9 \cdot 0.664$ bits
$\approx 79 \,\mathrm{MB}$.
\item Sync overhead: $m \cdot \delta_\infty = (10^9 / 31623) \cdot 1
\approx 3.16 \cdot 10^4$ bits $\approx 4 \,\mathrm{kB}$ (negligible).
\item AC slack: $m \cdot p = (10^9 / 31623) \cdot 32 \approx 1.01 \cdot
10^6$ bits $\approx 123 \,\mathrm{kB}$ (negligible).
\item Concentration --- rigorous Azuma upper bound. The boxed formula
gives $\sqrt{2N \log(2/\delta)} \cdot (W_N + 1 + \rho/(1-\rho)) \cdot p$;
for Chinchilla-70B with $W_N = 2048$ tokens of context, $\rho = 0.5$
(typical English-text mixing), this is $\sqrt{2 \cdot 10^9 \cdot 14.5}
\cdot 2050 \cdot 32 \approx 1.1 \cdot 10^{10}$ bits $\approx 1.4
\,\mathrm{GB}$ at $\delta = 10^{-6}$. This worst-case Azuma bound is
hopelessly loose for realistic transformers and exceeds the archive
size by an order of magnitude --- it serves only as a hard upper-
bound proof of finiteness, not as a practical sizing target.
\item Concentration --- practical CLT-based estimate. The
``Relation to CLT'' remark in Theorem 7.14 (citing McLeish 1974 martingale
CLT) gives a much tighter \emph{asymptotic} estimate: empirical
$\sigma^2_{L^{\mathrm{rep}}/N} \approx \mathrm{Var}(\ell_1)$ for stationary
sources, so total $\sigma_{L^{\mathrm{rep}}} \approx \sqrt{N} \cdot
\sigma_{\ell_1}$. Empirically, $\sigma_{\ell_1} \in [0, p]$ has
\emph{actual} variance typically much smaller than $p^2/12$
(uniform-on-$[0,p]$ worst case); for English text with
$H_{M'} \approx 0.664$ bpb, observed $\sigma_{\ell_1}^2 \approx 4$
giving raw $\sigma_{L^{\mathrm{rep}}} \approx \sqrt{4 \cdot 10^9} \approx
6.3 \cdot 10^4$ bits $\approx 8\,\mathrm{kB}$ (asymptotic CLT std).
For a $1-\delta$ Gaussian two-sided tail, multiply by
$\sqrt{2 \log(2/\delta)}$: at $\delta = 0.317$ ($1\sigma$, factor
$\approx 1.92$): $\approx 15\,\mathrm{kB}$; at $\delta = 10^{-6}$
(factor $\approx 5.4$): $\approx 43\,\mathrm{kB}$ ($\sim 0.05\%$
of archive). This is the
practical sizing estimate, well below the Azuma worst-case bound.
\item Predicted total: $\approx 79\,\mathrm{MB}$ at $\delta = 10^{-6}$,
with rigorous concentration bands of $\pm 1.4\,\mathrm{GB}$
(Azuma worst-case, T7.14) or $\pm 2.1$--$4.2\,\mathrm{MB}$
(variance-aware Bernstein, T7.25 with long-run variance
$v^2 = \sigma_\ell^2 (2W_N+1) + O(B^2)$ and
empirical $\sigma_\ell \approx 2$ bits for English text) and
practical concentration band $\pm 0.04\,\mathrm{MB}$ (CLT-based
empirical), giving $\sim 12.0\times$ compression vs $1\,\mathrm{GB}$
raw text, matching Del{\'e}tang et al.\ 2024 Table 1 measurement.
\item Per-query random-access cost: $O(\sqrt{N} \cdot \mathrm{cost}(M'))
\approx 31623 \cdot \mathrm{transformer\_eval} \approx$ tens of
milliseconds on a single GPU.
\end{itemize}

\emph{\textbf{Proof.}} Substitute the practical neural-LM predictor
$M'$ into Theorem 7.4 ($\varepsilon_{\mathrm{KL}} = H_{M'} - h(X)$),
giving $\mathbb{E}[L^{\mathrm{rep}}] \leq N \cdot H_{M'} + m \cdot
\delta_\infty + O(1) + \varepsilon_{\mathrm{ac}}^{\mathrm{total}}$.
Apply Theorem 7.5's optimal $K^* = \Theta(\sqrt{N})$ for balanced
storage-access workloads, giving $m \cdot \delta_\infty + \alpha Q
K^* = \Theta(\sqrt{N})$. Apply Theorem 7.14's concentration bound
with $\eta = 2^{-p}$, giving deviation
$\sqrt{2 N \log(2/\delta)} \cdot (W_N + 1 + \rho/(1-\rho)) \cdot p$
(predictor-context $W_N = O(\log N)$ plus source-mixing-tail
$\rho/(1-\rho) = O(1)$ under constant $\rho$; see Theorem 7.14's
$(W_N + 1 + \rho/(1-\rho))$-bound and its subsequent correction
remark). Combine all terms.

The $\varepsilon_{\mathrm{ac}}^{\mathrm{total}} = O(m \cdot p)$
follows from Theorem 7.2 convention C2 with per-sub-block stream
$\eta$-floor slack $O(p)$; sum over $m = N/K$ sub-blocks gives
$O((N/K) \cdot p) = O(\sqrt{N} \cdot p)$.

Per-query cost: Theorem 7.5 gives
$O(\log L(R) + (K + k) \cdot \mathrm{cost}(M'))$. For
$L(R) = \mathrm{poly}(N)$ and $K = \Theta(\sqrt{N})$:
$O(\log N + \sqrt{N} \cdot \mathrm{cost}(M')) = O(\sqrt{N} \cdot
\mathrm{cost}(M'))$ dominated by NN inference.
\hfill$\blacksquare$

\emph{Scope and significance.} Theorem 7.15 is the explicit end-to-end
synthesis for the paper's headline use case: \emph{neural lossless
compression with byte-granular random access}. It composes the four
main theorems (T7.2 Shannon-rate, T7.4 KL accounting, T7.5 K* optimal,
T7.14 concentration) into a single bound with concrete numerical
predictions calibrated to the empirical performance of modern
language models. The boxed equation can be used directly as a sizing
formula by implementers: choose $\delta$ for confidence,
$M'$ from available pre-trained LMs (with $H_{M'}$ measured on a
calibration corpus), and compute the predicted archive size with
$\pm$ confidence interval.

\emph{Comparison with existing neural compression systems.}
NNCP (Bellard 2019) and PixelCNN++-based compressors achieve similar
$N \cdot H_{M'}$ rates but lack byte-granular random access (sequential
arithmetic coding only). RNR's contribution is the byte-granular RA
preservation at $O(\sqrt{N})$ per-query cost, paid via the
$m \cdot \delta_\infty + m \cdot p$ sub-block overhead which is
$O(\sqrt{N})$ in total — negligible against the $N \cdot H_{M'}$
predictor cost for any non-trivial $H_{M'}$.

\emph{Implementation requirements (per §10 conformance).} The
predictor $M'$ must be evaluated in bit-exact integer arithmetic
(Theorem 10.6 conformance conditions (a)--(g)) across all encoder
and decoder devices. The $\eta$-floor arithmetic coder is a standard
fixed-point design (Said 2003 \emph{Arithmetic Coding for Data
Compression}, Hewlett-Packard Tech.\ Report; Witten--Neal--Cleary
1987 \emph{Comm.\ ACM} 30(6):520). The sub-block synchronisation
protocol is Theorem 10.3's byte-granular RA construction.

\emph{Side-information framing and model overhead (caveat).}
The $12.0\times$ ratio and $\approx 79\,\mathrm{MB}$ figure above
measure the \emph{archive bitstream} produced by the encoder; they
do \emph{not} include the size of the neural predictor $M'$ itself.
In Shannon-style or side-information-assisted compression
(Slepian--Wolf 1973, Wyner--Ziv 1976), the decoder is assumed to
possess $M'$ as shared side information, in which case the rate
$N \cdot H_{M'}$ is the operational quantity of interest. In strict
algorithmic/Kolmogorov-style framing, by contrast, the archive must
be self-contained: a Chinchilla-70B predictor stored in 16-bit
precision occupies $\approx 140\,\mathrm{GB}$, which dwarfs the
$1\,\mathrm{GB}$ enwik9 source --- yielding \emph{negative savings}
(i.e., the encoded archive expands the input by a factor of roughly
$140\times$) for a single archive of this size. The headline
$12\times$ rate therefore applies in the \emph{amortized} regime where
$M'$ is shared across many archives whose combined size exceeds the
model size (a property the rate-distortion literature calls
``universal'' or ``model-shared'' coding; see also the critique of
unadjusted LM-as-compressor metrics in Zhang et al.\ 2024
\emph{L3TC: Leveraging RWKV for Learned Lossless Low-Complexity Text
Compression}, arXiv:2412.16642). The break-even archive volume against
a fixed 140\,GB-class shared predictor satisfies
$V^* \cdot (\log_2|\Sigma| - H_{M'}) = 140\,\mathrm{GB} \cdot \log_2|\Sigma|$
(equating self-contained archive size with raw source size), giving
$V^* \approx 140\,\mathrm{GB} \cdot 8 / (8 - 0.664) \approx 152\,\mathrm{GB}$
of compressible byte-alphabet text. This break-even is comfortably
exceeded by cloud-backup, archival-storage, content-delivery, and
even individual large-corpus deployments (a single user's photo or
video library, a moderate-size Wikipedia mirror, etc.), but is
\emph{not} satisfied by single-document compression where the source
is megabyte-scale. Theorem
7.21 quantifies the predictor-archive size trade-off explicitly when the
model itself is encoded as part of the archive (knowledge distillation
to predictor $M''$ with bounded KL gap).

\emph{Remark 7.15a (Support selection for the neural-LM field: the §6.9
Gaussian-MESP picture transfers to the discrete autoregressive predictor,
cleaner in submodularity but with a new \#P-hard value-oracle obstacle that
has no Gaussian analogue).} The predictor $M'$ of Theorem 7.15 defines a
\emph{discrete autoregressive entropy field} --- the Type-III-C field of the
neural-LM deployment, with per-position conditionals
$Q_i(\cdot\mid x_{<i})$ in place of the Gaussian $\mathcal N(0,\Sigma)$ of
§6.9. We record how the §6.9 support-selection theory (Lemma 6.9 fixed-field
APX-hardness; Remark 6.9a free-field easiness; Corollary 6.9b greedy
additive-regret encoder) ports to this field. As in §6.9 the
\emph{fixed-field} lever is the operative one for Theorem 7.15, where $M'$
is a pre-trained \emph{shared} predictor: choose a coded set $S$, $|S|=s$,
of sub-block positions to store, the decoder reconstructing the complement
$X_{\bar S}$ from $M'$; the \emph{selection objective} is the model's exact
reconstruction residual for the complement, by the chain rule,
\[
C_{\mathrm{sel}}(S) \;=\; H(X_{\bar S}\mid X_S)
\;=\; H(X) - M(S), \qquad M(T):=H(X_T),
\]
so minimising it is the maximum-entropy-sampling (MESP) selection
$\max_{|S|=s} M(S)$ --- the Shannon-entropy mirror of Lemma 6.9's log-det
MESP (Lemma 6.9 writes the complementary form $H(X_S\mid X_{\bar S})=H(X)-
M(\bar S)$; under $S\leftrightarrow\bar S$ the two are the \emph{same} MESP
problem, and every claim below is identical for either direction --- the
total archive rate is $H(X)$ regardless of $S$, so $C_{\mathrm{sel}}$ is the
selection \emph{objective}, not the whole stored cost). Four facts (each
verified numerically in
\texttt{scripts/verify/remark\_7\_15a\_transformer\_support\_selection.py};
V1--V5):
\begin{enumerate}
\item \emph{(Submodularity is unconditional --- strictly cleaner than the
Gaussian case.)} $M(T)=H(X_T)$ is non-negative, monotone, and submodular
for \emph{every} discrete joint law, because entropy is a polymatroid rank
function (Fujishige, \emph{Information and Control} 39(1):55--72, 1978):
$M(T\cup\{i\})-M(T)=H(X_i\mid X_T)\ge 0$ (monotone) and is non-increasing in
$T$ (submodular, ``conditioning reduces entropy''). Unlike the Gaussian
differential entropy --- which can be negative and whose monotonicity needs
the calibration $\sigma^2\ge 1/(2\pi e)$ shown \emph{necessary} in Remark
6.9d --- discrete entropy needs \emph{no} calibration: that noise floor
simply vanishes (V1).
\item \emph{(Greedy retains the $(1-1/e)$ guarantee and the additive-regret
identity.)} Greedy maximisation of $M$ gives $M(T_{\mathrm{greedy}})\ge
(1-1/e)\,M(T_{\mathrm{opt}})$ (Nemhauser--Wolsey--Fisher, \emph{Math.\
Programming} 14:265--294, 1978), and the residual transfers \emph{additively}
exactly as in Corollary 6.9b: $C_{\mathrm{sel}}(\text{greedy}) -
C_{\mathrm{sel}}(\text{opt}) = M(T_{\mathrm{opt}}) - M(T_{\mathrm{greedy}})
\le (1/e)\,M(T_{\mathrm{opt}})$ --- not a multiplicative ratio, since the cost
is a difference $H(X)-M(S)$ (V2).
\item \emph{(The selection is NP-hard --- hardness is \emph{not} softened by
discreteness.)} Free choice of a size-$s$ coded set to optimise $M$ is
NP-hard already in the discrete setting: it contains \textsc{Max-Coverage}
(NP-hard, Feige, \emph{JACM} 45(4):634--652, 1998). Reduction: $U$ latent
independent fair coins $E_1,\dots,E_U$ (not selectable); one selectable
\emph{position} per candidate set $C_i$ with value $X_i:=(E_e)_{e\in C_i}$,
so $H(X_S)=\bigl|\bigcup_{i\in S} C_i\bigr|$ for any position-subset $S$ and
$\max_{|S|=s}H(X_S)$ is \textsc{Max-Coverage} --- a \emph{free} per-position
selection, faithful to Lemma 6.9's MESP rather than a candidate-restricted
variant (the positions $X_i$ are tuple-valued, within the general
discrete-field scope below). Here the value oracle $H(X_S)=|\bigcup C_i|$ is
poly-time, so the selection is hard \emph{even with an easy oracle} (V4).
\item \emph{(A new value-oracle complexity split with no Gaussian analogue.)}
Greedy needs an efficient value oracle for $H(X_T)$; for an AR field this is
marginal inference, and its complexity splits:
\begin{itemize}
\item[(a)] \emph{order-$1$ Markov, arbitrary $T$:} POLY-TIME. The induced
subset process is again order-$1$ (the observed previous kept symbol is the
full state, so no hidden-state belief branching), giving the chain-rule
entropy $H(X_T)=H(X_{t_1})+\sum_k H(X_{t_k}\mid X_{t_{k-1}})$ in $O(r\,
|\Sigma|^2)$ arithmetic for the entropy assembly plus $O(r\,|\Sigma|^3\log N)$
for the gap-power pairwise marginals (transfer matrix / Chapman--Kolmogorov)
--- polynomial, verified to equal the brute marginal entropy on
every subset, with no $|\Sigma|^{|T|}$ table built;
\item[(b)] \emph{bounded-order-$w$ Markov, \emph{contiguous} block
$T=\{a,\dots,b\}$:} POLY-TIME via the forward algorithm,
$O((b-a)\,|\Sigma|^{w+1})$ (poly for constant $w,|\Sigma|$). Contiguity is a
\emph{sufficient} condition here, not shown necessary: an \emph{arbitrary}
subset of an order-$w\ge 2$ chain induces a hidden-state process (a function of
a Markov chain), and computing $H(X_T)$ becomes the \emph{exact finite-block
Shannon entropy of an HMM} --- each $P(x_T)$ is poly-time (forward algorithm)
but the entropy sums $-P\log P$ over $|\Sigma|^{|T|}$ strings. This case is
\emph{genuinely open}: the order-$1$ method provably fails (the induced process
is not order-$1$; $I(X_{t_k};X_{t_{k-2}}\mid X_{t_{k-1}})>0$, V3b$'$, and a
numerical probe finds its forward belief set does not collapse), \emph{no}
poly-time exact algorithm is known, and \emph{no} \#P-/NP-hardness theorem is
known either --- the adjacent published results concern the entropy \emph{rate}
(Blackwell, \emph{Ann.\ Math.\ Statist.}\ 1957; no closed form), the most-likely
string (Lyng\o{}--Pedersen, \emph{J.\ Comput.\ Syst.\ Sci.}\ 65(3):545--569,
2002, NP-hard), or relative entropy between automata (Cortes--Mohri--Rastogi--
Riley, \emph{Int.\ J.\ Found.\ Comput.\ Sci.}\ 19(1):219--242, 2008,
PSPACE-complete), and the \emph{unambiguous}-PFA entropy is in fact tractable
--- none settles exact finite-block Shannon entropy. It sits, unresolved,
between the poly endpoints (a)/(b) and the \#P-hard general case (c);
\item[(c)] \emph{general AR / high-order arbitrary $T$:} the \emph{exact}
entropy value oracle is \#P-hard. Exact probabilistic inference is already NP-hard (Cooper,
\emph{Artif.\ Intell.}\ 42(2--3):393--405, 1990) and its counting / marginal
version is \#P-hard (Roth, \emph{Artif.\ Intell.}\ 82(1--2):273--302, 1996);
probability reduces to the entropy oracle explicitly: for an
event $E$ computed by a poly-size field coordinate, set $Y=E\wedge B$ with
$B\sim\mathrm{Bern}(1/2)$ independent; then $\Pr(Y{=}1)=\Pr(E)/2\in[0,\tfrac12]$
and $H(Y)=h_2(\Pr(E)/2)$, with $h_2$ strictly increasing on $[0,\tfrac12]$,
so an entropy oracle recovers the \#P-hard marginal $\Pr(E)=2\,h_2^{-1}(H(Y))$
(V5).
\end{itemize}
The Gaussian setting has \emph{neither} obstacle in (c): $\Sigma_T$ is read
off in $O(N^2)$ and its log-det computed in $O(s^3)$. Thus the discrete
neural-LM deployment carries \emph{two} independent obstacles --- NP-hard
selection (3) \emph{and} a \#P-hard value oracle (4c) --- where the Gaussian
theory of §6.9 carries only one (V3).
\end{enumerate}
\emph{Scope (honest, mirroring Lemma 6.9).} This is the \emph{fixed-field}
lever, the regime of Theorem 7.15 ($M'$ pre-trained and shared). In the
\emph{free-field amortised} regime the Remark 6.9a collapse applies verbatim:
a free amortised predictor codes every block at the support-invariant rate
$H(X)$, the chain rule $H(X_S)+H(X_{\bar S}\mid X_S)=H(X)$ erases the
combinatorial selection, and the problem is in $\mathrm P$. The NP-/\#P-hardness
of (3)--(4c) is established for \emph{general} discrete fields supplied by a
poly-size representation; whether it persists on the restricted class of
\emph{actual} transformer softmax fields is a separate question --- the
discrete counterpart of Lemma 6.9's item-(4) $\Sigma$-restriction caveat ---
and is left open. The transfer's two new contributions over §6.9 are
therefore (i) unconditional/cleaner submodularity for the discrete field, and
(ii) a representation-dependent value-oracle complexity split (poly endpoints
vs.\ a \#P-hard general case, with the noncontiguous bounded-order regime left
open) that locates the true practical bottleneck of the Theorem 7.15 encoder
when it performs support selection.

\emph{Remark 7.15b (Why Remark 7.15a's \#P-hardness does not reach the
Theorem 7.15 encoder: causality replaces marginalisation).} Remark 7.15a's
value oracle --- exact subset-entropy / marginal computation for a general
high-order AR field --- is \#P-hard, and a non-contiguous support selection
\emph{forces} it: reconstructing a non-prefix subset requires marginalising
the field over the gap positions (equivalently, the exact marginal of a
computed coordinate, as in the Remark 7.15a/V5 reduction). The Theorem 7.15
streaming coder never forms such a marginal. It codes each
sub-block left-to-right, coding $X_i$ under the predictor's conditional
$Q_i(\cdot\mid x_{<i})$ on the \emph{realised, decoder-reproducible} prefix;
the realised codelength of an input $x$ is
\[
L(x)\;=\;\sum_{i}\,-\log_2 Q_i(x_i\mid x_{<i})
\;=\;-\log_2\!\prod_i Q_i(x_i\mid x_{<i}),
\]
a sum of $K$ single-softmax evaluations --- \emph{one forward pass}, with no
marginalisation of any position. Two consequences:
\begin{itemize}
\item \emph{The causal cost oracle is polynomial for every AR field, of any
effective order.} Evaluating and paying the streaming-coder cost is
$O(K\,|\Sigma|\,\mathrm{cost}(M'))$ per sub-block: conditioning on a realised
prefix is a forward pass, never a marginalisation. This holds verbatim for a
full-attention transformer (effective order $=$ the whole prefix). The
\emph{identity} $\mathbb{E}[L]=\sum_i H(X_i\mid X_{<i})=H(X^N)$ holds (entropy
chain rule; the cross-entropy $N\,H_{M'}$ of Theorems 7.4/7.15 under a
mismatched $Q$), but the encoder never \emph{computes} this expectation ---
which may itself be \#P-hard from a compact high-order description; it
\emph{pays} the realised $L(x)$ (one forward pass) and \emph{estimates} the
rate as an empirical mean over a calibration corpus. So the predictor
cross-entropy term $N\,H_{M'}$ is poly-time \emph{achievable} and
\emph{estimable} with no exact entropy or marginal (the full archive length
adds the Theorem 7.15 sync and AC terms on top of this term).
\item \emph{The \#P-hardness is confined to optional, non-streaming
operations.} An exact marginal appears only in operations the deployed coder
does not perform --- chiefly the non-contiguous support selection of Remark
7.15a (coding a non-contiguous set and reconstructing its complement from a
non-prefix context), and exact expected-rate analysis. The streaming coder and
decoder, and the Theorem 7.15 synthesis, never invoke one; the \#P barrier is
therefore orthogonal to the deployed encoder's operation. \emph{In particular
this includes random access:} byte-granular RA (Theorems 7.2/7.5) restores the
predictor state from the nearest sync point and decodes the containing
sub-block \emph{causally} forward to the queried position --- conditioning only
on the realised in-block prefix, never reconstructing a non-prefix subset ---
so RA is itself a forward pass, $O(K\cdot\mathrm{cost}(M'))$, and inherits none
of the \#P-hardness. The non-contiguous reconstruction (which would require the
acausal conditional $P(x_i\mid x_S,\dots)$ with $S$ containing positions to the
right of $i$, hence a runtime marginal) is the \emph{separate, optional}
Type-III-C selection lever, not the random-access mechanism.
\end{itemize}
The structural moral is the autoregressive one: a model that only ever
conditions on realised data sidesteps the marginal-inference barrier its own
joint law carries --- exact high-order marginal inference being \#P-hard in
general (Remark 7.15a), though poly in the order-1 and contiguous cases --- and
the barrier re-materialises exactly when a coding decision forces a general
exact \emph{marginal} (non-contiguous reconstruction) rather than a
\emph{conditional} (causal coding). Verified in
\texttt{scripts/verify/remark\_7\_15b\_causal\_cost\_oracle.py} (C1--C4:
chain-rule identity; forward-only codelength $=-\log_2 P(x)$ with sample mean
$\to H$; non-contiguous residual requires a genuine gap marginal with
exponential op-count; order-independence at the full-attention regime).

\emph{Remark 7.15c (The open-middle oracle is poly-time $\varepsilon$-approximable
under filter stability: the \#P-hardness is an exact-computation artifact, dissolved
only in the approximate oracle).} Remark 7.15a
left the noncontiguous bounded-order oracle's \emph{exact} complexity open (it is
the exact finite-block entropy of the induced process $Y_k:=X_{t_k}$, a subsample
of $X$). Its \emph{operative} status is settled under a filter-stability hypothesis.
\emph{Hypothesis (FS).} The induced subsample process --- the order-$w$ hidden
context observed at $\{t_k\}$ through the deterministic readout $Y_k=X_{t_k}$, with
the chain evolving \emph{unobserved} across the gaps --- is \emph{uniformly
filter-stable}: its prediction filter (the belief over the hidden context given the
observed subsample) forgets its initialisation geometrically in total variation,
uniformly in the realisation. Equivalently, (FS) is exactly the displayed subsample
conditional-MI bound below. This is the standard exponential-filter-stability
conclusion, driven here by contraction of the unobserved gap-transitions; it holds
when the order-$w$ hidden chain is geometrically ergodic, under the usual
filter-stability conditions adapted to a function-of-Markov-chain observation
(Atar--Zeitouni 1997; Le Gland--M\'evel, \emph{Math.\ Control Signals Systems}
13:63--93, 2000; cf.\ Lemma 5.6e). Note (FS) is \emph{stronger} than the bare
conditional-MI mixing (5.8) of Theorem 7.2: $\sigma(Y)\subseteq\sigma(X)$ transfers
$\beta$-mixing but \emph{not} the conditional-MI bound below (conditional MI is not
monotone under dropping the skipped mediators), so (FS) is the right primitive.
Under (FS) the conditional mutual information
$I(Y_k;Y_{<k-d}\mid Y_{k-d:k-1})\le C\rho^d$ decays geometrically (the filter
contracts uniformly in total variation, and longer physical gaps $t_{k}-t_{k-1}$
only contract it faster). Hence the order-$d$ truncation
\[
A_d \;:=\; \sum_k H\!\left(Y_k\mid Y_{k-d:k-1}\right)\;\ge\; H(Y_1,\dots,Y_r),
\qquad
0\le A_d - H \;=\; \sum_k I\!\left(Y_k;Y_{<k-d}\mid Y_{k-d:k-1}\right)\;\le\; r\,C\,\rho^d,
\]
over-estimates the exact block entropy with \emph{geometric} error (conditioning
on less only raises entropy; the gap is the windowed-out tail dependence). $A_d$
uses only $(d{+}1)$-block marginals of $Y$, each computed by the forward algorithm
over the $|\Sigma|^w$ hidden state, so for
$d=O\!\big(\log(r/\varepsilon)/\log(1/\rho)\big)$ --- \emph{logarithmic} in
$1/\varepsilon$ --- the oracle is evaluated to additive error $\varepsilon$ in
$\mathrm{poly}(r,1/\varepsilon)$ time (table size
$|\Sigma|^{d+1}=(r/\varepsilon)^{O(\log|\Sigma|/\log(1/\rho))}$, not the
$|\Sigma|^{|T|}$ of exact brute marginalisation). The exponent
$O(\log|\Sigma|/\log(1/\rho))$ makes this polynomial in $(r,1/\varepsilon)$ \emph{for
any fixed alphabet and mixing rate}, but of \emph{degree} that itself grows as the
source mixes more slowly ($\rho\to1$) --- a slice-wise (``XP'') guarantee, not a
fixed-degree one. And greedy support selection
driven by the $\varepsilon$-accurate oracle attains $M(\text{greedy})\ge
(1-\tfrac1e)M(\text{opt})-O(s\varepsilon)$ (additive oracle error propagates
additively through the monotone-submodular greedy); taking $\varepsilon=
\varepsilon'/s$ --- adding $O(\log s)$ to the depth --- gives
$(1-\tfrac1e)M(\text{opt})-\varepsilon'$ at $\mathrm{poly}(r,1/\varepsilon')$ cost.
\emph{So the Remark 7.15a hardness is an artifact of \emph{exact} computation, not
of the mixing regime}: mixing does \emph{not} make the exact oracle polynomial ---
it is suspected to remain hard even on mixing sources --- but under (FS), and given
the bounded-order/finite-HMM representation, the operative
$\varepsilon$-\emph{approximate} oracle, and hence the whole greedy support
selection, is poly-time computable (in the slice-wise sense above). This complements
Remark 7.15b (causal coding and random access are already \emph{exactly} poly): the
optional non-contiguous \emph{selection} lever is operationally tractable too on a
filter-stable source; only the exact worst-case oracle carries the unresolved
exact-computation barrier (the suspected hardness) that Remark 7.15a leaves open.
Verified in
\texttt{scripts/verify/remark\_7\_15c\_mixing\_oracle\_approx.py} (M1--M5: monotone
geometric over-estimate $\to H$; decay rate tracking the mixing; logarithmic depth
budget; robust $(1-\tfrac1e)-O(\varepsilon)$ greedy).

\emph{Remark 7.15d (R\'enyi spectrum of the open-middle oracle: every integer order
is exactly poly, the min-entropy is NP-hard, the operative Shannon point is not
reached).} Remark 7.15a
left the \emph{exact} Shannon oracle $H(Y_1,\dots,Y_r)$ open and Remark 7.15c made it
$\varepsilon$-approximable under (FS); we now \emph{locate} the difficulty exactly,
via the R\'enyi spectrum $H_\alpha(Y)=\tfrac1{1-\alpha}\log_2\sum_y P(y)^\alpha$
($H_1$ Shannon, $H_\infty$ min-entropy). For \emph{integer} $\alpha\ge2$,
\[
\sum_y P(Y_T{=}y)^\alpha \;=\; \Pr\bigl(\alpha\text{ i.i.d.\ copies of }X
\text{ agree on }T\bigr),
\]
the partition function of $\alpha$ independent copies of the chain constrained to
agree at the subset positions and evolving freely across the gaps --- computed by a
\emph{forward} algorithm over the $\alpha$-fold product state ($|\Sigma|^{w\alpha}$),
poly in $(r,N)$ for fixed $\alpha,w$ (gaps by $O(\log)$ matrix powers), with
\emph{no} $\log$-of-a-sum and \emph{no} $|\Sigma|^{|T|}$ marginal. Hence \emph{$H_\alpha$
is poly-time exact for every integer $\alpha\ge2$ --- unconditionally, with no mixing
hypothesis}. At the other end $H_\infty=-\log_2\max_y P(y)$ is the most-likely subset
tuple, i.e.\ the HMM most-likely-\emph{output-string} / consensus problem (not the
poly Viterbi most-likely \emph{path}), \emph{NP-hard} (Lyng\o{}--Pedersen, \emph{J.\
Comput.\ Syst.\ Sci.}\ 65(3):545--569, 2002). The hardness is for our specific
subclass, not merely by membership in HMMs: any HMM is encodable as an order-$2$
chain alternating hidden-state and emitted-symbol positions with $T$ the emitted
positions, so $\max_y P(X_T{=}y)$ over noncontiguous subsets of bounded-order chains
already captures general HMM consensus. The Shannon point $H_1$ is \emph{not reached}
by this route: by R\'enyi monotonicity $H_2\le H_1\le H_0$ the exact poly $H_2$ is a
\emph{free, hypothesis-free poly lower bound} on the open oracle (and
$H_0=\log_2|\mathrm{supp}|$ bounds it above), but the moments do not \emph{recover}
$H_1$ --- analytic continuation to $\alpha=1$ would need more than finitely many
orders, which the per-order cost $|\Sigma|^{w\alpha}$ (constant for fixed $\alpha$
but growing with $\alpha$) precludes, and the non-integer orders $\alpha\in(1,\infty)$
themselves lack the finite-copy structure. So the spectrum exhibits a
\emph{representational discontinuity} rather than a clean isolation: the \emph{integer}
orders are finite-dimensional product-chain partition functions (poly), the
min-entropy is NP-hard, and the operative Shannon entropy --- the $\alpha\to1$
derivative, with no finite-moment representation --- remains open, now flanked by the
exact poly lower bound $H_2$ and the (FS) approximation of Remark 7.15c. This
complements Remark 7.15c with an \emph{exact, unconditional} poly handle ($H_2$, also
a lower bound) and a precise complexity-landscape map of the remaining open problem.
Verified in
\texttt{scripts/verify/remark\_7\_15d\_renyi\_spectrum.py} (R1--R4: the
product-forward DP equals $\sum_y P(y)^\alpha$ for $\alpha=2,3$ to machine precision,
no $|\Sigma|^{|T|}$ table; spectrum ordering; min-entropy maximiser).

\emph{Remark 7.15e (Mapping the wall: the exact oracle is in $\mathrm{FP}^{\#P}$ and
both natural attacks are obstructed --- the \#P-hardness route is self-defeating, the
poly route is Blackwell-blocked).} A genuine attempt to settle the exact Shannon
oracle of Remarks 7.15a/d. It does not break the wall, but pins it on both sides.
\emph{(E1) Upper bound: the oracle is in $\mathrm{FP}^{\#P}$.} Computing
$H(Y)=-\sum_{y\in\Sigma^r}P(y)\log_2 P(y)$ to $2^{-\mathrm{poly}}$ precision is a
weighted sum, over an exponential index set, of terms each poly-time computable (the
forward algorithm returns $P(y)$ as a poly-bit rational, and the transcendental
$\log_2$ to $2^{-\mathrm{poly}}$ precision in poly time); approximating each
$-P(y)\log_2 P(y)$ to guard bits, scaling to integers, and summing over $y$ with a
$\#P$ oracle places the value in $\mathrm{FP}^{\#P}$ --- computable by a poly-time
machine with a $\#P$ oracle. (The output is real/irrational, so this is an
$\mathrm{FP}^{\#P}$ approximation, not a $\#P$ function; the logs do not push the
class higher.) \emph{(E2) The
predicate-encoding \#P-hardness route is self-defeating for bounded-order chains.}
The general-AR \#P-hardness of the entropy oracle (Remark 7.15a, V5) encodes a
\#P-hard predicate $E$ into a coordinate $Y=E\wedge B$, so $H(Y)=h_2(\Pr(E)/2)$
reveals $\Pr(E)=\#E/|\Sigma|^n$. For a bounded-order-$w$ chain the coordinate $Y=E(a)$
must be computed \emph{as the chain generates} $a=a_1\cdots a_n$ \emph{once}, tracking
$E$ in its state of size $|\Sigma|^w$ --- i.e.\ $E$ is a width-$|\Sigma|^w$
\emph{read-once} branching program. But a general CNF needs \emph{exponential}
read-once width (it cannot be encoded), and for any read-once-encodable $E$,
model-counting $\#E$ is in $\mathrm P$ (a bottom-up DP over the BP nodes;
Van den Broeck--Suciu 2017). The gap positions do not provide a loophole: the
unobserved gaps marginalise by \emph{the same} forward DP (a matrix-product
walk-count, $\#\mathrm L\subseteq\mathrm{FP}$), the chain generates its sequence
\emph{once} (so it cannot re-read for multi-pass power), and a \emph{memoryless}
bounded-state chain cannot enforce the global (bijection / no-repeat) constraints
that make counting $\#P$-hard --- the predicate bit must be produced by the bounded
state regardless of where the assignment sits. So the \emph{direct} (inline)
predicate-encoding route reveals only $\mathrm{FP}$-computable counts: \emph{it can
exhibit only poly quantities, never \#P-hardness}. That is exactly the boundary
between the general-AR oracle (\#P-hard) and the bounded-order open middle (open).
(This bars the specific V5 mechanism, not every conceivable reduction.) \emph{(E3) The obvious
poly algorithm is Blackwell-blocked.} $H(Y)=\sum_k H(Y_k\mid Y_{<k})$ needs the filter
belief over the hidden state, whose reachable set does \emph{not} collapse (the
Blackwell measure is generically infinitely supported --- the 7.15a open-middle
probe), so the filtering sum branches exponentially; the finite-block entropy is a
Lyapunov-type functional of random matrix products, whose rate has no closed form
(Blackwell, \emph{The entropy of functions of finite-state Markov chains}, Trans.\
First Prague Conf.\ on Information Theory, 1957, pp.\ 13--20; the generically
infinite-dimensional predictive features are characterised by Jurgens--Crutchfield,
\emph{Chaos} 2021, arXiv:2008.12886). \emph{Net:} the wall
holds against both natural attacks --- the exact oracle sits in $[\,$poly (collision
$H_2$, Remark 7.15d)$\,,\ \mathrm{FP}^{\#P}\,]$, with the standard \#P-hardness
reduction obstructed (E2, read-once) and the standard algorithm obstructed (E3,
Blackwell). It is \emph{not broken}; but the map shows any resolution needs a
\emph{non-predicate-encoding} hardness construction or a \emph{non-filtering}
algorithm. Verified in \texttt{scripts/verify/remark\_7\_15e\_wall\_map.py} (W1:
read-once model-counting $=$ brute, poly; W2: the predicate-encoding reveals exactly
that poly count; W3 illustrates the weighted-sum form --- the
$\mathrm{FP}^{\#P}$ membership is the analytic argument above, not a numerical check).

\emph{Remark 7.15f (An exact-tractable island of the open middle: finite finite-prefix
predictive belief --- strictly beyond finite-order Markov --- and a belief-geometry
reading of the §7.15 map.)} The open middle (exact finite-block entropy of a
function-of-a-Markov-chain; Remarks 7.15a--7.15e), bracketed in $[\,\text{poly }H_2,\
\mathrm{FP}^{\#P}\,]$, has a clean \emph{exact}-tractable island. Let the
\emph{finite-prefix predictive belief} $b_t$ be the law of the hidden/causal state given
the finite observed past $X_1^{t-1}$, started from the stationary prior. \emph{When the
reachable set $\mathcal B$ of such finite-prefix beliefs is finite}, $b_t$ is a
finite-state Markov chain and, by causal sufficiency $\Pr(X_t\mid X_1^{t-1})=\Pr(X_t\mid
b_t)$,
\[
  H(X^N)\;=\;\sum_{t=1}^N\mathbb E_{\nu_t}\!\bigl[\,H\bigl(\Pr(X_t\mid b_t)\bigr)\,\bigr],
  \qquad \nu_t=\text{law of }b_t,\quad
  \nu_{t+1}(b')=\!\!\sum_b\nu_t(b)\!\!\sum_{x:\,b\xrightarrow{x}b'}\!\!\Pr(x\mid b),
\]
is computable in $O(N\cdot|\mathcal B|\cdot|\Sigma|)$ once $\mathcal B$ is known ---
\emph{polynomial}, with no $2^N$ enumeration. (The algorithm is the standard finite
$\varepsilon$-machine / computational-mechanics belief recursion; the point here is which
sources it covers and how they sit in the §7.15 map.) Three points:
\begin{enumerate}
\item \emph{\textbf{Strictly beyond finite-order Markov (7.15a).}} An order-$k$ Markov
source has $\le|\Sigma|^k$ point-mass beliefs, so 7.15a's order-$k$/contiguous cases are
a sub-class; the converse fails. The \emph{even process} (binary; $1$s occur only in
even-length runs --- the canonical \emph{sofic} source) has just \emph{four} reachable
finite-prefix beliefs, yet is not any finite-order Markov chain: for \emph{every} $k$,
two pasts ending in a $1$-run longer than $k$ of opposite parity share their last $k$
symbols but have different next-symbol laws ($\Pr(\text{next}=0)=\tfrac12$ after an even
run, $0$ after an odd one). So the island captures genuinely-hidden sofic sources no
fixed Markov order reaches.
\item \emph{\textbf{The condition is finite-prefix belief, \emph{not} finite Blackwell
support.}} Finiteness of $\mathcal B$ implies the Blackwell measure (the recurrent,
infinite-past belief law; Remark 7.15e) has finite support, but \emph{not conversely}:
a synchronising HMM (states $A,B$; symbols $a,b$ reveal the state, $1$ does not, with
self-loop probabilities $p\neq q$) has Blackwell support $\{\delta_A,\delta_B\}$, yet the
prefixes $1^n$ give pairwise-distinct beliefs $\propto(\pi_A p^n,\pi_B q^n)$ ---
\emph{infinitely many} finite-prefix beliefs. There they \emph{concentrate}
geometrically ($1^n\!\to\delta_A$), so that source is exactly the filter-stable /
Remark 7.15c \emph{approximable} case, not the exact island. The finite-block entropy is
driven by the finite-prefix (transient-inclusive) set, which is what must be finite.
\item \emph{\textbf{A belief-geometry reading of the §7.15 map.}} Organising sources by
the finite-prefix belief set: \emph{finite} $\Rightarrow$ exact poly (this island,
strictly containing finite-order Markov); \emph{infinite but geometrically concentrating}
(filter stability) $\Rightarrow$ poly-$\varepsilon$ (Remark 7.15c); \emph{proliferating
/ non-concentrating} (a generic noisy-emission HMM, whose belief set fills a continuum)
$\Rightarrow$ the obvious exact belief-DP is blocked, the generic $[\text{poly }H_2,
\mathrm{FP}^{\#P}]$ open case (Remarks 7.15d/e). The last is not asserted to be exactly
where $\#P$-hardness lives (7.15e leaves that open); only that the finite/concentrating
structure that makes the first two tractable is absent.
\end{enumerate}
Verified in \texttt{scripts/verify/remark\_7\_15f\_epsilon\_machine\_island.py} (V1 the
even process: finite belief set $|\mathcal B|{=}4$ and the poly belief-MC entropy matches
brute force to machine precision; V2 it is not order-$k$ Markov; V3 the finite-order
containment; V4 a noisy HMM with a continuum-filling belief set; V5 the belief-MC block
entropy at any $N$; V6 the synchronising HMM --- finite Blackwell support but infinitely
many concentrating finite-prefix beliefs, the FS case). The exact Shannon complexity of
the proliferating case is pinned in Remark 7.15g ($\#P$-hard); this remark identifies a standard exact-tractable
island (strictly beyond finite-order Markov) and organises the surrounding §7.15 cases by
the finite-prefix belief geometry.

\emph{Remark 7.15g (The §7.15 open middle is \#P-hard: a \#SAT reduction to the exact
coefficient-query HMM entropy).} Remarks 7.15d/e bracketed the exact finite-block entropy of a
function of a finite Markov chain in $[\,\text{poly }H_2,\ \mathrm{FP}^{\#P}\,]$ and left its
true complexity open, noting that the obvious predicate-encoding route self-defeats (a
bounded-order HMM is a read-once branching program, and counting its paths to a fixed output is
in $\mathrm{FP}$). We close the hardness side, in the natural exact-symbolic representation.
Since a finite-block entropy of a rational-parameter HMM is a rational combination
$H(X^N)=\sum_q\kappa_q\log q$ over primes $q$ (the $\kappa_q\in\mathbb Q$ unique, as prime
logarithms are $\mathbb Q$-linearly independent by unique factorisation), the natural exact query
is the \emph{coefficient problem}
\[
  \textsc{Coeff-HMM-Entropy}(M,N,q)\ :=\ \kappa_q,\ \text{the coefficient of }\log q\text{ in }H_M(X^N).
\]
\emph{This problem is \#P-hard, and in $\mathrm{FP}^{\#P}$, hence $\mathrm{FP}^{\#P}$-complete}
(Turing reductions). (We do \emph{not} claim hardness of returning the entire expansion, which
can be exponentially large, nor of a bounded-precision numerical value; see Scope.)

\emph{\#P-hardness (\#SAT $\le$ \textsc{Coeff-HMM-Entropy}).} Given a proper 3-CNF $\varphi$ with
$n$ variables and $m$ clauses (each clause on three distinct variables, so it has exactly
$2^{n-3}$ falsifying assignments; \#SAT for such formulas is \#P-complete; WLOG $\varphi$ is nonempty with
$m\ge1$, $n\ge3$, the $m{=}0$/$n{<}3$ cases being trivial), fix a prime
$P>\max(m,2)$ (so $P$ is odd and $1\le m<P$) and build a finite-state HMM emitting a length-$n$
block: the hidden chain first picks a mode --- \emph{dummy} with probability $8P/(m{+}8P)$, then
emitting a uniform assignment; or \emph{clause $i$} with probability $1/(m{+}8P)$ each, then
emitting a uniform assignment that \emph{falsifies} clause $i$ (its three variables fixed to the
falsifying bits, the other $n{-}3$ fair). The hidden state set is
$\{\text{dummy},\text{clause }1,\dots,m\}\times\{\text{position}\}$, size $O(mn)$ (the random
emissions absorb into a function-of-a-Markov-chain with constant-factor blow-up). With $u(x)$ the
number of clauses $x$ falsifies, the block marginal is $p(x)=(P+u(x))/Z$, $Z=2^{n-3}(8P+m)$
(indeed $\sum_x u(x)=m2^{n-3}$ gives $\sum_x p(x)=1$). Grouping by $u(x)=j$ with
$c_j=\#\{x:u(x)=j\}$ --- so $c_0=\#\mathrm{SAT}(\varphi)$ ---
\[
  H(X^n)=\log Z-\frac1Z\sum_{j=0}^{m}c_j\,(P+j)\log(P+j).
\]
Because $P$ is an \emph{odd} prime with $P\nmid Z=2^{n-3}(8P+m)$ (as $P\neq2$ and
$8P{+}m\equiv m\not\equiv0\bmod P$) and $P\nmid(P{+}j)$ for $1\le j\le m$ (as $0<j<P$), the prime
$P$ occurs only in the $j{=}0$ term, so
\[
  \kappa_P=-\frac{c_0\,P}{Z}=-\frac{\#\mathrm{SAT}(\varphi)\cdot P}{Z},
\]
and one query $\textsc{Coeff-HMM-Entropy}(M,n,P)$ returns $\#\mathrm{SAT}(\varphi)=-\kappa_P Z/P$.

\emph{$\mathrm{FP}^{\#P}$ upper bound.} Clearing denominators, $p(x)=w(x)/W$ with $w(x)$ a
poly-evaluable nonnegative integer (a weighted path count) and the common denominator
$W=\sum_x w(x)$. Then $\kappa_q=v_q(W)-\frac1W\sum_x w(x)\,v_q(w(x))$, where $v_q$ is the
$q$-adic valuation (convention $w(x)v_q(w(x))=0$ when $w(x)=0$); both $W$ and
$\sum_x w(x)v_q(w(x))$ are exponential sums of poly-evaluable nonnegative integers, hence
computable with a $\#P$ oracle. So
$\textsc{Coeff-HMM-Entropy}\in\mathrm{FP}^{\#P}$, giving $\mathrm{FP}^{\#P}$-completeness.

\emph{Corollary (multi-information).} The same reduction makes the exact (coefficient-query)
total correlation $\mathrm{TC}=\sum_i H(X_i)-H(X^N)$ (Theorem 5.5) \#P-hard: the single-symbol
marginals $H(X_i)$, hence their prime-log coefficients, are poly-computable, so
$\kappa_P(H(X^N))=\kappa_P\!\bigl(\sum_i H(X_i)\bigr)-\kappa_P(\mathrm{TC})$ and an exact-$\mathrm{TC}$
oracle again yields $\#\mathrm{SAT}$ (whether or not the marginals involve $P$; verified V7).

\emph{Corollary (the exact ideal RNR rate is \#P-hard --- the operational reading).} The
operationally central case is the ideal RNR repair length itself: by the chain rule the exact
expected ideal finite-block repair length under true causal arithmetic coding is
$L^\ast_{\mathrm{RNR}}(M,N)=\sum_{t=1}^N H_M(X_t\mid X_{<t})=H_M(X^N)$, so its coefficient-query
problem \emph{is} $\textsc{Coeff-HMM-Entropy}$ --- $\mathrm{FP}^{\#P}$-complete. Thus computing the
\emph{exact} ideal RNR archive size for a finite-state neural-surrogate source is \#P-hard, even
though the deployed encoder runs in one forward pass per symbol (Remark 7.15b): the intractability
is in pricing the optimum exactly, not in coding, which is precisely why RNR reports achievable
rates and bounds rather than the exact ideal.

\emph{Why this evades the Remark 7.15e obstruction.} The read-once/forward-DP objection of 7.15e
is about counting hidden paths to a \emph{fixed} word, which is in $\mathrm{FP}$. The entropy
instead \emph{groups} paths by their observed word and weights the grouped masses by $\log$; the
reduction exploits exactly that group-and-$\log$ structure (the collision multiplicities $c_j$
enter through $(P+j)\log(P+j)$, and the prime $P$ isolates $c_0$). It does not encode CNF
evaluation into one bounded-state predicate, so it is the resolution 7.15e flagged as missing,
not a contradiction of it.

\emph{Scope (honest).} (i) The hardness is of the \emph{coefficient query} (equivalently the
exact symbolic value), via the $\mathbb Q$-independence of prime logs; it does \emph{not} claim
hardness of a bounded-precision \emph{numerical} approximation, whose value folds all $c_j$ into
one real and need not expose $c_0$. This is complementary to Remark 7.15c, where the
\emph{approximate} open-middle oracle is poly-$\varepsilon$ under filter stability --- so the
open middle is \emph{exact-\#P-hard yet $\varepsilon$-approximable} (under FS), a clean
exact/approximate dichotomy. (ii) The HMM grows with $\varphi$ (the model-plus-$N$ input
convention of the 7.15d/e bracket); it is hardness in the model size, not for a fixed HMM as
$N\to\infty$. (iii) Adjacent known results are about different quantities: entropy
\emph{approximation} for circuit samplers is NISZK-complete (Goldreich--Sahai--Vadhan 1999) and
entropy \emph{difference} SZK-complete (Goldreich--Vadhan 1999), with a logspace/branching-program
analogue (Dvir--Goldwasser--Rothblum--Vadhan 2011); HMM \emph{min}-entropy / consensus-string is
NP-hard (Lyngs\o--Pedersen 2002); integer R\'enyi entropy is poly (Breitner--Skorski 2020); and
relative entropy of probabilistic automata is studied separately (Cortes--Mohri--Rastogi--Riley).
The exact \emph{Shannon block}-entropy \#P-hardness above is the direct statement for the §7.15
open middle and we did not find it published. Together with the finite-belief poly island
(Remark 7.15f, exact) and the FS poly-$\varepsilon$ oracle (Remark 7.15c, approximate), this pins
the open middle: the exact coefficient query is $\mathrm{FP}^{\#P}$-complete, the tractable
carve-outs being finite finite-prefix belief (exact) and filter stability (approximate).
Verified in \texttt{scripts/verify/remark\_7\_15g\_hmm\_entropy\_sharp\_p\_hard.py} (V1
$p(x)=(P{+}u(x))/Z$, $\sum_x p=1$; V2 the mode$\times$position HMM reproduces the marginal; V3
$c_0=\#\mathrm{SAT}$; V4 the entropy identity; V5 the reduction --- the exact $\log P$ coefficient
$=-c_0P/Z$ recovers $\#\mathrm{SAT}$ across SAT/mixed/UNSAT formulas, $P$ dividing no other term;
V6 the exact-vs-approximate separation; V7 the multi-information corollary --- exact $\mathrm{TC}$
recovers $\#\mathrm{SAT}$ via poly-computable marginals).

\emph{Remark 7.15h (The tractable boundary is poly belief-\emph{growth}, not finite belief: finite
belief is sufficient but not necessary).} Remark 7.15f frames a \emph{finite} finite-prefix
predictive-belief set as the exact-tractable island; that is \emph{sufficient but not necessary}. The
sharp sufficient condition is poly belief-\emph{growth}: let $B(N)$ be the number of distinct beliefs
reachable in $\le N$ steps from the stationary prior (the exact belief-DP costs
$\Theta(N\,B(N)\,|\Sigma|)$), so $B(N)=\mathrm{poly}(N)\Rightarrow$ exact $H(X^N)$ is poly-time, which is
strictly weaker than $|\mathcal B|<\infty$. \emph{The clean finite/infinite dichotomy is false:} the
discrete-time \emph{renewal} process (binary, with infinite-support inter-event law; Marzen--Crutchfield,
\emph{Inf.\ \& Causal Architecture of Discrete-Time Renewal Processes}) has a \emph{countably infinite}
belief set --- one distinct point-mass belief (hazard) per phase $=$ time since the last event --- yet
$B(N)=N{+}1$ and the exact finite-block entropy is computed in $O(N^2)$ by a phase DP
$H(X^N)=\sum_t\mathbb E_{\nu_t}[h_2(\mathrm{hazard})]$ (a countable \emph{unifilar} HMM; the spectral
$h_\mu(L)$ of Riechers--Crutchfield is poly exactly when the mixed-state operator is finite-rank). This
refines the §7.15 map into a belief-\emph{growth} trichotomy: (I) poly growth $B(N)=\mathrm{poly}(N)$
(finite belief 7.15f, and the infinite-belief renewal/countable-unifilar class) $\Rightarrow$
\emph{exact-poly}; (II) geometrically-concentrating infinite belief (filter stability, Remark 7.15c)
$\Rightarrow$ \emph{poly-$\varepsilon$} (exact complexity open); (III) exponential growth
$B(N)=2^{\Theta(N)}$ (continuum/fractal Blackwell measure) $\Rightarrow$ \emph{$\#P$-hard} --- the
regime of the 7.15g $\#\mathrm{SAT}$ construction, whose belief set we verify proliferates as $2^N$.
\emph{Conjecture:} exact $H(X^N)$ is poly-time iff $B(N)=\mathrm{poly}(N)$ --- the forward direction is
the above; the converse (exponential belief-growth $\Rightarrow$ not poly unless $\mathrm P=\#\mathrm P$,
beyond the single 7.15g family) is \emph{open}, as is the exact complexity of the filter-stable middle
(II). \emph{Validated} in \texttt{scripts/verify/remark\_7\_15h\_belief\_growth\_boundary.py} (the
infinite-belief renewal phase-DP matches brute-force $2^N$ enumeration to $10^{-13}$ for $N\le13$ and
reaches $N{=}1000$; the finite-support renewal recovers the finite-belief 7.15f island).

\emph{Remark 7.15i (Regime II is non-trivial: ergodic filter stability does not collapse exact
belief-growth, and the §7.15g hardness does not reach it).} Two facts pin the difficulty of the
filter-stable middle (regime II) of the 7.15h trichotomy. \emph{First, regime II genuinely contains
exponential exact belief-growth.} Generic stationary \emph{ergodic} non-unifilar finite-state HMMs have
a \emph{contracting} mixed-state (Blackwell-filter) IFS --- hence are filter-stable (geometric
forgetting; the 7.15c / Atar--Zeitouni hypothesis) --- \emph{and} an uncountable, positive-dimension
(fractal) mixed-state attractor (J\"urgens--Crutchfield, \emph{Chaos} 31:083114, 2021,
arXiv:2102.10487; companion arXiv:2008.12886). The two coexist because filter stability contracts the
\emph{distance between} beliefs from differing \emph{initialisations} but does not \emph{merge} the
\emph{distinct exact} beliefs reached along distinct \emph{words}: it bounds the $\varepsilon$-distinct
beliefs to $\mathrm{poly}$ (the 7.15c truncation) while the \emph{exact}-distinct count $B(N)$
proliferates as $2^N$ on the fractal attractor. We verify a minimal witness --- a generic ergodic
$2$-state binary-output non-unifilar HMM with TV-filter-contraction $\rho\approx0.24$ yet exact
$B(N)=2^{N+2}-1$, against a \emph{unifilar} control of the same forgetting rate with $B(N)\equiv3$
(exact-poly, regime I) --- so filter stability is \emph{orthogonal} to exact belief-growth and ergodic
forgetting does \emph{not} force exact-poly; the exact belief-DP (cost $\Theta(N\,B(N))$, 7.15h) is
blocked in (II) exactly as in (III). \emph{Second, the 7.15g $\#P$-hardness does not transfer to (II).}
That witness encodes $\#\mathrm{SAT}$ in a \emph{frozen-mode} mixture (a latent mode chosen once, never
transitioning) --- the antithesis of forgetting; the only adjacent exact-hardness for an HMM functional,
the $\#P$-hardness of HMM total-variation distance (Kiefer, ICALP 2018, arXiv:1804.06170), likewise
relies on non-forgetting reachability-tracking, and the consensus-string $\mathrm{NP}$-hardness
(Lyngs\o--Pedersen, \emph{JCSS} 65(3), 2002) and entropy-\emph{rate} Lyapunov/joint-spectral-radius
intractability (Jacquet--Seroussi--Szpankowski, \emph{TCS} 2008) are a different (rate, not finite-block)
object or non-forgetting. \emph{Verdict: regime II exact complexity is open.} The mixed-state operator is
infinite-rank when $d_\mu>0$ (so no finite-rank poly algorithm, Riechers--Crutchfield), yet no
forgetting-compatible $\#\mathrm{SAT}$ reduction is known. \emph{Conjecture:} exact $H(X^N)$ on
stationary ergodic filter-stable fractal-Blackwell sources ($d_\mu>0$) is $\#P$-hard, the missing
ingredient being a reduction that reads $\#\mathrm{SAT}$ from a prime-log coefficient of $H(X^N)$ carried
by the fractal Blackwell mass structure (rather than a persistent latent mode); equivalently, regime II
is exact-poly only if ergodic mixing provably ``forgets the count,'' for which there is no theorem.
\emph{What is known is a one-sided obstruction: the 7.15g prime-log mechanism does not transfer through
ergodic forgetting.} The most favourable adaptation is \emph{periodic injection} --- a period-$n$ hidden
clock that draws a \emph{fresh} 7.15g mode each block (the strongest forgetting: the mode is re-randomised
every period). With the block boundaries \emph{observable} the source synchronises, the belief set
collapses to $\mathrm{poly}$, and the count survives additively --- but that is exactly regime I, not the
fractal middle; with the boundaries \emph{hidden} (so the observer of the stationary stream must average
over the uniform phase, as in the cyclostationary stationarisation $y_n:=x_{n+\theta}$, Gardner) the
$\log P$ coefficient of $H(X^N)$ is no longer $-c_0P/Z$: the non-linear $-\sum_x p\,v_P(p)$ structure means
the phase-average does not pass through to the coefficient, and the cross-boundary windows replace the
single-aligned-assignment statistic $u(x)$ by a product over independent fresh assignments. We verify that
the resulting hidden-phase coefficient is not even a function of the collision multiset $\{c_j\}$ --- three
$3$-CNF formulas with \emph{identical} $\{c_j\}$ (hence identical $\#\mathrm{SAT}$) give three
\emph{different} coefficients --- let alone an affine function of $\#\mathrm{SAT}$, so no single coefficient
query recovers $\#\mathrm{SAT}$ by the 7.15g route. (This obstructs the \emph{mechanism}, not the
conjecture: it does not show regime II is sub-$\#P$, and is consistent with ergodic mixing being the very
tool that makes the \emph{approximate} count easy --- MCMC, and the 7.15c poly-$\varepsilon$ oracle under
filter stability.) \emph{Validated} in
\texttt{scripts/verify/remark\_7\_15i\_ergodic\_fs\_belief\_growth.py} (V1 ergodicity;
V2 TV-contraction $\rho\approx0.24$; V3 exact $B(N)=2^{N+2}-1$; V4 unifilar control $B(N)\equiv3$) and the
phase-blur obstruction in
\texttt{scripts/verify/remark\_7\_15i\_phase\_blur\_obstruction.py} (V1 aligned single block
$\kappa_P=-\#\mathrm{SAT}\,P/Z$; V2 additive over aligned blocks; V3 the hidden-phase coefficient $\neq$ the
aligned one and does not recover $\#\mathrm{SAT}$; V4 identical-collision formulas give distinct hidden-phase
coefficients).

\textbf{Theorem 7.16 (Adaptive sub-block sizing $K_n$ for piecewise-
stationary sources: strict improvement over fixed-$K$ Theorem 7.5
when local statistics vary).}
\emph{Let $X = (X_n)_{n \geq 1}$ be a piecewise-stationary source
partitioned into $R$ regions $\{\mathcal{R}_1, \ldots, \mathcal{R}_R\}$
of lengths $\{N_1, \ldots, N_R\}$ ($\sum_r N_r = N$), where each
region $\mathcal{R}_r$ is stationary exp-mixing with per-region entropy
rate $h_r$ and Crutchfield--Feldman excess $\delta_\infty^{(r)}$ (and
optionally per-region predictor cross-entropy $H_{M'}^{(r)}$, AC
precision $p_r$). Theorem 7.5's optimal sub-block size minimised over
\emph{adaptive} per-region $K_r$ gives}
\[
K_r^* \;=\; \sqrt{\gamma_r \cdot N_r / (\alpha_r \cdot Q_r)},
\quad\text{where } \gamma_r = \delta_\infty^{(r)} + \log(1/\eta_r),
\]
\emph{with corresponding minimum total cost}
\[
C(\{K_r^*\}) \;=\; \sum_r \Bigl[\,N_r \cdot h_r + 2 \sqrt{\gamma_r N_r \alpha_r Q_r}\,\Bigr] + O(1)
\;=\;
N \cdot \bar{h} + 2 \sum_r \sqrt{\gamma_r N_r \alpha_r Q_r} + O(1).
\]
\emph{Compared to fixed-$K$ uniform choice $K^* = \sqrt{N \bar{\gamma}/(\alpha Q)}$
(Theorem 7.5 with averaged $\bar{\gamma}, \alpha, Q$), the adaptive
optimum is \emph{strictly smaller} whenever the per-region quantities
$\gamma_r, \alpha_r, Q_r$ are not all equal, with relative savings}
\[
\Delta := \frac{C(K^*_{\text{fixed}}) - C(\{K_r^*\})}{C(K^*_{\text{fixed}})}
\;=\;
\frac{2 \sqrt{N \bar{\gamma} \bar{\alpha Q}} - 2 \sum_r \sqrt{\gamma_r N_r \alpha_r Q_r}}{C(K^*_{\text{fixed}})}
\;>\; 0
\]
\emph{by Cauchy--Schwarz, with equality iff all per-region parameters
coincide.}

\emph{\textbf{Proof.}} Direct calculus: minimize
$\sum_r [N_r h_r + \gamma_r N_r / K_r + \alpha_r Q_r K_r]$ over
$K_1, \ldots, K_R$ subject to no coupling between regions (each region
independently chooses its $K_r$ since sub-block sync points respect
region boundaries by construction). The per-region term is
$N_r h_r + \gamma_r N_r/K_r + \alpha_r Q_r K_r$; partial derivative
w.r.t.\ $K_r$ gives $K_r^* = \sqrt{\gamma_r N_r / (\alpha_r Q_r)}$
(Theorem 7.5 specialized to region $r$).

For the strict inequality versus uniform $K^*$: Cauchy--Schwarz applied
to vectors $(\sqrt{\gamma_r N_r})$ and $(\sqrt{\alpha_r Q_r})$ gives
$\sum_r \sqrt{\gamma_r N_r \alpha_r Q_r} \leq \sqrt{(\sum_r \gamma_r N_r)
(\sum_r \alpha_r Q_r)} = \sqrt{N \bar{\gamma} \cdot \bar{\alpha Q}}$,
with strict inequality unless $\gamma_r N_r / (\alpha_r Q_r)$ is
constant in $r$ (i.e., all regions equivalent).
\hfill$\blacksquare$

\emph{Practical applications for neural lossless compression.}
\emph{(a) HTML / Wikipedia text: } markup tags ($h_{\text{tag}} \approx
0.3$ bpb, low entropy) vs prose ($h_{\text{prose}} \approx 1.0$ bpb)
vs URLs ($h_{\text{url}} \approx 2.0$ bpb). Adaptive $K_r$ uses smaller
sub-blocks in high-entropy URL regions ($K_r \sim 256$) and larger in
low-entropy markup ($K_r \sim 4096$).

\emph{(b) Code repositories: } comments ($h \approx 1.0$ bpb), syntax
($h \approx 0.5$ bpb), whitespace ($h \approx 0.2$ bpb), strings
($h \approx 1.5$ bpb).

\emph{(c) Mixed-modality archives (text + images + tables):} different
per-region $h_r$ and $\delta_\infty^{(r)}$; adaptive $K$ exploits the
variation.

\emph{Magnitude of savings.} The savings are dominated by the
\emph{relative} variance of $\gamma_r = \delta_\infty^{(r)} +
\log(1/\eta_r)$ across regions. When the predictor precision $p =
\log(1/\eta)$ dominates over the per-region excess
$\delta_\infty^{(r)}$ (e.g., $p = 32$ bits while
$\delta_\infty^{(r)} \in [0.1, 2.0]$ bits), the relative variation of
$\gamma_r$ is small ($\Delta \sim 5\%$) and adaptive savings are
modest ($\sim 0.02$--$0.05\%$ on the sync-overhead component, $\sim 0$
on total). When $\delta_\infty^{(r)}$ dominates (low-precision AC
$p \leq 8$, or sources with widely varying $\delta_\infty$),
relative variation is larger and savings reach $\sim 5\%$ on
sync-overhead. The empirical magnitude depends on the
$\mathrm{Var}(\gamma_r) / \mathbb{E}[\gamma_r]^2$ ratio.

\emph{Implementation.} The region boundaries are pre-computed by a
segmentation pass (e.g., HTML parser, code lexer, content-type
sniffer); each region's $(h_r, \delta_\infty^{(r)})$ is estimated
from a calibration corpus or empirical sub-block measurements. The
$K_r$ values are recorded in the archive header $\Theta$ (one
$\log K$-bit integer per region, bounded total $R \log K = \mathrm{polylog}(N)$
contribution).

\emph{Relation to Theorem 7.5.} For $R = 1$ (single region = full
archive), Theorem 7.16 recovers Theorem 7.5 exactly. The adaptive
framework strictly subsumes the fixed-$K$ framework.

\textbf{Theorem 7.17 (Edit-tolerant RNR: localized re-encoding cost
under in-place updates, insertions, and deletions).}
\emph{Let $X^N$ be an archive compressed by RNR (Theorem 7.2 framework)
with sub-block size $K$. For each of the following edit operations on
the source, the re-encoding cost is bounded by a constant multiple of
the sub-block size $K$, independent of the archive size $N$:}

\emph{(a) [Single-byte replacement at position $i$]} \emph{Re-encode
the unique sub-block containing position $i$ ($j_i := \lfloor i/K
\rfloor$). Cost: $O(K \cdot \mathrm{cost}(M'))$ predictor evaluations,
$O(K)$ source-stream reads, $\Theta(K \cdot H_{M'})$ bits archive
delta. No other sub-block is touched. Header $\Theta$ updates: the
new sub-block's offset, hash, and (if adaptive) $K_r$ value.}

\emph{(b) [Range replacement of $M$ bytes starting at position $i$]}
\emph{Re-encode the $\lceil M/K \rceil + 1$ sub-blocks intersecting
the modified range $[i, i+M-1]$. Cost: $O(M \cdot \mathrm{cost}(M'))$,
$O(M \cdot H_{M'})$ bits archive delta. No re-encoding of unmodified
sub-blocks.}

\emph{(c) [Insertion of $M$ bytes at position $i$ under content-defined
chunking]} \emph{Assuming the encoder uses \emph{content-defined
chunking} (CDC) --- sub-block boundaries determined by rolling-hash
fingerprints of the source content (Rabin 1981 \emph{Fingerprinting by
Random Polynomials}, Center for Research in Computing Tech., Harvard;
Muthitacharoen--Chen--Mazi{\`e}res 2001 \emph{SOSP '01} ``A low-bandwidth
network file system''), rather than fixed-position $K$-aligned
boundaries --- insertion of $M$ bytes affects at most $O(1)$
sub-blocks straddling the insertion point. Re-encode those affected
sub-blocks (cost $O((M+K) \cdot \mathrm{cost}(M'))$) and shift
sync-metadata offsets for subsequent sub-blocks by $\Delta L$. The
header offset shift is $O(\log L(R))$ per sub-block, cumulative
$O(m \cdot \log L(R)) = \mathrm{polylog}(N)$ bits. \emph{Without CDC}
(fixed $K$-aligned boundaries), insertion realigns ALL subsequent
sub-blocks --- each now covers different source bytes than before, and
all $m - \lfloor i/K \rfloor$ subsequent sub-blocks must be re-encoded,
giving $O(N)$ cost; in this case insertion is no better than monolithic
re-encoding.}

\emph{(d) [Deletion of $M$ bytes at position $i$ under content-defined
chunking]} \emph{Symmetric to insertion: re-encode the affected
sub-blocks (CDC absorbs the deletion locally), shift offsets of
subsequent sub-blocks. Cost: $O((M+K) \cdot \mathrm{cost}(M'))$ +
header updates. Without CDC, all subsequent sub-blocks shift and
require re-encoding.}

\emph{Comparison with monolithic compressors.} Standard compressors
(gzip, zstd, Brotli without restart points) require \emph{full
re-encoding} of the archive after any edit, costing
$O(N \cdot \mathrm{cost}(\text{compressor}))$. RNR's sub-block
structure reduces this to $O(K)$ per edit operation, a speedup factor
of $N/K = m$. For $N = 1\,\mathrm{GB}$, $K = 1024$: $m = 10^6\times$
speedup on single-byte edits.

\emph{\textbf{Proof.}} \emph{(a) Single-byte replacement.}
Each sub-block in Theorem 7.2's RNR is encoded independently (sync
points at $K$-spaced boundaries reset the arithmetic coder and
predictor context). Hence changing one byte affects only the
arithmetic-coded payload of the sub-block containing it.
Re-encoding requires:
\begin{itemize}
\item Re-read the modified $K$-byte sub-block (memory access $O(K)$).
\item Re-run the predictor $M'$ on the new sub-block, generating $K$
positions of conditional distributions ($K \cdot \mathrm{cost}(M')$
time).
\item Re-encode via $\eta$-floor arithmetic coder ($K$ AC operations).
\item Write new payload bits to archive ($K \cdot H_{M'}$ bits in
expectation).
\item Update header: this sub-block's hash, offset (if length
changed), and (in adaptive $K$ regime) $K_r$ value.
\end{itemize}
Total cost: $O(K \cdot \mathrm{cost}(M'))$. No other sub-block touched
because Theorem 7.2's per-sub-block independence guarantees the
payload of sub-block $j$ depends only on $X_{[jK : (j+1)K)}$.

\emph{(b) Range replacement.} Affected sub-blocks are those with
indices $[\lfloor i/K \rfloor, \lfloor (i+M-1)/K \rfloor]$, totaling
at most $\lceil M/K \rceil + 1$ sub-blocks. Each re-encoded as in (a),
giving total cost $O(M \cdot \mathrm{cost}(M'))$.

\emph{(c, d) Insertion/deletion.} Re-encoding the affected sub-block
(which may grow/shrink by up to $M$ bytes) requires the same $O((M+K)
\cdot \mathrm{cost}(M'))$ work. Subsequent sub-blocks' encoded
payloads are UNCHANGED (their source content is unaffected), but
their \emph{offsets} in the archive shift by the bit-length change
$\Delta L$. The offset shift requires a sweep through the offset
table in the header, costing $O(m \cdot \log L(R)) = \mathrm{polylog}(N)$
bits and $O(m)$ time. The PAYLOAD bytes of unaffected sub-blocks are
not re-read, re-decoded, or re-encoded.
\hfill$\blacksquare$

\emph{Scope and significance.} Theorem 7.17 establishes RNR as a
practical compression scheme for \emph{mutable archives}: filesystems
(file edits), databases (record updates), version-controlled archives
(commit deltas), live document collaboration (text edits). The
$O(K)$ per-edit cost contrasts sharply with the $O(N)$ cost of
monolithic compressors, enabling responsive interactive applications
on multi-GB archives.

\emph{Comparison with related schemes.} Existing approaches for
mutable compressed archives:
\emph{(a) gzip / zstd / brotli + restart points (Z\_FULL\_FLUSH):}
require encoder-side restart-point spacing similar to RNR's sub-blocks,
but lack RNR's byte-granular RA without restart-to-restart decoding.
\emph{(b) bzip2 + block independence:} similar block-independent
structure to RNR, but block size $K = 900\,\mathrm{kB}$ (much larger
than RNR's typical $K$), and lacks adaptive predictor / byte-granular RA.
\emph{(c) git pack format:} delta-encoded objects, supports edits but
requires recomputing deltas between revisions; not random-access.
\emph{(d) Compressed B-trees / LSMs:} per-page compression with
$O(\text{page size})$ edit cost; similar to RNR but page-level
granularity, lacks neural-LM predictor integration.

\emph{Implementation requirements.} The edit operations require the
encoder to:
\begin{itemize}
\item Maintain a header offset table indexed by sub-block index,
updatable in $O(\log m)$ per edit (e.g., via a packed array or
order-statistic tree).
\item Re-validate sub-block hashes (Theorem 9.1 verified-only context
update) after re-encoding to maintain context-propagation invariants.
\item Handle archive-bit-length changes via an indirection table
(constant-time offset lookup after shift).
\end{itemize}

\textbf{Theorem 7.18 (Parallelizable RNR encoding: data-parallel
linear speedup up to $P = m$ workers).}
\emph{Let $X^N$ be encoded by RNR (Theorem 7.2 framework) with sub-block
size $K$ and $m = \lceil N/K \rceil$ sub-blocks. The per-sub-block
independence of Theorem 7.2 implies the encoding task decomposes into
$m$ independent jobs, each requiring $O(K \cdot \mathrm{cost}(M'))$
work, with no inter-job synchronisation during encoding. Hence for $P$
parallel workers:}
\[
T_{\mathrm{encode}}^{(P)}(X^N)
\;=\;
\Theta\!\Bigl(\,\frac{m}{P} \cdot K \cdot \mathrm{cost}(M')\,\Bigr) + T_{\mathrm{merge}}
\;=\;
\Theta\!\Bigl(\,\frac{N}{P} \cdot \mathrm{cost}(M')\,\Bigr) + O(m \cdot \log L(R))
\]
\emph{for any $P \leq m$, with $T_{\mathrm{merge}} = O(m \cdot \log L(R))$
the post-encoding header assembly cost (offset table + per-sub-block
hashes). Speedup over the sequential encoder ($P = 1$):
$T^{(1)} / T^{(P)} \to P$ as $N \to \infty$ (linear speedup up to
$P = m$, scattering bounded thereafter by Amdahl's law on the merge
step).}

\emph{Decoder parallelism.} The decoder admits the same parallelism
for batch decoding (decoding multiple sub-blocks concurrently for
multi-byte queries). Per-query single-byte decoding still requires
$O(K \cdot \mathrm{cost}(M'))$ on a single worker (the byte's
containing sub-block).

\emph{\textbf{Proof.}} Theorem 7.2's encoding partitions $X^N$ into $m$
disjoint sub-blocks $\{X_{[jK : (j+1)K)}\}_{j=0}^{m-1}$. For each $j$,
the encoder runs the $W_N$-bounded predictor with cold context
(no information from sub-blocks $j' \neq j$) and arithmetic-codes the
sub-block independently. The per-sub-block work
$O(K \cdot \mathrm{cost}(M'))$ depends only on $X_{[jK : (j+1)K)}$ and
the model $M'$ (shared, read-only across workers).

For $P$ workers: assign each worker $\lceil m/P \rceil$ consecutive
sub-blocks (or any balanced partition). Each worker processes its
sub-block subset in $O(\lceil m/P \rceil \cdot K \cdot \mathrm{cost}(M'))
= O(N/P \cdot \mathrm{cost}(M'))$ time, with no inter-worker
communication during encoding (only at the end, to merge headers).

Header merge: each worker produces its sub-blocks' encoded payloads
plus offset/hash metadata. Final assembly computes cumulative offsets
($O(m)$ time, $O(\log L(R))$ bits per offset) and concatenates
payloads. Total merge: $T_{\mathrm{merge}} = O(m \cdot \log L(R)) =
O((N/K) \cdot \log L(R))$. For $K = \mathrm{polylog}(N)$ this is
$O(N \cdot \log L(R) / \mathrm{polylog}(N))$ --- sub-linear in $N$
(strictly between $\mathrm{polylog}(N)$ and $N$), and asymptotically
dominated by the per-worker encoding cost $N \cdot \mathrm{cost}(M')/P$
whenever $P \cdot \log L(R) = o(K \cdot \mathrm{cost}(M'))$.

Hence
$T^{(P)} = N \cdot \mathrm{cost}(M') / P + O((N/K) \log L(R))$,
giving asymptotic linear speedup $T^{(1)}/T^{(P)} \to P$ as $N \to
\infty$ for any fixed $P \leq m$. For $P > m$, additional workers are
idle (Amdahl saturation at $P = m$).
\hfill$\blacksquare$

\emph{Practical realisation for neural lossless compression.}

\emph{(a) GPU batch encoding.} Modern GPUs (NVIDIA B100/B200 datacenter
or RTX 5090 Blackwell consumer, AMD MI325X) have $P \sim 10^4$--$10^5$
parallel cores and dedicated transformer-inference units (Tensor Cores
gen-5, AMD CDNA-3 Matrix Cores). For a $1\,\mathrm{GB}$ archive with
$K = 1024$, $m = 10^6$: full GPU utilisation ($P = 10^4$) gives
$10^4 \times$ speedup over single-core encoding. Wall-clock encoding
time on B200: $\sim 100\,\mathrm{ms}$ ($10^6$ sub-blocks of $1024$
bytes each, transformer inference at $\sim 1\,\mu\mathrm{s}$/sub-block
batched).

\emph{(b) Distributed multi-node encoding.} For archives at
$N = 10^{12}$ bytes (TB-scale, e.g., LLM training datasets):
$m = 10^9$ sub-blocks, $P = 10^3$ nodes give $10^3 \times$ speedup;
each node processes $10^6$ sub-blocks in parallel using its local GPU.

\emph{(c) Streaming encoding pipeline.} Even without parallel
workers, the per-sub-block independence enables \emph{pipelined}
encoding where one sub-block is being predicted while the previous is
being arithmetic-coded. Throughput speedup $\approx 2 \times$ at
no parallelism cost.

\emph{Limitations.} The predictor $M'$ itself may have internal
parallelism (e.g., transformer attention is parallelizable across
positions within a sub-block). Theorem 7.18 quantifies inter-sub-block
parallelism, on top of which intra-sub-block parallelism in $M'$ adds
a multiplicative factor. The full parallel cost is
$\mathrm{cost}(M') / P_{\mathrm{intra}}$ per position, giving
$T^{(P, P_{\mathrm{intra}})} = O(N \cdot \mathrm{cost}(M') /
(P \cdot P_{\mathrm{intra}}))$.

\emph{Comparison with sequential compressors.} gzip, zstd, brotli are
\emph{inherently sequential} due to causal arithmetic-coder state
dependencies on prior positions. Parallel variants require restart
points (Z\_FULL\_FLUSH), which effectively replicate RNR's sub-block
structure at coarser granularity. RNR's sub-block design natively
supports the parallelism without auxiliary tuning.

\textbf{Theorem 7.19 (Streaming RNR encoder: $O(K + \mathrm{size}(M'))$
memory and $K$-position bounded latency for memory-constrained devices,
under footer-layout convention).}
\emph{The Theorem 7.2 encoder admits a streaming implementation
processing the source $X^N$ as a sequential stream, with the following
resource bounds (under the convention that sub-block metadata --- offsets,
hashes, MAC tags --- is written as a \emph{footer} after the payload, in
the ZIP central-directory style, rather than a prefix header):}
\begin{itemize}
\item \emph{\textbf{Memory}: $O(K + \mathrm{size}(M'))$ bits encoder
state for the active sub-block plus predictor, \emph{plus} an
incrementally-accumulating footer buffer of $O(m \cdot \log L(R))
= O(N/K \cdot \log L(R))$ bits (or a streamed footer if the storage
medium supports random-access seek after payload emission).
For the on-line encoder: total memory is $O(K + \mathrm{size}(M') +
(N/K) \cdot \log L(R))$, with the last term sub-linear in $N$ for
$K = \mathrm{polylog}(N)$. Specifically: $K$-byte source buffer +
arithmetic-coder state $O(p)$ bits + predictor model parameters
$\mathrm{size}(M')$ bits + predictor hidden state
$O(W_N \cdot d_{\mathrm{hidden}})$ bits + footer accumulator
$O((N/K) \log L(R))$ bits.}
\item \emph{\textbf{Latency}: At most $K$ source positions of input
read before first output bit emitted (encoder must read full sub-block
before encoding). Steady-state: one sub-block of output emitted per
$K$ source positions consumed.}
\item \emph{\textbf{Throughput}: $\mathrm{cost}(M')^{-1}$ source
positions per unit time per single-thread encoder (assuming $M'$
inference is the bottleneck).}
\item \emph{\textbf{Bounded output rate}: average output bits per
source position $\bar{r} = H_{M'} \leq \log_2 |\Sigma|$, with
high-probability deviation per Theorem 7.14.}
\end{itemize}
\emph{The decoder admits the symmetric streaming pattern: $O(K +
\mathrm{size}(M'))$ memory state, $K$-position output latency from
seek to first decoded byte.}

\emph{\textbf{Proof.}} The Theorem 7.2 encoding partitions $X^N$ into
$m = \lceil N/K \rceil$ sub-blocks, each encoded independently.
The streaming pattern: maintain a $K$-byte ring buffer of recent
source positions; when buffer fills, run the $W_N$-bounded predictor
$M'$ over the sub-block (one forward pass per position), arithmetic-
code the sub-block using the $\eta$-floor distribution, emit the
encoded payload to the output stream, and reset for the next
sub-block. Predictor state $M'$ is loaded once at encoder startup
and shared across all sub-blocks (read-only); per-sub-block hidden
state $O(W_N \cdot d_{\mathrm{hidden}})$ is recomputed from cold
context.

Memory: only the current sub-block ($K$ bytes), the predictor model
($\mathrm{size}(M')$, e.g., a 1B-parameter LM at 16-bit quantization
is $2\,\mathrm{GB}$; smaller distilled LMs at $\sim 100\,\mathrm{M}$
parameters fit in $200\,\mathrm{MB}$), and the AC state ($O(p)$ bits)
are stored. Total: $O(K + \mathrm{size}(M'))$, independent of $N$.

Latency: the encoder must read $K$ source positions before producing
any output (sub-block-level granularity). For $K = 1024$ bytes:
$1\,\mathrm{kB}$ input latency, $\sim 1\,\mathrm{kB} / H_{M'}$ output
latency.

The decoder mirrors this: load model once, allocate $K$-byte sub-block
buffer, AC state $O(p)$ bits. To answer a byte query at position $i$,
seek to sub-block $j_i = \lfloor i/K \rfloor$, read its slice from the
archive ($O(\log L(R))$ for offset lookup), AC-decode the $K$ positions
(or up to $i \bmod K$ for early termination), return the requested byte.
\hfill$\blacksquare$

\emph{Practical realisation for memory-constrained devices.}

\emph{(a) Mobile / smartphone neural-LM compression.} A flagship
smartphone (e.g., 2025 Snapdragon X3 Elite or Apple A19) has
$\sim 16\,\mathrm{GB}$ DRAM; allocating $200\,\mathrm{MB}$ for a
distilled language model ($\sim 100\,\mathrm{M}$ parameters, INT8
quantized) plus $1\,\mathrm{kB}$ sub-block buffer fits comfortably.
Encoding throughput on Neural Engine: $\sim 100\,\mathrm{MB/s}$.

\emph{(b) Embedded / IoT (ESP32, Cortex-M).} For ultra-constrained
devices ($\sim 256\,\mathrm{kB}$ RAM), use a tiny predictor ($\sim
100\,\mathrm{kB}$ model, e.g., $4$-grams or small Markov chain instead
of NN), $K = 256$ bytes sub-block, total state $\sim 200\,\mathrm{kB}$.
The framework gracefully degrades from neural LM to
classical-predictor RNR while preserving sub-block structure.

\emph{(c) Streaming-data pipelines.} Continuous data sources (sensor
networks, log streams) can be compressed in real time with bounded
encoder lag of $K$ positions ($1$--$4$ kB typical). For $K = 1024$
and source rate $1\,\mathrm{MB/s}$: encoder lag $\sim 1\,\mathrm{ms}$.

\emph{Comparison with other streaming compressors.} gzip in streaming
mode buffers up to $32\,\mathrm{kB}$ window + DEFLATE state $\sim
64\,\mathrm{kB}$; zstd in streaming uses $\sim 1\,\mathrm{MB}$ block
buffer. RNR with $K = 1024$ + distilled LM achieves competitive
compression at $\sim 200\,\mathrm{MB}$ total memory (dominated by
predictor model), \emph{or} matches gzip's $\sim 100\,\mathrm{kB}$
memory footprint with a tiny classical predictor.

\emph{Relation to other §7 theorems.} Theorem 7.19 specifies an
\emph{implementation pattern} compatible with Theorems 7.2 (Shannon
rate), 7.4 (KL accounting), 7.14 (concentration), 7.17 (edit-tolerance),
and 7.18 (parallel encoding). The streaming encoder can be combined
with parallel encoding (Theorem 7.18) by streaming each sub-block
range to a different worker, recovering both parallelism and bounded
per-worker memory.

\textbf{Theorem 7.20 (Cache-aware RNR decoder: amortized per-query
cost under workload locality).}
\emph{Let RNR be deployed with sub-block size $K$, $m = \lceil N/K
\rceil$ sub-blocks, and let $\mathcal{Q} = (b_1, b_2, \ldots, b_T)$
be a sequence of $T$ byte-granular access queries against the archive.
Equip the decoder with an LRU cache of capacity $C$ \emph{decoded
sub-blocks} (storing the recovered $K$-byte source content and the
predictor's terminal hidden state for each cached sub-block). Define
the \emph{cache hit rate}}
\[
\rho_C(\mathcal{Q})
\;:=\;
\frac{1}{T} \sum_{t=1}^{T} \mathbf{1}\{\,\lfloor b_t/K \rfloor \in \mathrm{LRU}_{C}(b_1, \ldots, b_{t-1})\,\},
\]
\emph{the fraction of queries whose containing sub-block is in the LRU
cache at the moment of the query. Then the amortized per-query
decoding cost is}
\[
\boxed{
\;
\bar{T}_{\mathrm{query}}(\mathcal{Q}, C)
\;=\;
\rho_C \cdot O(1) + (1 - \rho_C) \cdot O\bigl(K \cdot \mathrm{cost}(M')\bigr).
\;
}
\]
\emph{For workloads with locality $\rho_C \to 1$ (sequential reads,
range scans, hot-key access), the amortized cost approaches the
cache-hit constant $O(1)$ — exponentially better than the cold-cache
$O(K \cdot \mathrm{cost}(M'))$.}

\emph{Workload analysis for typical patterns.}
\begin{itemize}
\item \emph{(Uniform random)}: $\rho_C = C/m$, no savings for $C \ll m$.
\item \emph{(Sequential scan)}: $\rho_C = 1 - 1/K$ (each new query
either hits same sub-block or fetches next). Amortized $O(1)$ per
query.
\item \emph{(Range scan over $L$ bytes)}: $\rho_C = 1 - K/L$. For
$L \gg K$, near-perfect locality.
\item \emph{(Zipf hot-key with skew $\alpha$)}: We compute the hit rate
of the Perfect-LFU / static-top-popularity cache (an upper bound on
LRU under the Independent Reference Model; Breslau et al.\ 1999
\emph{Proc.\ IEEE INFOCOM}, ``Web caching and Zipf-like distributions:
evidence and implications'').
The hit rate is the cumulative probability of the top-$C$ Zipf-ranked
items $\rho_C = \sum_{k=1}^C k^{-\alpha} / \sum_{k=1}^m k^{-\alpha}$,
with two regimes distinguished by normalizer convergence:
\emph{(divergent normalizer)} $0 < \alpha \leq 1$ gives
$\rho_C \approx (C/m)^{1-\alpha}$ for $\alpha < 1$
(equivalently, $\rho_C \to 0$ as $m \to \infty$ for fixed $C/m$ small),
and the harmonic limit $\rho_C \approx \log C / \log m$ at $\alpha = 1$;
\emph{(convergent normalizer)} $\alpha > 1$ gives the saturating
limit
$\rho_C \to 1 - \dfrac{C^{1-\alpha}}{(\alpha-1)\,\zeta(\alpha)}$
as $C \to \infty$ (using $\sum_{k>C} k^{-\alpha} \sim
C^{1-\alpha}/(\alpha-1)$ from Euler--Maclaurin).
Note: the Zipf distribution is power-law and hence \emph{heavy-tailed}
in both regimes (in the precise probability-theoretic sense of
sub-exponential tails); the $\alpha > 1$ vs $\alpha \leq 1$ split
concerns normalizer integrability on $\mathbb{N}$, not tail
heaviness. The LRU competitive analysis under IRM and general
arrival processes is in Ferragut, Rodriguez, and Paganini (2018)
\emph{Queueing Systems} 88(3):207--241 (online 2017, formal volume
2018), Theorem~3 (general timer-based-cache hit-rate, with LRU as
a limit).
\item \emph{(LRU under bursty workloads)}: $\rho_C$ tracks workload's
working-set size; cache performance characterised by Belady's
optimal-replacement bound and online competitive ratio
$\rho_{\mathrm{LRU}}/\rho_{\mathrm{Bel\acute{a}dy}} \leq C/(C - C_{\mathrm{wss}} + 1)$
(Sleator--Tarjan 1985 \emph{Comm.\ ACM} 28(2):202--208 ``Amortized
efficiency of list update and paging rules'').
\end{itemize}

\emph{\textbf{Proof.}} The decoder's per-query cost decomposes into:
\begin{itemize}
\item \emph{Cache hit} (sub-block already decoded): $O(1)$ to look up
the cached decoded bytes by index (assuming hash-table or direct-
indexed cache), $O(1)$ to extract byte $b \bmod K$ from the cached
$K$-byte buffer. Total: $O(1)$ per query.
\item \emph{Cache miss}: fetch encoded sub-block payload from archive
($O(\log L(R))$ for offset lookup), AC-decode the $K$ positions
($O(K \cdot \mathrm{cost}(M'))$), evict LRU entry, install new
sub-block in cache.
\end{itemize}
Amortizing over $T$ queries: total cost $T \cdot \bar{T}_{\mathrm{query}}$
where $\bar{T}_{\mathrm{query}} = \rho_C \cdot O(1) + (1 - \rho_C)
\cdot O(K \cdot \mathrm{cost}(M'))$.

The cache-hit rate $\rho_C$ depends on workload locality; classical
paging analysis applies (LRU is $C/(C - C_{\mathrm{wss}} + 1)$-
competitive against the optimal offline algorithm).
\hfill$\blacksquare$

\emph{Practical implications for neural lossless compression.}

\emph{(a) Database scan with neural-compressed pages.} A typical
database table-scan reads pages sequentially; cache hit rate $\rho_C
\approx 1$. Decoder amortizes $O(1)$ per page access despite each
page being a separately-encoded sub-block. Effective decoding
throughput: limited by L1/L2 cache memory bandwidth (\textit{not}
predictor inference cost).

\emph{(b) Hot-key access patterns.} For Zipf-distributed access with
$\alpha = 1.5$ (heavily skewed, convergent-normalizer regime) and
$m = 10^6$, $C = 100$: the exact static-top hit rate
$\rho_C = \sum_{k=1}^{100} k^{-1.5} / \sum_{k=1}^{10^6} k^{-1.5}
\approx 0.9243$ (92\% hit rate). For the typical Breslau measurement
range $\alpha \in [0.7, 0.85]$ (divergent-normalizer regime, web
traces): at $C=100, m= 10^6, \alpha=0.8$ the exact value is
$\rho_C = \sum_{k=1}^{100} k^{-0.8} / \sum_{k=1}^{10^6} k^{-0.8}
\approx 0.1087$ (11\% hit rate, modestly above uniform $C/m =
10^{-4}$); the leading asymptotic $(C/m)^{1-\alpha} = 0.158$ is too
crude at this $C$ (next-order corrections from harmonic sum
remainders are not negligible until $C \gg 1$).
Amortized cost reduced by factor $1/(1-\rho_C)$: $\sim 12\times$ for
$\alpha=1.5$, $\sim 1.1\times$ for the $\alpha=0.8$ web-trace regime.

\emph{(c) Random-access OLTP workloads.} Uniform random access has
$\rho_C \approx C/m$. For $C = 1024, m = 10^6$: $\rho_C \approx
10^{-3}$, cache barely helps. RNR's per-query cost remains $O(K \cdot
\mathrm{cost}(M'))$ — still polylog in $N$, just no cache amortization.

\emph{Combination with predictor hidden-state caching.} An advanced
cache stores not just decoded bytes but the predictor's \emph{terminal
hidden state} at sub-block boundary. For sequential access to an
adjacent sub-block, the warm-context from the cached state can
\emph{accelerate} the cold-start cost — though the encoded payload
was prepared with cold context, so the speedup applies only to the
predictor evaluation per position, not the AC slack. Empirical gain:
$\sim 1.5\times$ on sub-block decode under sequential locality.

\emph{Combination with Theorem 7.17 (edit-tolerant).} After an edit
to sub-block $j$, cache entries for $j$ are invalidated; other
sub-blocks' cache entries remain valid. The edit + cache-invalidation
cost is $O(K \cdot \mathrm{cost}(M'))$ per edit plus $O(1)$
cache-eviction.

\textbf{Theorem 7.21 (Predictor-archive size trade-off: optimal LM
parameter count $L^*(N)$ grows sub-linearly with archive size $N$
under empirical scaling laws).}
\emph{Let $X^N$ be compressed by RNR with neural LM predictor $M_L$
of $L$ parameters, where the empirical cross-entropy of $M_L$ on the
source follows a Kaplan-Hoffmann power-law scaling:}
\[
H_{M_L}(X) \;=\; h(X) + C \cdot L^{-\alpha}
\]
\emph{for empirical constants $C > 0$, $\alpha > 0$ (Kaplan et al.\ 2020
\emph{Scaling Laws for Neural Language Models}, arXiv:2001.08361;
Hoffmann et al.\ 2022 \emph{Training Compute-Optimal Large Language
Models}, arXiv:2203.15556). Under the C1 convention amortising the
predictor cost $L$ bits as part of the archive header, the total
storage cost is}
\[
\mathrm{TotalCost}(L, N) \;=\; L + N \cdot H_{M_L}(X) + (\text{RNR sub-block overhead})
\;=\; L + N \cdot h(X) + N \cdot C \cdot L^{-\alpha} + O(\sqrt{N}).
\]
\emph{Minimising over $L$ for fixed $N$:}
\[
\boxed{
\;
L^*(N) \;=\; (N \cdot C \cdot \alpha)^{1/(1+\alpha)},
\quad
\mathrm{TotalCost}(L^*, N) \;=\; N \cdot h(X) + \Theta\bigl(N^{1/(1+\alpha)}\bigr),
\;
}
\]
\emph{i.e., the optimal predictor size grows as $L^* = \Theta(N^{1/(1+\alpha)})$,
strictly sub-linear in $N$, and the per-archive amortised total overhead
is $(1 + 1/\alpha) \cdot L^* = \Theta(N^{1/(1+\alpha)})$ bits — also
sub-linear. The exponent $1/(1+\alpha) < 1$ for any $\alpha > 0$.}

\emph{Concrete scaling values.}
\begin{itemize}
\item \emph{$\alpha = 1$ (heuristic):} $L^* = \Theta(\sqrt{N})$,
overhead $= \Theta(\sqrt{N})$. For $N = 1\,\mathrm{GB}$: $L^* \sim
30\,\mathrm{k}$ parameters, overhead $\sim 60\,\mathrm{kB}$.
\item \emph{$\alpha = 0.5$ (intermediate):} $L^* = \Theta(N^{2/3})$,
overhead $= \Theta(N^{2/3})$. For $N = 1\,\mathrm{GB}$: $L^* \sim
3 \cdot 10^6$ params, overhead $\sim 9 \cdot 10^6$ bits = $1.1\,\mathrm{MB}$.
\item \emph{$\alpha = 0.34$ (Hoffmann/Chinchilla 2022 fit for compute-
optimal LLM training):} $L^* \sim N^{1/1.34} \approx N^{0.75}$, overhead
$= \Theta(L^*) \sim N^{0.75}$. Balanced regime.
\end{itemize}

\emph{\textbf{Proof.}} Calculus optimisation:
$\frac{d}{dL} \mathrm{TotalCost}(L, N) = 1 - N \cdot C \cdot \alpha
\cdot L^{-(\alpha+1)} = 0$
gives $L^{\alpha+1} = N \cdot C \cdot \alpha$, hence $L^* = (N \cdot
C \cdot \alpha)^{1/(\alpha+1)}$.
Second-order: $d^2/dL^2 = N \cdot C \cdot \alpha(\alpha+1) \cdot
L^{-(\alpha+2)} > 0$ at $L^*$, confirming minimum.
Substituting back:
$N \cdot C \cdot (L^*)^{-\alpha} = N \cdot C \cdot (N C \alpha)^{-\alpha/(1+\alpha)}
= L^* / \alpha$ (using $L^{1+\alpha} = NC\alpha$, so
$L^{-\alpha} = L/(NC\alpha)$, hence $N C L^{-\alpha} = L/\alpha$).
Therefore:
$\mathrm{TotalCost}(L^*, N) = L^* + N \cdot h(X) + L^*/\alpha
= (1 + 1/\alpha) L^* + N h(X) = \Theta(L^*) + N h(X)
= \Theta(N^{1/(1+\alpha)}) + N h(X)$.
The overhead grows as $L^* = \Theta(N^{1/(1+\alpha)})$, strictly
sub-linear in $N$.
\hfill$\blacksquare$

\emph{Scope and practical implications.}

\emph{(a) Knowledge distillation from large LMs.} For archives where
shipping a $1\,\mathrm{B}$-parameter pretrained LM (size $\sim
2\,\mathrm{GB}$) is wasteful, knowledge distillation (Hinton, Vinyals
\& Dean 2015 arXiv:1503.02531 ``Distilling the knowledge in a neural
network'') compresses the LM to its optimal $L^*(N)$ size while
retaining most of the predictive power.

\emph{(b) Archive-tier-specific predictors.} Different archive sizes
benefit from different predictors:
\begin{itemize}
\item Small personal archives ($N \sim 1\,\mathrm{MB}$): $L^* \sim
10\,\mathrm{k}$ params; tiny on-device LM (e.g., n-gram-NN hybrid).
\item Medium archives ($N \sim 1\,\mathrm{GB}$): $L^* \sim 10^6$ params;
distilled mobile-class LM.
\item Large archives ($N \sim 1\,\mathrm{TB}$): $L^* \sim 10^9$ params;
full datacenter LM.
\item Shared archives across many users (effective $N \to \infty$):
$L^*$ grows unboundedly, but practical caps at $L \sim 10^{12}$
(frontier model size).
\end{itemize}

\emph{(c) Cost-of-ownership analysis.} Total cost of ownership for
neural lossless compression on a user's archive includes both the
storage cost (paid per byte of archive) and the model-download cost
(paid once per LM version). Theorem 7.21 gives the cost-optimal LM
size as a function of archive size, enabling tier-based pricing /
deployment for compression services.

\emph{Limitations.} The Kaplan-Hoffmann scaling laws are empirical
fits for natural-language data and specific architectures; sources
with non-LM structure (genomics, sensor data, financial time series)
may follow different scaling exponents $\alpha$ or non-power-law
forms. The theorem's structure (optimisation over $L$, sub-linear
$L^*$) is general but the specific exponents require source-class
calibration.

\textbf{Theorem 7.22 (Fail-safe RNR with per-sub-block store-mode
fallback: archive size never exceeds the raw source).}
\emph{Augment the Theorem 7.2 RNR encoder with the following
\emph{per-sub-block store-mode fallback rule}: for each sub-block $j$,
compute both the RNR-encoded payload $L_j^{\mathrm{RNR}}$ (per-position
arithmetic coding with $\eta$-floor predictor $M'$) and the raw store
payload $L_j^{\mathrm{store}} = K \cdot \log_2 |\Sigma|$ bits; emit
the shorter of the two, prepended with a $1$-bit mode flag. Then,
\emph{for any source $X^N$ (including adversarial / random / worst-case
distributions) and any predictor $M'$ (including failed, mismatched,
or adversarial models)}:}
\[
L^{\mathrm{rep}}_{\mathrm{RNR}}(X^N; M')
\;\leq\;
N \cdot \log_2 |\Sigma| + m \cdot (1 + O(\log L(R)))
\;=\;
N \cdot \log_2 |\Sigma| + O(m \log L(R)).
\]
\emph{i.e., the encoded archive is at most $N$ bytes (for byte
alphabet $|\Sigma| = 256$) plus sub-linear overhead. Compression is
never worse than copying the raw source. Best/typical/worst cases:}
\begin{itemize}
\item \emph{Best (perfect prediction):} $L^{\mathrm{rep}} = N \cdot h(X) +
O(\sqrt{N})$ (Theorem 7.2).
\item \emph{Typical (mismatched predictor):} $L^{\mathrm{rep}} = N \cdot
(h(X) + \varepsilon_{\mathrm{KL}})$ (Theorem 7.4).
\item \emph{Worst (adversarial source/predictor):} $L^{\mathrm{rep}} =
N \cdot \log_2 |\Sigma|$ (store-mode triggered on all sub-blocks).
\end{itemize}

\emph{\textbf{Proof.}} Per-sub-block, the encoder runs both RNR and
store modes and selects the shorter encoding. For sub-block $j$:
\begin{itemize}
\item RNR-mode cost: $L_j^{\mathrm{RNR}} = -\sum_{i \in [jK, (j+1)K)}
\log_2 \max(P_{M'}(X_i \mid \text{context}_i), \eta) + O(p)$ AC slack.
\item Store-mode cost: $L_j^{\mathrm{store}} = K \cdot \log_2 |\Sigma|$
bits (raw bytes).
\item Mode flag: $1$ bit per sub-block.
\end{itemize}
Selected cost per sub-block: $\min(L_j^{\mathrm{RNR}}, L_j^{\mathrm{store}}) +
1 \leq L_j^{\mathrm{store}} + 1 = K \cdot \log_2 |\Sigma| + 1$.

Summing over $m$ sub-blocks: $L^{\mathrm{rep}} \leq m \cdot (K \cdot
\log_2 |\Sigma| + 1) + O(m \log L(R)) = N \cdot \log_2 |\Sigma| + m +
O(m \log L(R)) = N \cdot \log_2 |\Sigma| + O(m \log L(R))$.

For $K = \mathrm{polylog}(N)$: $m = N / \mathrm{polylog}(N)$, so the
$O(m \log L(R))$ overhead is $O(N / \mathrm{polylog}(N))$, strictly
sub-linear in $N$. Hence the archive size never exceeds $N \cdot
\log_2 |\Sigma|$ asymptotically.
\hfill$\blacksquare$

\emph{Decoder behaviour.} For each sub-block, the decoder reads the
mode flag (1 bit). If store-mode, read the $K$ raw bytes directly.
If RNR-mode, AC-decode the $K$ positions using the predictor.
Per-query cost: $O(1)$ for store-mode lookup, $O(K \cdot
\mathrm{cost}(M'))$ for RNR-mode (same as Theorem 7.2's per-query
cost).

\emph{Scope and significance.} Theorem 7.22 establishes RNR as a
\emph{fail-safe} compression scheme: the worst-case rate is bounded
by the raw source size, providing a hard guarantee against predictor
failures or adversarial inputs.

\emph{Comparison with existing compressors.}
\begin{itemize}
\item \emph{gzip / zlib (DEFLATE):} similar store-mode fallback for
incompressible blocks (`BTYPE = 00`), worst-case overhead $\sim 5$
bytes per $64\,\mathrm{kB}$ store-block. RNR's overhead is comparable
($O(m \log L(R))$ bits).
\item \emph{zstd:} \emph{frame format} includes uncompressed
``raw'' frames as fallback; worst-case overhead $\sim 5$ bytes per
$128\,\mathrm{kB}$ frame.
\item \emph{Modern neural compressors (NNCP, PixelCNN++):} typically
lack explicit store-mode; rely on predictor quality. Adversarial
source can inflate archive beyond raw size by factor $\log(1/\eta) /
\log_2|\Sigma|$ (e.g., $2\times$ for $\eta = 2^{-16}$, byte alphabet).
\end{itemize}

\emph{Practical implications.}

\emph{(a) Defence against malicious / adversarial inputs.} A storage
service compressing user-uploaded data faces potential adversarial
inputs (deliberately incompressible blobs designed to inflate
archives). Theorem 7.22 caps the inflation at $0$: every uploaded
file occupies at most its raw size on disk.

\emph{(b) Predictor failure modes.} Neural predictors may degrade
under distribution shift, training-data poisoning, or numerical
instability. Theorem 7.22 ensures graceful degradation: predictor
quality directly determines compression ratio, but never causes
\emph{negative} compression (archive larger than raw).

\emph{(c) Mixed-content archives.} For archives containing both
compressible (text, code) and incompressible (already-compressed
binaries, encrypted data) regions, store-mode triggers per-sub-block
in incompressible regions, preserving compression gains in
compressible regions. Theorem 7.16 (adaptive $K$) combined with
Theorem 7.22 (store-mode) gives optimal handling of mixed content.

\emph{Combination with Theorem 7.21 (predictor sizing).} The
fail-safe property allows aggressive predictor distillation:
smaller $L^*(N)$ predictor with higher $\varepsilon_{\mathrm{KL}}$ is
acceptable since store-mode caps the worst case. This loosens the
distillation accuracy requirement and lowers $L^*(N)$ further in
practice.

\textbf{Theorem 7.23 (Integrity-preserving RNR: cryptographic tamper
detection via per-sub-block authenticated encryption at $O(\tau \cdot
m)$ overhead).}
\emph{Augment the Theorem 7.2 RNR encoder with per-sub-block
authenticated encryption: for each sub-block $j$, compute a
message-authentication-code tag $\tau_j = \mathrm{HMAC}_K(\mathrm{payload}_j)$
of length $\tau$ bits using a per-archive symmetric key $K$ (HMAC-
SHA256 with $\tau = 256$; Poly1305 with $\tau = 128$; AES-GCM with
$\tau = 128$, etc.). Store the tag in the archive header alongside the
sub-block offset and hash. The decoder verifies the MAC before
decoding any sub-block; mismatch aborts the query with a tamper-detected
error. Then:}
\begin{itemize}
\item \emph{(Integrity overhead)}: $L^{\mathrm{rep}+\mathrm{MAC}}(X^N) =
L^{\mathrm{rep}}(X^N) + m \cdot \tau$ bits. For $\tau = 128$ and
$K = 1024$: $m \cdot \tau / N = 128 / (1024 \cdot 8) \approx 1.6\%$
overhead.
\item \emph{(Detection probability)}: Under a computationally bounded
adversary, modification of any sub-block payload is detected with
probability $\geq 1 - 2^{-\tau} - \mathrm{negl}(\lambda)$, where
$\lambda$ is the security parameter (HMAC key length). For
$\tau = 128$, the false-negative probability is $\leq 2^{-128}$
(cryptographically secure under the PRF assumption on the underlying
compression function; Bellare--Canetti--Krawczyk 1996 \emph{Advances
in Cryptology --- CRYPTO '96}, Springer LNCS 1109 ``Keying hash
functions for message authentication'', with the tighter PRF-only
characterisation in Bellare 2006 \emph{CRYPTO 2006},
LNCS 4117 ``New proofs for NMAC and HMAC: security without
collision-resistance''; for AEAD security (encrypt-then-MAC compositions),
Bellare \& Namprempre 2000 \emph{Advances in Cryptology --- ASIACRYPT
2000}, Springer LNCS 1976 ``Authenticated encryption: relations among
notions and analysis of the generic composition paradigm'').
\item \emph{(Per-query verification cost)}: $O(K + \tau)$ time
($O(K)$ for HMAC computation over the sub-block, $O(\tau)$ for tag
comparison). Polylog in $N$ for $K = \mathrm{polylog}(N)$.
\item \emph{(Byte-granular RA preserved)}: each query verifies only
the queried sub-block's MAC, no whole-archive scan required.
\end{itemize}

\emph{\textbf{Proof.}} The HMAC construction (Bellare--Canetti--
Krawczyk 1996) provides a pseudorandom function family $\mathrm{HMAC}_K
: \{0,1\}^* \to \{0,1\}^{\tau}$ such that, given oracle access to
$\mathrm{HMAC}_K(\cdot)$ without knowledge of $K$, the probability of
producing a valid $(m', \tau')$ pair with $m'$ not previously queried
and $\tau' = \mathrm{HMAC}_K(m')$ is at most $2^{-\tau} +
\mathrm{negl}(\lambda)$ for computationally bounded adversaries.

Per-sub-block MAC tag generation: compute $\tau_j = \mathrm{HMAC}_K(
\mathrm{payload}_j)$ at encoding time, store alongside sub-block
metadata in header $\Theta$. Header size grows by $m \cdot \tau$ bits.
For $m = N/K$ and $\tau = 128$: header overhead $\sim N \cdot \tau /
K = N / 64$ bits for $K = 1024$, i.e., $1.6\%$ of byte-archive
size.

Per-query: decoder computes $\mathrm{HMAC}_K(\mathrm{payload}_j)$ on
the fetched payload, compares to stored $\tau_j$. Mismatch implies
tampering; constant-time comparison prevents timing-channel attacks.

Sub-block independence (Theorem 7.2) preserved: each MAC is local to
its sub-block, no inter-sub-block cryptographic dependencies. Byte-
granular RA cost identical to plain RNR plus $O(\tau)$ per query.
\hfill$\blacksquare$

\emph{Implementation choices (practical).}
\begin{itemize}
\item \emph{HMAC-SHA256} ($\tau = 256$): conservative, FIPS-approved,
$3\%$ overhead at $K = 1024$.
\item \emph{Poly1305} ($\tau = 128$): faster ($\sim 5\times$ HMAC-
SHA256), $1.6\%$ overhead at $K = 1024$.
\item \emph{AES-GCM} ($\tau = 128$, $\nu = 96$-bit nonce): combined
encryption + authentication, $1.6\%$ MAC overhead, $1.2\%$ nonce
overhead, total $\sim 2.8\%$. Hardware-accelerated on modern CPUs
(AES-NI, ARM Crypto Extensions).
\item \emph{BLAKE3 keyed mode} ($\tau = 256$ or truncated): SIMD-
optimised, $\sim 3\times$ HMAC-SHA256 throughput.
\end{itemize}

\emph{Combined encryption + integrity (AEAD).} For
\emph{confidentiality} in addition to integrity, replace HMAC with an
authenticated encryption mode (AES-GCM, ChaCha20-Poly1305). The
encryption layer wraps each sub-block payload before storage; the
decoder decrypts and verifies before RNR decoding. RA and parallelism
properties (Theorems 7.18, 7.20) preserved since each sub-block has
independent (key, nonce) pair.

\emph{Scope and significance.} Theorem 7.23 brings RNR to parity with
modern secure-archive systems (LUKS encrypted volumes, encrypted ZFS
datasets, age encrypted files) on integrity / tamper-detection
properties, while preserving the RNR-specific byte-granular RA and
sub-block parallelism. Applications:

\emph{(a) Cloud backup with provider-untrusted integrity.} User uploads
RNR-compressed + MAC'd archive to cloud storage; provider cannot
modify any byte without detection on client-side query.

\emph{(b) Forensic archives with chain-of-custody.} Per-sub-block
hash + MAC chain provides evidence of integrity at any byte position
since archive creation.

\emph{(c) Software distribution.} Signed-and-compressed updates
deliver code via RNR + MAC; per-update RA enables differential
patching with cryptographic verification.

\emph{Comparison with related systems.}
\begin{itemize}
\item \emph{LUKS / dm-crypt}: AES-XTS encryption on $4\,\mathrm{kB}$
filesystem blocks. Comparable overhead, but no compression. RNR +
T7.23 adds compression at similar cost.
\item \emph{Tarsnap / restic / borg}: encrypted-compressed backup
systems. RNR achieves better compression via neural LMs while
preserving byte-granular RA (not supported in these systems).
\item \emph{Git pack with GPG sign}: per-pack signature, no per-byte
integrity. RNR + T7.23 provides finer-grained integrity at modest
overhead increase.
\end{itemize}

\emph{Combination with Theorem 9.1 (verified context update).}
Theorem 9.1 already provides per-sub-block hash verification for
context-update integrity. Theorem 7.23 extends this with
\emph{cryptographic} MAC (preimage- and forgery-resistant under
key-secrecy) versus 9.1's plain hash (collision-resistant but no key).
The cryptographic upgrade defends against malicious adversaries with
significant compute budgets.

\textbf{Theorem 7.24 (Decoder latency Hoeffding tail under IRM
workload and static top-popularity cache: minimal honest production
tail bound).}
\emph{Fix an RNR archive built under Theorem 7.15 hypotheses, source
$X^N$ held fixed (no source-randomness in this statement; the source-
side concentration of per-sub-block costs is supplied by Theorem 7.14
applied separately). Let $\mu$ be a probability distribution over
sub-block indices $\{1, \ldots, m\}$ and assume:}
\begin{itemize}
\item \emph{(IRM workload).} \emph{The query addresses
$b_1, \ldots, b_Q \in [N]$ are independent and identically distributed:
the sub-block index $J_i := \lceil b_i / K \rceil$ is drawn iid from
$\mu$ (and the byte offset $b_i \bmod K$ is drawn uniformly).}
\item \emph{(Static top-$C$ cache, ``Perfect-LFU'').} \emph{The
decoder caches exactly the $C$ most $\mu$-popular sub-blocks
$\mathcal{C}^* := \{\,j : \mu(j) \in \text{top-}C\,\}$ throughout the
workload (no eviction). Cache hit rate is
$\rho_C := \sum_{j \in \mathcal{C}^*} \mu(j)$.}
\item \emph{(Bounded per-miss cost).} \emph{Each cache-miss query has
cost $T_{\mathrm{miss}}(j) \in [0, T_{\max}]$ for some uniform upper
bound $T_{\max}$ (e.g., $T_{\max} = K \cdot c_{M'} + \Delta_{\mathrm{ac}}^{\max}$
with $\Delta_{\mathrm{ac}}^{\max}$ the maximum AC slack across sub-blocks
under Theorem 7.2 convention C2).}
\item \emph{(Workload $\perp$ source).} \emph{The workload $W$ is
independent of $X^N$.}
\end{itemize}
\emph{Then the total batched-decode latency
$T_{\mathrm{total}} := \sum_{i=1}^Q T_{\mathrm{query}}(b_i)$ satisfies,
for any $\delta > 0$:}
\[
\boxed{\;
\Pr\!\Bigl(\bigl|\,T_{\mathrm{total}} - Q \cdot \bar{T}\,\bigr|
\;>\;
\sqrt{\tfrac{Q}{2} \log(2/\delta)} \cdot T_{\max}\Bigr)
\;\leq\;
\delta\;
}
\]
\emph{where $\bar{T} := \mathbb{E}_\mu[T_{\mathrm{query}}] =
(1 - \rho_C) \cdot \mathbb{E}_{\mu \mid J \notin \mathcal{C}^*}
[T_{\mathrm{miss}}(J)] \in [0, (1-\rho_C) T_{\max}]$ is the per-query
expected cost under $\mu$.}

\emph{\textbf{Proof.}} Under IRM and static-cache assumptions, the
per-query costs $T_{\mathrm{query}}(b_1), \ldots, T_{\mathrm{query}}(b_Q)$
are i.i.d.\ random variables (a function of the i.i.d.\ $J_i$ alone,
since the cache content is determined a priori by $\mu$, not by query
history). Each is bounded: $T_{\mathrm{query}}(b_i) = 0$ if
$J_i \in \mathcal{C}^*$ (hit) and $T_{\mathrm{query}}(b_i) =
T_{\mathrm{miss}}(J_i) \in [0, T_{\max}]$ otherwise. Hoeffding's
inequality (Hoeffding 1963 \emph{J.\ Amer.\ Stat.\ Assoc.}\
58(301):13, Theorem 2) for iid bounded sums yields
$\Pr(|S_Q - \mathbb{E}[S_Q]| > t) \leq 2 \exp(-2t^2/(Q T_{\max}^2))$,
which gives the stated bound by setting $t = \sqrt{(Q/2)\log(2/\delta)}
\cdot T_{\max}$.
\hfill$\blacksquare$

\emph{Scope and significance.}
Theorem 7.24 is the \emph{minimal honest} tail bound for production
deployment. It deliberately makes restrictive hypotheses to remain
provably correct:
\begin{itemize}
\item \emph{Why IRM, not LRU?} Under LRU the cache state depends on
query history, so $T_{\mathrm{query}}(b_i)$ are not iid; the
counterexample workload $A, B, A, B, \ldots$ against an LRU cache of
size $C = 1$ produces $\Theta(Q)$ deviation in cache misses, defeating
$\sqrt{Q}$-style concentration. Perfect-LFU (static top-$C$) sidesteps
this by removing eviction dynamics: each query independently lands in
or out of the cache based on its sub-block identity alone. A correct
LRU+IRM concentration analysis (with non-trivial cache-state mixing
times) is left to future work (cf.\ Fagin 1977 \emph{J.\ Comput.\ Sys.\ Sci.}
14(2):222--250; Hendricks 1972; King 1971 \emph{Proc.\ IFIP Congress}, 485--490).
\item \emph{Why fix the source?} Per-sub-block cost
$T_{\mathrm{miss}}(j)$ depends on the log-loss of sub-block $j$, which
is random over the source $X^N$. Theorem 7.14 separately bounds the
per-sub-block log-loss concentration; the present theorem assumes a
uniform $T_{\max}$ upper bound that absorbs source-side worst-case
behavior. To compose with source randomness, apply Theorem 7.14 first
to bound $T_{\max}$ with prob.\ $\geq 1 - \delta'$, then Theorem 7.24
with $\delta$, giving total tail probability $\delta + \delta'$ by
union bound.
\item \emph{Why bounded $T_{\max}$, not variance-aware?} Hoeffding is
known to be loose when individual terms are concentrated. Bernstein's
inequality would give a tighter bound using
$\mathrm{Var}[T_{\mathrm{query}}]$, but requires more delicate
analysis. The Hoeffding form is chosen for transparency and
falsifiability.
\end{itemize}

\emph{Numerical worked example.}
$N = 10^9$, $K = \sqrt{N} \approx 31623$, $m = 31623$, $c_{M'} =
1\,\mu\mathrm{s}$, AC slack negligible relative to $K \cdot c_{M'}$,
so $T_{\max} \approx K \cdot c_{M'} = 32\,\mathrm{ms}$. Zipf $\mu$
with $\alpha = 0.8$, $C = 1000$: exact $\rho_C = \sum_{k=1}^{1000}
k^{-0.8} / \sum_{k=1}^{31623} k^{-0.8} \approx 0.438$. For
$Q = 10^6$, $\delta = 10^{-3}$: $\log(2/\delta) \approx 7.6$, so
$\sqrt{(Q/2)\log(2/\delta)} \cdot T_{\max} \approx \sqrt{3.8 \cdot 10^6}
\cdot 32\,\mathrm{ms} \approx 62.4\,\mathrm{s}$.
$\bar{T} \leq (1-\rho_C) T_{\max} \approx 0.562 \cdot 32\,\mathrm{ms}
\approx 18\,\mathrm{ms}$; expected total $\approx 18{,}000\,\mathrm{s}
\approx 5\,\mathrm{h}$; 99.9\%-tail addition $\sim 62\,\mathrm{s}$
($\sim 0.35\%$ tail). The bound is Hoeffding-loose (real
deviations under variance-aware bounds would be smaller); useful as
an upper-bound SLO target.

\emph{Relation to T7.20.}
Theorem 7.20 computes the average per-query latency
$\mathbb{E}[T_{\mathrm{query}}]$ for any workload-cache combination.
Theorem 7.24 adds, under the much-restricted IRM + Perfect-LFU
hypotheses, a worst-case Hoeffding tail. For LRU and non-IRM
workloads, neither theorem applies directly; empirical service-level
measurement is the operational fallback.

\textbf{Theorem 7.25 (Variance-aware Bernstein concentration of
$L^{\mathrm{rep}}$ under exp-mixing).} \emph{Let $X = (X_1, \ldots, X_N)$
be a stationary $\alpha$-mixing process with exponential mixing rate
$\alpha(k) \leq C_\alpha \exp(-c k)$. Assume the $\eta$-floor and
bounded-context hypotheses of Theorem 7.14 (per-position cost
$\ell_i := -\log_2 Q_i(X_i \mid \mathrm{ctx}_i)$ bounded above by
$B := \log_2(1/\eta)$, predictor context window $W_N$), and let
$\sigma_\ell^2 := \mathrm{Var}(\ell_1)$ denote the marginal
per-position-cost variance under the stationary law. Define the
long-run-variance proxy
\[
  v^2 \;:=\; \sigma_\ell^2 \cdot (2 W_N + 1) + 8 B^2 C_\alpha / (e^c - 1).
\]
Then for every $\delta \in (0, 1)$, with probability at least
$1 - \delta$,
\[
  \bigl| L^{\mathrm{rep}}(X) - \mathbb{E}[L^{\mathrm{rep}}(X)] \bigr|
  \;\leq\; C_1 \cdot \sqrt{N \cdot v^2 \cdot \log(2/\delta)}
  \;+\; C_2 \cdot B \cdot (\log N)^2 \cdot \log(2/\delta),
\]
where $C_1, C_2$ are absolute constants depending only on the mixing
rate $c$ (from Merlev\`ede-Peligrad-Rio 2009, IMS Collections 5,
Theorem 2). When $\sigma_\ell \cdot \sqrt{W_N+1} \gg B$, the variance
term $v^2 \approx \sigma_\ell^2 (2 W_N + 1)$ dominates and the bound
scales as $\sigma_\ell \sqrt{N(W_N+1)\log(2/\delta)}$; otherwise the
$B^2$ correction is of the same order.}

\emph{Scope and significance.} Theorem 7.25 replaces Azuma's
worst-case Lipschitz factor of $B = \log_2(1/\eta)$ in Theorem 7.14
by the long-run-variance proxy $v$, which under the empirical
assumption $\sigma_\ell \approx 2$ bits for English text (an
empirical estimate consistent with classical entropy-rate
measurements; Brown et al.\ 1992 and Shannon 1951) gives $v^2 \approx
4 \cdot 4097 + 8 \cdot 32^2 / (e^{1} - 1) \approx 16400 + 4760
\approx 2.1 \cdot 10^4$ bits${}^2$ for the typical regime ($c \approx 1$,
$C_\alpha \approx 1$). For a Chinchilla-class transformer with
$W_N = 2048$ on enwik9 at $\delta = 10^{-6}$, the sub-Gaussian
term yields $C_1 \sqrt{N v^2 \log(2/\delta)} \approx 2.1$--$4.2\,
\mathrm{MB}$ (constant $C_1 \in [1, 2]$ from the specific
Merlev\`ede-Peligrad-Rio formulation; Bernstein-tail term $C_2 B
(\log N)^2 \log(2/\delta) \approx 0.024\,\mathrm{MB}$ is sub-leading
by $\sim 100\times$), versus Theorem 7.14's Azuma worst-case of
$\approx 1.4\,\mathrm{GB}$ --- a $\approx 330\times$--$670\times$
improvement.
The bound remains $\sim 40$--$80\times$ loose against the practical
CLT estimate of $\approx 42\,\mathrm{kB}$; the residual gap stems
from the conservative $(2W_N+1)$ effective-mixing factor in $v^2$
versus the distribution-specific long-run covariance $\sum_k
\mathrm{Cov}(\ell_0, \ell_k)$ for the specific transformer-on-text
process. The bound applies under exp-mixing; the analogous
polynomial-mixing extension requires $\sum_k \alpha(k) < \infty$
(i.e.\ $\alpha(k) \leq C k^{-\beta}$ with $\beta > 1$) and the
Rio (2017) machinery instead of MPR.

\emph{\textbf{Proof.}} The per-position-cost sequence $(\ell_i)$
inherits the $\alpha$-mixing of $X$ at the shifted lag: for $k > W_N$,
$\ell_i$ and $\ell_{i+k}$ depend on disjoint coordinate windows of
$X$, so the standard inheritance argument (Bradley 2007, Vol.~1, §5)
gives $\alpha_{\ell}(k) \leq \alpha_X(k - W_N)$. For $k \leq W_N$,
the trivial bound $\alpha_\ell(k) \leq 1$ holds.

The long-run-variance proxy. For a stationary $\alpha$-mixing
bounded sequence with $|\ell_i| \leq B$, Davydov's covariance
inequality (Davydov 1968; cf.\ Doukhan 1994, Theorem~1.1; Bradley
2007, Vol.~1, Theorem~3.7) gives $|\mathrm{Cov}(\ell_0, \ell_k)|
\leq 4 B^2 \alpha_\ell(k)$, while Cauchy-Schwarz gives the
universal bound $|\mathrm{Cov}(\ell_0, \ell_k)| \leq \sigma_\ell^2$.
Combining, and splitting the sum at $k = W_N$:
\begin{align*}
  v^2 &:= \sigma_\ell^2 + 2 \sum_{k \geq 1} |\mathrm{Cov}(\ell_0, \ell_k)| \\
  &\leq \sigma_\ell^2 + 2 \sum_{k=1}^{W_N} \sigma_\ell^2
  + 2 \sum_{k > W_N} 4 B^2 C_\alpha \exp(-c(k - W_N)) \\
  &\leq \sigma_\ell^2 (2 W_N + 1) + 8 B^2 C_\alpha / (e^c - 1).
\end{align*}

Bernstein bound. Applying Merlev\`ede-Peligrad-Rio (2009, IMS
Collections 5:273--292, Theorem~2; the case of exponentially-mixing
bounded sequences) to the bounded centred sequence
$(\ell_i - \mathbb{E}[\ell_i])$ with per-term bound $M = B$, sample
size $N$, and the long-run variance $v^2$ derived above:
\[
  \Pr\bigl(\,|L^{\mathrm{rep}} - \mathbb{E}[L^{\mathrm{rep}}]| \geq \varepsilon\bigr)
  \;\leq\; \exp\!\Bigl(\,
    -\tfrac{C_3 \varepsilon^2}{N v^2 + B^2 + \varepsilon B (\log N)^2}
  \Bigr) ,
\]
where $C_3 > 0$ depends only on the mixing rate $c$.
Setting the right-hand side equal to $\delta$ and inverting in
$\varepsilon$ (the standard bilinear-Bernstein inversion: in the
sub-Gaussian regime $\varepsilon^2 \ll N v^2$, the first denominator
term dominates and $\varepsilon \leq C_1 \sqrt{N v^2 \log(2/\delta)}$;
in the tail regime, the third term dominates and $\varepsilon \leq
C_2 B (\log N)^2 \log(2/\delta)$; the $B^2$ constant in the
denominator is absorbed into $C_1$ for $N \geq 1$), we obtain
\[
  \varepsilon \;\leq\; C_1 \sqrt{N v^2 \log(2/\delta)}
  + C_2 B (\log N)^2 \log(2/\delta),
\]
which is the claim. \hfill$\blacksquare$

\emph{Remark 7.25a (regime of dominance).}
In the sub-Gaussian regime ($\varepsilon^2 \ll N v^2 / (B (\log N)^2)$),
the first term dominates and the bound scales like
$\sqrt{N v^2 \log(2/\delta)}$. In the Bernstein-tail regime
(very small $\delta$ or very large $W_N$), the second term dominates
and the bound scales like $B (\log N)^2 \log(2/\delta)$.
For the enwik9 example at $\delta = 10^{-6}$, the sub-Gaussian regime
holds and the bound is determined by $v$. Under
$\sigma_\ell \sqrt{W_N+1} \gg B$ (the regime where Theorem 7.25
strictly improves on Theorem 7.14), $v^2 \approx \sigma_\ell^2 (2W_N+1)$
so the bound effectively reads $\sigma_\ell \sqrt{2 N (W_N+1) \log(2/\delta)}$
--- this is the source of Theorem 7.25's improvement over Theorem 7.14.

\emph{Remark 7.25b (relation to Theorems 7.14 and 7.15).}
Theorem 7.14 uses Azuma's inequality on the Doob martingale of
$L^{\mathrm{rep}}$ with the worst-case per-step bound
$(W_N + 1 + \rho/(1-\rho)) B$, scale $A_{\rm Azuma}
\approx (W_N+1) B$ in the leading regime $\rho/(1-\rho) \ll W_N + 1$.
Theorem 7.25 replaces Azuma by Bernstein for the per-position
sequence $(\ell_i)$, yielding the variance-aware scaling
$A_{\rm Bern} \approx v = \sqrt{\sigma_\ell^2 (2W_N+1)
+ O(B^2/c)}$. The improvement factor is
$A_{\rm Azuma}/A_{\rm Bern} = (W_N+1) B / v$; under
$\sigma_\ell \sqrt{W_N+1} \gtrsim B$ this reduces to $\sqrt{W_N+1}
\cdot B/(\sigma_\ell \sqrt{2})$. For English text on Chinchilla
($\sigma_\ell \approx 2$, $W_N = 2048$, $B = 32$): $v^2 \approx
4 \cdot 4097 + 8 \cdot 1024 / (e-1) \approx 2.1 \cdot 10^4$,
$v \approx 145$; vs.\ $(W_N+1) B \approx 6.6 \cdot 10^4$ --- an
improvement factor of $A_{\rm Azuma}/A_{\rm Bern} \approx 450$
in the deviation scale (or equivalently $\approx 330\times$--$670\times$
in the absolute $\delta = 10^{-6}$ deviation bound, depending on
the choice of $C_1 \in [1, 2]$). The end-to-end
Chinchilla-class prediction of Theorem 7.15 thus admits the dual
concentration band $\pm 1.4\,\mathrm{GB}$ (Azuma, T7.14) or
$\pm 2.1$--$4.2\,\mathrm{MB}$ (Bernstein, T7.25; constants
$C_1 \in [1,2]$ from MPR 2009) at $\delta = 10^{-6}$ on enwik9.
Theorem 7.14 remains the operational bound for non-mixing sources
or when $\sigma_\ell$ cannot be estimated;
Theorem 7.25 is the preferred bound for exp-mixing sources where
the long-run-variance proxy $v$ can be measured or upper-bounded.

\textbf{Theorem 7.26 (LRU+IRM miss-count concentration via
Markov-chain Chernoff).} \emph{Let $(X_t)_{t \geq 1}$ be an IRM
request sequence on a finite alphabet $[N_{\rm dist}]$ with
popularity distribution $\pi$ satisfying full support
($\pi_{\rm min} := \min_i \pi_i > 0$). Equip the decoder with an LRU
cache of capacity $C \leq N_{\rm dist}$ (Theorem 7.20 setup). Let
$M_Q := \sum_{t=1}^{Q} \mathbf{1}\{X_t \notin \mathrm{Cache}_t\}$
denote the miss count over the first $Q$ queries, and $\mu_\infty
:= \lim_{t \to \infty} \Pr(X_t \notin \mathrm{Cache}_t)$ the
stationary miss rate (Fagin 1977; Che et al.\ 2002 characteristic-time
approximation, exact in the $C \to \infty$ limit). Then the
augmented cache-state Markov chain $(\mathrm{Cache}_t, X_t)_{t \geq
1}$ is finite, irreducible, and aperiodic, with a unique stationary
distribution and strictly positive spectral gap
$\gamma_{\rm LRU} > 0$ (with $\gamma_{\rm LRU}$ depending only on
$\pi$ and $C$, not on $Q$). By the Lezaud-Gillman Markov-chain
Chernoff bound (Lezaud 1998, Annals of Applied Probability
8(3):849--867; Gillman 1998, SIAM Journal on Computing
27(4):1203--1220; sharpened by Paulin 2015, Electronic Journal of
Probability 20(79)), for every $\varepsilon > 0$:
\[
  \Pr\!\bigl(|M_Q / Q - \mu_\infty| \geq \varepsilon\bigr)
  \;\leq\; C_{\rm Chern} \cdot \exp\!\Bigl( -\tfrac{\varepsilon^2
  \cdot Q \cdot \gamma_{\rm LRU}}{C_{\rm var}}\Bigr) ,
\]
where $C_{\rm Chern}$ and $C_{\rm var}$ are absolute constants
from the specific Markov-Chernoff form (Lezaud: $C_{\rm Chern} = 2$,
$C_{\rm var} = 5$; Paulin sharper: $C_{\rm Chern} = 2$, $C_{\rm var}
= 2$, requiring the pseudo-spectral gap rather than the additive
reversibilization gap).}

\emph{Scope and significance.} Theorem 7.26 closes the LRU+IRM
finite-batch concentration open problem of §13.2, complementing
Theorem 7.24's IRM + Perfect-LFU result (which sidesteps cache
dynamics by using a \emph{static} LFU cache and applying classical
Hoeffding to iid miss indicators). The novelty here is the
positive-spectral-gap argument: under IRM with full-support $\pi$,
the LRU cache-state chain is geometrically ergodic, and miss-count
deviations enjoy a $\sqrt{Q}$-Hoeffding tail (sub-Gaussian rate
$\Theta(Q \gamma_{\rm LRU})$ in the exponent). The bound combines
with Theorem 7.20's expected-cost formula to give a complete
characterization of the decoder-latency distribution: the mean
follows T7.20's amortization, the deviation follows T7.26's
$\sqrt{Q}/\gamma_{\rm LRU}$ scaling. The crude lower bound
$\gamma_{\rm LRU} \geq \pi_{\rm min}$ (proof in Remark 7.26a below)
suffices for the qualitative claim; tighter explicit constants
require LRU-specific spectral analysis (Aldous-Diaconis 1987
coupling-based bounds; Wolfer-Kontoyiannis 2019 non-reversible-chain
spectral gaps), open for future work.

\emph{\textbf{Proof.}}
Under IRM with the standard LRU update rule, the state is most
naturally encoded as the \emph{ordered LRU stack}
$\mathbf{s}_t = (s_t^{(1)}, \ldots, s_t^{(C)})$ where $s_t^{(1)}$
is the most-recently-requested item among those currently in cache
(King 1971; Kesidis et al.\ 2017, arXiv:1704.04849). Together with
the current query $X_t$, the joint state $(\mathbf{s}_t, X_t)$
lives in the finite space $\mathcal{S} := \{(\mathbf{s}, x) :
\mathbf{s} \in [N_{\rm dist}]^{(C)}, x \in [N_{\rm dist}]\}$ of
cardinality $P(N_{\rm dist}, C) \cdot N_{\rm dist} = N_{\rm dist}!
/ (N_{\rm dist} - C)! \cdot N_{\rm dist}$, where $[N_{\rm dist}]^{(C)}$
denotes ordered $C$-tuples of distinct elements. The transition
law is the standard LRU rule: $X_{t+1} \sim \pi$ independently;
if $X_{t+1} \in \mathbf{s}_t$, move $X_{t+1}$ to position 1 and
shift older items down; if $X_{t+1} \notin \mathbf{s}_t$, prepend
$X_{t+1}$ at position 1 and evict the position-$C$ item.

\emph{Irreducibility.} For any target ordered stack $\mathbf{s}'
= (a_1, \ldots, a_C)$ and any subsequent query $x' \in
[N_{\rm dist}]$, the deterministic query sequence $X_{t+1} = a_C,
X_{t+2} = a_{C-1}, \ldots, X_{t+C} = a_1, X_{t+C+1} = x'$ drives
the chain to state $(\mathbf{s}', x')$ regardless of the starting
state (each query moves its target item to position 1; after
$C$ such queries in reverse order, the stack becomes
$(a_1, \ldots, a_C) = \mathbf{s}'$). Each specific query has
probability $\geq \pi_{\rm min} > 0$ under IRM, so any state is
reached from any other in exactly $C + 1$ steps with probability
$\geq \pi_{\rm min}^{C+1} > 0$. Hence the chain is irreducible.

\emph{Aperiodicity.} For any state $(\mathbf{s}, x)$, querying
$X_{t+1} = s^{(1)}$ (the current top of stack) leaves $\mathbf{s}$
unchanged, giving the transition $(\mathbf{s}, x) \to
(\mathbf{s}, s^{(1)})$. A subsequent query $X_{t+2} = s^{(1)}$
yields $(\mathbf{s}, s^{(1)}) \to (\mathbf{s}, s^{(1)})$ --- a
self-loop with probability $\pi_{s^{(1)}} > 0$. Hence the chain
has period 1 (aperiodic).

\emph{Stationary distribution and spectral gap.} As a finite,
irreducible, aperiodic Markov chain, the joint chain admits a
unique stationary distribution $\pi^*$ (Levin-Peres-Wilmer 2017
``Markov Chains and Mixing Times'', Theorem 4.9), with
$\mu_\infty = \mathbb{E}_{\pi^*}[\mathbf{1}\{x \notin \mathbf{s}\}]$.
The chain is not reversible in general (LRU dynamics are
directional). However, the additive reversibilization
$\tilde{P} := (P + P^*)/2$ (where $P^*$ is the time-reversal of
$P$ w.r.t.\ $\pi^*$) is reversible by construction, and Aldous-Fill
(``Reversible Markov Chains and Random Walks on Graphs'', Chapter
9.4.1, ``Reversibilizations'') establishes that the spectral gap
$\gamma_{\rm LRU} := 1 - \lambda_2(\tilde{P})$ is strictly positive
for any finite irreducible aperiodic non-reversible chain
(equivalently: $\tilde{P}$ has eigenvalue 1 with multiplicity 1 by
Perron-Frobenius applied to the connected aperiodic reversibilization).

\emph{Application of Lezaud Chernoff.} Lezaud (1998, Theorem 1.1)
applies to any bounded $L^2$-functional $f$ of a finite-state
Markov chain with positive spectral gap. Taking
$f(S, x) := \mathbf{1}\{x \notin S\} - \mu_\infty$ (the centred
miss indicator, bounded by $\|f\|_\infty \leq 1$) and applying
Theorem 1.1 yields
\[
  \Pr\!\bigl(|S_Q^f| \geq Q \varepsilon \bigr)
  \;\leq\; C_{\rm Chern} \exp\!\Bigl( -\tfrac{\varepsilon^2 Q
  \gamma_{\rm LRU}}{C_{\rm var}}\Bigr) ,
\]
where $S_Q^f := \sum_{t=1}^{Q} f(\mathrm{Cache}_t, X_t) = M_Q - Q
\mu_\infty$, and the constants $C_{\rm Chern}, C_{\rm var}$ are
those of Lezaud's specific form (Lezaud's original $C_{\rm Chern}
= 2$, $C_{\rm var} = 5$ for bounded centred functionals with
$\|f\|_\infty \leq 1$; alternative sharper forms by Paulin 2015
via the pseudo-spectral gap give $C_{\rm var} = 2$). Rearranging
gives the claim. \hfill$\blacksquare$

\emph{Remark 7.26a (positivity of $\gamma_{\rm LRU}$ and explicit
lower bounds).} The positivity $\gamma_{\rm LRU} > 0$ is immediate
from the chain being finite, irreducible, and aperiodic (Aldous-Fill
RWG §9.4.1; cf.\ proof above). Explicit lower bounds on
$\gamma_{\rm LRU}$ are more subtle. A crude Dobrushin-style lower
bound: from the irreducibility argument above, the $(C+1)$-step
chain $P^{C+1}$ has Dobrushin coefficient $\beta_{C+1} := \max_{s,
s'} \|P^{C+1}(s, \cdot) - P^{C+1}(s', \cdot)\|_{\rm TV} \leq 1 -
\pi_{\rm min}^{C+1}$ (because the $C+1$-step deterministic-query
coupling above succeeds with probability $\geq \pi_{\rm min}^{C+1}$
from any pair of starting states). By the Dobrushin-spectral-gap
inequality (Diaconis-Stroock 1991 ``Geometric bounds for
eigenvalues of Markov chains''; LPW Theorem 13.1), $\gamma(P^{C+1})
\geq 1 - \beta_{C+1} \geq \pi_{\rm min}^{C+1}$, and the
multiplicative-step relation $\gamma(P^{C+1}) = 1 - (1 -
\gamma(P))^{C+1} \leq (C+1) \gamma(P)$ gives the one-step bound
$\gamma_{\rm LRU} := \gamma(P) \geq \pi_{\rm min}^{C+1} / (C+1)$.
For Zipf popularity $\pi_i \propto i^{-\alpha}$ on $[N_{\rm dist}]$,
$\pi_{\rm min} = N_{\rm dist}^{-\alpha} / \zeta_{N_{\rm dist}}(\alpha)$,
giving $\gamma_{\rm LRU} \geq (N_{\rm dist}^{-\alpha}/\zeta(\alpha))^{C+1}
/ (C+1)$ --- positive but exponentially small in $C$, hence not
practically useful for moderate cache sizes. Tighter bounds via
LRU-specific spectral structure (the Knudsen-Tsitsiklis 1990
``stack distance'' representation; coupling-from-the-past analyses
for the LRU permutation chain; Aldous-Diaconis 1987 cutoff
phenomena for related permutation chains) are open for future
work. The qualitative $\sqrt{Q}$-Hoeffding scaling of Theorem 7.26
holds for any positive $\gamma_{\rm LRU}$, regardless of the
explicit constant.

\emph{Remark 7.26b (relation to Theorems 7.20 and 7.24).}
Theorem 7.20 gives the \emph{expected} per-query decoder latency
under any cache + workload, including LRU + Zipf. Theorem 7.24
gives the \emph{Hoeffding tail} under the restricted hypothesis of
IRM + \emph{static} Perfect-LFU caching (no eviction dynamics).
Theorem 7.26 extends T7.24's tail-bound machinery to LRU + IRM
by replacing iid Hoeffding with Markov-chain Chernoff (Lezaud);
the cost is the $1/\gamma_{\rm LRU}$ factor in the exponent
(versus T7.24's iid rate). Together, T7.20 + T7.26 give the
expected-and-tail characterization of LRU+IRM decoder latency
that was open after T7.24's static-LFU restriction.

\textbf{Theorem 7.27 (Space-time product converse for sub-block
RNR random access).} \emph{Let $X$ be a length-$N$ source compressed
under any per-position sequential RNR mode (Type-I, per-byte
Type-III-B strategy (iii), or aligned Type-III-A as in
Theorem 10.3) with a $W$-bounded predictor $M$ (Definition 10.2,
per-evaluation cost $\mathrm{cost}(M)$) and uniform sub-block
sync-point spacing of size $K \geq W$ (Definition 10.3, uniform
variant). Assume the encoder achieves the Theorem 7.2 rate
$L^{\mathrm{rep}}_{\mathrm{RNR}} = N h(X) + o(N)$ on a stationary
exp-mixing source with non-degenerate entropy rate
$h(X) \geq h_{\min} > 0$. Let $C(K)$ denote the per-query
computational work (predictor evaluations $+$ bit operations) for a
\emph{first-access} single-byte random-access query --- i.e., a query
to a sub-block whose $M$-outputs have not been previously memoised
(the decoder cache is cold relative to the queried sub-block). Let
$S(K)$ denote the sync-overhead storage (the seek-index size, in bits,
excluding the compressed payload $R$ and excluding any decoder
side-cache of memoised $M$-outputs). Then:}
\[
  C(K) \cdot S(K) \;\geq\; N \cdot \mathrm{cost}(M) \cdot
  \bigl\lceil \log_2(N \cdot h_{\min}) \bigr\rceil
  \;\geq\; (1 - o(1)) \cdot N \cdot \mathrm{cost}(M) \cdot \log_2 N.
\]
\emph{The displayed bound holds uniformly in $K$ for the fixed-width
seek-index realisation of Definition 10.3 (each offset stored in
$\lceil\log_2 L(R)\rceil$ bits); for an index that uses succinct
monotone (Elias--Fano) encoding the same $\Theta(N\,\mathrm{cost}(M)
\log N)$ floor holds at the order-optimal balance $K=\Theta(\sqrt N)$
but degrades by up to a $\log N$ factor at $K=O(1)$ (Lemma 7.27b
Remark). Increasing $K$ trades $C(K)$ up linearly while reducing
$S(K)$; the geometric-mean balance point $K^{**} = \sqrt{N \cdot
\log_2 N / \mathrm{cost}(M)}$ (worst-case parameters, no operational
$\alpha Q$ weighting; to be distinguished from $K^*$ of T7.5 which
adds $\alpha Q$ and the Crutchfield-Feldman constant $\gamma$) is
one specific point on the iso-product curve
$C(K) \cdot S(K) = N \mathrm{cost}(M) \log_2 N$ where the two factors
are equal up to lower-order terms.}

\emph{Scope.} The theorem is a CONVERSE for the sub-block class of
random-access schemes characterised by Definition 10.3 (uniform or
non-uniform sync spacing with seek-index lookup, predictor-state /
entropy-coder-state reset at sync points, sequential decode from
sync to query position). Within this class, the bound is tight: it
matches Theorem 7.5's achievability formula $\gamma N / K + \alpha Q
K = 2\sqrt{\gamma N \alpha Q}$ structurally, with $\gamma$ playing
the role of $\log N \cdot |\mathrm{state}|$ and $\alpha Q$ playing
the role of $\mathrm{cost}(M)$. For non-sub-block schemes (e.g.,
hypothetical schemes that store the entire decoded archive
verbatim, or schemes based on succinct data structures with
non-sequential predictor invocation), the product bound does
\emph{not} apply directly --- such schemes may sidestep the
$K \cdot \mathrm{cost}(M)$ per-query lower bound by avoiding
sequential predictor evaluation entirely. Crucially, this can be done
\emph{without} abandoning the Shannon-near rate: Theorem 7.27$'$ below
exhibits the Ferragina--Venturini entropy-bounded structure, which
attains rate $N h(X) + o(N)$ with $O(1)$ query and no query-time
predictor calls for finite-order-Markov sources, refuting the
general-scheme converse. The sub-block restriction is therefore
essential; the §13.2 open problem is resolved negatively for
Markov-expressible sources (Theorem 7.27$'$) and remains open only for
genuinely neural predictors (Remark 7.27d).

\emph{\textbf{Proof.}} The proof combines two component lower bounds:
(a) per-query work, and (b) sync-overhead storage.

\emph{Lemma 7.27a (per-query predictor-evaluation lower bound,
first-access).} For any sub-block scheme satisfying Definition 10.3
with the encoder achieving $L^{\mathrm{rep}}_{\mathrm{RNR}} \leq N h(X)
+ o(N)$, a single-byte first-access random-access query at byte
position $p \in [0, N)$ --- i.e., a query whose target sub-block has
no $M$-outputs cached in any decoder side-storage --- requires the
decoder to evaluate the predictor $M$ at every byte position in
$[P, p]$, where $P = K \lfloor p / K \rfloor$ is the latest sync
point at or before $p$. In particular, $C(K) \geq K \cdot
\mathrm{cost}(M)$ in the worst case where $p \equiv K - 1
\pmod{K}$ (giving $K$ evaluations at positions $P, P+1, \ldots, P+K-1$).

\emph{Proof of Lemma 7.27a.} By Definition 10.3, the
predictor-state and entropy-coder state at position $P$ are
initialised to the same values as in the encoder's sync-point-aligned
full decode. By Theorem 10.1 (bit-exact deterministic recompute), the
decoder must reproduce the encoder's $M(C_i)$ at every position
$i \in [P, p]$ to keep its entropy-coder state synchronised with the
encoder's. Suppose, for contradiction, that the decoder skips
evaluation of $M$ at some position $i^* \in [P, p]$. Then either
(i)~the decoder uses a memoised value of $M(C_{i^*})$ from a previous
query --- but since the worst-case query is the first access to this
sub-block, no memoisation is available; or (ii)~the decoder predicts
$M(C_{i^*})$ from some other source --- but $M$ is a black-box
$W$-bounded function (Definition 10.2) and there is no
information-theoretic shortcut to its output (even an approximate
prediction would yield an entropy-coder-state divergence, by
Theorem 10.1 bit-exactness, propagating downstream and corrupting
the decoded byte at $p$). Hence the decoder must evaluate $M$
at every position in $[P, p]$, giving $(p - P + 1)$ evaluations.
In the worst case $p \equiv K - 1 \pmod{K}$, $p - P = K - 1$, so
$p - P + 1 = K$ and $C(K) \geq K \cdot \mathrm{cost}(M)$.
\hfill$\square$

\emph{Lemma 7.27b (sync-overhead storage lower bound).} Under
Definition 10.3 with uniform sync spacing and $m = \lceil N/K
\rceil$ sync points, the sync-overhead storage $S(K)$ satisfies the
information-theoretic floor
\[
  S(K) \;\geq\; \log_2 \binom{L(R)+1}{m}
  \;=\; m \log_2\!\frac{L(R)}{m} + \Theta(m)
  \;\geq\; \frac{N}{K}\,\log_2\!\bigl(h_{\min} K\bigr) ,
\]
and, \emph{for the fixed-width-offset realisation of the seek index}
(each entry an explicit $\lceil \log_2 L(R)\rceil$-bit pointer, the
form whose size the Definition 10.3 budget quotes), the stronger
per-entry bound
\[
  S_{\mathrm{fw}}(K) \;=\; m \cdot \bigl\lceil \log_2 L(R) \bigr\rceil
  \;\geq\; \frac{N}{K} \cdot \bigl\lceil \log_2(N \cdot h_{\min})
  \bigr\rceil .
\]

\emph{Proof of Lemma 7.27b.} By Definition 10.3 the seek index has
exactly $m = \lceil N/K \rceil$ entries, encoding the strictly
increasing offsets $0 \le o_1 < \cdots < o_m \le L(R)$ at which the
$m$ sync points begin in $R$ (the sync-point identifier $j$ maps
deterministically to the $j$-th smallest offset). Any such index
determines, and is determined by, an $m$-element subset of
$\{0,1,\ldots,L(R)\}$; by Shannon's noiseless source-coding lower
bound (Cover--Thomas 2006, Theorem 5.3.1) it therefore requires at
least $\log_2 \binom{L(R)+1}{m}$ bits, a floor attained up to
$O(m)$ by the Elias--Fano / succinct-monotone encoding (Munro--Raman
2001 \emph{SIAM J.\ Comput.}\ 31(3):762; Okanohara--Sadakane 2007,
``Practical entropy-compressed rank/select dictionary''). Using
$\binom{L(R)+1}{m}\ge (L(R)/m)^m$, the Theorem 7.2 rate hypothesis
$L(R)\ge N h(X)\ge N h_{\min}$, and $m=\lceil N/K\rceil \le N$, gives
$S(K)\ge m\log_2(L(R)/m)\ge (N/K)\log_2(N h_{\min}/m)
= (N/K)\log_2(h_{\min} K)$, the first display. If instead the index
is realised with \emph{fixed-width} entries --- each offset stored
verbatim in $\lceil\log_2 L(R)\rceil$ bits --- then trivially
$S_{\mathrm{fw}}(K)=m\lceil\log_2 L(R)\rceil$, and under the same
rate hypothesis $\lceil\log_2 L(R)\rceil \ge \lceil\log_2(N
h_{\min})\rceil = (1-o(1))\log_2 N$, the second display.
\hfill$\square$

\emph{Remark (the index floor is not the per-entry floor at small
$K$).} The two displays differ by the factor $\lceil\log_2
L(R)\rceil / \log_2(h_{\min}K)$, which is $\Theta(1)$ exactly when
$m=\lceil N/K\rceil = N^{\Theta(1)}$, e.g.\ at the balance
$K=\Theta(\sqrt N)$ where both read $\tfrac12(1+o(1))\log_2 N$ per
entry. For \emph{small} $K$ (many sync points, $m\to N$) the genuine
succinct floor degrades: at $K=1$, $\log_2(h_{\min}K)=\log_2
h_{\min}=O(1)$, so the storage floor is only $O(N)$ bits, a factor
$\log_2 N$ below the fixed-width count. Hence the $K$-invariant
product below holds rigorously for the fixed-width seek-index
realisation; for an index free to use succinct monotone encoding the
product floor is tight at the order-optimal balance
$K=\Theta(\sqrt N)$ but is \emph{not} uniform in $K$ --- it can drop
to $\Omega(N\,\mathrm{cost}(M))$ (a $\log N$ factor weaker) at
$K=O(1)$. The operationally relevant regime is the balance point,
where the converse is tight in either realisation.

\emph{Combining Lemmas 7.27a and 7.27b.} Use the per-query bound
$C(K)\ge K\,\mathrm{cost}(M)$ (Lemma 7.27a) with the storage bound
(Lemma 7.27b). For the \emph{fixed-width seek-index realisation}
($S_{\mathrm{fw}}(K)=m\lceil\log_2 L(R)\rceil$), multiplying the two
lower bounds gives
\begin{align*}
C(K) \cdot S_{\mathrm{fw}}(K) &\;\geq\; \bigl(K \cdot \mathrm{cost}(M)\bigr)
\cdot \bigl(\lceil N/K \rceil \cdot \lceil \log_2(N h_{\min}) \rceil\bigr) \\
&\geq\; N \cdot \mathrm{cost}(M) \cdot
\lceil \log_2(N h_{\min}) \rceil ,
\end{align*}
where the second inequality uses $K \cdot \lceil N/K \rceil \geq N$
for any integer $K \geq 1$. This product is $K$-invariant: the factor
of $K$ in the per-query bound cancels the $1/K$ in the sync-overhead
bound, leaving a quantity depending only on $N$, $\mathrm{cost}(M)$,
and $\log N$. For $h_{\min} \in (0, 1]$ fixed and $N \to \infty$,
$\lceil \log_2(N h_{\min}) \rceil = (1 + o(1)) \log_2 N$, so
$C(K) \cdot S_{\mathrm{fw}}(K) \geq (1 - o(1)) N \cdot \mathrm{cost}(M)
\cdot \log_2 N$. For an index free to use \emph{succinct monotone}
encoding the storage floor is the first display of Lemma 7.27b,
$S(K)\ge (N/K)\log_2(h_{\min}K)$, and the product becomes
$C(K)\cdot S(K) \ge N\,\mathrm{cost}(M)\,\log_2(h_{\min}K)$, which
equals $(1-o(1))\,\tfrac12\,N\,\mathrm{cost}(M)\log_2 N$ at the
balance $K=\Theta(\sqrt N)$ but degrades to
$\Omega(N\,\mathrm{cost}(M))$ as $K\to O(1)$ (Remark above). In
either realisation the converse is $\Theta(N\,\mathrm{cost}(M)\log N)$
at the order-optimal balance point; the uniform-in-$K$ form is
specific to fixed-width offsets.

\emph{Balance point.} Setting $C(K) = S(K)$ (the AM-GM minimum of
the sum $C + S$ under fixed product) gives $K \cdot \mathrm{cost}(M)
= (N/K) \cdot \log_2 N$, i.e., $K^{**} = \sqrt{N \log_2 N /
\mathrm{cost}(M)}$. For $\mathrm{cost}(M) = \mathrm{polylog}(N)$,
$K^{**} = \Theta(\sqrt{N})$ up to polylog factors, matching the
T7.5 balance point structurally. The T7.5 formula
$K^* = \sqrt{\gamma N / (\alpha Q)}$ adds the
operational $\alpha Q$ weighting and the Crutchfield-Feldman
constant $\gamma$, but the geometric-mean $\sqrt{N}$ scaling is
identical in both formulae. \hfill$\blacksquare$

\emph{Scope, significance, and limitations.} Theorem 7.27 closes the
matching-converse question for sub-block RNR random access in the
strong sense: the product $C(K) \cdot S(K)$ is at least
$N \cdot \mathrm{cost}(M) \cdot \lceil \log_2(N h_{\min}) \rceil =
\Theta(N \log N \cdot \mathrm{cost}(M))$ across the entire
sub-block parameter range, not just T7.5's natural $K = \sqrt{N}$
choice. This rules out the existence of a sub-block-based scheme
that strictly beats T7.5's trade-off curve via clever $K$-tuning. In
particular, no choice of $K$ inside the sub-block class can give
simultaneously $o(\sqrt{N})$ first-access work and $o(\sqrt{N})$ sync
overhead at balance: such a point would violate
$C \cdot S \geq N \log N \mathrm{cost}(M)$.

\emph{First-access vs.\ amortised work.} The bound applies to
\emph{first-access} per-query work: the cost to answer a query
against a sub-block whose $M$-outputs have not been previously
memoised. With a decoder cache (Theorem 7.20) hitting a hot
sub-block, the \emph{amortised} per-query work can be much smaller
than $K \cdot \mathrm{cost}(M)$ --- in the extreme, $O(K)$ cache
reads with zero predictor evaluations once the cache is full.
Theorem 7.27 bounds the worst-case first-touch cost, which is the
relevant quantity for archive cold-start and the upper bound from
Theorem 7.5's $\alpha Q K$ term (where $\alpha Q$ accounts for
expected query volume \emph{averaged} over first-access and
warm-cache outcomes). The cache-aware amortised version is the
subject of Theorems 7.20 / 7.24 / 7.26 (which give expected and
tail bounds for warm-cache cost), and Theorem 7.27 supplies the
matching cold-start converse.

\emph{Non-sub-block converse (resolved negatively; see Theorem 7.27$'$ below).} The theorem applies only to
schemes that fit Definition 10.3's sub-block template
(sync-point-aligned seek index containing bit-offsets only,
predictor-state reset at each sync point, sequential decode from
sync to query position). A hypothetical non-sub-block scheme can
sidestep the per-position $K$-evaluation lower bound by enlarging
the side storage beyond the seek index's bit-offsets and
pre-computing $M$-outputs there: for instance, caching $M(C_i)$ at
every position $i$ takes $N \cdot |\mathrm{output}(M)|$ bits of side
storage and reduces per-query work to $O(1)$ predictor evaluations,
giving a product $C \cdot S = O(N) \cdot O(|\mathrm{output}(M)|) =
O(N \cdot |\mathrm{output}(M)|)$. This trades against T7.27 by
exchanging the seek-index $\log N$ factor for an $|\mathrm{output}(M)|$
factor (typically $\geq 8$ bits per byte), and exits Definition 10.3's
architecture (the side storage is no longer just a seek index).
Other non-sub-block scheme candidates include succinct data
structures with random-access predictor invocation, or schemes that
store predictor-state snapshots at non-uniform information-theoretic
positions. The general converse across all random-access schemes is
addressed below by Theorem 7.27$'$, Remark 7.27d, and Corollary 7.27c
(and the §13.2 entry is updated accordingly). The short version: the
once-conjectured universal product $(\text{query work}) \cdot
(\text{total storage}) \geq \Omega(N \cdot \mathrm{cost}(M))$ for every
Shannon-near random-access scheme is \emph{false} for
finite-order-Markov-expressible sources --- the Ferragina--Venturini
(SODA 2007) entropy-bounded storage structure stores the realised string
in $N h(X) + o(N)$ bits with $O(1)$-time random access and \emph{no}
query-time predictor evaluations, giving $C \cdot S = O(N) = o(N\,
\mathrm{cost}(M)\log N)$ (Theorem 7.27$'$). The sub-block restriction of
Theorem 7.27 is therefore essential: the product bound is a property of
the predictor-sequential decode architecture, not of the random-access
task. What \emph{does} survive is the $\mathrm{cost}(M)$-bearing product
$C_{\mathrm{work}} \cdot S \geq (1-o(1)) N \cdot \mathrm{cost}(M) \cdot
\log N$ on the \emph{predictor-faithful} class (Corollary 7.27c), which
strictly contains both the sub-block schemes of Definition 10.3 and the
state-snapshot schemes of Theorem 10.4, and is the maximal class in
which dense predictor recomputation --- hence the $\mathrm{cost}(M)$
factor --- is forced. The $\log N$ factor is the seek-index
identifier-length floor and survives into Corollary 7.27c. For
genuinely neural (non-finite-order-Markov) predictors the
Ferragina--Venturini escape closes \emph{only} when the source's
effective Markov horizon grows super-logarithmically in $N$ (Corollary
7.27e, which corrects Remark 7.27d(ii): for every \emph{fixed} stationary
source, even infinite-order, FV still attains Shannon-near rate at build
order $o(\log_2 N/\log_2^2\sigma)$); in that growing-horizon regime the
existence of any model-free $O(1)$-access Shannon-near structure reverts
to open.

\emph{Comparison with classical cell-probe and succinct-data-structure
lower bounds.} The $\sqrt{N}$ time-space trade-off for sub-block
RNR is structurally analogous to the classical compressed
suffix-array sampling-rate trade-off (Sadakane 2003;
Grossi-Vitter 2005), where sampling rate $s$ gives storage $O(N/s)$
and query $O(s \log N)$, yielding $s = \Theta(\sqrt{N})$ at balance.
The general cell-probe lower bound for static membership in the
RAM model is $\Omega(\log N / \log\log N)$ per query under
linear-space (Andersson-Miltersen 1999, Pătraşcu-Demaine 2006
\emph{SIAM J.\ Comput.}\ 35(4):932); Larsen 2012
\emph{52nd FOCS}, pp.\ 293--301 (`The cell probe complexity of
dynamic range counting') established $\Omega(\log^{3/2} N)$ for
dynamic predecessor under polynomial space. These bounds are
weaker than T7.27 in absolute scale (they trade off probe count
versus storage in bits, whereas T7.27 trades $K \cdot \mathrm{cost}(M)$
versus seek-index storage), but they are of the same product-tradeoff
type and lend literature precedent to the sub-block-restricted
converse.

\textbf{Theorem 7.27$'$ (General-scheme converse is \emph{false}: a
rate-respecting scheme beats the product bound).} \emph{The
general-scheme analogue of Theorem 7.27 --- a product lower bound
$C \cdot S \geq \Omega(N \cdot \mathrm{cost}(M))$ (with or without the
$\log N$ factor) holding for every random-access scheme attaining the
Theorem 7.2 Shannon-near rate --- is false. Concretely, for the class of
stationary order-$k$ Markov sources (with $0 < h(X) < \log_2 \sigma$ and
the predictor $M$ outputting the true order-$W$ conditionals, $W \geq
k$, as in Lemma 5.1c), there is a random-access scheme with}
\[
  S_{\mathrm{tot}} \;=\; N \cdot h(X) + o(N) \quad\text{(total on-disk
  storage; Shannon-near rate, hypothesis of Theorem 7.2 met)},
\]
\[
  C \;=\; O(1) \text{ cell probes per single-byte query, with }
  \textbf{zero}\text{ predictor evaluations,}
\]
\emph{so the time--space product is $C \cdot S = O(N) = o\bigl(N \cdot
\mathrm{cost}(M) \cdot \log N\bigr)$ whenever $\mathrm{cost}(M) \cdot
\log N = \omega(1)$. The scheme is the Ferragina--Venturini (SODA 2007;
\emph{Theoret.\ Comput.\ Sci.}\ 372(1) (2007) 115--121) entropy-bounded
storage structure applied to the realised string $X^N$.}

\emph{Consequently the §13.2 ``general-scheme space-time converse'' is
\textbf{resolved negatively}: Theorem 7.27's sub-block restriction
(Definition 10.3) is not a technical artefact but \emph{essential} ---
the product bound is a property of the sequential-decode architecture,
not of the random-access task itself. The largest class for which a
$\mathrm{cost}(M)$-bearing product bound survives is identified
positively in Corollary 7.27c (the predictor-faithful class), and the
gap between that class and the unrestricted class is exactly the
Ferragina--Venturini escape.}

\emph{\textbf{Proof.}} Let $X$ be stationary order-$k$ Markov on a
$\sigma$-symbol alphabet with entropy rate $h(X) = H(X_{k+1} \mid
X_1, \ldots, X_k) \in (0, \log_2 \sigma)$. Fix a realisation
$x = X^N \in \Sigma^N$.

\emph{Step 1 (the realised string is constant-factor compressible and an
order-$k$ Markov predictor witnesses it).} The $k$-th order empirical
entropy of $x$ is $H_k(x) = \frac{1}{N} \sum_{q \in \Sigma^k}
|x_q|\, H_0(x_q)$, where $x_q$ is the subsequence of symbols immediately
following occurrences of context $q$ in $x$ and $H_0$ is the zeroth-order
empirical entropy (Manzini 2001, \emph{J.\ ACM} 48(3):407--430, the
$H_k$ definition). For a stationary ergodic order-$k$ Markov source, the
entropy rate is $h(X) = H(X_{k+1} \mid X_1^k)$ (Cover--Thomas 2006,
Chapter 4, ``Entropy Rates of a Stochastic Process''), and $N \cdot
H_k(x)$ equals the empirical negative log-likelihood
$-\log_2 \hat{p}(x)$ under the maximum-likelihood order-$k$ model; by
the law of large numbers applied to the per-context empirical
distributions $\hat{p}(\cdot \mid q) \to p(\cdot \mid q)$ (each context
$q$ recurs $\Theta(N)$ times a.s.\ by ergodicity), $H_k(x) \to
H(X_{k+1}\mid X_1^k) = h(X)$ almost surely; equivalently this is the
Shannon--McMillan--Breiman convergence $-\frac{1}{N}\log_2 p(x^N) \to
h(X)$ (Cover--Thomas 2006, \S16.8, the general AEP) specialised to the
order-$k$ Markov likelihood. Hence $N \cdot H_k(x) = N \cdot h(X) + o(N)$
with probability $1 - o(1)$. The same conditional
distributions $P(x_i \mid x_{i-k:i-1})$ are exactly what the order-$W$
($W \geq k$) predictor $M$ of Lemma 5.1c outputs; thus the RNR sub-block
encoder of Theorem 7.27 achieves rate $N h(X) + o(N)$ on this source
(its Theorem 7.2 hypothesis is met), and so any scheme we compare it
against is judged against the \emph{same} rate target.

\emph{Step 2 (Ferragina--Venturini structure: entropy storage with
$O(1)$ access, no predictor at query time).} Ferragina and Venturini
construct, for any string $x \in \Sigma^N$, a static data structure
of total size
\[
  N \cdot H_k(x) \;+\; O\!\left(\frac{N}{\log_\sigma N}
  \bigl(k \log \sigma + \log\log N\bigr)\right)
  \;=\; N \cdot H_k(x) + o(N \log \sigma)
\]
bits (valid for $k = o(\log_\sigma N)$, which holds for the fixed
order $k$ here), that retrieves any substring of $x$ of length $\ell$ in
$O(1 + \ell \log \sigma / N) = O(1)$ time for $\ell = O(N/\log\sigma)$,
in particular any \emph{single} symbol $x_p$ in $O(1)$ time. The
retrieval is a table lookup plus a constant number of cell probes; it
performs \textbf{no} predictor evaluation --- the predictor $M$ is used
\emph{only at construction time} to certify the entropy bound (Step 1)
and is not invoked when answering a query. Combining with Step 1,
$S_{\mathrm{tot}} = N H_k(x) + o(N\log\sigma) = N h(X) + o(N)$ almost
surely (for fixed $\sigma$ the $o(N\log\sigma)$ lower-order term is
$o(N)$, dominated by the $\Theta(N)$ leading term), and the worst-case
single-byte query cost is $C = O(1)$ cell probes.

\emph{Step 3 (the product is beaten).} The total storage in bits is
$S = S_{\mathrm{tot}} = \Theta(N)$ (since $h(X) > 0$); the per-query
work is $C = O(1)$. Hence $C \cdot S = O(N)$. Theorem 7.27's sub-block
product, on the \emph{same} source and rate target, is $C(K) \cdot S(K)
\geq (1-o(1)) N\, \mathrm{cost}(M) \log_2 N$. The ratio is
\[
  \frac{C \cdot S}{C(K) \cdot S(K)}
  \;=\; O\!\left(\frac{N}{N\, \mathrm{cost}(M) \log_2 N}\right)
  \;=\; O\!\left(\frac{1}{\mathrm{cost}(M) \log_2 N}\right)
  \;\to\; 0
\]
whenever $\mathrm{cost}(M)\log_2 N \to \infty$ (true for any
non-degenerate predictor as $N \to \infty$). Thus the Ferragina
--Venturini scheme strictly beats the product bound while meeting the
Shannon-near rate hypothesis, refuting the universal converse.
\hfill$\blacksquare$

\textbf{Remark 7.27d (why the escape does not contradict Theorem 7.27,
and what it costs).} The Ferragina--Venturini scheme does \emph{not}
contradict Theorem 7.27, because it is not a Definition 10.3 sub-block
scheme: it neither resets predictor state at sync points nor decodes
sequentially from a sync point. It replaces the predictor-sequential
decode by a \emph{precomputed, model-free} succinct index whose internal
tables encode the order-$k$ statistics once at build time. Three honest
observations sharpen the picture. \emph{(i) The escape is genuine, not
an accounting trick.} The structure's total on-disk size $N H_k(x) +
o(N)$ is counted in full (tables included); there is no hidden side
storage. It is the constant-factor compressibility of the RNR regime
($h(X) = \Theta(1) > 0$ a fixed positive constant, hence $S_{\mathrm{tot}} = \Theta(N)$) that
makes $O(1)$ entropy-compressed access possible --- in sharp contrast to
the \emph{exponentially} compressible regime ($L = 2^{\sqrt n}$, grammar
size $n$) where Chen--Verbin--Yu (\emph{Data structure lower bounds on
random access to grammar-compressed strings}, ICALP 2012; journal version
\emph{Algorithmica}) prove an $\Omega(n^{1/2 - \varepsilon})$ query lower
bound. The RNR product bound is thus architecture-specific, not
information-theoretic. \emph{(ii) The escape pays at build time and in
generality.} Ferragina--Venturini fixes the model to fixed-order
empirical statistics and bakes them into the index; it forfeits the
RNR framework's central feature of plugging in an \emph{arbitrary}
$W$-bounded neural predictor $M$ whose conditionals are not order-$k$
Markov for any fixed $k$. For a genuine neural $M$ (a transformer
LM) the order-$k$ surrogate $H_k$ that Ferragina--Venturini compresses
to satisfies $H_k(x) \to h(X)$ only as $k \to \infty$, while the
lower-order term $O(N k \log\sigma / \log_\sigma N)$ blows up as
$k$ grows; so for predictors that genuinely beat all finite-order Markov
models, the Ferragina--Venturini rate is \emph{not} Shannon-near at any
fixed $k$ and the escape closes. The refutation therefore holds for
finite-order-Markov-expressible sources but \emph{not} for the
neural-predictor sources that primarily motivate RNR; on the latter, no
$O(1)$-access escape is known and the question reverts to open (see
§13.2). \emph{(iii) The $\mathrm{cost}(M)$ factor is non-universal
regardless.} Even setting Ferragina--Venturini aside, a scheme may
\emph{partially} pre-cache: store $M(C_i)$ at the $o(N)$ positions of a
sparse set $\Pi$ with $|\Pi| \log_2\sigma = o(N)$ (rate-respecting).
Queries landing in $\Pi$ skip the predictor, violating Theorem 7.27's
per-query $K\,\mathrm{cost}(M)$ bound (Lemma 7.27a, which required
\emph{every} window position to be recomputed). Hence no bound with $C$
measured in predictor \emph{evaluations} can carry $\mathrm{cost}(M)$
universally; the factor survives only when dense recomputation is
forced, which is the predictor-faithful condition of Definition 7.27b.

\textbf{Definition 7.27b (Predictor-faithful decoder).} Fix a decode
window granularity $K$ (the analogue of sub-block spacing; for a general
scheme, $K$ is the number of consecutive source positions whose joint
decoding a single query triggers). A decoder $\mathcal{D}$ for a scheme
with predictor $M$ is \emph{$\epsilon$-predictor-faithful} if, on every
first-access single-byte query, at most an $\epsilon$ fraction of the
$K$ positions in the triggered window have their predictor output
$M(C_i)$ or recurrent state available verbatim in the archive (so that
the decoder may read it without an $M$-evaluation); equivalently, the
decoder must \emph{recompute} $M$ at $\geq (1-\epsilon) K$ of the
window positions. The archive may still store the seek index, the
compressed payload, and $o(N)$-density auxiliary data; what
$\epsilon$-faithfulness forbids is a \emph{dense} stored shadow of the
predictor's outputs/states across the window. Equivalently:
$\epsilon$-faithful decoders are those that retain the RNR
predictor-sequential decode architecture (up to an $\epsilon$-sparse
cache), as opposed to model-free succinct indices like Ferragina
--Venturini ($\epsilon = 1$: the predictor is never invoked at query
time, so the faithful bound below is vacuous --- correctly, since that
is exactly the escape scheme of Theorem 7.27$'$).

\textbf{Corollary 7.27c (Predictor-faithful product bound: $\mathrm{cost}
(M)$ survives, strictly beyond sub-blocks).} \emph{Let $(A,
\mathcal{D})$ be an $\epsilon$-predictor-faithful scheme (Definition
7.27b) with a $W$-bounded or non-$W$-bounded predictor $M$ of
per-evaluation cost $\mathrm{cost}(M)$, achieving the Theorem 7.2
Shannon-near rate, with seek-index-style storage of $m = \lceil N/K
\rceil$ bit-offset entries. Let $C_{\mathrm{work}}(K)$ be the worst-case
first-access per-query \emph{computational work} (predictor evaluations
$\times \mathrm{cost}(M)$ plus bit operations) and let $S(K)$ be the
seek-index storage in bits. Then}
\[
  C_{\mathrm{work}}(K) \;\geq\; (1 - \epsilon)\, K \cdot \mathrm{cost}(M),
  \qquad
  S(K) \;\geq\; \frac{N}{K} \cdot \bigl\lceil \log_2(N h_{\min})
  \bigr\rceil,
\]
\emph{and therefore}
\[
  C_{\mathrm{work}}(K) \cdot S(K) \;\geq\; (1 - \epsilon)\, N \cdot
  \mathrm{cost}(M) \cdot \bigl\lceil \log_2(N h_{\min}) \bigr\rceil
  \;=\; (1 - \epsilon)(1 - o(1))\, N \cdot \mathrm{cost}(M) \cdot \log_2 N.
\]
\emph{The predictor-faithful class strictly contains: (a) Definition
10.3 sub-block schemes ($\epsilon = 0$: nothing is stored beyond the
bit-offset index, so every window position is recomputed --- this
recovers Theorem 7.27 exactly); and (b) Theorem 10.4 state-snapshot
schemes, which store $\Theta(N/K) = o(N)$ recurrent-state snapshots,
hence have predictor state available at only the $m = N/K$ sync points
and must recompute $M$ at the other $K-1$ positions per window, giving
$\epsilon = O(1/K) \to 0$. State-snapshot schemes are \emph{not}
sub-block schemes (they store source-dependent recurrent state, exiting
Definition 10.3, cf.\ the open-extension note after Lemma 5.1e), so
Corollary 7.27c is a strict generalisation of Theorem 7.27 within the
predictor-faithful regime. The class \emph{excludes} the Ferragina
--Venturini escape ($\epsilon = 1$), consistent with Theorem 7.27$'$.}

\emph{\textbf{Proof.}} The storage bound on $S(K)$ is Lemma 7.27b
verbatim (it depends only on the seek index having $m = \lceil N/K
\rceil$ bit-offset entries into a repair stream of length $L(R) \geq N
h_{\min}$, which holds for any scheme with a sync-point bit-offset
index, sub-block or not). For the work bound, by $\epsilon$
-faithfulness the decoder reads stored $M$-data at $\leq \epsilon K$ of
the $K$ window positions; at the remaining $\geq (1-\epsilon)K$
positions it has no stored predictor output or state and, by the
black-box no-shortcut property of $M$ (Definition 10.2, exactly as in
the proof of Lemma 7.27a: any skipped or approximated $M$-evaluation
diverges the entropy-coder state by Theorem 10.1 bit-exactness and
corrupts the decoded byte), it must perform a fresh $M$-evaluation.
Hence $C_{\mathrm{work}}(K) \geq (1-\epsilon) K \cdot \mathrm{cost}(M)$.
Multiplying the two bounds and using $K \cdot \lceil N/K \rceil \geq N$
gives the product. For (a), $\epsilon = 0$ reduces to Theorem 7.27. For
(b), a state-snapshot scheme stores the recurrent state at the $m =
\lceil N/K \rceil$ sync points and only there; on a first-access query
to a window of $K$ positions, the snapshot provides the boundary state
at the single sync point $P = K\lfloor p/K\rfloor$, after which the
decoder must run $M$ forward at every position in $[P, p]$ (Theorem
10.4 proof: ``decode forward as in Theorem 10.3, cost $O((K+k)\,
\mathrm{cost}(M))$''), i.e.\ at $\geq K - 1$ of the $K$ positions; thus
the fraction with stored predictor \emph{output} available is $0$ and
the fraction with stored \emph{state} available is $1/K$, giving
$\epsilon \leq 1/K \to 0$ and $C_{\mathrm{work}}(K) \geq (1 - 1/K) K\,
\mathrm{cost}(M) = (K-1)\mathrm{cost}(M)$. \hfill$\blacksquare$

\emph{Consistency and significance.} Corollary 7.27c reproduces Theorem
7.27's product at $\epsilon = 0$ and extends it to the non-sub-block
state-snapshot architecture (Theorem 10.4) at $\epsilon = O(1/K)$, so
the $\sqrt N$ balance point $K^{**} = \sqrt{N \log_2 N / \mathrm{cost}
(M)}$ and the T7.5 operational balance $K^{*} = \sqrt{\gamma N /
(\alpha Q)}$ both lie on the predictor-faithful iso-product curve (under
the worst-case identification $\gamma \leftrightarrow \log_2 N$,
$\alpha Q \leftrightarrow \mathrm{cost}(M)$); the verify script
\texttt{thm\_7\_27\_prime\_general\_scheme\_converse.py} checks this
identity at machine precision and exhibits the Ferragina--Venturini
escape numerically (its $C \cdot S = \Theta(N)$ sits a factor
$\mathrm{cost}(M)\log N$ below the faithful product). The honest
resolution of the §13.2 open problem is therefore three-part:
\emph{(1)} the unrestricted general-scheme product bound is \emph{false}
for finite-order-Markov-expressible sources (Theorem 7.27$'$: the
Ferragina--Venturini structure achieves Shannon-near rate with $O(1)$
query and no query-time predictor calls); \emph{(2)} the
$\mathrm{cost}(M)$-bearing product $C_{\mathrm{work}} \cdot S =
\Omega(N\,\mathrm{cost}(M)\log N)$ holds exactly on the
predictor-faithful class (Corollary 7.27c), which strictly contains the
sub-block and state-snapshot architectures and is the maximal class for
which dense predictor recomputation is forced; \emph{(3)} for genuinely
neural predictors $M$ whose conditionals are not finite-order Markov,
the Ferragina--Venturini escape closes (its $H_k$ surrogate is not
Shannon-near; Remark 7.27d(ii)), so whether \emph{some other} model-free
succinct index can attain Shannon-near rate with $o(\sqrt{N})$ query for
such sources remains genuinely open --- now sharpened from ``is there a
general converse?'' to ``is there a model-free Shannon-near
random-access structure for non-Markov neural-compressible sources?''.
Corollary 7.27e below settles the Ferragina--Venturini half of that
question quantitatively (whether \emph{FV itself} closes) by computing
the redundancy-vs-order crossover; it corrects the over-broad reading of
Remark 7.27d(ii) and isolates the genuinely residual case.

\textbf{Corollary 7.27e (Ferragina--Venturini escape-closure crossover;
correction and sharpening of Remark 7.27d(ii)).} \emph{Fix an alphabet
size $\sigma \geq 2$ and write $L_\sigma(N) := \log_\sigma N = \log_2 N /
\log_2 \sigma$. For a stationary source let
$\Delta_k := H(X_{k+1} \mid X_1^k) - h(X) \geq 0$ be the order-$k$
conditional-entropy excess (decreasing, $\Delta_k \downarrow 0$ by the
entropy-rate theorem, Cover--Thomas 2006 Theorem 4.2.1), so that the
Crutchfield--Feldman excess entropy is $\mathbf{E} = \sum_{k \geq 0}
\Delta_k \in [0, \infty]$. Applying the Ferragina--Venturini (TCS 372(1)
(2007) 115--121) structure at empirical order $k$ to a realisation
$x = X^N$ yields, almost surely as $N \to \infty$ with $k = o(L_\sigma
(N))$, a total per-symbol rate}
\[
  r_k(N) \;=\; h(X) \;+\; \underbrace{\Delta_k + o(1)}_{\text{model gain}}
  \;+\; \underbrace{\Theta\!\Bigl(\tfrac{(k+1)\log_2\sigma + \log_2\log_2 N}
  {L_\sigma(N)}\Bigr)}_{\rho_k(N),\ \text{FV redundancy}} ,
\]
\emph{where the $o(1)$ inside the model gain is the per-context
estimation error (each order-$k$ context recurs $\Theta(N/\sigma^k)$
times a.s.\ for $k = o(L_\sigma(N))$, so $H_k(x) = h(X) + \Delta_k +
o(1)$). The FV redundancy obeys}
\[
  \rho_k(N) = o(1) \iff k = o\!\bigl(L_\sigma(N)/\log_2\sigma\bigr)
  = o\!\bigl(\log_2 N / \log_2^2 \sigma\bigr),
\]
\emph{so the build order at which FV's structural overhead reaches a
constant bits/symbol is $k_{\mathrm{cross}}(N) = \Theta(\log_2 N /
\log_2^2\sigma)$. Consequently:}
\begin{itemize}
\item[\emph{(a)}] \emph{(Escape stays OPEN for every \textbf{fixed}
stationary source --- correction of Remark 7.27d(ii).) For any fixed
source ($\Delta_k$ a fixed $N$-independent sequence with $\Delta_k
\downarrow 0$), the optimal FV order $k^\star(N) := \arg\min_k
[\Delta_k + \rho_k(N)]$ gives $\min_k [\Delta_k + \rho_k(N)] \to 0$, so
FV attains the Theorem 7.2 Shannon-near rate with $O(1)$ query and the
escape of Theorem 7.27$'$ \textbf{does not close}. This holds even for
\textbf{infinite Markov order} ($\Delta_k > 0$ for all $k$) and even for
\textbf{divergent excess entropy} ($\mathbf{E} = \infty$): e.g.\ for
$\Delta_k = \Theta(1/k)$ (harmonic, $\mathbf{E} = \infty$) the joint
minimum is $\min_k [a/k + \Theta(k/L_\sigma(N))] = \Theta(\sqrt{a /
L_\sigma(N)}) \to 0$, attained at $k^\star = \Theta(\sqrt{a\,
L_\sigma(N)}) = o(L_\sigma(N))$. Thus ``genuinely neural $=$
non-finite-order Markov'' is \textbf{not sufficient} to close the FV
escape; Remark 7.27d(ii)'s assertion that the escape closes for all
neural predictors is too strong.}
\item[\emph{(b)}] \emph{(Escape CLOSES iff the modelling horizon grows
super-logarithmically.) Let $k_\epsilon(N) := \min\{k : \Delta_k^{(N)}
\leq \epsilon\}$ be the order needed to reach excess $\epsilon$, allowing
a triangular-array source whose excess sequence $\Delta_k^{(N)}$ may
depend on $N$. The FV escape closes --- no fixed-order FV index attains
$r_k(N) = h(X) + o(1)$ --- if and only if}
\[
  \inf_{k = o(L_\sigma(N))} \bigl[\Delta_k^{(N)} + \rho_k(N)\bigr]
  \;=\; \Omega(1),
\]
\emph{equivalently iff for every $\epsilon > 0$ the order to reach excess
$\epsilon$ grows as $k_\epsilon(N) = \omega(\log_2 N / \log_2^2\sigma)$
(the source's effective Markov horizon outruns the FV crossover). A
canonical witness is the triangular array $\Delta_k^{(N)} =
c_0 L_\sigma(N)/k$ for $k \leq L_\sigma(N)$ (modelling horizon $\propto
L_\sigma(N)$): the unconstrained AM--GM optimum of $c_0 L_\sigma(N)/k +
\Theta(k/L_\sigma(N))$ sits at $k = \Theta(L_\sigma(N))$, \emph{outside}
FV's validity range $k = o(L_\sigma(N))$, so on the admissible range the
objective is still on its decreasing branch and $\inf_{k = o(L_\sigma
(N))}[\Delta_k^{(N)} + \rho_k(N)] = \Theta(c_0) = \Omega(1)$, bounded
away from $0$ --- the escape closes.}
\item[\emph{(c)}] \emph{(Byte-granular practical closure.) For
$\sigma = 256$, $k_{\mathrm{cross}}(N) = \log_2 N / 64$, and solving
$\rho_1(N) = (2\log_2\sigma + \log_2\log_2 N)/L_\sigma(N) = 1$ gives
$\rho_1(N) > 1$ bit/symbol for all $N \lesssim 2^{188} \approx 10^{57}$.
Hence on byte data at any physically realisable $N$ the FV redundancy
already exceeds a bit/symbol at $k = 1$, and the escape is
\textbf{practically closed} unless the source is Shannon-near already at
$k = 0$ (i.i.d.). The
$\Theta(1)$-bits/symbol FV overhead at small build orders is the
operative obstruction at deployment scale; the asymptotic $k_\epsilon$
vs.\ $k_{\mathrm{cross}}$ race of (a)--(b) is the obstruction in the
$N \to \infty$ limit.}
\end{itemize}
\emph{The verify script
\texttt{scripts/verify/thm\_7\_27e\_neural\_fv\_escape\_crossover.py}
locates the crossover (A), exhibits open joint-excess decay for fixed
finite-context and exp-mixing sources (B), closed $\Omega(1)$ joint
excess for growing-horizon sources (C), the fixed-vs-growing boundary
(D), and the byte-alphabet practical-closure regime (E).}

\emph{\textbf{Proof.}} The rate decomposition is the FV space bound
(Ferragina--Venturini 2007; equivalently Sadakane--Grossi, as recorded
at the FV structure citation above)
$|\mathrm{FV}_k| = N H_k(x) + O\bigl(\tfrac{N}{L_\sigma(N)}((k+1)\log_2
\sigma + \log_2\log_2 N)\bigr)$ divided by $N$, with $H_k(x) = h(X) +
\Delta_k + o(1)$ a.s.\ (Manzini 2001 $H_k$ definition together with the
order-$k$ Markov AEP of Theorem 7.27$'$ Step~1, the $o(1)$ being the
empirical-to-true conditional gap, which vanishes when each context
recurs $\Theta(N/\sigma^k) \to \infty$ times, i.e.\ for $k = o(L_\sigma
(N))$). The redundancy equivalence is algebra: $\rho_k(N) = o(1)$ iff
$(k+1)\log_2\sigma = o(L_\sigma(N)) = o(\log_2 N/\log_2\sigma)$, i.e.\
$k = o(\log_2 N/\log_2^2\sigma)$; the $\log_2\log_2 N / L_\sigma(N) =
o(1)$ term is always lower order. For (a), with $\Delta_k$ fixed and
$\downarrow 0$: given any $\eta > 0$ pick $k_0$ with $\Delta_{k_0} <
\eta/2$ (possible since $\Delta_k \to 0$); then for $N$ large enough that
$\rho_{k_0}(N) < \eta/2$ (possible since $k_0$ is fixed and $\rho_{k_0}
(N) \to 0$), $\min_k[\Delta_k + \rho_k(N)] \leq \Delta_{k_0} +
\rho_{k_0}(N) < \eta$; as $\eta$ was arbitrary the minimum $\to 0$. The
harmonic computation is the AM--GM minimisation of $a/k + bk$ with $b =
\Theta(1/L_\sigma(N))$, giving minimum $2\sqrt{ab} = \Theta(\sqrt{a/
L_\sigma(N)})$ at $k^\star = \sqrt{a/b} = \Theta(\sqrt{a L_\sigma(N)})$,
which is $o(L_\sigma(N))$ since $a$ is constant; both the model-gain and
redundancy conditions are met, so FV is Shannon-near. For (b), if the
displayed infimum is $\Omega(1)$ then by definition no $k = o(L_\sigma
(N))$ achieves $r_k(N) = h(X) + o(1)$ (and $k = \Omega(L_\sigma(N))$ is
outside FV's validity and anyway makes $\rho_k = \Omega(1)$), so FV
cannot be Shannon-near; conversely if the infimum is $o(1)$ the
minimising $k$ gives a Shannon-near FV index. The $k_\epsilon$
reformulation follows because $\Delta_k + \rho_k(N)$ is small only where
both summands are small, and $\rho_k(N) = o(1)$ caps the admissible $k$
at $o(k_{\mathrm{cross}}(N))$; the witness $\Delta_k^{(N)} = c_0 L_\sigma
(N)/k$ has $a = c_0 L_\sigma(N)$, so the unconstrained AM--GM optimum
$k^\star = \sqrt{a/b} = \Theta(L_\sigma(N))$ lies on the boundary of FV
validity. Restricting to the admissible $k = o(L_\sigma(N))$ leaves the
objective on its decreasing branch, where $\Delta_k^{(N)} + \rho_k(N) \geq
\Delta_{cL_\sigma(N)}^{(N)} = c_0/c = \Theta(1)$ for any $k = cL_\sigma(N)
+ o(L_\sigma(N))$; taking the infimum over the admissible range gives
$\Theta(c_0)$, a positive constant independent of $N$. (The verify
script's (C)/(D) cap $k$ at $\tfrac12 L_\sigma(N)$ and report the
matching $\approx 2c_0$ value, not the smaller unconstrained $\sqrt
{c_0}$.) Claim (c) substitutes $\sigma = 256$, $\log_2\sigma = 8$ into
$k_{\mathrm{cross}}(N) = \log_2 N/\log_2^2\sigma = \log_2 N/64$ and
solves $\rho_1(N) = (2\log_2\sigma + \log_2\log_2 N)/L_\sigma(N) = 1$,
i.e.\ $16 + \log_2\log_2 N = (\log_2 N)/8$, whose root is $\log_2 N
\approx 188$ (so $\rho_1(N) > 1$ for $N \lesssim 2^{188} \approx
10^{57}$). \hfill$\blacksquare$

\emph{Net effect on the open problem.} Corollary 7.27e \textbf{narrows}
the residual open question of Corollary 7.27c piece (3) but does
\textbf{not} fully close it. It establishes: \emph{(i)} the
Ferragina--Venturini escape itself does not close for fixed neural
(infinite-order, even divergent-excess) sources --- so against such
sources a general converse is \emph{false} via FV, exactly as for
finite-order Markov; the honest scope of Theorem 7.27$'$'s refutation is
therefore \emph{every fixed stationary source}, not merely finite-order
Markov. \emph{(ii)} The escape closes precisely for triangular-array
sources whose effective Markov horizon grows as $\omega(\log_2 N /
\log_2^2\sigma)$ (and, at deployment scale, for byte data essentially
always, by (c)). What remains genuinely open is the \emph{same} residual
as before but now correctly delimited: for an escape-\emph{closed} source
(horizon $\omega(\log_2 N)$, e.g.\ a model whose context grows with the
corpus), is there \emph{some other} (non-FV) model-free succinct
random-access structure attaining Shannon-near rate with $o(\sqrt{N})$
query, or does a converse force $\Omega(\sqrt{N})$ there? Corollary 7.27e
removes the false premise (``neural $\Rightarrow$ FV closes'') and pins
the open case to the growing-horizon regime, where the question is
genuinely hard and unresolved.

\textbf{Corollary 7.27f (Growing-horizon block-table floor; and the
unconditional $\sqrt N$ converse is \emph{not} provable).} \emph{Fix
$\sigma \geq 2$, write $L := L_\sigma(N) = \log_2 N/\log_2\sigma$, and let
$X^{(N)}$ be a growing-horizon source (Corollary 7.27e(b)): a triangular
array whose order-$k$ excess obeys $\Delta_k^{(N)} = c_0 L/k$ for $k \leq
L$ (effective Markov horizon $\propto L$, so $k_\epsilon(N) = \Theta(L/
\epsilon) = \omega(\log_2 N/\log_2^2\sigma)$ and the FV escape closes by
7.27e(b)). Call a random-access structure \emph{block-context-table} if it
partitions $x \in \Sigma^N$ into blocks of $b = O(\log_\sigma N)$ symbols
and answers a single-symbol query by a constant number of cell probes into
(i)~a block-local code and (ii)~shared order-$\leq k$ context-decode
tables, attaining per-symbol space $H_k(x) + \rho_k(N)$ with the
canonical redundancy}
\[
  \rho_k(N) \;=\; \Theta\!\Bigl(\tfrac{(k+1)\log_2\sigma + \log_2\log_2 N}
  {L}\Bigr) .
\]
\emph{This is the functional form proved for \textbf{every known}
$O(1)$-access high-order-entropy structure: Ferragina--Venturini (TCS
372(1) (2007) 115); Gonz\'alez--Navarro, ``Statistical encoding of
succinct data structures'' (CPM 2006) / ``Random access to high-order
entropy compressed text''; Grossi--Orlandi--Raman and Belazzougui
--Navarro (\emph{ACM Trans.\ Algorithms} 11(4) (2015), Art.\ 31), whose
access operation runs in $O(1)$ time within this compressed space. Then:}
\begin{itemize}
\item[\emph{(i)}] \emph{(\textbf{POSITIVE, conditional --- new content
beyond 7.27e}.) No block-context-table structure attains the Theorem 7.2
Shannon-near rate on a growing-horizon source. Quantitatively, for the
witness above,}
\[
  \inf_{1 \leq k \leq L}\bigl[\Delta_k^{(N)} + \rho_k(N)\bigr]
  \;=\; \Theta(1) \quad\text{(bounded away from $0$, $N$-independent).}
\]
\emph{The obstruction is a \textbf{model-order / redundancy
incompatibility}: Shannon-near rate forces $\Delta_k^{(N)} = c_0 L/k =
o(1)$, i.e.\ build order $k = \omega(L)$; but $\rho_k(N) = o(1)$ forces
$k = o(L/\log_2\sigma) = o(L)$. The smallest order capturing the context
($k_{\mathrm{model}} = \omega(L)$) strictly exceeds the largest order the
redundancy budget allows ($k_{\mathrm{red}} = o(L)$), so no single block
order is simultaneously low-redundancy and low-model-residual. This
closes the \emph{entire known family} of $O(1)$-access structures (not
just Ferragina--Venturini, the only member 7.27e ruled out) on
growing-horizon sources.}
\item[\emph{(ii)}] \emph{(\textbf{NEGATIVE, unconditional}.) The target
\emph{unconditional} converse --- ``any random-access structure at rate
$h(X)+o(1)$ needs query time $\Omega(\sqrt N)$ or storage $\omega(N
h(X))$, for all schemes'' --- is \textbf{false for every fixed source and
unprovable in general}, because pure single-symbol \emph{access} is not a
hard cell-probe problem. (a)~Belazzougui--Navarro's lower bound is a
\emph{rank}/\emph{select} bound, $\rho\cdot t_{\mathrm{acc}}\cdot
t_{\mathrm{sel}} = \Omega(\log_2^2\sigma/w)$; it is vacuous for pure
access ($t_{\mathrm{sel}}$ unconstrained), and the same authors show
access alone runs in $O(1)$ time at the entropy bound. (b)~The generic
static cell-probe \emph{access} floor at linear space is only $\Omega(\log
N/\log\log N)$ (Andersson--Miltersen 1999; P\u{a}tra\c{s}cu--Demaine 2006
\emph{SIAM J.\ Comput.}\ 35(4):932), asymptotically far below $\sqrt N$.
(c)~A genuine $\Omega(\sqrt N)$-type access bound is known only
\textbf{conditionally}: through the BWT/MTF representation
(Sinha--Weinstein, ``Local decodability of the Burrows--Wheeler
transform,'' STOC 2019: $\Omega(N/t^2)$ for symmetric structures, i.e.\
$t = \Omega(\sqrt{N/r})$ at redundancy $r$), or for \emph{exponentially}
compressible inputs (Chen--Verbin--Yu, ICALP 2012: $\Omega(n^{1/2
-\varepsilon})$ at grammar size $n$ with $L(x) = 2^{\sqrt n}$). The latter
is inapplicable in the RNR regime, where $h(X) > 0$ forces $S = \Theta(N)$
bits (constant-factor, not exponential, compressibility; cf.\ Remark
7.27d(i)).}
\end{itemize}
\emph{Hence the honest status: the growing-horizon escape closes against
the whole \textbf{block-context-table architecture} (the only known route
to $O(1)$-access high-order entropy), but ruling out \textbf{all} succinct
structures would require an unconditional $\Omega(\sqrt N)$ access lower
bound that contradicts the $O(1)$-access achievability above. The residual
of Corollary 7.27e therefore stays \textbf{genuinely open}, now sharpened
to: \emph{is there a model-free $O(1)$-access Shannon-near structure for
growing-horizon sources that is \textbf{not} block-context-table?} A
positive answer needs an architecture escaping the per-block context
table; a negative answer needs a cell-probe access bound stronger than
any currently known (no such unconditional $\Omega(\sqrt N)$ bound exists
for access). The verify script
\texttt{scripts/verify/cor\_7\_27f\_growing\_horizon\_block\_table\_floor.py}
checks (i) the $\Theta(1)$ joint-excess floor on growing-horizon sources,
its vanishing on fixed sources (consistency with 7.27e(a)), the
$k_{\mathrm{model}}/k_{\mathrm{red}} \to$ const $>1$ incompatibility, and
(ii) the $\log N/\log\log N \ll \sqrt N$ gap that blocks an unconditional
converse.}

\emph{\textbf{Proof.}} \emph{(i)} The redundancy $\rho_k(N)$ is the FV
space bound of Corollary 7.27e (its proof), shared verbatim --- up to the
$\Theta(\cdot)$ constant --- by the Gonz\'alez--Navarro and
Grossi--Orlandi--Raman / Belazzougui--Navarro structures, all of which
block the string into $b = \Theta(\log_\sigma N)$-symbol blocks and store
order-$\leq k$ context-decode tables of total size $\Theta(\tfrac{N}{L}
(k\log_2\sigma + \log_2\log_2 N))$ bits; this is the only mechanism known
to deliver $O(1)$-time access at high-order empirical entropy (the
block-local code is read in $O(1)$ probes, the table lookup in $O(1)$).
The floor is on this \emph{context-modeling} layer --- the bits needed to
identify, per block, which order-$k$ conditional model applies and to
store the $\Theta(\sigma^{\leq k})$ active model parameters at the
$\Theta(N/L)$ block granularity --- and is therefore independent of how
near the entropy bound the inner zeroth-order code is (e.g.\ Belazzougui
--Navarro's $o(n)(H_0+1)$ zeroth-order redundancy does not reduce it: that
improvement is to the per-block code, not to the high-order context
overhead, which is $\rho_k(N)$).
Thus on any such structure the per-symbol rate is $h(X) + \Delta_k^{(N)} +
\rho_k(N) + o(1)$ exactly as in 7.27e, with the model term now the
growing-horizon excess. Substituting $\Delta_k^{(N)} = c_0 L/k$ and
$\rho_k(N) = \Theta(k\log_2\sigma/L)$ and minimising over the FV-admissible
range $k \in [1, L]$: the unconstrained AM--GM optimum of $c_0 L/k +
\Theta(k\log_2\sigma/L)$ sits at $k^\star = \Theta(L/\sqrt{\log_2\sigma})
= \Theta(L)$, on the boundary of the admissible range, where the value is
$\Theta(\sqrt{c_0\log_2\sigma}) = \Theta(1)$; and on the admissible
interior the objective is $\geq \Delta_{cL}^{(N)} = c_0/c = \Theta(1)$ for
$k = cL$. Either way $\inf_{k \leq L}[\Delta_k^{(N)} + \rho_k(N)] =
\Theta(1)$, bounded away from $0$ independently of $N$, so no admissible
block order achieves rate $h(X) + o(1)$; orders $k = \omega(L)$ exceed FV
validity and have $\rho_k(N) = \omega(1)$ anyway. The incompatibility is
the stated order separation: $\Delta_k^{(N)} \leq \epsilon \Rightarrow k
\geq c_0 L/\epsilon =: k_{\mathrm{model}}$, while $\rho_k(N) \leq \epsilon
\Rightarrow k \leq \epsilon L/(C\log_2\sigma) =: k_{\mathrm{red}}$, and
$k_{\mathrm{model}}/k_{\mathrm{red}} = \Theta(C\log_2\sigma/\epsilon^2) >
1$ for $\epsilon$ small --- no $k$ meets both. \emph{(ii)} The achievability
of $O(1)$-time pure access at the high-order entropy bound is
Belazzougui--Navarro (op.\ cit., abstract: ``within this compressed space,
operations access and select can be solved in constant or almost-constant
time''); their product lower bound $\rho\,t_{\mathrm{acc}}t_{\mathrm{sel}}
= \Omega(\log_2^2\sigma/w)$ constrains rank/select, not access in
isolation. The generic cell-probe access floor $\Omega(\log N/\log\log N)$
is Andersson--Miltersen / P\u{a}tra\c{s}cu--Demaine. Since $\log N/\log
\log N = o(\sqrt N)$, no unconditional cell-probe argument yields
$\Omega(\sqrt N)$ for access; the only $\sqrt{\cdot}$-type access bounds in
the literature are conditional on a specific representation (BWT/MTF,
Sinha--Weinstein) or on exponential compressibility (Chen--Verbin--Yu),
neither of which holds unconditionally in the constant-factor-compressible
RNR regime. Hence an unconditional $\sqrt N$ converse, which would
contradict the $O(1)$-access achievability for fixed sources (7.27e(a)),
cannot be established. \hfill$\blacksquare$

\emph{Net effect (final delimitation of the 7.27e residual).} Corollary
7.27f resolves the growing-horizon residual \emph{conditionally and
honestly}: \emph{(1)} positively, the converse extends from FV alone
(7.27e(b)) to the \textbf{entire known block-context-table family} of
$O(1)$-access high-order-entropy structures --- on growing-horizon
sources every such structure pays $\Omega(1)$ bits/symbol over $h(X)$, by
the model-order/redundancy incompatibility; \emph{(2)} negatively, a
\textbf{fully architecture-free} $\Omega(\sqrt N)$ converse is unprovable,
since pure access is $O(1)$-achievable at the entropy bound and the
generic cell-probe access floor is only $\log N/\log\log N$. The genuinely
open question is thereby pinned to its irreducible core: the existence (or
not) of a \emph{non-block-context-table} model-free succinct structure
attaining Shannon-near rate with $o(\sqrt N)$ access on growing-horizon
sources. This neither relabels 7.27e (which ruled out only FV) nor
overclaims an unconditional converse; it is the maximal honest statement
the present techniques support.


\textbf{Theorem 7.28 (Non-asymptotic CLT for $L^{\mathrm{rep}}$:
Berry--Esseen rate under exp-mixing).} \emph{Adopt the hypotheses of
Theorem 7.25: $X = (X_1, \ldots, X_N)$ stationary $\alpha$-mixing with
$\alpha_X(k) \leq C_\alpha e^{-c k}$; per-position cost
$\ell_i := -\log_2 Q_i(X_i \mid \mathrm{ctx}_i)$ bounded,
$0 \leq \ell_i \leq B = \log_2(1/\eta)$, depending on a window of $W_N$
coordinates of $X$; with marginal variance
$\sigma_\ell^2 := \mathrm{Var}(\ell_1)$ and long-run variance
\[
  v^2 \;:=\; \sigma_\ell^2 + 2 \sum_{k \geq 1}
  \mathrm{Cov}(\ell_0, \ell_k) ,
\]
which is finite and strictly positive (the series converges absolutely
by Davydov's inequality under exp-mixing; positivity is the standard
non-degeneracy assumption, satisfied whenever $(\ell_i)$ is not an
a.s.\ coboundary). Let $\Phi$ be the standard normal CDF and
$\kappa := m_3 / (B\, v^2)$ the dimensionless third-moment ratio, where
$m_3 := \mathbb{E}|\ell_1 - \mathbb{E}\ell_1|^3 \leq B \sigma_\ell^2$
(so $\kappa \leq \sigma_\ell^2 / v^2$). Then there is a constant
$C_{\mathrm{BE}}$, depending only on $(c, C_\alpha)$ and on $\kappa$,
such that for all $N \geq 2$,
\[
  \sup_{x \in \mathbb{R}}
  \left|
    \Pr\!\left(
      \frac{L^{\mathrm{rep}}(X) - \mathbb{E}[L^{\mathrm{rep}}(X)]}
           {\sqrt{N\, v^2}}
      \leq x
    \right)
    - \Phi(x)
  \right|
  \;\leq\;
  C_{\mathrm{BE}} \cdot
  \frac{B \,(\log N)^2}{v \,\sqrt{N}} .
\]
Equivalently, the Kolmogorov distance between the law of the normalised
repair length and the standard Gaussian is $O\!\bigl(N^{-1/2}
(\log N)^2\bigr)$, the implied constant scaling linearly in the
dependence ratio $\kappa = \Theta(\sigma_\ell^2/v^2)$.}

\emph{Scope, significance, and the log factor.} Theorem 7.28
\emph{complements} Theorem 7.25: where T7.25 gives a one-sided
\emph{tail bound} (an upper bound on $\Pr(|L^{\mathrm{rep}} - \mu|
\geq \varepsilon)$), T7.28 gives a two-sided \emph{distributional
approximation} --- the full CDF of $L^{\mathrm{rep}}$ is uniformly
close to a Gaussian. The operational payoff is \emph{exact} two-sided
confidence intervals: at confidence $1 - \delta$,
$L^{\mathrm{rep}} \in \mathbb{E}[L^{\mathrm{rep}}] \pm
z_{1-\delta/2}\sqrt{N v^2}$ where $z_{1-\delta/2} = \Phi^{-1}(1 -
\delta/2)$ is the Gaussian quantile, valid up to the additive CDF
error above. For moderate $\delta$ this is strictly tighter than the
Bernstein band of T7.25, because the Gaussian quantile $z_{1-\delta/2}$
replaces the Bernstein factor $\sqrt{2\log(2/\delta)}$ (at
$\delta = 10^{-6}$: $z \approx 4.89$ versus $\sqrt{2\log(2/\delta)}
\approx 5.27$, a $\approx 7\%$ reduction in the deviation scale before
constants).

The crucial point --- and the reason this is \emph{not} a routine
``CLT holds'' statement --- is the \emph{explicit log factor}. For
\emph{independent} summands Berry--Esseen gives the log-free rate
$O(N^{-1/2})$, but for strongly mixing summands the best known (and,
for the Kolmogorov metric, generally unavoidable) rate carries a
$(\log N)^2$ correction. This is Tikhomirov's (1980) theorem: under
$\alpha(k) \leq K e^{-\beta k}$ and a finite $(2+\delta)$ absolute
moment, the rate is $O\bigl(N^{-\delta/2}(\log N)^{1+\delta}\bigr)$ for
$0 < \delta \leq 1$; the bounded case $|\ell_i| \leq B$ has all moments,
so we take $\delta = 1$ and obtain $O(N^{-1/2}(\log N)^2)$. (The same
leading term $A N^{-1/2}(\log N)^2$ appears in the recent mixing
Berry--Esseen synthesis; cf.\ the survey rate in
\texttt{arXiv:2604.03712}, Remark~1. Sunklodas 2007 obtains the milder
$O(N^{-1/2}\log N)$ but only for \emph{smooth} / Lipschitz test
functions; passing to indicator test functions --- i.e.\ Kolmogorov
distance --- costs the extra logarithm.) A statement of T7.28 with a
bare $N^{-1/2}$ rate would be \emph{incorrect}: the verify script shows
the $(\log N)^2 \approx 429$ factor at $N = 10^9$ inflates the naive
estimate by two orders of magnitude.

\emph{Reliability regime (honest limitation).} With a literature-order
constant $C_{\mathrm{BE}} \in [1, 5]$, the certified CDF error at the
enwik9 scale ($N = 10^9$, $v^2 \approx 2.1 \cdot 10^4$, $v \approx
145$, $B = 32$) is $D_N = C_{\mathrm{BE}} B (\log N)^2 / (v\sqrt{N})
\approx 3 \cdot 10^{-3}$ to $1.5 \cdot 10^{-2}$. Under the usual rule
that a Gaussian tail estimate at level $\delta$ is trustworthy only
when $\delta \gtrsim 10\,D_N$, the \emph{rigorous} worst-case bound
certifies the Gaussian CI only for $\delta \gtrsim 1.5 \cdot 10^{-1}$
--- so for the extreme tail $\delta = 10^{-6}$, T7.28's certificate is
\emph{vacuous} and Theorem 7.25's Bernstein bound remains the rigorous
tool. (This estimate already treats $C_{\mathrm{BE}}$ as order-unity;
the worst-case $W_N$-dependence of the Tikhomirov constant noted in the
proof, Step 1, could inflate it further, only reinforcing the deferral
to T7.25 for rigorous extreme tails.) We state this limitation plainly.
The redeeming empirical fact
(verified below) is that the \emph{realised} constant is far smaller
than the worst-case $C_{\mathrm{BE}}$ (consistent with $\kappa \leq
\sigma_\ell^2/v^2 \approx 1$ and $B/v \approx 0.22$ at enwik9): the
simulated KS statistic obeys $\mathrm{KS} \cdot \sqrt{N} / (\log N)^2
\approx 0.01$, giving an \emph{actual} CDF error $\approx 1.4 \cdot
10^{-4}$ at enwik9 --- roughly $100\times$ below the certified worst
case --- so the Gaussian CI is \emph{practically} accurate for moderate
$\delta \gtrsim 10^{-3}$, the regime T7.28 is intended for. T7.28 is
thus the sharp inner estimate (exact CIs for moderate $\delta$); T7.25
is the rigorous outer envelope (certified tails for extreme $\delta$);
the two are numerically consistent (below).

\emph{\textbf{Proof.}} The argument has three steps: (1) the cost
sequence inherits exp-mixing; (2) it satisfies the hypotheses of a
Berry--Esseen theorem for strongly mixing sequences; (3) the
moving-function window only shifts the mixing lag and does not alter
the rate.

\emph{Step 1 (mixing inheritance).} Exactly as in the proof of
Theorem 7.25: for $k > W_N$, $\ell_i$ and $\ell_{i+k}$ are functions of
disjoint coordinate windows of $X$, so $\alpha_\ell(k) \leq
\alpha_X(k - W_N) \leq C_\alpha e^{-c(k - W_N)} = (C_\alpha e^{c W_N})
e^{-c k}$ (Bradley 2007, Vol.~1, §5, measurable-function inheritance);
for $k \leq W_N$ the trivial bound $\alpha_\ell(k) \leq 1$ is dominated
by $(\max(1,C_\alpha) e^{c W_N}) e^{-c k}$, since the latter is
$\geq e^{c(W_N - k)} \geq 1$ for $k \leq W_N$. Hence $(\ell_i)$ is
stationary $\alpha$-mixing with exponential rate $\alpha_\ell(k) \leq
C'_\alpha e^{-c k}$, $C'_\alpha := \max(1, C_\alpha) e^{c W_N}$. The
mixing constant is inflated by the window-dependent factor $e^{c W_N}$
but the \emph{rate} $c$ is preserved; this is the only place the
moving-function structure enters, and it affects $C_{\mathrm{BE}}$
(through $C'_\alpha$), not the $N$-dependence.

A caution on this repackaging. Writing $\alpha_\ell(k) \leq C'_\alpha
e^{-ck}$ with $C'_\alpha = e^{c W_N}$ is legitimate but \emph{lossy}:
the BE constant of Tikhomirov's theorem (and of the
\texttt{arXiv:2604.03712} synthesis) depends polynomially on the
mixing prefactor $K = C'_\alpha$, so this packaging would make the
\emph{rigorous} worst-case constant scale like $e^{O(c W_N)}$ ---
catastrophic at $W_N = 2048$. The non-lossy quantity is the
\emph{cumulative} mixing budget $\sum_{k \geq 1} \alpha_\ell(k) \leq
W_N + \sum_{j \geq 1}\alpha_X(j) \leq W_N + C_\alpha/(e^c - 1)$, which
is \emph{linear} in $W_N$; dependence-measure formulations of
Berry--Esseen (Jirak 2016, \emph{Ann.\ Probab.}\ 44(3):2024--2063) are
controlled by such cumulative sums, suggesting the constant grows only
polynomially in $W_N$. We do not pin down the exact $W_N$-dependence
of $C_{\mathrm{BE}}$ here; for the operational claim this does not
matter, because (a) the \emph{rate in $N$} is unaffected, and (b) the
$\delta$-regime in which T7.28 is recommended (moderate $\delta$, where
the \emph{empirically realised} constant governs) is exactly the regime
where the worst-case $W_N$-blowup is irrelevant. For rigorous
extreme-tail guarantees we defer to T7.25 regardless (see the
reliability paragraph and Remark 7.28b). The exact $W_N$-scaling of the
optimal BE constant for moving-window functionals is an open
refinement.

\emph{Step 2 (Berry--Esseen for mixing).} The centred sequence
$Y_i := \ell_i - \mathbb{E}[\ell_i]$ is stationary, bounded
($|Y_i| \leq B$, hence has a finite third absolute moment,
$\delta = 1$), $\alpha$-mixing with $\alpha_\ell(k) \leq C'_\alpha
e^{-ck}$, and has strictly positive long-run variance $v^2 =
\lim_{N} N^{-1}\mathrm{Var}(L^{\mathrm{rep}}) > 0$ (the limit exists
and equals $\sigma_\ell^2 + 2\sum_{k} \mathrm{Cov}(\ell_0,\ell_k)$ by
stationarity and absolute summability of the covariances;
Ibragimov--Linnik 1971, Theorem 18.5.3). These are precisely the
hypotheses of Tikhomirov's Berry--Esseen theorem (Tikhomirov 1980,
\emph{Theory Probab.\ Appl.}\ 25(4):800--818, the main theorem for
$\alpha$-mixing sequences): for an exponentially $\alpha$-mixing
stationary sequence with finite $(2+\delta)$ absolute moment,
$0 < \delta \leq 1$, and $\mathrm{Var}(S_N) \sim N v^2$ with $v^2 > 0$,
\[
  \sup_x \left| \Pr\!\Bigl( \tfrac{S_N - \mathbb{E} S_N}{\sqrt{N v^2}}
  \leq x \Bigr) - \Phi(x) \right|
  \;\leq\; A_{\delta}\, \frac{m_{2+\delta}}{v^{2+\delta}}\,
  N^{-\delta/2} (\log N)^{1+\delta} ,
\]
where $m_{2+\delta} := \mathbb{E}|Y_1|^{2+\delta}$ and $A_\delta$
depends only on $(\delta, c, C'_\alpha)$ (the explicit form of the
dependence is in Tikhomirov's Theorem; see also the exposition in Rio
2017, \emph{Asymptotic Theory of Weakly Dependent Random Processes},
Ch.~7, and the mixing synthesis in \texttt{arXiv:2604.03712}).
Specialising to $\delta = 1$ and $S_N = L^{\mathrm{rep}}$: the third
absolute moment satisfies $m_3 = \mathbb{E}|Y_1|^3 = \mathbb{E}\bigl[
|Y_1|\, Y_1^2 \bigr] \leq B\, \mathbb{E}[Y_1^2] = B\, \sigma_\ell^2$
(using $|Y_1| \le B$, valid for the nonnegative cost
$0 \le \ell_i \le B$ since then $|Y_1| = |\ell_1 - \mathbb{E}\ell_1|
\le B$); hence the moment ratio is
\[
  \frac{m_3}{v^3} \;\leq\; \frac{B\, \sigma_\ell^2}{v^3}
  \;=\; \frac{B}{v}\cdot \frac{\sigma_\ell^2}{v^2}
  \;=\; \frac{B}{v}\cdot \kappa' , \qquad
  \kappa' := \frac{\sigma_\ell^2}{v^2} ,
\]
a \emph{dimensionless} factor. Writing $C_{\mathrm{BE}} := A_1 \kappa'$
(an absolute constant once the mixing rate $(c, C'_\alpha)$ and the
dependence ratio $\kappa'$ are fixed) yields
\[
  \sup_x \left| \Pr\!\Bigl( \tfrac{L^{\mathrm{rep}} - \mathbb{E}
  L^{\mathrm{rep}}}{\sqrt{N v^2}} \leq x \Bigr) - \Phi(x) \right|
  \;\leq\; A_1 \kappa'\, \frac{B (\log N)^2}{v \sqrt{N}}
  \;=\; C_{\mathrm{BE}}\, \frac{B (\log N)^2}{v \sqrt{N}} ,
\]
which is the claim. The constant $C_{\mathrm{BE}}$ depends on
$(c, C'_\alpha)$ and on the dimensionless dependence ratio
$\kappa' = \sigma_\ell^2/v^2$ (and hence, via Step 1, on $W_N$ through
$C'_\alpha$, but \emph{not} on $N$). In the positive-dependence regime
typical of log-loss costs ($\sigma_\ell^2 \leq v^2$, i.e.\
$\kappa' \leq 1$) the ratio is $O(1)$; the verify script measures
$\kappa' \approx 1.0$ for the two-state Markov source.

\emph{Step 3 (no rate change from the moving window).} The only effect
of the $W_N$-window moving-function structure is, via Step 1, the
inflation $C_\alpha \mapsto C'_\alpha = \max(1,C_\alpha) e^{c W_N}$ of
the mixing prefactor; the exponential \emph{rate} $c$ is unchanged, so
$\sum_k \alpha_\ell(k) < \infty$ with the same geometric decay and the
Tikhomirov hypotheses hold verbatim. Consequently the
$N^{-1/2}(\log N)^2$ rate is identical to the non-moving (independent
window) case; the window enters only through the constant. \hfill
$\blacksquare$

\emph{Remark 7.28a (numerical confirmation; consistency with T7.25).}
The verify script \texttt{thm\_7\_28\_berry\_esseen\_l\_rep.py}
simulates $L^{\mathrm{rep}}$ for the two-state Markov source of T7.25
(matched predictor, $\eta = 2^{-32}$, $B = 32$) and (i) computes the
one-sample Kolmogorov--Smirnov distance between the empirical CDF of
$(L^{\mathrm{rep}} - \hat\mu)/\hat\sigma$ and $\Phi$ across
$N \in \{125, \ldots, 8000\}$, fitting $\mathrm{KS} \sim N^{a}$ with
$a \approx -0.42$ (consistent with $-1/2$ mildly flattened by the
$(\log N)^2$ factor; $\mathrm{KS}\cdot\sqrt{N}/(\log N)^2 \approx 0.01$
is roughly flat, confirming the \emph{shape} of the rate); (ii)
verifies empirical quantiles at $p \in \{0.01, 0.5, 0.99\}$ match the
Gaussian predictions $\hat\mu + \hat\sigma\, \Phi^{-1}(p)$ to within
$0.5$ bits at $N = 8000$. The Gaussian quantile $\Phi^{-1}$ is computed
by the Beasley--Springer--Moro / Acklam rational approximation
\emph{with both numerator and denominator polynomials} and a Halley
refinement step, verified to give $\Phi^{-1}(1 - 5\cdot10^{-7}) =
4.8916$ (round-trip error $3\cdot10^{-15}$); a prior implementation
that dropped the denominator returned the spurious value $729$ and is
explicitly guarded against. At the enwik9 scale the script reports the
rigorous worst-case CDF error $D_N \approx 3\cdot10^{-3}$ to
$1.5\cdot10^{-2}$ (constant $C_{\mathrm{BE}} \in [1,5]$) --- documenting
that the bound is \emph{not} below $\delta = 10^{-6}$ --- alongside the
empirical error $\approx 1.4\cdot10^{-4}$, and checks that the Gaussian
$\delta = 10^{-6}$ band $z\sqrt{N v^2} \approx 2.68\,\mathrm{MB}$
($z = 4.8916$) is the \emph{inner} estimate of T7.25's Bernstein band
$2.1$--$4.2\,\mathrm{MB}$, the gap being the factor
$z / \sqrt{2\log(2/\delta)} \approx 0.93$ together with the
$C_1 \in [1,2]$ Bernstein constant.

\emph{Remark 7.28b (which bound to use).} For \emph{point estimates and
moderate confidence} ($\delta \gtrsim 10^{-3}$): use T7.28's Gaussian
CI $\mathbb{E}[L^{\mathrm{rep}}] \pm z_{1-\delta/2}\sqrt{N v^2}$ --- it
is exact up to the (empirically tiny) Berry--Esseen error and tighter
than Bernstein. For \emph{rigorous extreme-tail guarantees}
($\delta \leq 10^{-6}$, e.g.\ a hard storage-budget certificate): use
T7.25's Bernstein bound, whose validity does not degrade with $\delta$.
For \emph{worst-case / non-mixing} sources: use T7.14's Azuma bound.
The three results are nested in sharpness (Azuma $\subseteq$ Bernstein
$\subseteq$ Gaussian-CLT, loosest to tightest deviation scale) and in
rigor-robustness (Gaussian-CLT $\subseteq$ Bernstein $\subseteq$ Azuma,
by breadth of validity), so the operator selects by the $\delta$ regime
and the available knowledge of the source.

\textbf{Theorem 7.29 (Second-order coding rate of RNR; the random-access /
dispersion interplay).} \emph{Adopt the hypotheses of Theorem 7.28 ---
\emph{stationary $\alpha$-mixing with exponentially decaying coefficients} ---
with the ideal predictor $Q_i=P(\cdot\mid X_{i-W_N:i-1})$ on a context window
$W_N=\Theta(\log N)$, so $\ell_i\to-\log_2 P(X_i\mid X_{<i})$. Under exp-mixing
the windowed per-position information converges to the true information
exponentially in $W_N$, so the long-run variance $v^2$ of Theorem 7.28 equals
the source \emph{varentropy} $V:=\lim_N \mathrm{Var}\bigl(-\log_2 P(X^N)\bigr)/N$
(finite, positive) up to a $O(\rho^{W_N})=N^{-\Theta(1)}$ relative error; the
$m=N/K$ per-sub-block context resets add at most $O(m)=O(\sqrt N)$ to
$\mathrm{Var}(L^{\mathrm{rep}})=NV(1+o(1))$ (an $o(\sqrt N)$ effect on the
standard deviation in the balanced regime), so $\sqrt{N v^2}=\sqrt{N V}+
o(\sqrt N)$. Let the
sub-block spacing be $K$ and let $\delta_\infty>0$ be the per-sync-point excess
overhead (Crutchfield--Feldman), so the sync cost is $m\,\delta_\infty=
(N/K)\delta_\infty$. Write $Q^{-1}(\varepsilon):=\Phi^{-1}(1-\varepsilon)$ for the
Gaussian \emph{tail}-quantile (the upper-$(1{-}\varepsilon)$ point; positive for
$\varepsilon<\tfrac12$). Then:}
\begin{enumerate}
\item \emph{(Achievability.) For every $\varepsilon\in(0,1)$, with probability at
least $1-\varepsilon$,}
\[
L^{\mathrm{rep}}_{\mathrm{RNR}}(X^N)\;\le\;
N h \;+\; \tfrac{N}{K}\,\delta_\infty \;+\; \sqrt{N V}\,Q^{-1}(\varepsilon)
\;+\; o(\sqrt N),
\]
\emph{the $\sqrt{NV}\,Q^{-1}(\varepsilon)$ term coming from Theorem 7.28's
uniform Berry--Esseen CLT --- its Kolmogorov (CDF) error $O((\log N)^2/\sqrt N)$
translates, at the Gaussian density $\varphi(Q^{-1}(\varepsilon))$, into a
\emph{quantile} error $O((\log N)^2)=o(\sqrt N)$ --- and the mean
$Nh+(N/K)\delta_\infty$ from Theorem 7.15.}
\item \emph{(Dispersion floor --- converse.) For a memoryless or finite-state
irreducible aperiodic Markov source, the varentropy $V$ is the source dispersion:
under the excess-length criterion ($\Pr[\,\mathrm{length}>L\,]\le\varepsilon$),
the optimal fixed-to-variable codelength obeys $L^\star(N,\varepsilon)\ge
Nh+\sqrt{NV}\,Q^{-1}(\varepsilon)-o(\sqrt N)$ (Kontoyiannis--Verd\'u, \emph{IEEE
Trans.\ Inf.\ Theory}\ 60(2):777--795, 2014, DOI
10.1109/TIT.2013.2291007). So $\sqrt{NV}\,Q^{-1}(\varepsilon)$ cannot be removed
(for a general exp-mixing source the same floor holds wherever the
information-density CLT of K-V's broader conditions applies).}
\item \emph{(The interplay.) At the balanced spacing $K=c\sqrt N$ (Theorem 7.5)
the random-access overhead and the dispersion are \emph{both} $\Theta(\sqrt N)$,}
\[
L^{\mathrm{rep}}_{\mathrm{RNR}}
\;=\; N h \;+\; \sqrt N\Bigl[\,\tfrac{\delta_\infty}{c}
\;+\; \sqrt V\,Q^{-1}(\varepsilon)\,\Bigr]\;+\;o(\sqrt N),
\]
\emph{with their ratio $\dfrac{\delta_\infty/c}{\sqrt V\,Q^{-1}(\varepsilon)}$
\emph{independent of $N$}. Thus byte-granular random access \emph{as realised by
this sub-block scheme} exacts a second-order rate penalty of the \emph{same
order} as the intrinsic stochastic variability; the $\sqrt N$-coefficient
decreases in the access budget $c=K/\sqrt N$ toward the irreducible dispersion
floor $\sqrt V\,Q^{-1}(\varepsilon)$. \emph{(Scope: item 2 is a global source
converse --- it does not assert the $(N/K)\delta_\infty$ overhead is mandatory.
The $(\delta_\infty/c)\sqrt N$ penalty is what \emph{this} scheme pays
(achievability); whether some non-sub-block random-access architecture could
reach the dispersion floor without it --- a tight second-order converse for
local decodability --- is open, cf.\ the space-time product of Theorem 7.27.)}}
\end{enumerate}

\emph{Significance.} Theorem 7.29 upgrades the first-order rate $Nh+o(N)$
(Theorem 7.2) and the distributional CLT (Theorem 7.28) to an \emph{operational}
second-order coding rate, and isolates an RNR-specific phenomenon: the price of
byte-granular random access, $(N/K)\delta_\infty$, lives at exactly the $\sqrt N$
scale of the fundamental source dispersion $\sqrt{NV}\,Q^{-1}(\varepsilon)$, so at
the order-optimal $K=\Theta(\sqrt N)$ the two redundancies are comparable and add.
Larger $K$ shrinks the access overhead toward the dispersion floor at the cost of
$O(K)$ per-query work (Theorem 7.5); the dispersion floor itself is irreducible
(item 2). \emph{Scope:} the identity $v^2=V$ uses the ideal predictor; a
mismatched $Q$ shifts the mean by $N\varepsilon_{\mathrm{KL}}$ (Theorem 7.4) and
replaces $V$ by the long-run variance of the cross-entropy density. The
per-sync-point constant is more precisely the Theorem 7.5 sync constant
$\gamma=\delta_\infty+\log_2(1/\eta)$ (the $\eta$-floor AC slack adds to the
Crutchfield--Feldman excess $\delta_\infty$); the displayed $\Theta(\sqrt N)$
interplay is unchanged with $\gamma$ in place of $\delta_\infty$. It assumes a
\emph{fixed} source of finite, summable excess (it does not cover source families
with non-summable excess). Verified in
\texttt{scripts/verify/thm\_7\_29\_dispersion\_ra\_interplay.py} (V1: analytic
varentropy $=$ empirical $\mathrm{Var}(\sum\ell_i)/N$; V2: the
$(1-\varepsilon)$-quantile matches $Nh+\sqrt{NV}\,Q^{-1}(\varepsilon)$ to
$o(\sqrt N)$; V3: the $\Theta(\sqrt N)$ overhead/dispersion ratio is
$N$-independent).

\emph{Remark 7.29a (The $(N/K)\delta_\infty$ overhead is tight for the
strictly-charged sub-block model: a matching second-order converse).} Theorem 7.29
stated the $(N/K)\delta_\infty$ penalty only as achievability. It is also necessary
in the \emph{strictly-charged} sub-block model: each block $i$ is decoded from its
own \emph{disjoint} pair $(\mathrm{payload}_i,\mathrm{sync}_i)$ alone; all shared
metadata (the predictor / model) is \emph{source-independent}; and the file size is
$\sum_i(|\mathrm{payload}_i|+|\mathrm{sync}_i|)$ (the Lemma 5.1e charging convention).
A \emph{source-dependent} global header is excluded --- it can collapse the bound, as
the $X_n\equiv Y$ counterexample (\S2, the source-dependent-metadata caveat) shows
($Y$ stored once, empty payloads). Then, since $X_{\mathrm{block}_i}$ is determined
by $(\mathrm{payload}_i,\mathrm{sync}_i)$, a source-coding (Shannon) lower bound gives
\[
\mathbb E\bigl[|\mathrm{payload}_i|+|\mathrm{sync}_i|\bigr]\;\ge\;
H(X_{\mathrm{block}_i})\;=\;Kh+\delta_K, \qquad
\delta_K:=H(X^K)-Kh\;\xrightarrow[K\to\infty]{}\;\delta_\infty,
\]
the block-$K$ excess, converging under exp-mixing to the excess entropy
$\delta_\infty$. \emph{Summing the disjoint per-block bounds} (no independence
assumption is needed --- the charging is disjoint) over $m=N/K$ blocks,
\[
\mathbb E[L_{\mathrm{RNR}}]\;\ge\; Nh + \tfrac{N}{K}\,\delta_K
\;=\; Nh+\tfrac{N}{K}\delta_\infty - o(\sqrt N) \qquad(K=\Theta(\sqrt N)),
\]
\emph{matching} the Theorem 7.29 achievability. The intuition is exact: every
independently-decodable block re-pays the block excess $\delta_K\to\delta_\infty$
--- the past--future information a \emph{global} code carries once --- so random
access pays it $m$ times. (Grouping $r$ blocks under one sync is \emph{not} a
loophole: it is the spacing $K\to rK$, moving along the same
$(N/K)\delta_\infty$-overhead vs.\ $O(K)$-access curve.) \emph{Scope:} the
strictly-charged sub-block model (disjoint per-block objects, source-independent
shared metadata, $K=\omega(1)$); it does \emph{not} cover source-dependent global
structures, non-sub-block codes (e.g.\ the Ferragina--Venturini $O(1)$-access
structure refuting the general-scheme space-time bound, Theorem 7.27$'$), or
non-mixing / global-latent sources --- and whether \emph{any} random-access code
reaches the dispersion floor without an $\omega(1)$-per-block penalty (the
redundancy-scale analogue of the Theorem 7.27 space-time product) remains open.
Verified in \texttt{scripts/verify/remark\_7\_29a\_subblock\_ra\_converse.py} (C1:
$\delta_\infty=H(X)-h=I(X_0;X_1)$ for order-$1$; C2: a cold block costs $Kh+\delta_K$;
C3/C4 illustrate the per-block bound and the $(N/K)\delta_\infty=\Theta(\sqrt N)$
scaling).

\emph{Remark 7.30 (Fill-in-the-middle / bisection coding: a neural realisation of the
Theorem 7.27$'$ escape, and its limits).} A causal predictor supplies $P(X_i\mid X_{<i})$;
a \emph{direction-flagged} / fill-in-the-middle (FIM) predictor --- as deployed code models
already are (Bavarian et al., \emph{Efficient Training of Language Models to Fill in the
Middle}, arXiv:2207.14255, 2022; Fried et al., \emph{InCoder}, arXiv:2204.05999, 2022) ---
additionally supplies the two-sided \emph{bridge} $P(X_i\mid X_L,X_R)$. This enables
\textbf{bisection-order} coding: store the endpoints, then code each interval's midpoint
conditioned on its two endpoints, recursing. Three facts, stated honestly --- this is a
\emph{Remark} (a neural \emph{framing} of an escape already established combinatorially in
Theorem 7.27$'$), not a new bound.

\emph{(1) Ideal rate $=H(X)$ for observable finite-order sources.} Bisection order is a
coordinate permutation, so $\sum_{\text{nodes}}-\log_2 P(X_{\text{node}}\mid
\text{coded-so-far})=-\log_2 P(X^N)$ exactly (chain rule). For an \emph{observable}
order-$m$ Markov source with $m$-wide anchors the two anchors $d$-separate the midpoint, so
the bridge equals the chain-rule conditional and the ideal rate is $H(X^N)=Nh+O(1)$ (a
benchmark, not a computable target: the exact $H(X^N)$ is $\mathrm{FP}^{\#P}$-complete,
Remark 7.15g, $\varepsilon$-approximable in poly time only under filter stability, Remark
7.15c --- the deployed coder pays the realized one-forward-pass length, Remark 7.15b); a
model $M'$ adds the \emph{bisection-query} divergence $\sum_{\text{nodes}}\mathbb E\,
D_{\mathrm{KL}}\!\big(P(\cdot\mid\text{coded})\,\Vert\,M'(\cdot\mid\text{anchors})\big)$
(\emph{not} the causal $N\,D_{\mathrm{KL}}(P\Vert M')$). Verified to machine precision
(V1--V2).

\emph{(2) It is the \emph{neural} realisation of the Theorem 7.27$'$ escape --- removing
the sub-block $\delta_\infty$ by \emph{sharing} anchors, not a new bound.} Adjacent blocks
\emph{share} their boundary anchors: each boundary anchor (the $m$-wide boundary context
for an order-$m$ source --- a single symbol only at $m=1$) is coded \emph{once} (as data),
not redundantly re-stored in two disjoint sync states, so the scheme is
\emph{non-sub-block} and the disjoint-charging converse of Remark 7.29a does not bind. This
is the \emph{same} conclusion already reached by Theorem 7.27$'$ --- the
Ferragina--Venturini structure attains $Nh+o(N)$ with fast access for finite-order-Markov
sources, refuting the sub-block penalty --- so the FIM/bisection construction is not a new
asymptote but the realisation of that escape with a \emph{learned bidirectional predictor}
in place of a model-free succinct index, a candidate for the Remark 7.27d(ii) question of
whether a \emph{genuine neural} predictor can take it. Compressed random access at entropy
is itself classical (Grossi--Gupta--Vitter wavelet trees, SODA 2003; Ferragina--Venturini,
\emph{Theoret.\ Comput.\ Sci.}\ 372:115--121, 2007); the contribution here is the
neural/FIM framing and the reading of $\delta_\infty$ as disjoint-sync \emph{redundancy}
that anchor-sharing removes.

\emph{(3) The limits are real --- two adversarial counterexamples.} (i) \emph{Observable
order only, \textbf{not} finite-state.} For a \emph{hidden}-state source (an HMM: hidden
$Y_t$ with $P(Y_t{=}Y_{t-1}){=}0.9$, noisy emission $P(X_t{=}Y_t){=}0.8$) the sufficient
statistic is the \emph{filter} state, not any fixed window of symbols, so $m$-wide anchors
do \emph{not} screen: $P(X_1{=}1\mid X_0{=}1,X_2{=}1)=0.734\neq0.757=P(X_1{=}1\mid
X_0{=}1,X_2{=}1,X_4{=}1)$, and bisection pays a strictly positive excess ($0.028$ bits at
$N{=}5$). Exact bounded-memory screening of a function-of-a-Markov-chain is precisely the
§7.15 ``open middle'' wall; the rate-$H$ claim is therefore scoped to \emph{observable}
finite-order sources, and longer-memory / neural laws pay a bisection excess
$\sum_{\text{nodes}}I(X_{\text{mid}};X_{\text{far}}\mid\text{anchors})$ that widening the
anchors only suppresses. (ii) \emph{$O(\log N)$ is logical depth, not bitstream access.} A
single arithmetic-coded stream is sequential: in pre-order a far-right query's second
ancestor sits at rank $\Theta(N)$, so naive access is $\Theta(N)$, not $O(\log N)$
(arithmetic-coded streams need reset / entry points for random access). Physical random
access uses independently-decodable units at block granularity --- $O(K)$ access with
$o(N)$ reset and seek-index overhead, the \emph{same} access/index regime as the sub-block
scheme but \emph{without} its $(N/K)\delta_\infty$ rate penalty. Verified in
\texttt{scripts/verify/thm\_7\_30\_bidirectional\_bisection\_ra.py} (V1: bridge exact \&
normalised; V2: bisection rate $=-\log_2 P(X^N)$ to machine precision for an observable
order-$1$ source; V3: tree depth $\lceil\log_2(N-1)\rceil$ \emph{and} the $\Theta(N)$
pre-order rank of a far ancestor --- depth is not access; V4: anchor width must match the
memory order; V5: the HMM screening failure above; V6: a pruned bisection with
\emph{two-sided} interior blocks $=H(X^N)$ while the \emph{same} leaves coded one-sided
pay a positive penalty --- so anchor-sharing genuinely removes the $(N/K)\delta_\infty$,
the leaves do \emph{not} ``become causal sub-blocks'').

\emph{Open direction (toward Remark 7.27d).} For an exponentially $\beta$-mixing source one
expects the per-node excess $I(X_{\text{mid}};X_{\text{far}}\mid\text{anchors of width }w)$
to decay in $w$ --- but this is a \emph{filter-stability} question (cf.\ Remark 7.15c;
Atar--Zeitouni), not settled by $\beta$-mixing alone, since conditioning on anchor strings
can amplify dependence. Whether $w=\Theta(\log N)$ anchors then drive the total excess to
$o(N)$ at polylog access is open, and is \emph{not} claimed here. The point would be the
\emph{mechanism} --- a \emph{learned bidirectional predictor} rather than growing-order
empirical tables --- not the achievable rate: Corollary 7.27e shows the
Ferragina--Venturini escape \emph{itself} stays Shannon-near on such mixing sources via
growing empirical order, so on rate alone there is nothing left to win.

\emph{Remark 7.30a (Genuinely-neural FV escape under two-sided filter stability:
closing the Remark 7.30 open direction for hidden-state sources).} The open direction
of Remark 7.30 --- whether a learned bidirectional / fill-in-the-middle predictor
attains the Theorem 7.27$'$ escape (rate $Nh+o(N)$ at polylog anchor context, $O(K)$
physical access) on a \emph{hidden}-state, non-observable-finite-order mixing source
--- was left open because $\beta$-mixing alone does not control the per-node bisection
excess $\Delta(w):=I(X_{\mathrm{mid}};X_{\mathrm{far}}\mid X_{\mathrm{anchor}}(w))$
(conditioning on the anchor strings can \emph{amplify} dependence, so the marginal
mixing rate is the wrong quantity; cf.\ Remark 7.15c). We close it under a
\emph{two-sided filter-stability} hypothesis, exactly as Remark 7.15c closed the
open-middle oracle under one-sided filter stability.

\emph{(TS-FS).} Assume the hidden chain is geometrically ergodic \textnormal{(GE)}
with positive emissions \textnormal{(EP)} --- the Lemma 5.6e / Remark 7.15c hypotheses,
under which the forward (and, by stationarity, the time-reversed) filter is
exponentially stable (Atar--Zeitouni 1997, \emph{Ann.\ IHP}\ 33(6):697--725; van~Handel
2008, \emph{PTRF}\ 145:35--74). A \emph{one-sided} screening argument then suffices, with
no separate smoother result: the $w$-wide left endpoint screens $X_{\mathrm{mid}}$ (and
indeed everything to its right) from the far-left up to $O(\rho^w)$ in total variation ---
the forward filter at the left boundary, run from the $w$ endpoint symbols, differs from
the filter run from the entire far-left past by $O(\rho^w)$ (filter stability) --- and
symmetrically the $w$-wide right endpoint screens the far-right via the reversed filter;
summing the two one-sided contributions,
\[
  \Delta(w)\;=\;I\bigl(X_{\mathrm{mid}};X_{\mathrm{far}}\mid X_{\mathrm{anchor}}(w)\bigr)
  \;\le\;C\,\rho^{\,w},\qquad \rho\in(0,1)\text{ the filter-forgetting rate,}
\]
where $X_{\mathrm{anchor}}(w)$ is the $w$-symbol interval \emph{endpoint} bracketing the
midpoint on each side (the bisection anchors of Remark 7.30, widened from $w=1$ to $w$),
$X_{\mathrm{far}}$ \emph{all} the coded points beyond them (the nearest exterior
dominating, the farther coded points screened geometrically more, so they add only a
$\sum_j\rho^{w+j}=O(\rho^w)$ tail), and the bound is governed by the endpoint
\emph{width} $w$ (the boundary filter is pinned to $O(\rho^w)$ by its $w$ observations)
\emph{independently of the endpoint distance / interior gap} --- the far is screened at
the boundary, not along the interior, so the rate $\rho$ is the same whatever the
interval size (verified gap-robust, V6; far-window saturation, V3). Then:
\begin{enumerate}
\item \emph{\textbf{(The escape closes.)} Choosing $w=\lceil\log_{1/\rho}N\rceil=
\Theta(\log N)$ gives $\rho^w\le1/N$, hence $\Delta(w)\le C/N$, so the total bisection
excess over the $m=O(N)$ interior nodes is $\sum_{\mathrm{nodes}}\Delta(w)\le m\cdot C/N
=O(1)$ --- a bounded total, a fortiori $o(N)$. Hence the FIM/bisection RNR coder attains
$Nh+O(1)$ with $w=\Theta(\log N)$ anchor context and $O(K)$ block-granular physical
access (Remark 7.30(ii)) --- the genuinely-neural
Ferragina--Venturini escape (Remark 7.27d(ii)) now for \emph{hidden}-state
filter-stable sources, not only the observable finite-order sources of Remark 7.30(1).}
\item \emph{\textbf{(The rate is the forgetting rate.)} The required anchor width
scales as $\log N/\log(1/\rho)$: a near-deterministic hidden chain (slow forgetting,
$\rho\to1$) needs proportionally wider anchors. This is precisely why $\beta$-mixing
alone was insufficient (Remark 7.30): the governing quantity is the filter/smoother
forgetting rate $\rho$ (controlled by TS-FS), not the marginal mixing rate.}
\end{enumerate}

\emph{Scope (honest).} TS-FS ($=$ GE${}+{}$EP $\Rightarrow$ smoother stability) is the
two-sided analogue of the Remark 7.15c / Lemma 5.6e filter-stability hypothesis;
it is \emph{sufficient, not necessary}, and --- as in Remark 7.15c --- it makes the CMI
bound $\Delta(w)\le C\rho^w$ the operative primitive, sidestepping the non-monotonicity
of conditional MI under anchor-conditioning that defeats a $\beta$-mixing-only argument.
It establishes only that the excess is \emph{geometrically small} (all the rate needs);
it does not give the exact optimal $w$, nor does it touch the §7.15 ``open middle'' ---
\emph{computing} $\Delta(w)$ exactly is still the exact-finite-block hidden-Markov
problem (Remarks 7.15a--7.15e), here only \emph{bounded}, not evaluated. The non-FS /
non-geometric case (e.g.\ heavy-tailed forgetting) remains open. Verified in
\texttt{scripts/verify/remark\_7\_30a\_filter\_stable\_bisection.py} (V1 the Remark 7.30
HMM excess $\Delta(w)$ decays in $w$, single-symbol anchors failing as in 7.30 V5; V2
the decay is geometric, $\rho\approx0.36$; V3 robust to the far-window; V4 $w=\Theta(\log
N)\Rightarrow$ total excess $o(N)$ at polylog access; V5 the rate is the forgetting rate
--- a near-deterministic chain decays slower, so the escape needs FS, not mixing alone;
V6 the \emph{actual} bisection excess with $w$-wide interval endpoints at a gap $g$ decays
geometrically in the width $w$ at the same rate $\rho$, gap-robust, confirming the
screening is at the boundary).

\emph{Significance.} Remark 7.30a completes the genuinely-neural FV-escape arc ---
Theorem 7.27$'$ (model-free succinct escape) $\to$ Remark 7.30 (neural FIM realisation,
observable finite-order) $\to$ Remark 7.30a (hidden-state, the open direction closed
under TS-FS) --- showing that a learned bidirectional predictor attains the Shannon-near
random-access escape on exactly the filter-stable hidden-Markov sources that model real
neural-LM data, the case Remark 7.30 explicitly left open. The residual is the §7.15
open-middle (exact excess) and the genuinely non-forgetting case.

\textbf{Theorem 7.31 (Large-deviations reliability function of the RNR archive; the
large-deviation scale of the deviation hierarchy and the two-scale random-access
picture).}
\emph{Adopt the per-position notation of Theorems 7.25/7.28/7.29:
$\ell_i:=-\log_2 Q_i(X_i\mid\mathrm{ctx}_i)$ is the realised coded length of position
$i$ and $L^{\mathrm{rep}}=\sum_{i=1}^N\ell_i$. Let $X$ be either i.i.d.\ or a
finite-state irreducible aperiodic Markov chain over a finite alphabet $\Sigma$,
coded under a (possibly mismatched) $W$-bounded model $M$ with
$\mathrm{supp}(P)\subseteq\mathrm{supp}(M)$, so each step cost
$g=-\log_2 M(\cdot\mid\mathrm{window})$ is a bounded functional of a finite
irreducible aperiodic chain. Write $h_{\mathrm{cross}}:=\mathbb E_P[-\log_2
M(X\mid\mathrm{ctx})]=h+D_{\mathrm{KL}}(P\Vert M)$ for the (cross-)entropy rate
(Theorem 7.4) and define the base-$2$ scaled cumulant generating function}
\[
  \Lambda(t)\;:=\;\lim_{N\to\infty}\tfrac1N\log_2\mathbb E\bigl[2^{\,t\,L^{\mathrm{rep}}}\bigr]
  \;=\;
  \begin{cases}
    \log_2\sum_{x}P(x)\,M(x)^{-t}, & \text{(memoryless source \& model)}\\[2pt]
    \log_2\rho\bigl(A_t\bigr),\quad A_t(x,x'):=P(x'\mid x)\,M(x'\mid x)^{-t}, & \text{(Markov),}
  \end{cases}
\]
\emph{$\rho(\cdot)$ the Perron--Frobenius eigenvalue of the tilted transfer matrix
(the memoryless display needs a context-free model $M(x)$; for a window-$W$ model, a
higher-order source, \emph{or an i.i.d.\ source under a contextual ($W\ge1$) model}
--- whose per-step costs are then dependent --- $x$ denotes the finite context state
and $A_t$ the corresponding augmented tilted matrix on $\mathcal Y\times\Sigma^W$,
cf.\ scope (i)).
$\Lambda$ is finite, convex, and real-analytic on all of $\mathbb R$. Then:}
\begin{enumerate}
\item \emph{\textbf{(Reliability function.)} Write $r_+:=\lim_{t\to+\infty}\Lambda'(t)$
for the maximal \emph{sustainable} long-run average cost: the maximum mean over
directed cycles of the support graph (Karp's maximum cycle mean), equal to the
single-step maximum $\ell_{\max}:=\max\{-\log_2 M(x'\mid x):P(x'\mid x)>0\}$ in the
memoryless case but in general strictly below it --- a high-cost edge lying off every
high-mean cycle cannot be sustained. For every overflow rate $R$ in the interior
effective domain $h_{\mathrm{cross}}<R<r_+$,}
\[
  \lim_{N\to\infty}-\tfrac1N\log_2\Pr\!\bigl[L^{\mathrm{rep}}\ge NR\bigr]
  \;=\; E(R)\;:=\;\sup_{t\ge0}\,\bigl[\,tR-\Lambda(t)\,\bigr]\in(0,\infty),
\]
\emph{a continuous, strictly increasing, convex rate with $E(h_{\mathrm{cross}})=0$
(assuming the non-degeneracy $v^2>0$); i.e.\
$\Pr[L^{\mathrm{rep}}\ge NR]=2^{-N(E(R)+o(1))}$. Crucially this includes the
matching \emph{lower} bound $\Pr[L^{\mathrm{rep}}\ge NR]\ge 2^{-N(E(R)+o(1))}$ ---
the overflow is no rarer than $E(R)$ states --- which the one-sided tail bounds
of Theorems 7.14 (Azuma), 7.25 (Bernstein) and 7.26 (Lezaud) do \emph{not}
supply. (At the right endpoint $R=r_+$ the supremum is the $t\to\infty$ limit and
$E$ stays finite --- realisations concentrating on the maximum-mean cycle; for
$R>r_+$ no realisation sustains that average, the event is impossible, and
$E\equiv+\infty$. The lower tail $r_-<R<h_{\mathrm{cross}}$, with
$r_-:=\lim_{t\to-\infty}\Lambda'(t)$ the minimum cycle mean, has the symmetric
exponent $\sup_{t\le0}[tR-\Lambda(t)]$.)}
\item \emph{\textbf{(One density, three regimes.)} The same $\ell$-density governs
every deviation scale with consistent constants:}
\[
  \Lambda'(0)=h_{\mathrm{cross}},\qquad
  \Lambda''(0)=\ln2\cdot v^2,\qquad
  E(h_{\mathrm{cross}}+a)=\frac{a^2}{2\ln2\cdot v^2}+O(a^3),
\]
\emph{where $v^2$ is the long-run variance of Theorems 7.25/7.28 (and $v^2=V$, the
varentropy, in the matched case of Theorem 7.29). The quadratic curvature of $E$
at its zero \emph{reproduces} the dispersion floor $\sqrt{NV}\,Q^{-1}(\varepsilon)$
(Theorem 7.29) and the Berry--Esseen Gaussian (Theorem 7.28). This is the unifying
point: one per-position density $\ell$ and its long-run variance $v^2$ govern the
concentration bounds (Theorems 7.14, 7.25), the central-limit/dispersion regime
(Theorems 7.28, 7.29), and the large-deviation exponent here, the three regimes
meeting at the mean where $E$'s leading curvature $a^2/(2\ln2\,v^2)$ reproduces the
dispersion. The exponents of the manuscript's tail bounds (Theorems 7.14, 7.25) are
each lower bounds on the exact rate $E$, which dominates them; among themselves they
do \emph{not} order uniformly in $a$ --- the variance-aware Bernstein bound is sharp
near the mean (matching the leading $a^2$ constant) while the worst-case Azuma range
bound can be the tighter of the two in the tail --- so only the exact exponent $E$ is
the envelope.}
\item \emph{\textbf{(Matched form is the R\'enyi dual.)} If $M=P$ then (memoryless)
$\Lambda(t)=\log_2\sum_x P(x)^{1-t}=t\,H_{1-t}(X)$, the order-$(1{-}t)$ R\'enyi
entropy, so}
\[
  E(R)\;=\;\sup_{t\ge0}\,t\,\bigl(R-H_{1-t}(X)\bigr):
\]
\emph{the archive-overflow reliability function is the Legendre dual of the source
R\'enyi-entropy spectrum (Remark 7.15d). For the Markov case $\Lambda(t)$ is the
corresponding tilted-matrix R\'enyi rate.}
\item \emph{\textbf{(Two-scale random access.)} The byte-granular random-access
overhead --- the cold-context penalty $(N/K)\delta_\infty=\Theta(\sqrt N)$
(Theorem 7.29) plus the seek-index storage $S(K)=\Theta\!\bigl((N/K)\log K\bigr)=
\Theta(\sqrt N\,\log N)$ (Lemma 7.27b) at the order-optimal $K=\Theta(\sqrt N)$ ---
is $o(N)$, hence sub-exponential. It shifts $R$ by $o(1)$ and leaves $E(R)$
\emph{unchanged}: byte-granular random access is asymptotically \emph{free at the
large-deviations (exponential-rate, $-\tfrac1N\log_2\Pr$) scale}. By contrast at the
central-limit scale it is \emph{not} free --- the cold-context term is
$\Theta(\sqrt N)$, exactly the order of the source dispersion
$\sqrt{NV}\,Q^{-1}(\varepsilon)$ (Theorem 7.29), and the seek index is
$\Theta(\sqrt N\log N)$, a logarithmic factor \emph{above} the dispersion. The
random-access price is thus \emph{confined to the $\sqrt N$ (and $\sqrt N\log N$)
scale}: present in the second-order rate, invisible in the large-deviation
exponent.}
\end{enumerate}

\emph{\textbf{Proof.}} \emph{Existence and form of $\Lambda$.} For i.i.d.\ $X$,
independence gives $\mathbb E[2^{tL^{\mathrm{rep}}}]=\prod_i\mathbb E[2^{t\ell_i}]
=\bigl(\sum_x P(x)M(x)^{-t}\bigr)^N$, so $\frac1N\log_2(\cdot)$ equals
$\log_2\sum_x P(x)M(x)^{-t}$ identically. For the Markov case, with state $x$ and
step cost $g(x,x')=-\log_2 M(x'\mid x)$, the tilted transfer matrix
$A_t(x,x')=P(x'\mid x)2^{tg(x,x')}=P(x'\mid x)M(x'\mid x)^{-t}$ satisfies
$\mathbb E[2^{tS_N}]=\pi^{\!\top}A_t^{\,N}\mathbf 1$ (sum over paths, $\pi$ the
stationary law); since $P$ is irreducible aperiodic and $M(x'\mid x)>0$ on the
support, $A_t$ is a nonnegative primitive matrix, so by Perron--Frobenius it has a
simple dominant eigenvalue $\rho(A_t)>0$ analytic in $t$, and
$\frac1N\log_2(\pi^{\!\top}A_t^{\,N}\mathbf 1)\to\log_2\rho(A_t)=\Lambda(t)$.
\emph{LDP.} On a finite alphabet $\Lambda$ is finite and differentiable on all of
$\mathbb R$ (a fortiori essentially smooth and steep), so the G\"artner--Ellis
theorem (Dembo--Zeitouni, \emph{Large Deviations Techniques and Applications},
2nd ed.\ 1998, Theorem 2.3.6; Cram\'er's theorem, Theorem 2.2.3, is the i.i.d.\
special case) yields the full large-deviation principle for $\frac1N
L^{\mathrm{rep}}$ with good convex rate $E=\Lambda^\ast$. Convex duality gives
continuity, monotonicity and $E(h_{\mathrm{cross}})=0$ (as $\Lambda'(0)=
h_{\mathrm{cross}}$). Since the costs are bounded and
the chain irreducible, $\Lambda'$ is increasing and maps $\mathbb R$ onto the open
interval $(r_-,r_+)$ of achievable long-run averages (its endpoints the min/max
cycle means); for $R$ in this interior the supremum is attained at the finite
$t^\ast$ solving $\Lambda'(t^\ast)=R$ (with $t^\ast>0$ for $R>h_{\mathrm{cross}}$), so
$\sup_{t\ge0}=\sup_{t\in\mathbb R}$ there. At the endpoint $R=r_+$ the maximiser
recedes to $t^\ast\to\infty$ and $E(r_+)$ is the (finite) rate of concentrating on the
maximum-mean cycle; for $R>r_+$ no realisation sustains that average, $\Pr=0$,
$E\equiv+\infty$ (in the memoryless case $r_+=\ell_{\max}$, the all-max-cost realisation). \emph{Consistency (item 2).} $\Lambda'(0)=\mathbb E[\ell]=h_{\mathrm{cross}}$;
in the memoryless case, writing $\psi(s)=\mathbb E[e^{s\ell}]$ and
$\Lambda(t)=(\ln2)^{-1}\ln\psi(t\ln2)$, gives $\Lambda''(0)=(\ln\psi)''(0)\cdot\ln2
=\ln2\cdot\mathrm{Var}(\ell)$; in the Markov case the curvature of the
Perron-eigenvalue log at the origin is exactly the asymptotic (long-run) variance of
the additive functional, $\Lambda''(0)=\ln2\cdot v^2$ with $v^2=\lim_N
\mathrm{Var}(L^{\mathrm{rep}})/N$ the Davydov-summable long-run variance of
Theorem 7.28 (standard for additive functionals of finite irreducible chains).
A second-order Taylor expansion of the Legendre dual at its zero gives
$E(h_{\mathrm{cross}}+a)=a^2/(2\Lambda''(0))+O(a^3)=a^2/(2\ln2\,v^2)$; the Gaussian
tail of $\mathcal N(h_{\mathrm{cross}},v^2/N)$ has base-$2$ exponent
$a^2/(2v^2\ln2)$, identical, so the LDP curvature matches the dispersion. The
Bernstein bound $\Pr[\cdot\ge N(h_{\mathrm{cross}}+a)]\le\exp(-Na^2/(2(v^2+Ba/3)))$
is an \emph{upper} bound on the probability, hence its base-$2$ exponent
$a^2/(2(v^2+Ba/3))/\ln2$ is a \emph{lower} bound on $E$, with the same $a\to0$
leading term and a strictly smaller value for $a>0$. \emph{Item 3.} The R\'enyi
identity $H_\alpha=(1-\alpha)^{-1}\log_2\sum_x P(x)^\alpha$ gives $\sum_x P(x)^{1-t}
=2^{tH_{1-t}}$, whence $\Lambda(t)=tH_{1-t}$. \emph{Item 4.} Let $\Delta_{\mathrm{RA}}
=(N/K)\delta_\infty+S(K)=O(\sqrt N\log N)=o(N)$. Then $\Pr[L^{\mathrm{rep}}+
\Delta_{\mathrm{RA}}\ge NR]=\Pr[L^{\mathrm{rep}}\ge N(R-\Delta_{\mathrm{RA}}/N)]$
and $\Delta_{\mathrm{RA}}/N\to0$, so by continuity of $E$ on the interior the exponent
converges to $E(R)$; at the central-limit scale the cold-context part is $\Theta(\sqrt
N)$ and the seek index $\Theta(\sqrt N\log N)$, both non-vanishing relative to the
dispersion $\Theta(\sqrt N)$, hence visible there (Theorem 7.29). \hfill$\blacksquare$

\emph{Corollary 7.31a (Incompressibility exponent; probabilistic refinement of
Theorem 8.1).} \emph{The probability that the archive fails to compress ---
$L^{\mathrm{rep}}\ge N\log_2|\Sigma|$, the uniform/incompressible rate --- decays as
$2^{-N\,E(\log_2|\Sigma|)}$ with $E(\log_2|\Sigma|)\in(0,\infty)$ finite and positive
\emph{whenever} $\log_2|\Sigma|$ lies in the interior effective domain --- both
$h_{\mathrm{cross}}<\log_2|\Sigma|$ (the model compresses in the mean) \emph{and}
$\log_2|\Sigma|<r_+$ (the incompressible rate is a \emph{sustainable} average, below the
maximum cycle mean; in the memoryless case $r_+=-\log_2 P_{\min}$, so this reads
$h<\log_2|\Sigma|<-\log_2 P_{\min}$). The strict interior is \emph{sufficient}; the
boundary $\log_2|\Sigma|=r_+$ may also give a finite positive exponent but is
periodicity/path-length delicate. Neither sub-condition is dispensable:
$h_{\mathrm{cross}}<\log_2|\Sigma|$ alone fails for a uniform $M$
($\ell\equiv\log_2|\Sigma|$ deterministically, $\Pr=1$, $E=0$), while
$\log_2|\Sigma|<r_+$ can fail \emph{even for non-uniform $M$} on a sparse-support chain
(e.g.\ $P=\bigl(\begin{smallmatrix}.9&.1\\1&0\end{smallmatrix}\bigr)$,
$M=\bigl(\begin{smallmatrix}.6&.4\\.95&.05\end{smallmatrix}\bigr)$ give
$r_+=0.737<1=\log_2|\Sigma|$) --- there the over-raw-rate event has probability zero
for all large $N$ ($E\equiv+\infty$ at $R=\log_2|\Sigma|$), a strong \emph{asymptotic}
no-expansion guarantee (finite-$N$ transients may still exceed it) rather than a defect. When both hold, near-incompressible realisations are exponentially
rare but do occur with a computable rate --- the operational, probabilistic form of the
no-universal-dominance Theorem 8.1: not only does no code shorten \emph{every}
input, the over-budget event has an explicit positive exponent.}

\emph{Scope and limitations.} (i) The finite tilted-matrix $\Lambda=\log_2\rho(A_t)$
extends beyond observable sources: because $M$ is $W$-bounded, the per-step cost is a
finite-window functional, so for \emph{any} source whose (hidden-state, $W$-window)
augmentation is a finite Markov chain --- including hidden-Markov / function-of-Markov
sources --- $\Lambda(t)$ remains a Perron root of the tilted \emph{augmented} transfer
matrix (a larger but still finite state space $\mathcal Y\times\Sigma^W$). The §7.15
``open middle'' obstruction is a \emph{different} object: it concerns the exact entropy
\emph{rate} $h=\Lambda'(0)$ under the \emph{ideal filter} predictor, whose sufficient
statistic is the continuous Blackwell measure rather than a finite window and which has
no finite closed form (Remarks 7.15c--7.15e) --- the finite-window CGF here does not
touch it. A genuinely non-finite-state source (continuous hidden state, non-summable
memory) instead needs the general G\"artner--Ellis existence hypotheses (CGF exists and
is essentially smooth) in place of the Perron form. (ii) Mismatch requires
$\mathrm{supp}(P)\subseteq\mathrm{supp}(M)$ (else a codeword has infinite length);
on a finite alphabet $\Lambda$ is then finite for all $t$. (iii) The exponent is
for the idealised realised length $\sum_i\ell_i$; the $o(N)$ sync/seek/arithmetic-
flush terms do not change it (item 4), whereas a $\Theta(N)$ model drift
($N\varepsilon_{\mathrm{KL}}$, Theorem 7.4) is absorbed into $h_{\mathrm{cross}}$ and
$\Lambda$ via $M$. (iv) This is the reliability function for a \emph{known}
source/model; the universal (unknown-source) overflow exponent --- combining with
the $d/2$ minimax regret of Theorem 7.6$'$ --- is not treated here. \emph{Verified
in} \texttt{scripts/verify/thm\_7\_31\_archive\_ldp\_reliability.py} (V1 i.i.d.\
exact multinomial tail $\to$ Cram\'er $E(R)$ with Bahadur--Rao prefactor; V2 Markov
exact transfer-matrix CGF $\to\log_2\rho(A_t)$ plus tail-trend and long-run-variance
cross-check; V3 mean/dispersion/Bernstein consistency; V4 mismatch; V5 two-scale
random access; V6 incompressibility exponent).

\emph{Significance.} Theorem 7.31 extends the §7 deviation hierarchy --- Azuma
(7.14) $\to$ Bernstein (7.25) $\to$ Berry--Esseen / dispersion (7.28--7.29) $\to$
large deviations (7.31) --- to the exponential-rate scale for known finite-state
sources, supplying the exact archive-overflow \emph{rate} and, for the first time in
the chain, the matching \emph{lower} bound on the overflow probability (a storage
engine cannot be luckier than $E(R)$ promises). It
identifies that rate as the Legendre dual of the source's R\'enyi spectrum
(Remark 7.15d), unifying two otherwise separate parts of the paper, and it pins the
cost of byte-granular random access to the dispersion scale: present at order
$\sqrt N$ (Theorem 7.29) but asymptotically free at the reliability-function scale.

\textbf{Theorem 7.32 (Optimal archive-overflow exponent, and the universality
dividend: a fixed model is tail-suboptimal).} \emph{Let $X^N$ be i.i.d.\ from
$P_0$ over a finite alphabet $\Sigma$ ($|\Sigma|=A$), with $P_0$ of \emph{full
support and non-uniform} (the interior of the simplex, so $D(Q\Vert P_0)<\infty$ for
every $Q$ and $h:=H(P_0)<\log_2 A$ --- a degenerate boundary $P_0$ makes both
exponents below $+\infty$ only on a measure-zero set and is excluded), and the
parametric class the $d=(A{-}1)$-dimensional simplex (extension: order-$k$ Markov,
$d=A^k(A{-}1)$).
Compare three RNR repair-stream codes: the \emph{fixed/ideal} code under the true
model (Theorem 7.31 with $M=P_0$; length $L_{\mathrm{fix}}=-\log_2 P_0(X^N)$), the
\emph{universal} Bayes-mixture code ($P_{\mathrm{mix}}=\int P_\theta(X^N)\,w(\theta)\,
d\theta$ with a smooth full-support prior --- e.g.\ Krichevsky--Trofimov,
$w=\mathrm{Dir}(\tfrac12)$; length $L_{\mathrm{mix}}=-\log_2 P_{\mathrm{mix}}(X^N)$),
and the \emph{overflow-optimal} (length-ranked) code. For an overflow rate
$R\in(h,\log_2 A)$ define the entropy- and cross-entropy-constrained Sanov rates}
\[
  E^\ast(R):=\!\!\inf_{Q:\,H(Q)\ge R}\!\! D(Q\Vert P_0),
  \qquad
  E_{\mathrm{arith}}(R):=\!\!\inf_{Q:\,H(Q)+D(Q\Vert P_0)\ge R}\!\! D(Q\Vert P_0).
\]
\emph{Then:}
\begin{enumerate}
\item \emph{\textbf{(Optimal exponent; converse.)} For every uniquely-decodable code
and every $R\in(h,\log_2 A)$,}
\[
  \Pr\bigl[L\ge NR\bigr]\;\ge\;2^{-N(E^\ast(R)+o(1))},
\]
\emph{and the length-ranked code attains it, so $E^\ast(R)$ is the \emph{optimal}
archive-overflow exponent for $P_0$. The minimiser is the entropy-flattened tilt
$Q_R\propto P_0^{\,s}$ ($s\in(0,1)$ chosen so $H(Q_R)=R$), the same R\'enyi-tilt
family as Theorem 7.31.}
\item \emph{\textbf{(A fixed model --- even the true one --- is tail-suboptimal.)}
The ideal arithmetic code of Theorem 7.31 ($M=P_0$) achieves only}
\[
  \lim_N -\tfrac1N\log_2\Pr[L_{\mathrm{fix}}\ge NR]\;=\;E_{\mathrm{arith}}(R)\;\le\;E^\ast(R),
\]
\emph{with \emph{strict} inequality for $R>h$. Reason: $\{H(Q)\ge R\}\subseteq
\{H(Q)+D(Q\Vert P_0)\ge R\}$, and the $E_{\mathrm{arith}}$-minimiser raises the
cross-entropy through $D$ (an empirical law concentrated on $P_0$-\emph{rare}
symbols, with $H<R$) rather than through entropy --- so the fixed code overflows on
low-entropy-but-atypical realisations that an adaptive code compresses to
$\approx N\widehat H<NR$.}
\item \emph{\textbf{(Universality dividend.)} For empiricals in the interior --- in
particular near the entropy-constrained minimiser $Q_R$ (full support), which
\emph{dominates the Sanov exponent} of $\{\widehat H\ge R\}$ --- the mixture code has}
\[
  L_{\mathrm{mix}}(X^N)\;=\;N\,\widehat H(X^N)\;+\;\tfrac d2\log_2 N\;+\;O(1)
\]
\emph{($\widehat H$ the empirical entropy; Laplace / Clarke--Barron 1990 --- the
leading $\tfrac d2\log_2 N$ is common to all smooth full-support priors for interior
empiricals, the $O(1)$ prior-dependent and minimised by the Jeffreys/KT prior).
Boundary empiricals (some symbol absent) \emph{can} also enter $\{\widehat H\ge R\}$
when $A\ge3$, carrying a prior-dependent boundary $\log_2 N$ coefficient --- a type
with $k$ positive symbols gets $[(A-k)\alpha+\tfrac{k-1}2]\log_2 N$ under a
$\mathrm{Dir}(\alpha)$ mixture, which equals the interior $\tfrac d2\log_2 N$ for
KT/Jeffreys ($\alpha=\tfrac12$) and is larger for $\alpha>\tfrac12$. Either way their
type has $D(Q\Vert P_0)>E^\ast$ strictly, so they are exponentially negligible and
leave the exponent unchanged. The mixture code therefore attains the \emph{optimal} exponent
$E^\ast(R)$ for \emph{every} interior $P_0$ in the class simultaneously (the mixture
length is source-independent; Csisz\'ar's method of types). Its only
price over the fixed code is the universal regret: for the KT/Jeffreys mixture (and
NML) this is the \emph{minimax} $\tfrac d2\log_2 N+O(1)$ (Theorem 7.6$'$), while for
an arbitrary smooth full-support prior the \emph{interior/source-typical} redundancy
is $\tfrac d2\log_2 N+O(1)$ (its boundary worst case may be larger, but
exponent-negligibly so). In every case it is a third-order term, $o(\sqrt N)$ ---
invisible at the dispersion scale (Theorem 7.29) --- and $o(N)$ --- invisible at the
large-deviation scale (Theorem 7.31). So universality, ordinarily counted a \emph{cost}, is asymptotically
a pure overflow-tail \emph{improvement} $E_{\mathrm{arith}}\to E^\ast$: a third-order
mean cost, invisible at the dispersion and large-deviation scales, in exchange for a
strict gain exactly where reliability is measured --- the exponent, for nonuniform
full-support $P_0$ and fixed $R\in(h,\log_2 A)$.}
\end{enumerate}

\emph{\textbf{Proof.}} \emph{(1) Converse by the method of types.} By the
Kraft--McMillan inequality any uniquely-decodable code has fewer than $2^{NR}$
codewords of length $<NR$, so at most $2^{NR}$ sequences are short-coded. But a type
$Q$ with $H(Q)>R$ contains $2^{N H(Q)+o(N)}>2^{NR}$ sequences, so at most the
vanishing fraction $2^{-N(H(Q)-R)}$ of them can be short-coded --- essentially every
sequence of every $H(Q)>R$ type overflows. Each such sequence has $P_0$-mass
$2^{-N(H(Q)+D(Q\Vert P_0))}$, so the overflow mass is at least
$\sum_{Q:H(Q)>R}2^{N H(Q)}2^{-N(H(Q)+D(Q\Vert P_0))}(1-o(1))=\sum_{Q:H(Q)>R}
2^{-N D(Q\Vert P_0)}\doteq 2^{-N\inf\{D(Q\Vert P_0):H(Q)\ge R\}}$ (the number of
types is polynomial; Csisz\'ar--K\"orner / Sanov). Hence $\Pr[L\ge NR]\ge
2^{-N(E^\ast+o(1))}$ for \emph{every} uniquely-decodable code. The length-ranked
code (shortest codewords to the $2^{NR}$ most-probable sequences) leaves, up to the
exponent, the union of all $H(Q)>R$ types, of mass $2^{-N(E^\ast+o(1))}$, attaining
it (ranking by probability rather than by entropy can differ on low-entropy rare
types but not at exponent scale).
The program $\min\{D(Q\Vert P_0):H(Q)\ge R\}$ is convex ($D$ convex, $\{H\ge R\}$
convex); its stationarity $\nabla_Q[D(Q\Vert P_0)-\lambda H(Q)]=\mathrm{const}$
gives $\log(Q/P_0)\propto -\lambda\log Q$, i.e.\ $Q_R\propto P_0^{1/(1+\lambda)}$,
$\lambda\ge0$ fixed by $H(Q_R)=R$. \emph{(2)} $E_{\mathrm{arith}}$ is the matched
Theorem 7.31 exponent: $\sup_{t}[tR-\Lambda_{P_0}(t)]=\inf\{D(Q\Vert P_0):\mathbb
E_Q[-\log_2 P_0]\ge R\}$ (Cram\'er $=$ Sanov-contraction), and $\mathbb E_Q[-\log_2
P_0]=H(Q)+D(Q\Vert P_0)$. The inclusion of feasible sets gives $E^\ast\ge
E_{\mathrm{arith}}$; strictness for $R>h$ because the $E_{\mathrm{arith}}$-minimiser
$Q$ has $H(Q)<R$ (it is the cross-entropy shell, reached at lower entropy by tilting
toward rare symbols), so it is infeasible for $E^\ast$, which must pay strictly more
$D$. \emph{(3)} $L_{\mathrm{mix}}=-\log_2\int P_\theta(X^N)w(\theta)d\theta$; Laplace
about the interior MLE $\hat\theta$ (the empirical law, positive prior density)
gives $-\log_2 P_{\hat\theta}(X^N)+\tfrac d2\log_2\tfrac{N}{2\pi}+\tfrac12\log_2\det
\hat I-\log_2 w(\hat\theta)+o(1)$, and $-\log_2 P_{\hat\theta}(X^N)=N\widehat
H(X^N)$, giving the display on interior types (Clarke--Barron 1990, \emph{IEEE
Trans.\ Inf.\ Theory} 36(3):453--471). The mixture length depends only on the data,
\emph{not} on $P_0$. By Sanov under $P_0$, $\Pr[L_{\mathrm{mix}}\ge NR]$ is dominated
by the type minimising $D(Q\Vert P_0)$ over $\{Q:L_{\mathrm{mix}}/N\to H(Q)\ \&\
H(Q)\ge R\}$; the minimiser is the interior $Q_R$, where $L_{\mathrm{mix}}=N\widehat
H+\tfrac d2\log_2 N+O(1)$ and the threshold $R-\tfrac d2(\log_2 N)/N\to R$, giving
$-\tfrac1N\log_2\Pr[L_{\mathrm{mix}}\ge NR]\to\inf\{D(Q\Vert P_0):H(Q)\ge R\}=
E^\ast(R)$. Boundary types in $\{\widehat H\ge R\}$ (possible for $A\ge3$) have
$D(Q\Vert P_0)>E^\ast$ and are exponentially negligible whatever their $\log_2 N$
coefficient, so they do not change the exponent. Being source-independent, the single
universal code attains $E^\ast$ for every interior $P_0$ (the streaming/type form of
Csisz\'ar's method of types, \emph{IEEE Trans.\ Inf.\ Theory} 44(6):2505--2523,
1998). Finally $\tfrac d2\log_2 N=o(\sqrt N)=o(N)$. \hfill$\blacksquare$

\emph{Scope and significance.} The result is for nonuniform full-support
i.i.d./finite-order Markov sources with a smooth full-support prior (a
\emph{sufficient} regularity condition --- KT/Jeffreys is the canonical minimax
choice --- not a necessary one); the converse is the standard Kraft/type bound, the
$\tfrac d2\log_2 N$ regret is the Clarke--Barron / Theorem 7.6$'$ minimax-leading
constant (up to an $O(1)$), and the per-source optimality of the universal code is
method-of-types (Csisz\'ar--K\"orner). For hidden-state / continuous sources the
empirical-entropy functional and the regret constant change --- the §7.15 ``open
middle'' returns for the \emph{exact} constant --- and the qualitative dividend is
\emph{expected} to persist under the analogous regularity (an adaptive sufficient
statistic carrying $o(\sqrt N)$ regret \emph{and} obeying a large-deviation
principle), but $o(\sqrt N)$ regret alone does not suffice and we do not establish it
there. The significance is a reversal of the usual
accounting: the $\tfrac d2\log_2 N$ price of \emph{not knowing the source} buys a
\emph{strictly better} archive-reliability exponent, and the fixed-true-model RNR
code of Theorem 7.31 is itself overflow-suboptimal --- it charges the full
cross-entropy on atypical data that an adaptive coder re-fits. Theorem 7.32 thus
completes the universal half of the §7 reliability picture (Theorem 7.31 being the
known-source half). Verified in
\texttt{scripts/verify/thm\_7\_32\_universality\_overflow\_dividend.py} (V1
$L_{\mathrm{mix}}=N\widehat H+\tfrac d2\log_2 N+O(1)$; V2 $E^\ast>E_{\mathrm{arith}}$
strict for $R>h$; V3 the length-ranked converse $\to E^\ast$; V4 Monte-Carlo, the
universal-vs-fixed overflow ratio compounding as $2^{N(E^\ast-E_{\mathrm{arith}})}$;
V5 the redundancy sign-flip on the gap event; V6 the $o(\sqrt N)$/$o(N)$ invisibility
with $E^\ast>E_{\mathrm{arith}}$).

\textbf{Remark 7.32a (Universal and distributed RNR at second order: the parameter is
a cheap header, the side information is not).} The clean $\sqrt N$ dispersion of
Theorem 7.29 is stated for a single \emph{known} source. Lifting that assumption in
the two canonical directions --- \emph{unknown} parameter and \emph{distributed} side
information --- has \emph{opposite} second-order consequences for random access, and
the contrast is the point.

\emph{(i) Universal: the $\sqrt N$ floor survives random access via a cheap header.}
For a $d$-dimensional smooth full-support parametric family at an interior $\theta$
(varentropy $V_\theta>0$), the global Bayes/NML mixture attains
$N h_\theta+\tfrac d2\log_2 N+\sqrt{N V_\theta}\,Q^{-1}(\varepsilon)+o(\sqrt N)$
(Clarke--Barron 1990 regret atop the Kontoyiannis--Verd\'u 2014 dispersion), but it is
\emph{sequential} --- block $i$'s code depends on the running posterior, so it is not
locally decodable. Random access is recovered by a \emph{two-pass} scheme: estimate
$\hat\theta$ from the whole stream, store it as a global header of
$\tfrac d2\log_2 N=o(\sqrt N)$ bits, and code each sub-block independently against the
\emph{fixed} $P_{\hat\theta}$. The aggregate repair length is then
\[
  L^{\mathrm{rep}}_{\mathrm{univ}}(N,\varepsilon,K)
  = N h_\theta + \tfrac NK\,\delta_\infty + \tfrac d2\log_2 N
    + \sqrt{N V_\theta}\,Q^{-1}(\varepsilon) + o(\sqrt N),
\]
\emph{locally decodable}, with the Kontoyiannis--Verd\'u dispersion
$\sqrt{N V_\theta}\,Q^{-1}(\varepsilon)$ \emph{intact}: the header
$\tfrac d2\log_2 N$ is dispersion-invisible. So the price of not knowing the source is
second-order-free \emph{even with random access} --- because the unknown is only $d$
real numbers, amortisable as one cheap global header. This header $\hat\theta$ is itself
\emph{source-dependent} (a function of the stream), so the two-pass scheme lies formally
\emph{outside} the strictly-charged sub-block model of Lemma 5.1e / Remark 7.29a, whose
per-block converse admits only source-\emph{independent} shared metadata. It nonetheless
does not contradict that converse's conclusion: collapsing the $\Theta(\sqrt N)$ floor
the converse establishes would take a $\Theta(N)$-bit source-dependent header (the
$X_n\equiv Y$ degeneracy of Remark 7.29a), whereas $\hat\theta$ carries only
$o(\sqrt N)$ bits and no per-block content --- so achievability (i) and the Remark 7.29a
converse coexist (the floor stands; the header is simply too small to defeat it). This
is the CLT companion to the large-deviation dividend of Theorem 7.32.

\emph{(ii) The header earns its keep.} A naive single-pass alternative whose
sufficient statistics \emph{reset} at each of the $m=N/K$ sync boundaries re-pays the
Clarke--Barron regret per sub-block:
\[
  \tfrac NK\cdot\tfrac d2\log_2 K \;=\; \Theta(\sqrt N\log N)
  \qquad\text{at } K=\Theta(\sqrt N),
\]
a $\log N$ factor \emph{above} the dispersion. This is not a tax of universality but of
\emph{discarding} the cheap header: (i) shows it is avoidable. (An achievability
comparison of two schemes, not a converse.)

\emph{(iii) Distributed (RNR Slepian--Wolf, Theorem 7.8): the genuinely open tax.}
Here the analogue of the ``unknown'' is the side information $Y$, which is $\Theta(N)$
\emph{source data}, not a $d$-parameter --- there is no cheap global header to amortise
it. The global SW second order at an interior rate point is the $O(\sqrt N)$ term
governed by the entropy-density dispersion matrix of
$(-\log P(X\mid Y),-\log P(Y\mid X),-\log P(X,Y))$ (Tan--Kosut 2014), but the RNR
sub-block construction decodes each block independently, and the per-block union bound
holding all $m$ corner constraints at confidence $\varepsilon/m$ pays
\[
  m\sqrt K\,Q^{-1}(\varepsilon/m) \;=\; \Theta\!\big(N^{3/4}\sqrt{\log N}\big)
  \qquad\text{at } K=\Theta(\sqrt N),
\]
strictly above $\sqrt N$. Whether a locally decodable sub-block Slepian--Wolf code can
restore the $\sqrt N$ floor (sharing typicality slack across blocks instead of a naive
$\varepsilon/m$ union bound) is \emph{open} --- the genuine second-order random-access
question, in sharp contrast to (i), where the cheap parameter header settles it.

\emph{The contrast.} For the natural sub-block schemes the achievable second orders are
\[
  \underbrace{\sqrt N}_{\substack{\text{known source}\\ \text{\emph{and} universal (header)}}}
  \;\;\lneq\;\;
  \underbrace{\sqrt N\log N}_{\substack{\text{universal,}\\ \text{naive reset (avoidable)}}}
  \;\;\lneq\;\;
  \underbrace{N^{3/4}\sqrt{\log N}}_{\substack{\text{Slepian--Wolf}\\ \text{sub-block (open)}}} ,
\]
and the lesson is the dichotomy: a $d$-dimensional \emph{parameter} is a cheap header
that keeps random-access universality at the $\sqrt N$ floor, whereas $\Theta(N)$ bits
of \emph{side information} admit no such header and leave the random-access second order
genuinely open. Verified in
\texttt{scripts/verify/remark\_7\_32a\_ra\_secondorder\_tax.py} (V1 the header
$\tfrac d2\log_2 N=o(\sqrt N)$, so (i) keeps the $\sqrt N$ floor; V2 the avoidable
naive-reset cost $\Theta(\sqrt N\log N)$; V3 the SW sub-block union-bound cost
$\Theta(N^{3/4}\sqrt{\log N})$; V4 the scale ordering).

\textbf{Theorem 7.33 (Moderate-deviation rate of the RNR archive; the random-access
penalty fades from the exponent but not from the probability).} \emph{Adopt the
hypotheses of Theorem 7.29 with $X$ \emph{i.i.d.\ or a finite-state irreducible
aperiodic Markov} source (the clean Cram\'er--Petrov regime; the general exp-mixing
case holds on a restricted moderate range, see Scope): ideal predictor, bounded
per-position cost $\ell_i$, entropy rate $h$, varentropy $V>0$, and
$L^{\mathrm{rep}}=\sum_i\ell_i$ the \emph{repair stream} (payload $+$ sync, as in
Theorem 7.29; the seek index is treated separately below). With the
\emph{moderate-deviation} scaling --- a budget $Nh+a_N\sqrt{NV}$ with $a_N\to\infty$,
$a_N=o(\sqrt N)$ --- then:}
\begin{enumerate}
\item \emph{\textbf{(Moderate-deviation rate.)}}
\[
  \Pr\!\bigl[L^{\mathrm{rep}}\ge Nh+a_N\sqrt{NV}\bigr]
  \;=\;\exp\!\Bigl(-(1+o(1))\,\tfrac{a_N^2}{2}\Bigr),
  \quad\text{i.e.}\quad
  -\tfrac1{a_N^2}\log_2\Pr[\cdot]\to\tfrac1{2\ln2}.
\]
\item \emph{\textbf{(The random-access penalty is a constant standardized shift.)} With
the sub-block overhead $(N/K)\delta_\infty$ at the order-optimal $K=c\sqrt N$, the RNR
archive obeys}
\[
  \Pr\!\bigl[L^{\mathrm{rep}}_{\mathrm{RNR}}\ge Nh+a_N\sqrt{NV}\bigr]
  \;=\;\exp\!\Bigl(-(1+o(1))\,\tfrac{(a_N-\beta)^2}{2}\Bigr),
  \qquad \beta:=\frac{\delta_\infty}{c\sqrt V},
\]
\emph{the cold-context payload penalty entering only as the $N$-independent shift
$\beta$ of the standardised deviation. (The seek index of Theorem 7.31, being
$\Theta(\sqrt N\log N)$, would add a \emph{further} standardized shift of order
$\log N/\sqrt V$ --- diverging, not constant --- inflating the overflow probability
even more; the constant-$\beta$ form is for the payload$+$sync stream, excluding the
index, consistent with the $L^{\mathrm{rep}}$ of Theorems 7.29/7.31.)}
\item \emph{\textbf{(The fade --- unifying Theorems 7.29 and 7.31.)} As $a_N$ ranges
over the moderate scale,}
\[
  \underbrace{\frac{(a_N-\beta)^2}{a_N^2}=\Bigl(1-\tfrac{\beta}{a_N}\Bigr)^2\to1}
  _{\text{exponent: RA \emph{invisible} (}\to\text{Thm 7.31)}}
  \qquad\text{while}\qquad
  \underbrace{\frac1{a_N}\log\frac{\Pr[L^{\mathrm{rep}}_{\mathrm{RNR}}\ge\cdot]}
  {\Pr[L^{\mathrm{rep}}\ge\cdot]}\;\to\;\beta>0}_{\text{probability: RA \emph{visible}}},
\]
\emph{i.e.\ the random-access overflow probability is inflated by the diverging factor
$e^{(1+o(1))\beta a_N}\to\infty$. At the central-limit end ($a_N=\Theta(1)$) the shift
$\beta$ is comparable to $a_N$ and the penalty is fully visible, the exact ratio being
the Gaussian tail ratio $Q(a_N-\beta)/Q(a_N)$ (Theorem 7.29; its large-$a_N$ Mills form
is $e^{\beta a_N-\beta^2/2}$); as $a_N\to\infty$ it vanishes from the exponent
(Theorem 7.31) yet keeps multiplying the probability by $e^{(1+o(1))\beta a_N}$. So
across the moderate-deviation bridge the random-access penalty \emph{fades from the
exponent but not from the probability}.}
\end{enumerate}

\emph{\textbf{Proof.}} \emph{(1)} For i.i.d./finite-state-Markov $(\ell_i)$ (bounded,
so Cram\'er's condition holds) the normalised partial sum
$S^\ast_N:=(L^{\mathrm{rep}}-Nh)/\sqrt{NV}$ obeys the sharp Cram\'er--Petrov expansion
$\log\Pr[S^\ast_N\ge x]=\log Q(x)+\tfrac{x^3}{\sqrt N}\lambda\!\bigl(\tfrac{x}{\sqrt
N}\bigr)+o(1)$ uniformly for $x=o(\sqrt N)$, where $Q=1-\Phi$ and $\lambda$ is the
Cram\'er series (Petrov, \emph{Sums of Independent Random Variables}, 1975,
Ch.~VIII). At $x=a_N\to\infty$ this gives $\Pr[S^\ast_N\ge a_N]=\exp(-(1+o(1))a_N^2/2)$
(the rate $a_N^2/2$, the moderate-deviation principle); the $m=N/K=O(\sqrt N)$ context
resets add $o(N)$ to the variance (Theorem 7.29 scope), absorbed in the $o(1)$.
\emph{(2)} By Theorem 7.29,
$L^{\mathrm{rep}}_{\mathrm{RNR}}=L^{\mathrm{rep}}+(N/K)\delta_\infty+o(\sqrt N)$
deterministically up to an $o(\sqrt N)$ fluctuation negligible at this scale, so the
event becomes $\{L^{\mathrm{rep}}\ge Nh+a_N\sqrt{NV}-(N/K)\delta_\infty\}=
\{(L^{\mathrm{rep}}-Nh)/\sqrt{NV}\ge a_N-\beta\}$ with $\beta=(N/K)\delta_\infty/
\sqrt{NV}=\sqrt N\,\delta_\infty/(c\sqrt{NV})=\delta_\infty/(c\sqrt V)$; apply (1) at
$a_N-\beta$. \emph{(3)} The exponent ratio $(a_N-\beta)^2/a_N^2=(1-\beta/a_N)^2\to1$
($\beta/a_N\to0$), so $-\tfrac1{a_N^2}\log\Pr$ is unchanged. For the probability ratio
the bare moderate-deviation principle is \emph{insufficient} (its $o(a_N^2)$ error,
divided by $a_N$, would swamp the $O(a_N)$ effect); we use the Cram\'er--Petrov
expansion of (1). With $Q=1-\Phi$ and $\lambda_0=\lambda(0)$,
\[
  \log\frac{\Pr[S^\ast_N\ge a_N-\beta]}{\Pr[S^\ast_N\ge a_N]}
  =\underbrace{\log\tfrac{Q(a_N-\beta)}{Q(a_N)}}_{=\,\beta a_N-\beta^2/2+O(\beta/a_N)\ (\text{Mills})}
  +\underbrace{\tfrac{1}{\sqrt N}\bigl[(a_N-\beta)^3\lambda(\cdot)-a_N^3\lambda(\cdot)\bigr]}
  _{=\,-3\beta\lambda_0 a_N^2/\sqrt N\,(1+o(1))}+o(1).
\]
Dividing by $a_N$, both corrections --- $\beta/a_N$ and $\lambda_0 a_N/\sqrt N$ ---
vanish for $a_N=o(\sqrt N)$, giving $\tfrac1{a_N}\log(\text{ratio})\to\beta$. At
$a_N=\Theta(1)$ the exact ratio is $Q(a_N-\beta)/Q(a_N)$ (the Theorem 7.29 /
Berry--Esseen regime, Mills form $e^{\beta a_N-\beta^2/2}$); at $a_N=\Theta(\sqrt N)$,
$\beta/a_N=\Theta(1/\sqrt N)\to0$ (the Theorem 7.31 regime, the $O(\sqrt N\log N)$
index overhead being $o(N)$).
\hfill$\blacksquare$

\emph{Scope and significance.} The clean statement is for i.i.d.\ / finite-state
irreducible aperiodic Markov sources (bounded $\ell_i$, $V>0$, ideal predictor), where
the Cram\'er--Petrov expansion holds uniformly over the full moderate range
$a_N=o(\sqrt N)$ and controls the probability ratio. A mismatched $M$ replaces $h$ by
$h_{\mathrm{cross}}$ and $V$ by the cross-entropy long-run variance (Theorem 7.31). For
general exp-mixing sources only the cruder moderate-deviation \emph{rate} is available,
and published $\alpha$-mixing moderate-deviation principles restrict the speed (roughly
$a_N=o(\sqrt N/(\log N\,\log\log N))$ for $\alpha$-mixing; the full $o(\sqrt N)$ range
needs stronger $\phi$-mixing / projective conditions, Merlev\`ede--Peligrad--Rio,
Dedecker--Merlev\`ede--Peligrad--Utev); the ratio refinement there is correspondingly
weaker. The seek index of Theorem 7.31, if counted, adds a diverging $\Theta(\log N)$
standardized shift (item 2), so the constant-$\beta$ statement is for the payload$+$sync
stream. Theorem 7.33 closes the last scale of
the §7 deviation hierarchy --- concentration (7.14/7.25), central limit / dispersion
(7.28/7.29), \emph{moderate deviations} (7.33), large deviations (7.31/7.32) --- and
unifies its two random-access verdicts: the single constant shift $\beta=\delta_\infty/
(c\sqrt V)$ is \emph{both} the Theorem 7.29 dispersion-scale penalty (visible when the
deviation is $\Theta(\sqrt N)$) \emph{and} the Theorem 7.31 large-deviation
non-penalty (invisible when the deviation is $\Theta(N)$); the moderate regime
interpolates exactly, showing the penalty fades from the exponent at the rate
$1-(1-\beta/a_N)^2=2\beta/a_N+O(a_N^{-2})$ while its multiplicative imprint on the
overflow probability, $e^{\beta a_N}$, grows without bound. Verified in
\texttt{scripts/verify/thm\_7\_33\_moderate\_deviation\_ra\_fade.py} (V1 the MDP rate
$\to\tfrac1{2\ln2}$ along $a_N\sim N^{1/4}$; V2 the RA penalty as the constant shift
$\beta$; V3 the fade --- exponent ratio $\to1$ while the probability factor $\to\infty$;
V4 the LDP rate's quadratic-at-mean equals the MDP rate).

\textbf{Theorem 7.35 (Extreme-value law for the worst-case random-access query cost; the
read-buffer provisioning margin).} \emph{Let $\{X_t\}$ be stationary with entropy rate $h$
and varentropy rate $V=\lim_N\frac1N\mathrm{Var}(-\log_2 P(X^N))\in(0,\infty)$. Partition
$X^N$ into $m=N/K$ contiguous blocks and let $L_i=-\log_2 P(X_{(i)}\mid\mathrm{ctx}_i)+O(1)$
be the codelength of block $i$ under the matched RNR sub-block coder --- the repair-stream
bits a random-access query landing in block $i$ must read. Let $M_m:=\max_{i\le m}L_i$ be
the worst-case query cost (the read buffer that serves any single byte-granular query), and
let $K\to\infty$, $m=N/K\to\infty$.}

\emph{\textbf{(a) (Leading-order provisioning law.)} If $X$ is i.i.d.\ or a finite-state
irreducible aperiodic Markov source and $\ln m=o(K)$, then}
\[
  M_m\;=\;Kh\;+\;\sqrt{2KV\ln m}\,\bigl(1+o_P(1)\bigr).
\]
\emph{(The leading-order law extends to stationary exp-mixing $\alpha$ sources --- the
Theorem 7.28/7.29 class --- on the restricted range $\ln m=o\bigl(K/(\log K\,\log\log
K)^2\bigr)$, where the published $\alpha$-mixing moderate-deviation principle reaches the
extremal level $x\asymp\sqrt{2\ln m}$, the same MDP-speed caveat as Theorem 7.33; i.i.d.\ and
finite-state Markov admit the full $\ln m=o(K)$.)}
\emph{\textbf{(b) (Refined Gumbel law.)} If moreover $X$ is i.i.d.\ or a finite-state
irreducible aperiodic Markov source (the sharp Cram\'er--Petrov regime of Theorem 7.33) and
$\ln m=o(K^{1/3})$ --- equivalently $K\gg(\ln m)^3$, i.e.\ $(\ln m)^{3/2}=o(\sqrt K)$ ---
then, with $a_m=\sqrt{2\ln m}$ and $b_m=a_m-\dfrac{\ln\ln m+\ln4\pi}{2a_m}$,}
\[
  a_m\!\left(\frac{M_m-Kh}{\sqrt{KV}}-b_m\right)\;\Longrightarrow\;\Lambda,\quad
  \Pr(\Lambda\le y)=e^{-e^{-y}}\ (\text{Gumbel}),
\]
\emph{and $\ \mathbb E\,M_m=Kh+\sqrt{KV}\bigl(b_m+\tfrac{\gamma}{a_m}\bigr)
+o\!\bigl(\tfrac{\sqrt{KV}}{a_m}\bigr)$, $\gamma$ the Euler--Mascheroni constant. A
mismatched coder $Q$ with $Q>0$ on the source support (so the per-symbol log-loss
$-\log_2 Q$ is finite, i.e.\ $P\ll Q$) replaces $(h,V)$ by the cross-entropy rate and
cross-information varentropy $(h_\times,v^2)$ of Theorem 7.29 (provided $v^2>0$ and the same
tail/dependence hypotheses hold for the cross-information process).}

\emph{\textbf{Proof.}} \emph{(1) Reduction.} $M_m$ is the maximum of the $m$ block
codelengths $L_1,\dots,L_m$, each a sum of $K$ per-symbol information densities; the
per-block CLT (Theorem 7.28) gives $(L_i-Kh)/\sqrt{KV}\Rightarrow\mathcal N(0,1)$.
\emph{(2) Tail control --- the two scale conditions.} For i.i.d.\ / finite-state Markov
$L_i$ (bounded densities, so Cram\'er's condition holds) the sharp Cram\'er--Petrov
expansion (Petrov, \emph{Sums of Independent Random Variables}, 1975, Ch.~VIII) gives
\[
  \Pr\!\bigl(L_i-Kh>x\sqrt{KV}\bigr)=Q(x)\,\exp\!\Bigl(\tfrac{x^3}{\sqrt K}\,
  \lambda\bigl(\tfrac{x}{\sqrt K}\bigr)\Bigr)\,(1+o(1))\quad\text{uniformly for }x=o(\sqrt K),
\]
$Q$ the standard-normal tail, $\lambda$ the Cram\'er series, $\lambda(0)=\mu_3/(6\sigma^3)$
($\mu_3,\sigma^2$ the per-symbol surprisal's third central moment and variance; for a Markov
functional the long-run analogues). At the extremal level $x\asymp\sqrt{2\ln m}$: for the
\emph{leading order} (a) the exponent is $-x^2/2+(x^3/\sqrt K)\lambda+o(1)$ with the
correction $o(x^2)\iff x=o(\sqrt K)\iff\ln m=o(K)$, so the level solving $m\Pr(\cdot)\asymp1$
is $x=\sqrt{2\ln m}(1+o(1))$ and $M_m=Kh+\sqrt{2KV\ln m}(1+o_P(1))$. For the \emph{refined
centering} (b) the multiplicative correction must satisfy $\exp((x^3/\sqrt K)\lambda)\to1$ at
$x\asymp\sqrt{2\ln m}$, i.e.\ $x^3/\sqrt K\to0\iff(\ln m)^{3/2}=o(\sqrt K)\iff\ln m=o(K^{1/3})$.
This is necessary, not merely sufficient: for $\ln m=o(K)$ but $\gg K^{1/3}$ the centre
$b_m$ acquires a non-vanishing skewness shift and the displayed Gumbel form fails (e.g.\
i.i.d.\ $p=(.85,.08,.05,.02)$ with $m=\lfloor e^{\sqrt K}\rfloor$: then $\ln m=o(K)$ yet
$x^3/\sqrt K\asymp K^{1/4}\to\infty$, so $m\Pr(\cdot)\to\infty$, not $e^{-y}$). If $\mu_3=0$
the next Cram\'er term is fourth-order ($x^4/K$), a sufficient relaxation to $\ln m=o(\sqrt K)$.
\emph{(3) I.i.d.\ Gumbel.} For i.i.d.\ $X$ the disjoint blocks are independent; with
$u_m(y)=Kh+\sqrt{KV}(b_m+y/a_m)$, step (2) and Mills' ratio $Q(x)\sim\phi(x)/x$ give
$m\Pr(L_1>u_m(y))\to e^{-y}$ ($b_m$ chosen so $mQ(b_m)\to1$), and independence yields
$\Pr(M_m\le u_m(y))=(1-\Pr(L_1>u_m(y)))^m\to\exp(-e^{-y})$, i.e.\ $a_m(Z_m-b_m)\Rightarrow
\Lambda$ with $Z_m=(M_m-Kh)/\sqrt{KV}$. \emph{(4) Finite-state Markov dependence.} The block
sequence $\{L_i\}$ is a block-level functional of the chain; finite-state irreducibility and
aperiodicity give \emph{geometric} $\alpha$-mixing, under which the array satisfies
Leadbetter's long-range condition $D(u_m)$ (from the geometric mixing) and the
anti-clustering condition $D'(u_m)$. For this \emph{triangular} block array ($K=K(m)$),
$D'(u_m)$ requires $m\sum_{j=2}^{\lfloor m/k\rfloor}\Pr(L_1>u_m,L_j>u_m)\to0$: non-adjacent
blocks ($j\ge3$) factorise up to the geometric $\alpha(\cdot)$ correction, and \emph{adjacent}
blocks, sharing only their $O(1)$-mixing boundary out of $K$ symbols, obey the bivariate
moderate-tail bound $\Pr(L_1>u_m,L_2>u_m)\le\Pr(L_1>u_m)\Pr(L_2>u_m)(1+o(1))$, so
$m\,\Pr(L_1>u_m,L_2>u_m)=O(1/m)\to0$; hence the extremal index is
$\theta=1$ and the Gumbel limit of (3) holds (Leadbetter--Lindgren--Rootz\'en,
\emph{Extremes and Related Properties}, 1983, Thm.~3.4.1/3.5.2). \emph{(Berman's covariance
condition is for stationary \emph{Gaussian} sequences; the $L_i$ are only asymptotically
Gaussian, so $D/D'$, not Berman, is the correct device.)} \emph{(5) Lower bound.} Use a
big-block/small-block split: a $q_K$-spaced sub-family of alternate blocks with
$q_K\to\infty$, $q_K=o(m)$, and $\ln q_K=o(\ln m)$ (e.g.\ $q_K=e^{\sqrt{\ln m}}$); geometric
mixing bounds its total-variation distance from independence by
$m\,\alpha(q_KK)=O(m\rho^{q_KK})\to0$ (since $\ln m=o(K)$). The sub-family's $m/q_K$
near-independent blocks have maximum $\sqrt{2KV\ln(m/q_K)}=\sqrt{2KV\ln m}(1-o(1))$ (using
$\ln q_K=o(\ln m)$), so $M_m$ is no smaller; with (3)--(4),
$M_m=Kh+\sqrt{2KV\ln m}(1+o_P(1))$. \hfill$\blacksquare$

\emph{Significance --- the extreme-value entry of the deviation hierarchy, and the
read-buffer law.} The §7 hierarchy until now described only the \emph{total} repair length
$\sum_iL_i$: concentration (7.14/7.25), CLT/dispersion (7.28/7.29), moderate deviations
(7.33), large deviations (7.31/7.32). Theorem 7.35 governs the per-block \emph{maximum}
$\max_iL_i$ --- the quantity that decides random access: the read buffer that must be
provisioned so that \emph{any} single byte-granular query can be served. A buffer of the
mean block size $Kh$ overflows on the worst of $m=N/K$ blocks with probability $\to1$; the
tight provisioning is $Kh+\sqrt{2KV\ln(N/K)}$ (part (a), for i.i.d./finite-state Markov, and
exp-mixing $\alpha$ on the restricted range above), a $\sqrt{\ln(N/K)}$ extreme-value
inflation over the mean driven by the \emph{same}
varentropy $V$ that sets the archive dispersion (Theorem 7.29); part (b) pins the residual
$O(1/a_m)$ fluctuation to a Gumbel law in the sharp-tail regime $K\gg(\ln m)^3$. Thus one
source functional $V$ controls both the central fluctuation of the total length and the
worst-case query latency. At the balanced $K=\Theta(\sqrt N)$ (where 7.29's random-access
overhead and dispersion are both $\Theta(\sqrt N)$) the mean block is $Kh=\Theta(\sqrt N)$
and the margin $\sqrt{2KV\ln(N/K)}=\Theta(N^{1/4}\sqrt{\log N})$ is of lower order --- so
the worst block is $(1+o(1))Kh$, the margin being the precise second-order headroom (and
$K=\Theta(\sqrt N)$ comfortably meets $K\gg(\ln m)^3$, so the Gumbel refinement applies).
The phenomenon needs $V>0$: a uniform-i.i.d.\ source has $V=0$, every block has identical
length, the maximum equals the mean, and there is no provisioning headroom --- the EVT
spread is a strict-varentropy effect. The nontrivial analytic ingredient is the uniform
moderate-deviation tail at $x\asymp\sqrt{2\ln m}$ (step 2), which also fixes the exact
applicability boundary $\ln m=o(K^{1/3})$. \emph{A third resource axis.} The read buffer is
distinct from the two resources in the space--time product of Theorem 7.27 --- per-query
\emph{compute} $C(K)=\Theta(K\,\mathrm{cost}(M))$ and seek-index \emph{storage} $S(K)$ ---
measuring the \emph{bits read} per query, $\Theta(Kh)$, rather than predictor evaluations or
index bits. The four quantities pull against each other in $K$: the index storage $S(K)$ and
the sync overhead $(N/K)\delta_\infty$ (Theorem 7.29) fall with $K$, while the per-query
compute (Theorem 7.27) and the read buffer (this theorem) rise. At the balanced
$K=\Theta(\sqrt N)$ the sync overhead, the per-query compute, and the per-query read buffer
are \emph{simultaneously} $\Theta(\sqrt N)$ --- one block size balancing three distinct
resources. Like the compute floor, the $\Theta(Kh)$ read buffer is a \emph{sub-block-class}
property: the Ferragina--Venturini escape (Theorem 7.27$'$) reads only $O(\mathrm{polylog}\,N)$
bits per query, sidestepping the read-buffer axis just as it does the compute floor. Verified in
\texttt{scripts/verify/thm\_7\_35\_worst\_case\_query\_evt.py} (V1 the constants $Kh,KV$;
V2 the refined Gumbel centring $b_m+\gamma/a_m$ in the sharp regime $K\gg(\ln m)^3$, to
$<\!1\%$; V3 the leading-order ratio $\to1$ from below; V4 a finite-state Markov source with
the fundamental-matrix varentropy cross-checked against the empirical block variance; V5 the
\emph{sharp} scale boundary --- the refined-Gumbel error tracks $x^3/\sqrt K$ and is small
iff $K\gg(\ln m)^3$, growing from $\sim\!2\%$ at $K=3.6(\ln m)^3$ to $\sim\!23\%$ at
$K=0.03(\ln m)^3$, confirming the $o(K^{1/3})$ condition and the i.i.d.\ counterexample of
step 2; V6 the matching lower bound from a separated sub-family).

\textbf{Theorem 7.36 (Functional CLT for the RNR repair process; the streaming-buffer
high-water-mark).} \emph{Let $X$ be a finite-state irreducible aperiodic Markov source (so
the realised conditional surprisals $\ell_i=-\log_2 P(X_i\mid X_{i-1})$ are bounded,
$|\ell_i|\le-\log_2 p_{\min}$, and form a geometrically $\alpha$-mixing finite-window
functional), with entropy rate $h$ and varentropy rate
$V=\lim_N\frac1N\mathrm{Var}(\sum_{i\le N}\ell_i)\in(0,\infty)$ (the Theorem 7.28/7.29
constant). Let $D_t=\sum_{i\le t}(\ell_i-h)$ be the centred repair process.}

\emph{\textbf{(a) (Functional CLT, ideal stream.)} For the ideal repair process (full,
order-$1$ context; no sub-block sync resets) the rescaled path
$W_N(t):=D_{\lfloor Nt\rfloor}/\sqrt{NV}$ converges in distribution on the Skorokhod space
$(D[0,1],J_1)$ to standard Brownian motion, $W_N(\cdot)\Rightarrow W(\cdot)$.}

\emph{\textbf{(b) (Buffer high-water-marks.)} By the continuous mapping theorem and uniform
integrability, the \emph{net cumulative surplus} (the buffer of an encoder pre-credited
against its rate reservation, allowed to run a temporary deficit) obeys}
\[
  \frac{\max_{t\le N}D_t}{\sqrt{NV}}\;\Longrightarrow\;\sup_{0\le s\le1}W(s)\;\overset{d}{=}\;|Z|,
  \ Z\sim\mathcal N(0,1)\ (\text{half-normal}),\qquad
  \mathbb E\Bigl[\tfrac{\max_t D_t}{\sqrt{NV}}\Bigr]\to\sqrt{\tfrac2\pi},
\]
\emph{$\Pr(\max_{t\le N}D_t>a\sqrt{NV})\to2(1-\Phi(a))$. The \emph{physical} empty-start
queue (Lindley occupancy $Q_t=\max(0,Q_{t-1}+\ell_t-h)$, no negative occupancy) has
high-water-mark equal to the maximal draw-up $\max_{j\le t\le N}(D_t-D_j)$, which by L\'evy's
theorem satisfies}
\[
  \frac{\max_{t\le N}Q_t}{\sqrt{NV}}\;\Longrightarrow\;\sup_{0\le s\le1}|W(s)|\quad
  (\text{L\'evy theta-function law, mean }\sqrt{\pi/2}\approx1.2533),
\]
\emph{not half-normal. At $t=1$, $W_N(1)\Rightarrow\mathcal N(0,1)$ recovers the Theorem 7.28
marginal CLT.}

\emph{\textbf{(c) (Sub-block implementation: a drifted buffer.)} The Theorem 7.5 sub-block
coder, draining at the matched ideal rate $h$, resets context at the $m=N/K$ sync points,
adding the deterministic cold-context overhead $(N/K)\delta_\infty$ (Theorem 7.29): a
per-symbol fill drift $\delta_K/K\to\delta_\infty/K$. At $K=c\sqrt N$ the rescaled net
surplus converges to Brownian motion \emph{with drift},}
\[
  \frac{D^{\mathrm{sb}}_{\lfloor Nt\rfloor}}{\sqrt{NV}}\;\Longrightarrow\;W(t)+\beta t,
  \qquad \beta=\frac{\delta_\infty}{c\sqrt V}\ \ (\text{the Theorem 7.33 shift}),
\]
\emph{so the high-water-mark is $\sqrt{NV}\,\sup_{[0,1]}(W+\beta s)$ (net surplus) or
$\sqrt{NV}\,\sup_{0\le s\le t\le1}\bigl[(W_t+\beta t)-(W_s+\beta s)\bigr]$ (physical queue
draw-up), reducing to (b) at $\beta\to0$ (large $K$). The drain rate $h$ is critical for the ideal stream; the sub-block
mean rate is $h+\delta_K/K$, so critical draining is at $h+\delta_K/K$ ($\delta_K\to
\delta_\infty$). A mismatched coder replaces $(h,V)$ by $(h_\times,v^2)$ of Theorem 7.29 and
$\delta_\infty$ by the corresponding cross cold-start excess (under the usual
support/nondegeneracy conditions).}

\emph{\textbf{Proof.}} \emph{(a)} For a finite-state irreducible aperiodic chain, $\xi_i:=
\ell_i-h=-\log_2 P(X_i\mid X_{i-1})-h$ is a bounded ($|\xi_i|\le-\log_2 p_{\min}+h$) function
of the consecutive pair $(X_{i-1},X_i)$, hence a stationary, mean-zero, geometrically
$\alpha$-mixing sequence with long-run variance $\sigma^2=\lim_N\mathrm{Var}(\sum_{i\le N}
\xi_i)/N=V>0$ (the fundamental-matrix variance of Theorem 7.28). The functional CLT for
additive functionals of finite irreducible aperiodic Markov chains (the associated
martingale/regeneration FCLT; equivalently Billingsley 1968, \S20, for the geometrically
mixing case) gives $D_{\lfloor N\cdot\rfloor}/\sqrt{N\sigma^2}\Rightarrow W$ on $(D[0,1],J_1)$;
with $\sigma^2=V$ this is (a). \emph{(b)} $\sup$ is continuous on $(D[0,1],J_1)$ and $W$ has
continuous paths a.s., so the continuous mapping theorem gives $\max_{t\le N}D_t/\sqrt{NV}
\Rightarrow\sup_sW(s)$ (with $D_0=0$); the reflection principle gives $\Pr(\sup_{[0,1]}W\ge a)
=2(1-\Phi(a))$, i.e.\ $\sup W\overset{d}{=}|Z|$, mean $\sqrt{2/\pi}$. Convergence of the mean
uses uniform integrability: the bounded-increment geometric-mixing maximal moment inequality
(M\'oricz 1976; Serfling 1970) gives $\sup_N\mathbb E[(\max_{t\le N}D_t/\sqrt{NV})^2]<\infty$.
For the physical queue, the Lindley recursion unrolls to $\max_{t\le N}Q_t=\max_{0\le j\le
t\le N}(D_t-D_j)$, the maximal draw-up; the draw-up process $D_t-\min_{j\le t}D_j$ converges
(continuous mapping) to $W_t-\min_{s\le t}W_s\overset{d}{=}|W_t|$ (L\'evy's theorem), whose
running maximum is $\sup_{[0,1]}|W|$ (theta-function law, mean $\sqrt{\pi/2}$). \emph{(c)}
Writing $D^{\mathrm{sb}}_t=D_t+\sum_{b\le t/K}\delta_K+o(\cdot)$ (the per-sync cold-start
excess $\delta_K\to\delta_\infty$, concentrated in the first post-reset symbol for order-$1$;
Lemma 5.1d), the deterministic part rescales to $\beta t$ at $K=c\sqrt N$, and its $O_P(\sqrt
m)=O_P(N^{1/4})$ fluctuation is negligible after $\sqrt N$ scaling; adding to (a) and applying
continuous mapping gives the drifted limit and its supremum. \hfill$\blacksquare$

\emph{Scope.} The clean statement is for finite-state irreducible aperiodic Markov sources,
where $\ell_i$ is a bounded finite-window functional with geometric $\alpha$-mixing (every
FCLT hypothesis met). For a general stationary exp-$\alpha$-mixing source the
\emph{infinite-past} surprisal $-\log_2 P(X_i\mid X_{<i})$ neither is bounded nor inherits
mixing automatically; one works instead with the windowed predictor $\ell_i^{(W_N)}=-\log_2
P(X_i\mid X_{i-W_N:i-1})$ of Theorem 7.29 --- a finite $W_N=\Theta(\log N)$-window functional
(so $\alpha$-mixing is inherited at the shifted lag) bounded by the $\eta$-floor of
Theorem 7.5 ($|\ell_i^{(W_N)}|\le\log_2(1/\eta)$) --- and passes to the limit via the
near-epoch-dependence/continuity of $\ell_i$ on the past (the exp-mixing windowed convergence
of Theorem 7.28); the FCLT then holds under the Merlev\`ede--Peligrad (2000, \emph{Ann.\
Probab.}\ 28:1336--1352) / Doukhan--Massart--Rio (1994, \emph{Ann.\ IHP}\ 30:63--82)
strong-mixing invariance principle. We state the Markov case to keep the hypotheses verified.

\emph{Significance --- one varentropy, three buffers (all carrying the $\beta$ penalty).}
Theorem 7.36 is the \emph{process-level} strengthening of the marginal CLT (Theorem 7.28),
which it recovers at $t=1$, and the encoder-side companion of the decoder-side Theorem 7.35.
The three operational buffer/latency quantities of RNR are governed by the \emph{single}
source functional $V$ (the varentropy): the archive dispersion $\sqrt{NV}\,Q^{-1}(\varepsilon)$
(Theorem 7.29, total size), the worst-case random-access \emph{read} buffer
$Kh+\sqrt{2KV\ln(N/K)}$ (Theorem 7.35, decoder side), and the streaming-encoder \emph{write}
buffer high-water-mark (this theorem, $\sqrt{NV}\,|Z|$ net-surplus or $\sqrt{NV}\,\sup|W|$
physical queue) --- the last carrying, like the archive dispersion (7.29) and the
moderate-deviation overflow (7.33), the \emph{same} cold-context shift $\beta=\delta_\infty/
(c\sqrt V)$ at the balanced $K=c\sqrt N$. Operationally, a streaming RNR encoder (Theorem 7.19)
draining at $h$ over a horizon-$N$ block needs a physical buffer $b\sqrt{NV}$ to bound the
overflow probability, with $b$ set by the $\sup|W|$ (theta-function) quantile (net-surplus
$\sup W$ if pre-credited); the sub-block implementation adds the $\beta$ drift, critical drain
$h+\delta_K/K$. The content is a genuine process limit (Donsker), not a marginal one: all path
functionals (running maximum, draw-up, range, first-passage, occupation) inherit Brownian
limits. Verified in \texttt{scripts/verify/thm\_7\_36\_functional\_clt\_streaming\_buffer.py}
(V1 the finite-dim Brownian distributions $\mathrm{Var}\,W_N(t)\to t$, $\mathrm{Cov}\to
\min(s,t)$; V2 the net-surplus high-water-mark $\to$ half-normal, mean $\sqrt{2/\pi}$, tail
$2(1-\Phi(a))$; V3 the two-sided $\to\sup|W|$; V4 a finite-state Markov source with the
long-run varentropy; V5 distributional convergence by quantile match; V6 the same-$V$
unification; V7 the sub-block drift $\to\sqrt{NV}\sup(W+\beta s)$, $\beta=\delta_\infty/
(c\sqrt V)$, not half-normal; V8 the physical Lindley queue high-water-mark $\to\sup|W|$,
mean $\sqrt{\pi/2}$, distinct from the net-surplus half-normal).

\textbf{10.7 Random access}

Theorems 10.1 and 10.2 concern bit-exact decoder reproduction. A
separate, operationally important question is whether a decoder can
extract a small portion of $X$ without decoding all of it --- the
seek-and-read functionality of bgzip and zstd seekable mode.
Existing neural compressors lack this capability for two distinct
reasons: autoregressive byte/pixel models (NNCP, PixelCNN++,
PixelSNAIL) require sequential evaluation of all $p$ previous bytes
before byte $p$ is available (cost $\Theta(p)$ per random read);
invertible-flow and bits-back block codecs (IDF++, Bit-Swap,
HiLLoC) operate on whole blocks of statically determined size and
state, so even a one-byte read forces decoding of the enclosing
block and any state on which it depends. In either case the
effective minimum random-read granularity is large and growing with
archive size (the first case linearly in $p$; the second
block-by-block with no sub-block seek). This is a deployment
blocker for archive-style use cases where partial reads are common.
RNR archives that use per-position sequential coding modes (e.g.,
Type-I, the per-byte sequential Type-III-B strategy, or Type-III-A
composition over per-position sequential per-tile modes) support
byte-granular random access in per-query time that is
polylogarithmic in $|X|$ (or independent of $|X|$ in the
uniform-spacing direct-lookup case when $\log L(R)$ is also
held fixed), provided the predictor has
bounded context, in a parameter regime where $K$, $k$,
$\mathrm{cost}(M)$, and $\log L(R)$ are held fixed; the resident
seek index is $\Theta(|X|/K)$. The factor $\mathrm{cost}(M)$ is one
per-position neural-predictor evaluation --- non-universal
(architecture-dependent, Lemma 5.1b) and the dominant term of the
per-query cost --- so the polylog claim bounds only the dependence on
$|X|$ with $\mathrm{cost}(M)$ held fixed. Whole-block Type-II partition-rank
modes and Type-III-B partition modes give only sub-block-granular
random access at a separate $O(K \cdot \mathrm{unrank}(K))$ per
requested sub-block (Remark 10.3a). We make the per-position
sequential case precise.

\textbf{Definition 10.2 ($W$-bounded predictor).} A predictor $M$ is
$W$-bounded if its output $Y_i = M(C_i)$ at position $i$ depends only
on the previous $W$ bytes of public context:
$Y_i = M(X_{i-W}, \ldots, X_{i-1})$ when $C$ is built incrementally
from decoded subblocks. The bound $W$ is a structural property of
$M$'s architecture, not a tunable parameter at inference time.

Transformers with fixed context window $W$ are $W$-bounded, as are
$n$-gram and finite-state predictors of order $W$ and any feedforward
architecture with bounded receptive field. RNNs, LSTMs, GRUs, and
state-space models (Mamba, S4) are not $W$-bounded in general; their
state can depend on arbitrarily long history. Theorem 10.4 below
covers the non-$W$-bounded case via state snapshots.

\textbf{Definition 10.3 (Sync point).} A sync point at position $P$ in
the source is a position where (a)~the predictor's context state is
reset to a known initial value $Q_{\text{init}}$ independent of bytes
$[0, P)$, and (b)~the entropy coder's internal state (arithmetic-coder
interval, ANS register, or analogous) is flushed at the encoder and
reset to a known initial value at the decoder, so that the next bit
in the repair stream begins a fresh entropy-coder block. The encoder
inserts sync points at fixed intervals $K$ so that the byte offset
in $X$ of the $j$-th sync point is the deterministic value $j K$
(no byte-offset storage required), and includes a seek index in the
archive mapping each sync point identifier to its bit offset in $R$
only. The seek index has size at most
$\lceil |X|/K \rceil \cdot \lceil \log_2 L(R) \rceil$ bits.
Definition 10.3 admits two variants: \emph{uniform spacing} (sync
points at byte positions $j K$, byte offsets implicit, used by
default) and \emph{non-uniform spacing} (sync points at
encoder-chosen byte positions $\Delta_0 = 0 < \Delta_1 < \cdots <
\Delta_{m-1} < |X|$ subject to four invariants: (i) the first
sync point $\Delta_0 = 0$ (so every query position $p \in [0, |X|)$
has some prior sync point); (ii) maximum gap between consecutive
sync points $\Delta_{j+1} - \Delta_j \leq K$ for $j = 0, \dots,
m-2$; (iii) tail gap $|X| - \Delta_{m-1} \leq K$ (so every query
position is within $K$ bytes of its prior sync point); (iv)
minimum gap $\Delta_{j+1} - \Delta_j \geq K/2$ for $j = 0, \dots,
m-2$. Byte offsets are stored explicitly per entry; this variant
is used when token-edge alignment or other encoder constraints
require it. The max-gap invariant ensures that the latest sync
point before any query position $p$ is within $K$ bytes of $p$, so
the $K + k$ decode-window cost of Theorem 10.3 is preserved; the
min-gap invariant bounds the entry count to at most
$\lceil 2 |X|/K \rceil$ (factor of 2 over uniform spacing). Under
non-uniform spacing each entry additionally stores a
$\lceil \log_2 |X| \rceil$-bit source position; the
seek-index size is at most
$\lceil 2 |X|/K \rceil \cdot (\lceil \log_2 L(R) \rceil + \lceil
\log_2 |X| \rceil)$ bits. Both variants are referenced in
subsequent theorems and remarks.
The entropy-coder flush adds $O(1)$ bits of termination per sync
point and contributes $O(|X|/K)$ bits of additional repair-stream
overhead in total; the predictor and entropy-coder reset together
make each sync-point bit offset a self-contained restart point for
decoding.

Each sync point forfeits some compression efficiency because the
predictor restarts cold and cannot exploit the prior $W$ bytes of
context. \emph{Bounded-cost assumption.} We assume the entropy
coder's per-position symbol distribution satisfies a probability
floor $\eta > 0$ (the floor is on whichever symbol the coder
actually emits per position: for per-position sequential Type-I
this is the class+payload symbol distribution; for per-byte
sequential Type-III-B this is the high-nibble distribution and
then the low-nibble conditional distribution per byte; for direct
byte arithmetic coding it is $\Pr_M(b \mid c)$). In practice this
is enforced by mixing the model's distribution with a small
uniform escape at the coder level, $\eta \geq 2^{-12}$ being
standard. Under this floor, the per-position cold-context
arithmetic-code length is bounded above by $S \cdot \log_2(1/\eta)$,
where $S$ is the number of floored symbols the coder emits per
source byte: $S = 1$ for Type-I (one class+payload symbol per
byte) and direct byte arithmetic, $S = 2$ for per-byte sequential
Type-III-B (high nibble + low nibble), and more generally
$S = $ (symbols-emitted-per-source-byte) for the chosen mode. The
per-sync-point excess cost over warm-context encoding of the same
$W$ source bytes is therefore bounded by
\[
\Delta_{\mathrm{sync}} \leq \sum_{i=P}^{P+W-1}
S \cdot \bigl[\log_2(1/\eta) - (-\log_2 \Pr_M(\text{symbols}_i \mid C_i^{\mathrm{warm}}))/S\bigr]
\leq S \cdot W \cdot \log_2(1/\eta),
\]
where $C_i^{\mathrm{warm}}$ denotes the warm $W$-byte context the
encoder would have used in a sync-point-free encoding, and the
inner term is averaged over the $S$ symbols emitted at position
$i$. Under the probability floor, both cold and warm per-position
symbol costs lie in $[0, \log_2(1/\eta)]$, so both per-byte costs
lie in $[0, S \cdot \log_2(1/\eta)]$, and each summand of
$\Delta_{\mathrm{sync}}$ lies in $[-S \cdot \log_2(1/\eta), S
\cdot \log_2(1/\eta)]$. (A pointwise negative summand is possible
when the cold/initial context happens to predict $X_i$ better
than the warm $W$-byte context --- empirically rare but
permitted by the floor alone; the floor caps both costs above
without enforcing cold $\geq$ warm.) The trivial upper bound
$S \cdot W \cdot \log_2(1/\eta)$ is therefore an upper bound on
$\Delta_{\mathrm{sync}}$ from above, even if some summands are
negative. It is loose and does not depend on the warm rate. Sharper upper
bounds require additional assumptions: e.g., a uniform per-byte
warm-context \emph{lower} bound $r^{\mathrm{warm}}_{\min}$ on the
per-byte warm cost valid at all positions, yielding
$\Delta_{\mathrm{sync}} \leq W \cdot (S \cdot \log_2(1/\eta) - r^{\mathrm{warm}}_{\min})$.
(An upper bound on warm cost would give a lower, not an upper,
bound on $\Delta_{\mathrm{sync}}$, so the warm-rate assumption must
be a lower bound to yield a useful overhead ceiling.) Without the
probability-floor assumption (e.g., a predictor that can place
arbitrarily small mass on the true byte after reset), the
per-sync-point cost is unbounded. The total compression cost of
sync points under the floor assumption is at most
$m \cdot S \cdot W \cdot \log_2(1/\eta)$ bits, where $m$ is the
total number of sync points: $m = \lceil |X|/K \rceil$ for uniform
spacing and $m \leq \lceil 2|X|/K \rceil$ for non-uniform spacing
(Definition 10.3); using $m$ rather than $|X|/K$ keeps the bound
correct when $|X|$ is not a multiple of $K$ (e.g., when $|X| < K$
the archive still has one sync point and $m = 1$). With $S$ as
defined above ($S = 1$ for Type-I and direct byte arithmetic,
$S = 2$ for per-byte sequential Type-III-B), the bound is loose
by $S \cdot W \cdot \log_2(1/\eta)$ at most.
Concrete-number estimates in the abstract and §11 (e.g., the
``approximately one percent'' figure) are heuristic conjectures
based on empirical observations of cold-vs-warm rates for small
distilled transformers; they are not theorems and depend on the
sharper assumption above.

\textbf{Corollary 10.2a (Assumption-graded per-sync overhead
ladder).} \emph{The three per-sync-point cold-context ceilings
established above and in Lemmas 5.1a, 5.1c stand in a strict
assumption-vs-sharpness order: each rung below buys a tighter
ceiling at the cost of a stronger hypothesis on the source/predictor,
and every entry is a restatement of an already-proved bound (no new
result is asserted). The fully general tight per-sync overhead --- a
ceiling matching the realized cold-context excess for an arbitrary
source under the floor assumption alone --- remains open; rungs~2--3
are tight (or near-tight) only under their named hypotheses.}

\begingroup\footnotesize\sloppy
\setlength{\tabcolsep}{4pt}
\noindent\begin{tabular}{@{}p{2.0cm}p{3.7cm}p{3.6cm}p{2.6cm}@{}}
\hline
\textbf{rung (assumption strength)} & \textbf{hypothesis} &
\textbf{per-sync ceiling $\Delta_{\mathrm{sync}}$ / total over $m$ syncs} &
\textbf{source}\\
\hline
1. loose (floor only)
& predictor probability floor $\eta>0$ on the emitted symbol
distribution (nothing on the warm rate or source law)
& $\Delta_{\mathrm{sync}}\le S\cdot W\cdot\log_2(1/\eta)$; total
$\le m\cdot S\cdot W\cdot\log_2(1/\eta)$ ($S$ floored symbols/byte,
$m$ sync count, Def.\ 10.3)
& this §10.7 (eqns above); Lemma 5.1a\\
2. warm-rate refined
& floor $\eta>0$ \emph{and} a uniform per-byte warm-context
\emph{lower} bound $r^{\mathrm{warm}}_{\min}$ valid at all positions
& $\Delta_{\mathrm{sync}}\le W\cdot(S\cdot\log_2(1/\eta)-r^{\mathrm{warm}}_{\min})$;
total $\le m\cdot W\cdot(S\cdot\log_2(1/\eta)-r^{\mathrm{warm}}_{\min})$
& this §10.7 (warm-rate refinement; a lower bound on warm cost is
required to yield an upper bound on overhead)\\
3. exact-Markov ($\delta_{W'}$)
& strictly stationary order-$W'$ Markov source, $W$-bounded predictor
with $W\ge W'$ outputting the true conditional, $K\ge W'$, full
sub-blocks $N=mK$
& $\Delta_{\mathrm{sync}}=\delta_{W'}=H(X^{W'})-W'h(X)\ge 0$
(\emph{exact}, not a bound); total excess over $H(D_N)$ is
$(m-1)\cdot\delta_{W'}$
& Lemma 5.1c (eqn 5.9)\\
\hline
\end{tabular}
\par\endgroup

\emph{Reading guide.} Rung~1 holds whenever the coder runs under the
floor and is the worst case used in Lemma 5.1a and Theorem 10.3; it
does not depend on the warm rate and is loose by construction. Rung~2
subtracts a guaranteed warm-context rate floor $r^{\mathrm{warm}}_{\min}$
from the same $S\log_2(1/\eta)$ per-byte cap --- it is sharper exactly
when the predictor is provably never worse than $r^{\mathrm{warm}}_{\min}$
bits/byte in warm context, and (as flagged above) the warm-rate
assumption must be a \emph{lower} bound on warm cost to yield an
\emph{upper} bound on overhead. Rung~3 replaces the bound by the
\emph{exact} finite source-specific constant $\delta_{W'}$ for
finite-order Markov sources (independent of $W$ for $W\ge W'$ and of
$K$ for $K\ge W'$); for natural-data Markov sources $\delta_{W'}\ll
S\cdot W\cdot\log_2(1/\eta)$ (Lemma 5.1c gives a $\sim340\times$ gap
for the order-1 English bigram instance), which is why the realized
overhead the abstract conjectures ($\approx$ one percent) lives far
below rung~1 and near rung~3 when the source is approximately Markov.
No rung discharges the general case: for an arbitrary source under the
floor alone we have only rung~1, and a tight ceiling matching the
realized excess in full generality is not proved here. A synthetic
heuristic check that the realized overhead-as-fraction respects rung~1
and lands in the conjectured band on Markov sources (and can fail on a
non-Markov stress source) is
\texttt{scripts/verify/bug\_003\_syncpoint\_overhead\_fraction.py}.

\textbf{Theorem 10.3 (Random access, output-sensitive, for
per-position sequential coding).} \emph{Let $X$ be compressed under
RNR using a \emph{per-position sequential} coding mode: Type-I, the
per-byte sequential Type-III-B strategy of §6.3 (strategy iii),
Type-III-A composition over per-position sequential per-tile modes
where additionally sync points are placed at tile boundaries (so
no sync point lands inside a tile's framed repair stream), or any
other coding scheme whose repair stream is read one position at a
time. The theorem explicitly excludes modes whose repair-stream
encoding spans an entire block as a single combinatorial object,
such as Type-II's multinomial/partition enumerative rank coding
without sub-block sync (Remark 10.3a) and Type-III-B's direct or
hierarchical partition strategies (i)/(ii) of §6.3. For Type-III-A,
the encoder must therefore choose tile and sync-point sizes such
that $K$ is an integer multiple of (or aligned to) the tile size,
ensuring sync-point boundaries coincide with tile boundaries. Let $M$ be a $W$-bounded predictor satisfying Theorem 10.1
and let sync points be placed
at uniformly spaced source positions $j K$ for $j = 0, 1, \dots$
(Definition 10.3). Assume the entropy coder operates under a
predictor probability floor $\eta > 0$ (Section 10.7) on its
per-position emitted symbol distribution, so each source byte
consumes at most $S \cdot \log_2(1/\eta) = O(1)$ repair-stream
bits, where $S$ is the number of floored entropy-coder symbols per
source byte for the chosen mode ($S = 1$ for Type-I and direct byte
arithmetic, $S = 2$ for per-byte sequential Type-III-B). Let $|X|$ denote the archive's uncompressed source length.
Random access to $k$ consecutive bytes starting at any position
$p \in [0, |X|)$ with $p + k \leq |X|$ (or, equivalently, requesting
$k' = \min(k, |X| - p)$ bytes if $p + k$ would exceed $|X|$, as a
short read) requires: (i)~one direct seek-index lookup at
entry $\lfloor p/K \rfloor$, reading $O(\log L(R))$ bits to obtain
the repair-stream bit offset of the sync point at byte position
$P = K \lfloor p/K \rfloor$; (ii)~decoding of at most $K + k$
positions of $R$, of which at most $K$ positions are warmup
(discarded output); (iii)~$K + k$ predictor evaluations,
entropy-coder updates, and $O((K+k) \cdot \log_2(1/\eta))$
repair-stream bits read. Total work is
$O(\log L(R) + (K + k) \cdot \mathrm{cost}(M))$. (For non-uniform
sync-point placement, each seek-index entry stores both a
$\lceil \log_2 |X| \rceil$-bit source position and a
$\lceil \log_2 L(R) \rceil$-bit repair-stream offset, so the
binary-search lookup costs
$O(\log m \cdot (\log |X| + \log L(R)))$ bits, where $m$ is the
number of sync-point entries; for $|X| \geq K$ this gives
$O(\log(|X|/K) \cdot (\log |X| + \log L(R)))$ bits, and for
$|X| < K$ there is a single entry and the lookup is $O(1)$. The
rest of the analysis is unchanged.) \emph{This cost is polylogarithmic in archive size
$|X|$ only when $K$, $k$, and $\mathrm{cost}(M)$ are held fixed or
grow at most polylogarithmically in $|X|$ and
$L(R) = \mathrm{poly}(|X|)$; for
$K = \Theta(|X|)$ or $k = \Theta(|X|)$ the cost is
$\Theta(|X| \cdot \mathrm{cost}(M))$, no better than full sequential
decode}. When $(K+k) \cdot
\mathrm{cost}(M) \gg \log L(R)$ the second term dominates and the
cost is $O((K+k) \cdot \mathrm{cost}(M))$, independent of $|X|$ in
that parameter regime. \emph{The factor $\mathrm{cost}(M)$ is one
per-position evaluation of the neural predictor $M$; it is
non-universal --- architecture- and context-dependent (Lemma 5.1b
records $\Theta(W)$ per token for a windowed transformer of context
$W$, and $\Theta(\sigma^W)$ model storage with $O(1)$ per-evaluation
lookup for an explicit table) --- and, since the decode window
dominates $\log L(R)$ in any realistic regime, it is the dominant
term of the per-query cost. The polylog conclusion is therefore a
statement about the dependence on $|X|$ with $\mathrm{cost}(M)$ held
fixed, not a claim that the per-query cost is small in absolute
(wall-clock) terms.}}

\emph{\textbf{Proof.}} Under the uniform-spacing convention of
Definition 10.3, the byte position of the latest sync point
$P \leq p$ is the deterministic value $P = K \cdot \lfloor p/K \rfloor$
(a constant-time computation, no search required). Look up the
corresponding bit offset in $R$ from the seek index: the index has
exactly $\lceil |X|/K \rceil$ entries, each of width
$\lceil \log_2 L(R)\rceil + O(1)$ bits, indexed by $\lfloor p/K \rfloor$
in $O(\log L(R))$ bits per direct probe. (For non-uniform spacing
with $m$ sync-point entries, each entry is
$\lceil \log_2|X|\rceil + \lceil \log_2 L(R)\rceil + O(1)$
bits wide, and binary search over the index costs
$O(\log m \cdot (\log |X| + \log L(R)))$ bits in total
(equivalently $O(\log(|X|/K) \cdot (\log |X| + \log L(R)))$ when
$|X| \geq K$, since $m \leq \lceil 2|X|/K \rceil$); the
uniform case avoids both the $\log m$ factor and the
$\log |X|$ per-entry term at the cost of one extra constant-time
arithmetic operation.) By
Definition 10.3 (a, b), decoding from $R$'s bit offset at $P$
proceeds with both the predictor state initialized to
$Q_{\text{init}}$ and the entropy-coder state initialized to its
post-flush starting value; no prior bits of $R$ are required to
parse the bits following the sync-point offset. Decode forward one
position at a time: let $k' = \min(k, |X| - p)$ (the actual number
of output bytes, equal to $k$ when $p + k \leq |X|$, less when
$p + k > |X|$). At each position $i \in \{P, P+1, \ldots, p+k'-1\}$,
compute $Y_i = M(C_i)$ where $C_i$ is the context built from
previously decoded positions $[P, i)$. By the $W$-bounded property
(Definition 10.2), $Y_i$ depends only on bytes $[i-W, i)$, which are
all available within the decoded prefix once $i$ exceeds $P + W$. The
first $W$ positions after $P$ may have shortened context (positions
$[P, P+W)$ see fewer than $W$ bytes), which the encoder accommodates
by encoding those positions with reduced-context predictions; this is
exactly the cold-context cost analyzed before. Continue decoding
until position $p + k'$; output bytes $[p, p+k')$. By Theorem 10.1,
all predictor evaluations are deterministic and bit-identical to
the encoder's sync-pointed full decode (i.e., the encoding produced
under the same sync-point structure starting from position 0 and
resetting predictor/entropy-coder state at each $jK$), so the
output matches the original $X$ on $[p, p+k')$. Total decoded
positions: $p + k' - P \leq K + k$ since $P > p - K$.
\hfill$\blacksquare$

\textbf{Remark 10.3a (Type-II partition modes give only
sub-block-granular random access).} Type-II's enumerative-rank
coding encodes a whole-block partition (Theorems 5.2--5.3) as a
single combinatorial rank, not as a per-position sequential repair
stream. As stated, Theorem 10.3 does not apply to a Type-II archive
that uses a single global rank for the full block. The standard
workaround is to subdivide the block at sync-point boundaries into
independently partition-coded sub-blocks of length $K$ each
(recovering only the sub-block-level Type-II separation, as already
noted in §10.7's ``Type-II semi-non-factorized coding'' trade-off
bullet). Each sub-block is then a standalone Type-II rank, and
random access has \emph{sub-block granularity}. For a read
$[p, p+k)$, let $j_0 = \lfloor p/K \rfloor$ and
$j_1 = \lfloor (p+k-1)/K \rfloor$; the request spans
$j_1 - j_0 + 1 = \lceil k/K \rceil + O(1)$ sub-blocks (one more
when $p$ is not aligned to a sub-block boundary). The driver looks
up each of the affected sub-blocks' bit offsets (cost
$O((\lceil k/K \rceil + 1) \cdot \log L(R))$ under uniform spacing)
and decodes each in full (cost $O((\lceil k/K \rceil + 1) \cdot K
\cdot \mathrm{unrank}(K))$), producing $(j_1 - j_0 + 1) \cdot K$
output bytes of which only $k$ are actually requested. For
$k \ll K$ the read costs $O(\log L(R) + K \cdot \mathrm{unrank}(K))$
and produces $K + O(K)$ bytes (decoding overhead is one full
sub-block per request); for $k \gg K$ the cost is
$O((k/K) \cdot K \cdot \mathrm{unrank}(K)) = O(k \cdot
\mathrm{unrank}(K))$, near-linear in $k$. Byte-granular access
within a sub-block requires decoding the whole sub-block, not just
the prefix up to $p+k$ as in Theorem 10.3's per-position sequential
decode. This is a strictly weaker random-access property than the
per-position sequential case.

Theorem 10.3 establishes that random access in RNR archives using
\emph{per-position sequential} coding modes runs in time with only
polylogarithmic dependence on archive size in the parameter regime
where $K$, $k$, $\mathrm{cost}(M)$, and $\log L(R)$ are polylog in
$|X|$ (a $\log L(R)$ direct seek-index cost under uniform spacing,
or a $\log(|X|/K) \cdot (\log |X| + \log L(R))$ binary-search cost
under non-uniform spacing, plus a $(K+k) \cdot \mathrm{cost}(M)$
decode-window cost).
Whole-block Type-II/III-B partition-rank modes give only
sub-block-granular access at a separately analyzed
$O(K \cdot \mathrm{unrank}(K))$ per requested sub-block (Remark
10.3a). For $K = 64$ KiB and $k = 1$ KiB, the cost is roughly
$K + k = 65$ KiB worth of predictor evaluations --- typically on
the order of $650\,\text{ms}$ on a single CPU core for a small
distilled transformer (one forward pass at $\sim 10\,\mu\text{s}$
per byte; faster on GPU/NPU). This is asymptotically a worst-case (end-of-archive read)
factor of $|X|/(K + k)$ faster than full sequential decoding, or
about $|X| / (2(K + k))$ for an average random read at position
$|X|/2$. For $|X| = 1$ GiB and $K = 64$ KiB the worst-case
speedup is roughly 16{,}000-fold and the average is roughly
8{,}000-fold.

\emph{Per-mode access-granularity and per-query cost (collation of
Theorem 10.3, Remark 10.3a, §6.3, and the §10.7 trade-off bullet).} The
random-access granularity is not uniform across coding modes: it is the
per-position-sequential property of Theorem 10.3 that delivers
byte-granular access, while the whole-block enumerative-rank modes of
Remark 10.3a give only sub-block granularity. The table below
consolidates, per coding mode, the granularity actually established, the
per-query cost formula in that mode, the wasted decode incurred to serve
one read, and the source of each entry. No row states a new result; each
is a restatement of the cited theorem, remark, or trade-off bullet.
In all rows $\mathrm{cost}(M)$ is the per-position predictor evaluation,
the non-universal, architecture-dependent factor that dominates the
per-query cost in any realistic regime (Lemma 5.1b); ``polylog in $|X|$''
qualifies only the \emph{dependence on archive size} with
$K,k,\mathrm{cost}(M),\log L(R)$ held fixed or polylog.

\begingroup\footnotesize\sloppy
\setlength{\tabcolsep}{4pt}
\noindent\begin{tabular}{@{}p{2.5cm}p{1.9cm}p{3.0cm}p{2.6cm}p{2.4cm}@{}}
\hline
\textbf{coding mode} & \textbf{access granularity} &
\textbf{per-query cost} & \textbf{wasted decode per query} &
\textbf{source}\\
\hline
Type-I (direct byte arithmetic; $S\!=\!1$)
& byte
& $O(\log L(R) + (K\!+\!k)\,\mathrm{cost}(M))$
& $\le K$ warmup positions (discarded); decode $\le K\!+\!k$ total
& Thm 10.3 (incl.); §10.7 bullet (``not affected'')\\
Type-II enumerative / partition-rank (whole block)
& sub-block ($K$)
& $O(\log L(R) + K\!\cdot\!\mathrm{unrank}(K))$ per requested sub-block; $\lceil k/K\rceil\!+\!O(1)$ sub-blocks for a $k$-read
& up to one full sub-block ($K\!+\!O(K)$ bytes) per request; $k$ requested
& Rem 10.3a; Thm 10.3 (excl.); §10.7 ``Type-II semi-non-factorized'' bullet\\
Type-III-A composition over per-position-sequential per-tile modes
& byte, \emph{iff} $K$ tile-aligned (sync points at tile boundaries)
& $O(\log L(R) + (K\!+\!k)\,\mathrm{cost}(M))$
& $\le K$ warmup positions; tile/sync misalignment voids byte granularity
& Thm 10.3 (incl., tile-alignment caveat)\\
Type-III-B (byte tensor; §6.3): per-byte sequential strategy (iii); $S\!=\!2$
& byte
& $O(\log L(R) + (K\!+\!k)\,\mathrm{cost}(M))$
& $\le K$ warmup positions (discarded); decode $\le K\!+\!k$
& Thm 10.3 (incl., strategy (iii)); §6.3 (iii); §10.7 bullet\\
Type-III-B direct / hierarchical partition strategies (i)/(ii) of §6.3
& sub-block ($K$)
& $O(\log L(R) + K\!\cdot\!\mathrm{unrank}(K))$ per sub-block (as Type-II)
& up to one full sub-block per request
& Thm 10.3 (excl.); Rem 10.3a; §6.3 (i)/(ii); §10.7 bullet\\
Type-III-C dictionary / grammar (flat terminal cover)
& byte (per-position-sequential carrier), under uniform sync only if no token crosses $jK$
& $O(\log L(R) + (K\!+\!k)\,\mathrm{cost}(M))$; non-uniform sync: a binary-search seek at polylog cost (\S10.7 non-uniform branch)
& $\le K$ warmup; tokens crossing $jK$ forfeit $\le 8|\tau|$ bits each (split), bounded $8|X|\overline{|\tau|}/K$
& §10.7 bullet (Type-III-C); Thm 10.3 (non-uniform branch)\\
\hline
\end{tabular}
\par\endgroup

\emph{Reading guide.} Rows 1, 3, 4, and 6 are the per-position
sequential modes in Theorem 10.3's scope: each is byte-granular and pays
the $O(\log L(R) + (K\!+\!k)\,\mathrm{cost}(M))$ window cost, with the
only wasted work the $\le K$ warmup positions discarded before position
$p$ (Theorem 10.3 step (ii); total decoded positions $p+k'-P\le K+k$).
Rows 3 and 6 carry an alignment caveat: Type-III-A is byte-granular only
when $K$ is tile-aligned, and Type-III-C only when no dictionary token
crosses a sync boundary $jK$ (otherwise the crossed token is split and
forfeits its saving, §10.7's third trade-off bullet; a non-uniform
sync-point variant aligns boundaries to token edges at the cost of the
binary-search seek term). Rows 2 and 5 are the whole-block
enumerative-rank modes excluded from Theorem 10.3 (Remark 10.3a): they
are only sub-block-granular, must decode and discard up to a full
$K$-byte sub-block to serve one byte, and replace the
$(K\!+\!k)\,\mathrm{cost}(M)$ window term with $K\!\cdot\!\mathrm{unrank}(K)$
per affected sub-block. No quantity in this table is new; each cell is
cross-referenced to Theorem 10.3, Remark 10.3a, §6.3, or the §10.7
trade-off bullet from which it is taken.

\textbf{Theorem 10.4 (Random access for non-$W$-bounded predictors).}
\emph{Under the same scope as Theorem 10.3 (per-position sequential
RNR coding modes; whole-block Type-II partition-rank modes are out
of scope), for a non-$W$-bounded predictor (RNN, LSTM, state-space
model), the predictor-reset condition (a) of Definition 10.3 cannot
be satisfied by cold-resetting context alone: the predictor's
recurrent hidden state depends on arbitrarily long history. Replace
condition (a) by an \emph{augmented sync point}: instead of
resetting the predictor's hidden state, store its full state at
position $P$ in the seek-index entry (so the decoder, when seeking
to $P$, loads the stored state and proceeds with bit-exact
continuation). Condition (b) on entropy-coder flush/reset remains
unchanged. The augmented snapshot stores the predictor's recurrent
hidden state, any adaptive parameter snapshot (if Theorem 10.2 is
in effect), and the entropy-coder boundary state (range/state for
arithmetic coding or ANS state, if a different flush convention is
used). Let $L(\text{state})$ be the bit length of this combined
boundary snapshot. Under the uniform-spacing convention of
Definition 10.3, the augmented seek index has exactly
$\lceil |X|/K \rceil$ entries (one per uniform sync point), with
total size $O((|X|/K) \cdot (\log L(R) + L(\text{state})))$ bits.
The random-access cost per query (to $k$ bytes at position $p$) is
$O\bigl(\log L(R) + L(\text{state}) + (K + k) \cdot
\mathrm{cost}(M)\bigr)$ under uniform spacing (direct index lookup),
or $O\bigl(\log(|X|/K) \cdot (\log |X| + \log L(R) + L(\text{state})) + (K + k)
\cdot \mathrm{cost}(M)\bigr)$ under non-uniform spacing (binary
search, with each entry containing a $\log |X|$ source-position
key, a $\log L(R)$ repair-stream offset, and the $L(\text{state})$
snapshot), where the index-related terms cover the entry lookup
plus loading and deserializing the snapshot before decoding can
resume.}

\emph{\textbf{Proof.}} Augment each sync point's index entry with
all boundary state at position $P$, serialized as a fixed-bit-width
integer tuple as required by Theorem 10.1. The combined snapshot adds
$L(\text{state})$ bits per sync point, multiplying the index size by
$(1 + L(\text{state}) / \log L(R))$. Under uniform sync-point
spacing (Definition 10.3), to random-access bytes $[p, p+k)$, the
latest sync point $P \leq p$ is the deterministic value
$P = K \lfloor p/K \rfloor$ (constant-time computation); direct
lookup at entry $\lfloor p/K \rfloor$ reads
$O(\log L(R) + L(\text{state}))$ bits to obtain the bit offset and
snapshot. Under non-uniform spacing, each entry is
$\lceil \log |X| \rceil + \lceil \log L(R) \rceil + L(\text{state})
+ O(1)$ bits wide (source-position key plus repair offset plus
snapshot), and the lookup is a binary search over the augmented
index, costing
$O(\log(|X|/K) \cdot (\log |X| + \log L(R) + L(\text{state})))$
bits in total.
In either case, deserialize the boundary snapshot, restore the
decoder's predictor, adaptive-parameter, and entropy-coder boundary
states, and decode forward as in Theorem 10.3 (cost
$O((K+k) \cdot \mathrm{cost}(M))$).
The decoded output is bit-identical by Theorem 10.1 applied to the
restored state. \hfill$\blacksquare$

For LSTM predictors with $d$-dimensional hidden state and
$d$-dimensional cell state (the standard LSTM has both), the
snapshot is $L(\text{state}) = 2 d \cdot b$ bits where $b$ is the
per-component bit-width. For modern state-space models with
structured state-transition matrices, $L(\text{state})$ can be
$O(d^2)$. With $d = 1024$ and $b = 16$, an LSTM state snapshot is
$2 \cdot 1024 \cdot 16 = 32{,}768$ bits $= 4$ KiB; placing one
snapshot per $K = 64$ KiB sync point adds $4 / 64 = 6.25\%$ of the
source size as augmented seek-index overhead (a $\sim$1000$\times$
increase over the offset-only index, which is only $\sim$8 bytes
per entry). State-space models with $d = 256$ and structured-state
representation can require 64 KiB snapshots per sync point, giving
$64 / 64 = 100\%$ of source size as augmented-index overhead at
$K = 64$ KiB. Since the underlying compressed archive is typically
$\sim$10--50\% of source size for the data classes RNR targets, an
augmented index this large adds at least $2 \times$--$10 \times$
to the on-disk archive size, which is deployment-prohibitive
without raising $K$ to amortize the snapshot cost. This motivates the $W$-bounded
design choice for deployable RNR codecs, particularly for archive
formats where the seek index is co-located with the data and
contributes to its size.

\emph{\textbf{Trade-offs}}

Three trade-offs govern random-access design in RNR archives.

\begin{itemize}
\item
  Sync-point spacing $K$. Smaller $K$ means faster seek and smaller
  warmup cost, but more sync points and therefore more cold-context
  compression overhead. Typical choices: $K = 4$ KiB (aggressive
  seek), $K = 64$ KiB (balanced), $K = 1$ MiB (compression-favored).
\item
  Type-II semi-non-factorized coding. Theorem 5.4's $\Theta(N)$
  separation requires joint structure across all $N$ positions of a
  block. Inserting sync points every $K$ positions partitions the
  block into independent $K$-length subblocks, recovering only
  $\Theta(K)$ separation per subblock. For $K = 64$ KiB and $N = 1$
  GiB, the per-subblock advantage is preserved but the cross-subblock
  joint structure is forfeited. Type-I is not affected by this
  trade-off because its construction is per-position sequential
  (Theorem 10.3 applies directly). Type-III-B \emph{is} affected
  when it uses the direct or hierarchical partition coding modes of
  §6.3 (which are Type-II at the 256-symbol alphabet); under those
  modes Type-III-B requires the same sub-block sync structure as
  Type-II and inherits the sub-block-granular random-access
  property of Remark 10.3a. A per-byte sequential Type-III-B coder
  --- one that codes each byte's high-then-low nibble at its own
  position using the joint $Q_{h,l}(i)$ --- does fall under Theorem
  10.3's per-position sequential scope and gets byte-granular
  random access. Type-III-C
  is affected when dictionary tokens span longer fragments or
  grammar expansions. We restrict attention to the standard case
  where the top-level parse is a flat non-overlapping cover of $X$
  by terminal-level token instances (each occurrence is a maximal
  expansion at source level); recursive grammar productions are
  expanded for the purpose of counting source-byte crossings. A
  token that would cross a uniform Def-10.3 sync-point boundary
  $jK$ must be split, with the affected occurrence forfeiting its
  dictionary saving (at most $8 \cdot |\tau|$ bits per crossed
  token, since the raw-literal baseline costs $8\ell$ bits for an
  $\ell$-byte token). Under the flat-cover restriction, at most one
  top-level token instance crosses any given source boundary
  (recursive expansions inside that one instance are all confined to
  the same $\ell$-byte span and contribute no additional crossings
  beyond their containing token's split). Over $|X|/K$ sync points
  the worst-case loss is bounded by $8 \cdot |X| \cdot
  \overline{|\tau|_{\text{crossed}}} / K$ bits, where
  $\overline{|\tau|_{\text{crossed}}}$ is the average length of
  crossed tokens (bounded above by $\max_\tau |\tau|$). For
  $K \gg \max_\tau |\tau|$ this is negligible; for $K$ comparable to
  $\max_\tau |\tau|$ the worst-case fractional overhead approaches
  $8 \cdot \max_\tau|\tau| / K$ bits per source byte, which can be
  substantial. Specifically, $K = 4$ KiB and a worst-case
  $\max_\tau|\tau| = 1$ KiB give a worst-case overhead bound of
  $2$ bits per source byte (25\% of an 8 bits/byte raw source);
  in practice the average crossed token is shorter than the maximum,
  and the typical overhead is a few per-cent on natural data. A
  non-uniform sync-point variant that aligns boundaries to token
  edges instead is also available, at the cost of storing byte
  offsets in the seek index (Definition 10.3's non-uniform spacing
  case, adding $\lceil \log_2 |X|\rceil$ bits per entry); the
  random-access cost then follows the non-uniform branch of
  Theorem 10.3 ($O(\log(|X|/K) \cdot (\log |X| + \log L(R)))$
  binary-search bits, including the per-entry source-position
  width).
\item
  Verification under random access. Theorem 9.1's chained-hash
  verification covers sequential decoding; a random-access reader
  cannot verify bytes $[0, p)$ it never decoded. Three solutions
  exist: (a)~per-sync-point hashes that allow local verification,
  costing $O((|X|/K) \cdot |\text{hash}|)$ bits; (b)~a Merkle tree
  over \emph{all} archive-mounted bytes --- the repair stream,
  seek index, sync-point metadata, predictor parameters, and
  per-subblock source hashes --- so that any tampering with index
  entries (which would otherwise let a corrupted seek-index entry
  point a valid repair-stream region to the wrong source range)
  is also detected; this still gives $O(\log |X|)$ hash overhead
  per read but requires that the Merkle root be supplied via a
  trusted channel; (c)~explicitly unverified random access for
  trusted-source archives. The Merkle-tree option is standard in
  content-addressable storage and is the recommended default for
  archive deployments.
\end{itemize}

\emph{\textbf{Comparison with existing systems}}

Standard random-access lossless compressors:

\begin{itemize}
\item
  bgzip (block-gzip): independent gzip blocks, ratio cost
  approximately 3 to 7 percent versus monolithic gzip, random access
  in $O(K + k)$ time. Widely used in bioinformatics for indexed FASTA
  and BAM files.
\item
  zstd seekable mode: independent zstd frames with a footer index,
  ratio cost approximately 2 to 10 percent versus monolithic zstd,
  random access in $O(K + k)$ time.
\item
  LZMA seekable wrappers (ad hoc xz multi-stream archives): ratio
  cost approximately 5 to 15 percent, random access in $O(K + k)$
  time.
\end{itemize}

Existing neural lossless compressors lack sub-archive random access,
with two distinct underlying obstructions. Autoregressive byte and
pixel models (NNCP, PixelCNN++, PixelSNAIL) require sequential
evaluation of bytes $[0, p)$ before byte $p$ can be decoded; cost is
$\Theta(p)$ predictor evaluations. Block-codec bits-back and
invertible-flow models (Bit-Swap, IDF++, HiLLoC) decode in fixed
block/patch units, so even a one-byte read forces decoding of the
enclosing block; effective seek granularity is the block size, not
the byte (typical block sizes are kilobytes to megabytes). Neither
family supports byte-resolution random access as currently designed.
For a 1 GiB archive the autoregressive case is hours of wall-clock
time per access; the block-codec case decodes the enclosing block
on every read.

RNR with a $W$-bounded transformer predictor, sync points every
$K$ positions, and a per-position sequential coding mode (Type-I,
per-byte sequential Type-III-B, or Type-III-A composition over
per-position sequential per-tile modes): per-query random access
in $O(\log L(R) + (K + k) \cdot \mathrm{cost}(M))$ time for uniform
spacing (or with an extra $\log(|X|/K)$ factor for non-uniform),
independent of $|X|$ in the predictor-cost term and polylog in
$|X|$ in the seek-index-lookup term, at a per-sync-point
worst-case excess of $S \cdot W \cdot \log_2(1/\eta)$ bits
(Section 10.7), giving a worst-case total of $m \cdot S \cdot W
\cdot \log_2(1/\eta)$ bits over the archive where $m$ is the
sync-point count ($m = \lceil |X|/K \rceil$ uniform, $m \leq \lceil
2|X|/K \rceil$ non-uniform per Definition 10.3); $S$ is
symbols-emitted-per-source-byte ($S = 1$ for Type-I/direct byte
arithmetic, $S = 2$ for per-byte sequential Type-III-B). For
$K = 64$ KiB, $W = 4$ KiB, $\eta = 2^{-12}$ (standard small
uniform escape mixed with the predictor's main distribution;
$\eta = 2^{-8}$ would force a uniform predictor and contradict
any cold-vs-warm gap below), $S = 1$ (Type-I), and a 1 GiB source,
the worst-case bound evaluates to $16384 \cdot 4096 \cdot 12$
bits $= 96$ MiB, or roughly $9.4\%$ of the source size; for $S = 2$
(per-byte Type-III-B) the worst-case bound doubles to $\approx
19\%$. The realistic overhead is smaller because the
per-byte excess in practice is the cold-vs-warm \emph{rate gap}
$\delta$ (typically a few tenths of a bit per byte for natural
data) rather than the full $\log_2(1/\eta) = 12$ bits. With
$\delta \approx 0.5$ bit/byte, $(|X|/K) \cdot W \cdot \delta =
16384 \cdot 4096 \cdot 0.5$ bits $\approx 32$ Mibits $\approx 4$
MiB, or roughly $0.4\%$ of the source size; this estimate is
heuristic and is not a theorem. The ratio comparison
against bgzip and zstd seekable (typically 2--7 percent overhead)
is empirical and is the subject of §11.

This is the framework\textquotesingle s deployment-grade distinguishing
feature relative to other neural compressors. The bounded-context
architecture, the deterministic reproducibility guarantee (Theorem
10.1), and the sync-point construction (Definition 10.3) jointly
enable seek-and-read access at a parameter-dependent ratio
overhead. For per-position sequential modes (Type-I, the per-byte
sequential Type-III-B strategy, and the per-tile sequential
branches of Type-III-A) the overhead is on the order of a few
percent for natural data at typical $K$ under the
probability-floor assumption (a heuristic estimate; not a theorem),
and random access is byte-granular. For whole-block rank-coded
Type-II and partition-coded Type-III-B modes, access is only
sub-block-granular per Remark 10.3a, with a separate $O(K \cdot
\mathrm{unrank}(K))$ per requested sub-block. For Type-III-C the
worst-case bound depends additionally on $K$ and $\max_\tau
|\tau|$ as detailed above. This is a capability that
LSTM-based or state-space-based
neural compressors can only provide with the expensive state
snapshots of Theorem 10.4 (4 KiB or more per sync point); that
NNCP-style autoregressive transformer compressors lack only because
their encoders do not insert sync points or maintain a seek index;
and that bits-back / invertible-flow block codecs (Bit-Swap, IDF++,
HiLLoC) lack because their block/ANS/latent decoding requires
recovering the enclosing block (or vectorized batch of patches) as
a unit, so the effective seek granularity is the block, not the
byte; sub-block-granularity seek is not provided by the current
designs without modifying their block structure.

\textbf{10.8 Block-device and filesystem deployment}

A natural application of the random-access result is to expose an
RNR-compressed archive as a read-only block device or filesystem through
an operating-system driver, allowing files or sectors to be accessed on
demand without ever decompressing the whole archive. This is the
deployment pattern of squashfs, EROFS, btrfs with zstd compression, and
Apple Disk Images: read-only compressed system images mounted by the
kernel and served to user-space applications transparently. Existing
neural compressors cannot serve this use case because their
$O(|X|)$ sequential decoding requirement makes any access to byte
$p$ cost the same as decompressing all preceding bytes; for archive
sizes typical of disk images (gigabytes to terabytes), this is
prohibitive. RNR removes this obstruction. The two theorems below
formalize the application.

\textbf{Theorem 10.5 (Block-device interface, parameter-conditional).}
\emph{Under the same scope as Theorem 10.3 (per-position sequential
RNR coding modes; Type-II rank-coded modes give only
sub-block-granular access per Remark 10.3a and require a separate
unranking-cost analysis), let $X$ be compressed under RNR with a
$W$-bounded predictor $M$ satisfying Theorem 10.1 and sync points
placed every $K$ positions; let $|X|$ denote the uncompressed
source length. A block-device driver that maps logical-block-
address-indexed read requests (LBA $\mathcal{L}$, length $k$ bytes)
to RNR random-access calls can serve any read request as follows: let $p = \mathcal{L}
\cdot \text{sector\_size}$. If $p \geq |X|$, return an empty read
($0$ bytes) and the standard end-of-device indication. Otherwise
$p \in [0, |X|)$ and the driver returns $k' = \min(k, |X| - p)
\geq 1$ bytes from byte range $[p, p+k')$ (a full read when
$p+k \leq |X|$, a short read when $p+k > |X|$, mirroring POSIX
\texttt{pread}). The boundary condition is inherited from Theorem
10.3. The read is served in time
$O(\log L(R) + (K + k) \cdot \mathrm{cost}(M))$ under uniform
sync-point spacing (Theorem 10.3; non-uniform spacing with $m$
sync-point entries replaces the direct lookup with binary search
at $O(\log m \cdot (\log |X| + \log L(R)))$ bits) and
working memory $O((K + k) \cdot 8 + L(M) + L(\mathrm{workspace}(M))
+ L(\mathrm{seek\_index}))$ \emph{bits}, equivalently
$O(K + k + L(M)/8 + L(\mathrm{workspace}(M))/8 +
L(\mathrm{seek\_index})/8)$ \emph{bytes}, where $L(M)$ is the
predictor description size in bits (kept resident in the driver),
$L(\mathrm{workspace}(M))$ is the per-inference activation/KV/scratch
workspace required by the inference engine on its target hardware
(architecture-specific; e.g., transformer KV-cache for the last
$W$ positions, layer-wise activation buffers), and
$L(\mathrm{seek\_index})$ is the seek-index size in bits (typically
$O((|X|/K) \log L(R))$ under uniform spacing). The cost is polylogarithmic in
archive size $|X|$ only in the parameter regime where $K$, $k$,
$\mathrm{cost}(M)$, and $\log L(R)$ are held fixed or grow at most
polylogarithmically in $|X|$; outside that regime the
output-sensitive cost of Theorem 10.3 applies.}

\emph{\textbf{Proof.}} Each LBA-indexed read request $(\mathcal{L},
k)$ corresponds to a byte-range request on the underlying source $X$
starting at byte position
$p = \mathcal{L} \cdot \text{sector\_size}$ and continuing for $k$
bytes. By Theorem 10.3, RNR random access to $k$ bytes at position
$p$ runs in $O(\log L(R) + (K + k) \cdot \mathrm{cost}(M))$ time
(under uniform spacing) and
requires $O(K + k)$ bytes (equivalently $O((K+k) \cdot 8)$ bits) of
decode-window working memory (warmup buffer plus output buffer).
The driver additionally keeps the predictor weights ($L(M)$ bits),
the predictor inference workspace ($L(\mathrm{workspace}(M))$ bits;
e.g., a $W$-position KV-cache for a transformer plus layer-wise
activation buffers), and the seek index ($L(\mathrm{seek\_index})$
bits) resident; these are fixed at archive-mount time, not
per-query. The translation from $(\mathcal{L}, k)$ to
$(p, k)$ is a constant-time arithmetic operation. The total
per-request cost is therefore as stated, with the parameter regime
caveat inherited from Theorem 10.3. \hfill$\blacksquare$

The $\mathrm{cost}(M)$ factor governs whether the driver is fast
enough for interactive use. For a small distilled byte-level
transformer (10 to 50 million parameters, int8 weights), one
single-byte forward pass costs approximately $10\,\mu\text{s}$ on a
single CPU core, $1\,\mu\text{s}$ on a server GPU, and
$0.1\,\mu\text{s}$ on a high-end NPU (engineering estimates; actual
numbers vary by quantization scheme and hardware generation).
Treating one forward pass as one decoded byte, per-byte inference
throughput is thus approximately $100$ KB/s on CPU, $1$ MB/s on
GPU, and $10$ MB/s on NPU. For $K + k = 64$ KiB, the random-access
latency for cold reads is therefore approximately
$64 \text{ KiB} \cdot 10\,\mu\text{s/byte} \approx 650\,\text{ms}$
(CPU), $\approx 65\,\text{ms}$ (GPU), and
$\approx 6.5\,\text{ms}$ (NPU). The NPU end of this range is in the
same order of magnitude as a slow random-read of a multi-MB region
from local NVMe (single-block NVMe random-read latency is
$\sim 100\,\mu\text{s}$, so RNR-NPU is about $65 \times$ slower per
request, comparable to a server-side decompressor); the CPU end
($\sim 650$ ms per cold read) is two orders of magnitude slower
than HDD seek latency (single-digit ms) and is in the range of
disk-network or multi-second offline-access workloads. Hardware-batching multiple
forward passes per kernel invocation can improve effective
throughput by a factor of 10--100 on GPU/NPU when access patterns
are sequential, narrowing the cold-read latency proportionally.
Production deployments will target NPU or GPU paths when low latency
is required and fall back to CPU only for archive inspection or
compatibility.

\textbf{Theorem 10.6 (Hardware-target portability, explicit
conformance).} \emph{Let $M$ satisfy the integer-arithmetic,
lookup-table, and fixed-reduction-order conditions of Theorem 10.1.
A hardware target $T$ is said to \emph{conform} for $M$ if it pins,
in a target-specific test vector matched to the engineering spec,
all of the following: (a)~signed integer representation (two's
complement) at the bit-widths declared by the model description;
(b)~the result of every left/right shift, division, and modulo at
those bit-widths (including round-toward-zero vs floor and signed
vs unsigned semantics); (c)~the saturation-vs-wrap behaviour at each
arithmetic operation; (d)~the byte-level serialization of every
non-linearity lookup table; (e)~the reduction-tree shape and order;
(f)~the per-target serialization of all approximation functions
(integer Newton-Raphson, table-lookup interpolation, etc.) used in
the forward pass; (g)~the byte-level encoding of all model
parameters and intermediate tensors, including endianness, signed
vs unsigned interpretation of each declared bit-width, the packing
order of sub-byte quantizations (e.g., int4-pair-per-byte: low
nibble first vs high nibble first), zero-point and scale storage
layout, and tensor index-to-memory order for any layout-sensitive
tensor (e.g., NHWC vs NCHW for convolution weights; row-major vs
column-major for matrix multiplications). Under (a)-(g), the output
$M(C)$ computed on $T$ is bit-identical to that computed on any
other conforming target. The choice of target affects throughput
and latency but not the produced bytes; the engineering spec lists
the conformance tests that establish (a)-(g) for each named target
class.}

\emph{\textbf{Proof.}} Direct corollary of Theorem 10.1, with the
explicit observation that conformance is a property of the
target\textquotesingle s integer-arithmetic, lookup-table, and
data-layout semantics rather than of its parallelism strategy or
vendor identity. Two's-complement signed integer addition and
multiplication at a fixed bit-width are associative and commutative
\emph{within the representable range and under a pinned wrap-vs-saturate
convention}; condition (a) fixes the representation, (c) fixes the
wrap-vs-saturate behaviour for each operation, and the per-layer
bit-width bound of Theorem 10.1 keeps every accumulator within the
representable range so neither wrap nor saturation triggers in a
conforming forward pass --- under those conditions the produced
bits are identical across platforms. Conditions (b), (d), (e), (f),
and (g) pin the remaining sources of cross-platform divergence
(shift/division/modulo semantics, lookup-table layout, reduction
order, approximation functions, and weight/input byte
serialization); each is enforced by the engineering spec's
target-specific test vectors. Without (a)--(g), bit-identity is
not guaranteed: e.g., the same algorithm with different sub-byte
packing orders produces different weights. Fixed reduction order,
specified by the model description, is honored by any sequential
implementation and by any parallel implementation that uses a
tree-reduction with the specified schedule. Composing these properties
layer by layer through $M$ yields bit-exact output. \hfill$\blacksquare$

Candidate hardware-target classes that an implementation could
potentially make conforming under the (a)-(g) requirements include:
scalar CPU on any architecture supporting the required integer
bit-widths (x86-64, ARM64, RISC-V, POWER); CPU SIMD using integer
SIMD instructions (AVX2, AVX-512 with VNNI and IFMA extensions, ARM
NEON, RISC-V vector extension); GPU using integer compute kernels
and integer Tensor Cores (NVIDIA Turing or later with int8 and int4
Tensor Core paths, AMD CDNA with matrix accelerators); dedicated
NPUs targeting integer inference (Apple Neural Engine, Google Edge
TPU, Qualcomm Hexagon NPU, Huawei Ascend, AWS Inferentia, AWS
Trainium in inference mode). Theorem 10.6 is conditional on a
target satisfying (a)-(g): whether any specific named class
\emph{actually} conforms is an engineering claim about that
target's instruction set, kernel library, and serialization
conventions, not a corollary of Theorem 10.6. The conformance tests
that establish (a)--(g) for each named target class are specified in
the companion engineering specification (\emph{RNR Coding --- Reference
Implementation Conformance Specification}, §4 per-target paths and §12
conformance test suite, with cross-target byte-identity as acceptance
criterion R-12.2.3); a target's certification against that suite is the
empirical Hypothesis H14 of the experimental-design companion. With
such a certification of each target, the same RNR archive can be
encoded on a server GPU, served from cloud storage, and decoded on a
phone NPU with byte-identity guaranteed by Theorem 10.6 --- a
conditional guarantee, not a demonstration carried out here.

\emph{\textbf{Cache amortization}}

A read-cache of size $C$ placed between the driver and the OS page
cache permits steady-state throughput exceeding the per-access bound
of Theorem 10.5. For sequential reads over a region of size $R \gg
K$ (so a single $K$-byte warmup is amortized over many output
bytes) up to cache size $C$, the amortized cost per byte approaches
$O(\mathrm{cost}(M))$: one predictor evaluation per byte once the
warmup pays back, with $(K + R)/R$ amortization factor (asymptotic
in $R/K$). For small sequential reads $R \lesssim K$ the
per-access $K + R$ warmup still dominates and the per-byte cost is
$O((K + R) \cdot \mathrm{cost}(M) / R)$. For random read \emph{re-access patterns} where the same regions
of size $\lesssim C$ are revisited many times (the working-set fits
in cache), the cache hit rate approaches one on repeated accesses,
and the amortized cost per byte after warmup approaches
$O(\mathrm{memcpy})$ speed, dominated by main memory bandwidth.
A single forward scan over a region without re-access amortizes
its $K$-byte warmup over the scan as covered above (one warmup
per access, not one per byte), but does not benefit from the
re-access caching effect; the relevant benefit is the
$R/K$-amortized cost of Theorem 10.3, not a hit-rate near one. For
uniformly random reads at distance larger than $C$, the amortized
cost per byte equals the Theorem 10.5 bound divided by $k$, that is,
$O((K + k) \cdot \mathrm{cost}(M) / k)$, which approaches
$\mathrm{cost}(M)$ as $k$ grows. In typical deployments the
operating-system page cache provides an effective cache of several
gigabytes, absorbing read locality so that the worst-case Theorem
10.5 bound is the latency budget for cold reads only.

\emph{\textbf{Read-write semantics}}

RNR archives are read-only by construction: the predictor $M$ and the
sync-point structure assume a fixed source $X$. Three patterns
support read-write deployments.

\begin{itemize}
\item
  Copy-on-write overlay. The compressed archive is mounted as a
  read-only base; writes are directed to a separate writable layer
  (Linux overlayfs, unionfs, AUFS). Reads consult the overlay first
  and fall through to the base only on miss. This is the
  squashfs-plus-overlayfs pattern that container runtimes (Docker,
  Podman, containerd) and live Linux distributions use today.
\item
  Region rewrite (fixed-predictor case only). A modification within
  sync-point region $[P, P + K)$ re-encodes only that region:
  re-runs the predictor on the modified bytes starting from $P$,
  re-encodes the repair stream for the region, and updates the seek
  index entry for the affected sync point. This is well-defined
  \emph{only when the predictor's behaviour at every subsequent
  position is independent of the bytes in $[P, P+K)$}, i.e., for a
  fixed-parameter predictor with no online adaptation. For adaptive
  predictors (Theorem 10.2) the parameter trajectory after $P + K$
  depends on the bytes in $[P, P+K)$; for non-$W$-bounded
  predictors (Theorem 10.4) the state snapshot at the \emph{next}
  sync point also depends on those bytes. In those cases the
  region rewrite must cascade through all sync points that read
  modified-region bytes into their context, and the cost is
  $O(|X| \cdot \mathrm{cost}(M))$ in the worst case (essentially a
  full re-encoding), not $O(K \cdot \mathrm{cost}(M))$. The
  fixed-predictor case below is the only one that supports cheap
  local rewrite. Cost (fixed predictor): $O(K \cdot \mathrm{cost}(M))$
  per modification.
  The naive append-in-place form is well-defined only if the
  re-encoded region's bit length is bounded by a pre-reserved slot
  size; otherwise the rewritten region shifts every subsequent
  sync-point bit offset, and the archive needs either (i)~per-region
  fixed-size slots with reserved padding, (ii)~an indirection layer
  that decouples logical sync-point identifiers from physical bit
  offsets (e.g., a separately mutable offset table), or
  (iii)~rewriting all later seek-index entries on each write,
  raising the per-modification cost to
  $O(K \cdot \mathrm{cost}(M) + (N-P)/K)$ for the index update.
  Suitable for low-frequency writes on archives that benefit from
  in-place mutation, with the slot/indirection choice fixed at
  archive format design time.
\item
  Full re-encoding. Periodic offline compaction re-encodes the whole
  archive, merging accumulated overlay writes back into the base.
  Standard practice for read-mostly snapshot stores; the compaction
  cost is $O(N \cdot \mathrm{cost}(M))$ but is amortized over many
  writes.
\end{itemize}

For the disk-image use case named at the start of this section (compress
an ISO or block-device image, mount through a driver, use without full
decompression), the read-only base plus copy-on-write overlay pattern is
the natural deployment. It matches the architecture used today for
compressed container images and live system images. RNR would
replace the dictionary-based compressed filesystem in the base
layer; whether the resulting archive achieves a higher compression
ratio than the dictionary-based baseline at identical access
semantics is conjectured but is not derived in this paper (the
ratio claim is empirical and conditional on the predictor
calibration; see §11 for conjectural predictions and §13.3 for the
empirical validation deferral).

\emph{\textbf{Comparison with existing compressed filesystems}}

Compressed read-only filesystems are widely deployed:

\begin{itemize}
\item
  squashfs (Linux): gzip, xz, lzma, lz4, or zstd compression in 4 KiB to
  1 MiB blocks. Typical ratio 1.5 to 3 times. Random-access cost O(K +
  k). Used by live Linux distributions, snap packages, and container
  image layers.
\item
  EROFS (Linux): lz4, lzma, or deflate in 4 to 64 KiB blocks. Typical
  ratio 1.5 to 2.5 times. Random-access cost O(K + k). Used by Android
  system partitions and Huawei storage products.
\item
  btrfs with zstd (Linux): zstd compression in 4 to 128 KiB extents.
  Typical ratio 1.5 to 2 times. Random-access cost O(K + k).
\item
  APFS compressed files (macOS) and NTFS compressed files (Windows):
  LZ4-class compression at fixed sector boundaries. Typical ratio 1.2 to
  1.8 times. Random-access cost O(sector size).
\end{itemize}

RNR with a well-fitted transformer predictor is conjectured to
achieve compression ratios materially better than the dictionary-
based systems above on structured data classes (system images,
source archives, log files, structured binaries), at the same
random-access cost class and at the same conformance overhead.
Specific numerical ratios are not derived here and are not theorems;
order-of-magnitude predictions are stated as conjectures in §11
(``Type-I, Code12 on text and source code'', and the corresponding
predictions for other Types) and require empirical validation
against NNCP, PixelCNN++, LZMA, and zstd baselines. For
read-mostly use cases with rare writes --- container images, system
snapshots, archival storage, software distribution images --- the
qualitative deployment claim is that RNR matches the access pattern
of these systems while potentially offering a compression advantage
that empirical validation must establish. The cost is the predictor
inference compute, which is feasible at NPU or GPU speeds and acceptable
at CPU speeds for non-interactive workloads.

\emph{\textbf{Reference driver architecture}}

A reference RNR block-device driver consists of the following
components, mirroring the structure of existing compressed-filesystem
drivers:

\begin{itemize}
\item
  Header parser. Reads the archive header into memory: predictor
  weights, sync-point seek index, verification metadata, type
  metadata.
\item
  Predictor inference engine. Implements $M$ on one or more hardware
  paths (CPU scalar, CPU SIMD, GPU integer, NPU integer), each path
  conforming to Theorem 10.6. The path is selected at driver load
  time based on hardware availability.
\item
  Range-decode function. Given $(p, k)$ with $p \in [0, |X|)$,
  computes $k' = \min(k, |X| - p)$, looks up the latest sync point
  $P \leq p$ in the seek index, then dispatches by the archive's
  coding mode: per-position sequential modes use the Theorem 10.3
  decode path (decode from $P$ to $p + k'$, return bytes $[p, p + k')$);
  whole-block partition-rank modes use the Remark 10.3a decode
  path (for each of the $\lceil k'/K \rceil + O(1)$ sub-blocks
  intersecting $[p, p + k')$, look up the sub-block's bit offset
  and unrank the full $K$-byte sub-block; assemble the requested
  $k'$ bytes from the concatenated sub-block outputs). The mode is
  recorded in the archive header. The function returns empty when
  $p \geq |X|$.
\item
  Cache. LRU or ARC cache of recently decoded regions, with eviction
  policy tuned to expected access pattern. Sits between the
  range-decode function and the OS page cache.
\item
  OS interface. Translates kernel-level block-device or filesystem
  read requests into range-decode calls. Options include Linux FUSE
  for user-space deployment, Linux NBD or device-mapper for
  kernel-level deployment, the macOS FSKit interface, or a Windows
  file-system minifilter driver.
\end{itemize}

Every component except the predictor inference engine has direct
counterparts in existing compressed-filesystem drivers; the integration
work is to replace the dictionary-based decompressor with the neural
inference engine. Production implementations of small transformers on
each target hardware path already exist; combining them into a
Theorem-10.6-conforming inference engine is engineering work and is
deferred to a companion artifact (§13.3).

\textbf{Open problems and status.} The open-problem ledger for the RNR
series is maintained in [RNR-I, Section 13]; the items bearing on this Part
--- the residual gaps in the Section 7.15 exact-entropy arc (the open
middle between the polynomial finite-belief island and the
coefficient-query hardness wall), the sub-block-restriction question for
Theorem 7.27, and the LRU / cache-locality item for random access --- are
recorded there and are not reproduced here to keep this Part self-contained.
The operational rate--distortion dispersion programme (the Section 7.34
family) is developed separately in [RNR-III].
\textbf{References}

\textbf{{[}1{]}} Banner, R., Hubara, I., Hoffer, E., and Soudry, D.
(2018). Scalable Methods for 8-bit Training of Neural Networks. NeurIPS.
\textbf{{[}2{]}} Bellard, F. (2019). Lossless Data Compression with
Neural Networks. Technical report and software,
https://bellard.org/nncp/.
\textbf{{[}3{]}} Bernstein, J., Wang, Y.-X., Azizzadenesheli, K., and
Anandkumar, A. (2018). signSGD: Compressed Optimisation for Non-Convex
Problems. ICML.
\textbf{{[}4{]}} Burrows, M., and Wheeler, D. J. (1994). A Block-Sorting
Lossless Data Compression Algorithm. SRC Research Report 124, Digital
Equipment Corporation.
\textbf{{[}5{]}} Carlsson, G., Ishkhanov, T., de Silva, V., and
Zomorodian, A. (2008). On the Local Behavior of Spaces of Natural
Images. International Journal of Computer Vision, 76(1), 1-12.
\textbf{{[}6{]}} Chen, X., Mishra, N., Rohaninejad, M., and Abbeel, P.
(2018). PixelSNAIL: An Improved Autoregressive Generative Model.
International Conference on Machine Learning.
\textbf{{[}7{]}} Charikar, M., Lehman, E., Liu, D., Panigrahy, R.,
Prabhakaran, M., Sahai, A., and Shelat, A. (2005). The Smallest Grammar
Problem. IEEE Transactions on Information Theory, 51(7), 2554-2576.

\textbf{{[}8a{]}} Casel, K., Fernau, H., Gaspers, S., Gras, B., and
Schmid, M. L. (2021). On the Complexity of the Smallest Grammar Problem
over Fixed Alphabets. Theory of Computing Systems, 65(1), 158-203.

\textbf{{[}8b{]}} Ga\'nczorz, M. (2018). Entropy bounds for grammar
compression. arXiv:1804.08547 (also presented at DCC 2018/2020).

\textbf{{[}8c{]}} Kozma, L., and Voderholzer, J. (2024). The
Computational Complexity of Tokenisation. arXiv:2411.08671.

\textbf{{[}8d{]}} Je\.z, A. (2015). A really simple approximation of
smallest grammar. Theoretical Computer Science, 616, 141-150.

\textbf{{[}8e{]}} Karzanov, A. V. (1998). Minimum 0-Extensions of Graph
Metrics. European Journal of Combinatorics, 19(1), 71-101.

\textbf{{[}8f{]}} Schwartz, R., and Tur, S. (2023). Improved hardness
of approximation for graph problems via Boolean PCPs.
arXiv:2104.11670 (also conference version, ICALP 2023).
\textbf{{[}8{]}} Cover, T. M. (1973). Enumerative Source Encoding. IEEE
Transactions on Information Theory, 19(1), 73-77.
\textbf{{[}9{]}} Cover, T. M., and Thomas, J. A. (2006). Elements of
Information Theory, 2nd edition. Wiley.
\textbf{{[}10{]}} Donoho, D. L. (2000). High-Dimensional Data Analysis:
The Curses and Blessings of Dimensionality. AMS Math Challenges Lecture.
\textbf{{[}11{]}} Duda, J. (2009). Asymmetric Numeral Systems.
arXiv:0902.0271.
\textbf{{[}12{]}} Frey, B. J. (1997). Bayesian Networks for Pattern
Classification, Data Compression, and Channel Coding. Ph.D. thesis,
University of Toronto.
\textbf{{[}13{]}} Frey, B. J. (1998). Graphical Models for Machine
Learning and Digital Communication. MIT Press.
\textbf{{[}14{]}} Hinton, G. E., and van Camp, D. (1993). Keeping Neural
Networks Simple by Minimizing the Description Length of the Weights.
Conference on Learning Theory.
\textbf{{[}15{]}} Hoogeboom, E., Peters, J. W. T., van den Berg, R., and
Welling, M. (2019). Integer Discrete Flows and Lossless Compression.
NeurIPS.
\textbf{{[}16{]}} Huffman, D. A. (1952). A Method for the Construction
of Minimum-Redundancy Codes. Proceedings of the IRE, 40(9), 1098-1101.
\textbf{{[}17{]}} Kingma, F. H., Abbeel, P., and Ho, J. (2019).
Bit-Swap: Recursive Bits-Back Coding for Lossless Compression with
Hierarchical Latent Variables. ICML.
\textbf{{[}18{]}} Knoll, B. (ongoing). Cmix: a lossless data compression
program. https://www.byronknoll.com/cmix.html.
\textbf{{[}19{]}} Levina, E., and Bickel, P. J. (2004). Maximum
Likelihood Estimation of Intrinsic Dimension. NIPS.
\textbf{{[}20{]}} Liveris, A. D., Xiong, Z., and Georghiades, C. N.
(2002). Compression of Binary Sources with Side Information at the
Decoder Using LDPC Codes. IEEE Communications Letters, 6(10), 440-442.
\textbf{{[}21{]}} Mentzer, F., Agustsson, E., Tschannen, M., Timofte,
R., and Van Gool, L. (2019). Practical Full Resolution Learned Lossless
Image Compression. CVPR.
\textbf{{[}22{]}} Pradhan, S. S., and Ramchandran, K. (2003).
Distributed Source Coding Using Syndromes (DISCUS): Design and
Construction. IEEE Transactions on Information Theory, 49(3), 626-643.
\textbf{{[}23{]}} Reed, I. S., and Solomon, G. (1960). Polynomial Codes
over Certain Finite Fields. Journal of SIAM, 8(2), 300-304.
\textbf{{[}24{]}} Rissanen, J. (1978). Modeling by Shortest Data
Description. Automatica, 14(5), 465-471.
\textbf{{[}25{]}} Rissanen, J., and Langdon, G. G. (1979). Arithmetic
Coding. IBM Journal of Research and Development, 23(2), 149-162.
\textbf{{[}26{]}} Salimans, T., Karpathy, A., Chen, X., and Kingma, D.
P. (2017). PixelCNN++: Improving the PixelCNN with Discretized Logistic
Mixture Likelihood and Other Modifications. ICLR.

\textbf{{[}29a{]}} Crutchfield, J. P., and Feldman, D. P. (2003).
Regularities Unseen, Randomness Observed: Levels of Entropy
Convergence. Chaos, 13(1), 25-54. DOI: 10.1063/1.1530990.
\textbf{{[}27{]}} Sahni, S., and Gonzalez, T. (1976). P-Complete
Approximation Problems. Journal of the ACM, 23(3), 555-565.
\textbf{{[}28{]}} Shannon, C. E. (1948). A Mathematical Theory of
Communication. Bell System Technical Journal, 27, 379-423 and 623-656.

\textbf{{[}31a{]}} Shannon, C. E. (1951). Prediction and Entropy of
Printed English. Bell System Technical Journal, 30(1), 50-64.
\textbf{{[}29{]}} Slepian, D., and Wolf, J. K. (1973). Noiseless Coding
of Correlated Information Sources. IEEE Transactions on Information
Theory, 19(4), 471-480.
\textbf{{[}30{]}} Tenenbaum, J. B., de Silva, V., and Langford, J. C.
(2000). A Global Geometric Framework for Nonlinear Dimensionality
Reduction. Science, 290(5500), 2319-2323.
\textbf{{[}31{]}} Townsend, J., Bird, T., and Barber, D. (2019).
Practical Lossless Compression with Latent Variables Using Bits Back
Coding. ICLR.
\textbf{{[}32{]}} Townsend, J., Bird, T., Kunze, J., and Barber, D.
(2020). HiLLoC: Lossless Image Compression with Hierarchical Latent
Variable Models. ICLR.
\textbf{{[}33{]}} van den Berg, R., Gritsenko, A. A., Dehghani, M.,
Sonderby, C. K., and Salimans, T. (2021). IDF++: Analyzing and Improving
Integer Discrete Flows for Lossless Compression. ICLR.
\textbf{{[}34{]}} van den Oord, A., Kalchbrenner, N., and Kavukcuoglu,
K. (2016). Pixel Recurrent Neural Networks. ICML.
\textbf{{[}35{]}} van den Oord, A., Kalchbrenner, N., Vinyals, O.,
Espeholt, L., Graves, A., and Kavukcuoglu, K. (2016). Conditional Image
Generation with PixelCNN Decoders. NeurIPS.
\textbf{{[}36{]}} Rhee, H., Jang, Y. I., Kim, S., and Cho, N. I. (2022).
LC-FDNet: Learned Lossless Image Compression with Frequency
Decomposition Network. CVPR, 6023-6032.
\textbf{{[}37{]}} Witten, I. H., Neal, R. M., and Cleary, J. G. (1987).
Arithmetic Coding for Data Compression. Communications of the ACM,
30(6), 520-540.
\textbf{{[}38{]}} Wu, S., Yan, J., Wang, J., Maxson, A., Xu, K., and He,
K. (2018). Training and Inference with Integers in Deep Neural Networks.
ICLR.
\textbf{{[}39{]}} Wyner, A. D., and Ziv, J. (1976). The Rate-Distortion
Function for Source Coding with Side Information at the Decoder. IEEE
Transactions on Information Theory, 22(1), 1-10.
\textbf{{[}40{]}} Ziv, J., and Lempel, A. (1977). A Universal Algorithm
for Sequential Data Compression. IEEE Transactions on Information
Theory, 23(3), 337-343.
\textbf{{[}41{]}} Ziv, J., and Lempel, A. (1978). Compression of
Individual Sequences via Variable-Rate Coding. IEEE Transactions on
Information Theory, 24(5), 530-536.
\textbf{{[}42{]}} Watanabe, S. (1960). Information Theoretical
Analysis of Multivariate Correlation. IBM Journal of Research and
Development, 4(1), 66-82.
\textbf{{[}43{]}} Studen\'y, M., and Vejnarov\'a, J. (1998). The
Multiinformation Function as a Tool for Measuring Stochastic
Dependence. In: \emph{Learning in Graphical Models}, Jordan, M. I.
(ed.), MIT Press, pp. 261-297.
\textbf{{[}44{]}} Lund, C., and Yannakakis, M. (1994). On the
Hardness of Approximating Minimization Problems. Journal of the ACM,
41(5), 960-981.
\textbf{{[}45{]}} Lehman, E., and Shelat, A. (2002). Approximation
Algorithms for Grammar-Based Compression. Proceedings of the
Thirteenth Annual ACM-SIAM Symposium on Discrete Algorithms (SODA),
205-212.
\textbf{{[}46{]}} Berman, P., and Karpinski, M. (1999). On Some
Tighter Inapproximability Results. Proceedings of the 26th
International Colloquium on Automata, Languages and Programming
(ICALP), Lecture Notes in Computer Science 1644, Springer, 200-209.
\textbf{{[}47{]}} Storer, J. A., and Szymanski, T. G. (1982). Data
Compression via Textual Substitution. Journal of the ACM, 29(4),
928-951.
\textbf{{[}48{]}} Ar{\i}kan, E. (2009). Channel Polarization: A Method
for Constructing Capacity-Achieving Codes for Symmetric Binary-Input
Memoryless Channels. IEEE Transactions on Information Theory, 55(7),
3051-3073.
\textbf{{[}49{]}} Azuma, K. (1967). Weighted Sums of Certain Dependent
Random Variables. T\^ohoku Mathematical Journal, 19(3), 357-367.
\textbf{{[}50{]}} Bellare, M. (2006). New Proofs for NMAC and HMAC:
Security Without Collision-Resistance. In: Advances in Cryptology ---
CRYPTO 2006, Lecture Notes in Computer Science 4117, Springer, 602-619.
\textbf{{[}51{]}} Bellare, M., Canetti, R., and Krawczyk, H. (1996).
Keying Hash Functions for Message Authentication. In: Advances in
Cryptology --- CRYPTO '96, Lecture Notes in Computer Science 1109,
Springer, 1-15.
\textbf{{[}52{]}} Bellare, M., and Namprempre, C. (2000). Authenticated
Encryption: Relations Among Notions and Analysis of the Generic
Composition Paradigm. In: Advances in Cryptology --- ASIACRYPT 2000,
Lecture Notes in Computer Science 1976, Springer, 531-545.
\textbf{{[}53{]}} Berger, T. (1971). Rate Distortion Theory: A
Mathematical Basis for Data Compression. Prentice-Hall, Englewood
Cliffs, NJ.
\textbf{{[}54{]}} Breslau, L., Cao, P., Fan, L., Phillips, G., and
Shenker, S. (1999). Web Caching and Zipf-like Distributions: Evidence
and Implications. Proceedings of IEEE INFOCOM, 126-134.
\textbf{{[}55{]}} Cover, T. M. (1975). A Proof of the Data Compression
Theorem of Slepian and Wolf for Ergodic Sources. IEEE Transactions on
Information Theory, 21(2), 226-228.
\textbf{{[}56{]}} Cover, T. M., and King, R. C. (1978). A Convergent
Gambling Estimate of the Entropy of English. IEEE Transactions on
Information Theory, 24(4), 413-421.
\textbf{{[}57{]}} Csisz\'ar, I., and K\"orner, J. (2011). Information
Theory: Coding Theorems for Discrete Memoryless Systems, 2nd edition.
Cambridge University Press. (First edition: Akad\'emiai Kiad\'o, 1981.)
\textbf{{[}58{]}} Cybenko, G. (1989). Approximation by Superpositions
of a Sigmoidal Function. Mathematics of Control, Signals, and Systems,
2(4), 303-314.
\textbf{{[}59{]}} Del\'etang, G., Ruoss, A., Duquenne, P.-A., Catt, E.,
Genewein, T., Mattern, C., Grau-Moya, J., Wenliang, L. K.,
Aitchison, M., Orseau, L., Legg, S., and Veness, J. (2024). Language
Modeling Is Compression. Proceedings of ICLR. arXiv:2309.10668.
\textbf{{[}60{]}} Doshi, V., Shah, D., and M\'edard, M. (2007). Source
Coding with Distortion through Graph Coloring. Proceedings of the IEEE
International Symposium on Information Theory (ISIT), 1501-1505.
\textbf{{[}61{]}} Doshi, V., Shah, D., M\'edard, M., and Effros, M.
(2010). Functional Compression through Graph Coloring. IEEE
Transactions on Information Theory, 56(8), 3901-3917.
\textbf{{[}62{]}} Equitz, W. H. R., and Cover, T. M. (1991). Successive
Refinement of Information. IEEE Transactions on Information Theory,
37(2), 269-275.
\textbf{{[}63{]}} Fagin, R. (1977). Generalized First-Order Spectra and
Polynomial-Time Recognizable Sets. In: Complexity of Computation,
SIAM-AMS Proceedings 7, 43-73.
\textbf{{[}64{]}} Ferragut, A., Rodr{\'\i}guez, I., and Paganini, F.
(2018). Optimizing TTL Caches under Heavy-Tailed Demand. Queueing
Systems, 88(3-4), 207-241.
\textbf{{[}65{]}} Grossi, R., and Vitter, J. S. (2005). Compressed
Suffix Arrays and Suffix Trees with Applications to Text Indexing and
String Matching. SIAM Journal on Computing, 35(2), 378-407.
\textbf{{[}66{]}} Hinton, G., Vinyals, O., and Dean, J. (2015).
Distilling the Knowledge in a Neural Network. arXiv:1503.02531.
\textbf{{[}67{]}} Hoeffding, W. (1963). Probability Inequalities for
Sums of Bounded Random Variables. Journal of the American Statistical
Association, 58(301), 13-30.
\textbf{{[}68{]}} Hoffmann, J., Borgeaud, S., Mensch, A., Buchatskaya,
E., Cai, T., Rutherford, E., et al. (2022). Training Compute-Optimal
Large Language Models. arXiv:2203.15556.
\textbf{{[}69{]}} Hornik, K., Stinchcombe, M., and White, H. (1989).
Multilayer Feedforward Networks Are Universal Approximators. Neural
Networks, 2(5), 359-366.
\textbf{{[}70{]}} Kaplan, J., McCandlish, S., Henighan, T., Brown,
T. B., Chess, B., Child, R., Gray, S., Radford, A., Wu, J., and
Amodei, D. (2020). Scaling Laws for Neural Language Models.
arXiv:2001.08361.
\textbf{{[}71{]}} Korada, S. B., and Urbanke, R. L. (2010). Polar
Codes Are Optimal for Lossy Source Coding. IEEE Transactions on
Information Theory, 56(4), 1751-1768.
\textbf{{[}72{]}} McDiarmid, C. (1989). On the Method of Bounded
Differences. In: Surveys in Combinatorics 1989, London Mathematical
Society Lecture Notes 141, Cambridge University Press, 148-188.
\textbf{{[}73{]}} McLeish, D. L. (1974). Dependent Central Limit
Theorems and Invariance Principles. Annals of Probability, 2(4),
620-628.
\textbf{{[}74{]}} Muthitacharoen, A., Chen, B., and Mazi\`eres, D.
(2001). A Low-Bandwidth Network File System. Proceedings of the 18th
ACM Symposium on Operating Systems Principles (SOSP), 174-187.
\textbf{{[}75{]}} Orlitsky, A., and Roche, J. R. (2001). Coding for
Computing. IEEE Transactions on Information Theory, 47(3), 903-917.
\textbf{{[}76{]}} Polyanskiy, Y., Poor, H. V., and Verd\'u, S. (2010).
Channel Coding Rate in the Finite Blocklength Regime. IEEE Transactions
on Information Theory, 56(5), 2307-2359.
\textbf{{[}77{]}} Rabin, M. O. (1981). Fingerprinting by Random
Polynomials. Technical Report TR-15-81, Center for Research in Computing
Technology, Harvard University.
\textbf{{[}78{]}} Richardson, T., and Urbanke, R. (2008). Modern Coding
Theory. Cambridge University Press.
\textbf{{[}79{]}} Rimoldi, B. (1994). Successive Refinement of
Information: Characterization of the Achievable Rates. IEEE Transactions
on Information Theory, 40(1), 253-259.
\textbf{{[}80{]}} Rio, E. (2000). Th\'eorie Asymptotique des Processus
Al\'eatoires Faiblement D\'ependants. Math\'ematiques \& Applications 31,
Springer.
\textbf{{[}81{]}} Rissanen, J. (1996). Fisher Information and
Stochastic Complexity. IEEE Transactions on Information Theory, 42(1),
40-47.
\textbf{{[}82{]}} Sadakane, K. (2003). New Text Indexing Functionalities
of the Compressed Suffix Arrays. Journal of Algorithms, 48(2), 294-313.
\textbf{{[}83{]}} Said, A. (2003). Arithmetic Coding for Data
Compression. Technical Report, Hewlett-Packard Laboratories.
\textbf{{[}84{]}} Samson, P.-M. (2000). Concentration of Measure
Inequalities for Markov Chains and $\Phi$-Mixing Processes. Annals of
Probability, 28(1), 416-461.
\textbf{{[}85{]}} Shannon, C. E. (1959). Coding Theorems for a Discrete
Source with a Fidelity Criterion. IRE National Convention Record, Part 4,
142-163.
\textbf{{[}86{]}} Sleator, D. D., and Tarjan, R. E. (1985). Amortized
Efficiency of List Update and Paging Rules. Communications of the ACM,
28(2), 202-208.
\textbf{{[}87{]}} Steinberg, Y., and Merhav, N. (2004). On Successive
Refinement for the Wyner-Ziv Problem. IEEE Transactions on Information
Theory, 50(8), 1636-1654.
\textbf{{[}88{]}} Witsenhausen, H. S. (1976). The Zero-Error Side
Information Problem and Chromatic Numbers. IEEE Transactions on
Information Theory, 22(5), 592-593.
\textbf{{[}89{]}} Zhang, Y., Huang, Z., Chen, K., Tao, D., et al.
(2024). L3TC: Leveraging RWKV for Learned Lossless Low-Complexity Text
Compression. arXiv:2412.16642.
\textbf{{[}90{]}} Krichevsky, R. E., and Trofimov, V. K. (1981). The
Performance of Universal Encoding. IEEE Transactions on Information
Theory, 27(2), 199-207.
\textbf{{[}91{]}} Baum, L. E., and Petrie, T. (1966). Statistical
Inference for Probabilistic Functions of Finite State Markov Chains.
Annals of Mathematical Statistics, 37(6), 1554-1563.
\textbf{{[}92{]}} Lauritzen, S. L. (1996). Graphical Models. Oxford
University Press.
\textbf{{[}93{]}} Lauritzen, S. L., and Spiegelhalter, D. J. (1988).
Local Computations with Probabilities on Graphical Structures and
Their Application to Expert Systems. Journal of the Royal Statistical
Society, Series B, 50(2), 157-194.
\textbf{{[}94{]}} Dinur, I. (2007). The PCP Theorem by Gap
Amplification. Journal of the ACM, 54(3), Article 12.
\textbf{{[}95{]}} Dinur, I., Kindler, G., Raz, R., and Safra, S.
(2003). Approximating CVP to Within Almost-Polynomial Factors is
NP-Hard. Combinatorica, 23(2), 205-243.
\textbf{{[}96{]}} Barron, A. R. (1985). The Strong Ergodic Theorem
for Densities: Generalized Shannon-McMillan-Breiman Theorem. Annals
of Probability, 13(4), 1292-1303.
\textbf{{[}97{]}} Gray, R. M. (1990). Entropy and Information
Theory. Springer.
\textbf{{[}98{]}} Pradhan, S. S., and Ramchandran, K. (1999).
Distributed Source Coding Using Syndromes (DISCUS): Design and
Construction. Proceedings of the Data Compression Conference (DCC),
158-167.
\textbf{{[}99{]}} Stratonovich, R. L. (1960). Conditional Markov
Processes. Theory of Probability and Its Applications, 5(2), 156-178.
\textbf{{[}100{]}} Wainwright, M. J., and Jordan, M. I. (2008).
Graphical Models, Exponential Families, and Variational Inference.
Foundations and Trends in Machine Learning, 1(1-2), 1-305.
\textbf{{[}101{]}} Arora, S., Babai, L., Stern, J., and Sweedyk, Z.
(1997). The Hardness of Approximate Optima in Lattices, Codes, and
Systems of Linear Equations. Journal of Computer and System Sciences,
54(2), 317-331.
\textbf{{[}102{]}} Birch, B. J. (1962). Approximations for the Entropy
for Functions of Markov Chains. Annals of Mathematical Statistics,
33(3), 930-938.
\textbf{{[}103{]}} Manurangsi, P. (2017). Almost-Polynomial Ratio
ETH-Hardness of Approximating Densest $k$-Subgraph. Proceedings of the
49th Annual ACM Symposium on Theory of Computing (STOC), 19-23.
\textbf{{[}104{]}} Pakzad, P., and Anantharam, V. (2002). Estimation
and Marginalization Using Kikuchi Approximation Methods. Neural
Computation, 17(8), 1836-1873.
\textbf{{[}105{]}} Kieffer, J. C., and Yang, E.-H. (2000). Grammar-Based
Codes: A New Class of Universal Lossless Source Codes. IEEE Transactions
on Information Theory, 46(3), 737-754.
\textbf{{[}106{]}} Merlev\`ede, F., Peligrad, M., and Rio, E. (2009).
Bernstein Inequality and Moderate Deviations under Strong Mixing
Conditions. In: High Dimensional Probability V: The Luminy Volume,
IMS Collections 5, 273-292.
\textbf{{[}107{]}} Bradley, R. C. (2005). Basic Properties of Strong
Mixing Conditions: A Survey and Some Open Questions. Probability
Surveys, 2, 107-144.
\textbf{{[}108{]}} Rio, E. (2017). Asymptotic Theory of Weakly
Dependent Random Processes. Probability Theory and Stochastic
Modelling 80, Springer.
\textbf{{[}109{]}} Kontorovich, L., and Ramanan, K. (2008).
Concentration Inequalities for Dependent Random Variables via the
Martingale Method. Annals of Probability, 36(6), 2126-2158.
\textbf{{[}110{]}} Davydov, Yu.\ A. (1968). Convergence of
Distributions Generated by Stationary Stochastic Processes. Theory
of Probability and Its Applications, 13(4), 691-696.
\textbf{{[}111{]}} Bradley, R. C. (2007). Introduction to Strong
Mixing Conditions, Volumes 1--3. Kendrick Press.
\textbf{{[}112{]}} Doukhan, P. (1994). Mixing: Properties and
Examples. Lecture Notes in Statistics 85, Springer.
\textbf{{[}113{]}} Brown, P. F., Della Pietra, V. J., Mercer,
R. L., Della Pietra, S. A., and Lai, J. C. (1992). An Estimate of
an Upper Bound for the Entropy of English. Computational Linguistics,
18(1), 31-40.
\textbf{{[}114{]}} Lezaud, P. (1998). Chernoff-Type Bound for Finite
Markov Chains. Annals of Applied Probability, 8(3), 849-867.
\textbf{{[}115{]}} Che, H., Wang, Z., and Tung, Y. (2002). Analysis
and Design of Hierarchical Web Caching Systems. Proceedings of IEEE
INFOCOM, 1416-1424.
\textbf{{[}116{]}} Aldous, D., and Diaconis, P. (1987). Shuffling
Cards and Stopping Times. American Mathematical Monthly, 93(5),
333-348.
\textbf{{[}117{]}} Knudsen, J. B., and Tsitsiklis, J. N. (1990).
On the Time Required to Drain a Network. Mathematics of Operations
Research, 15(4), 685-708.
\textbf{{[}118{]}} Levin, D. A., Peres, Y., and Wilmer, E. L. (2017).
Markov Chains and Mixing Times, 2nd edition. American Mathematical
Society.
\textbf{{[}119{]}} Gillman, D. (1998). A Chernoff Bound for Random
Walks on Expander Graphs. SIAM Journal on Computing, 27(4),
1203-1220.
\textbf{{[}120{]}} Paulin, D. (2015). Concentration Inequalities
for Markov Chains by Marton Couplings and Spectral Methods.
Electronic Journal of Probability, 20(79), 1-32.
\textbf{{[}121{]}} Aldous, D., and Fill, J. A. (2002).
Reversible Markov Chains and Random Walks on Graphs.
Unfinished monograph, available at
https://www.stat.berkeley.edu/\textasciitilde aldous/RWG/book.html.
\textbf{{[}122{]}} Diaconis, P., and Stroock, D. (1991). Geometric
Bounds for Eigenvalues of Markov Chains. Annals of Applied
Probability, 1(1), 36-61.
\textbf{{[}123{]}} King, W. F. (1971). Analysis of Demand Paging
Algorithms. Proceedings of the IFIP Congress, 485-490.
\textbf{{[}124{]}} Kesidis, G., Shan, Y., Fumagalli, A., Gulati,
A., and Singh, A. (2017). An Analysis of LRU Replacement Caches
under Independent Reference Models. arXiv:1704.04849.
\textbf{{[}125{]}} Niyogi, P., Smale, S., and Weinberger, S. (2008).
Finding the Homology of Submanifolds with High Confidence from
Random Samples. Discrete and Computational Geometry, 39(1-3),
419-441. (Introduces the ``reach'' parameter of a sub-manifold and
gives quantitative manifold-learning sampling bounds.)
\textbf{{[}126{]}} Muirhead, R. J. (1982). Aspects of Multivariate
Statistical Theory. Wiley.
\textbf{{[}127{]}} Diaconis, P., and Freedman, D. (1987).
A dozen de~Finetti-style results in search of a theory. Annales de
l'Institut Henri Poincar\'e (B) Probabilities et Statistiques, 23(S2),
397-423.
\textbf{{[}128{]}} Gersho, A., and Gray, R. M. (1992). Vector
Quantization and Signal Compression. Kluwer Academic Publishers.
(Bennett high-resolution quantization theory, Chapter 5.)
\textbf{{[}129{]}} Talagrand, M. (1995). Concentration of Measure
and Isoperimetric Inequalities in Product Spaces. Publications
Math\'ematiques de l'IH\'ES, 81, 73-205.
\textbf{{[}130{]}} Bennett, W. R. (1948). Spectra of Quantized
Signals. Bell System Technical Journal, 27(3), 446-472.
\textbf{{[}131{]}} Ryabko, B. Ya. (1984). Twice-Universal Coding.
Problems of Information Transmission, 20(3), 173-177.
\textbf{{[}132{]}} Willems, F. M. J., Shtarkov, Y. M., and Tjalkens,
T. J. (1995). The Context-Tree Weighting Method: Basic Properties.
IEEE Transactions on Information Theory, 41(3), 653-664.

\textbf{[RNR-I]} Ratushnyi, P. (2026). Ratushnyi Neural Repair Coding,
Part I: Framework, Coder Types, and the Achievability--Converse Matrix.
Companion paper (Part I of the RNR series).

\textbf{[RNR-III]} Ratushnyi, P. (2026). Ratushnyi Neural Repair Coding,
Part III: Operational Rate--Distortion Dispersion. Companion paper (Part
III of the RNR series).

\textbf{[C-CB90]} Clarke, B. S., and Barron, A. R. (1990). Information-
Theoretic Asymptotics of Bayes Methods. IEEE Transactions on Information
Theory, 36(3), 453-471.

\textbf{[C-KV14]} Kontoyiannis, I., and Verd\'u, S. (2014). Optimal
Lossless Data Compression: Non-Asymptotics and Asymptotics. IEEE
Transactions on Information Theory, 60(2), 777-795.

\textbf{[C-KV12]} Kostina, V., and Verd\'u, S. (2012). Fixed-Length Lossy
Compression in the Finite Blocklength Regime. IEEE Transactions on
Information Theory, 58(6), 3309-3338.

\textbf{[C-TK14]} Tan, V. Y. F., and Kosut, O. (2014). On the Dispersions
of Three Network Information Theory Problems. IEEE Transactions on
Information Theory, 60(2), 881-903.

\textbf{[C-Petrov75]} Petrov, V. V. (1975). Sums of Independent Random
Variables. Springer-Verlag, Berlin (Ergebnisse der Mathematik und ihrer
Grenzgebiete, Band 82).

\textbf{[C-LLR83]} Leadbetter, M. R., Lindgren, G., and Rootz\'en, H.
(1983). Extremes and Related Properties of Random Sequences and Processes.
Springer-Verlag, New York.

\textbf{[C-FV07]} Ferragina, P., and Venturini, R. (2007). A Simple
Storage Scheme for Strings Achieving Entropy Bounds. Theoretical Computer
Science, 372(1), 115-121 (also SODA 2007).

\textbf{[C-MR-succ]} Munro, J. I., and Raman, R. (2001). Succinct
Representation of Balanced Parentheses and Static Trees. SIAM Journal on
Computing, 31(3), 762-776.

\textbf{[C-WK19]} Wolfer, G., and Kontoyiannis, I. (2019). Minimax
Learning of Ergodic Markov Chains. Proceedings of the 30th International
Conference on Algorithmic Learning Theory (ALT), PMLR 98, 904-930.

\textbf{[C-Feige98]} Feige, U. (1998). A Threshold of $\ln n$ for
Approximating Set Cover. Journal of the ACM, 45(4), 634-652.

\textbf{[C-Moricz76]} M\'oricz, F. (1976). Moment Inequalities and the
Strong Laws of Large Numbers. Zeitschrift f\"ur Wahrscheinlichkeits-
theorie und Verwandte Gebiete, 35(4), 299-314.

\textbf{[C-Serfling70]} Serfling, R. J. (1970). Convergence Properties of
$S_n$ Under Moment Restrictions. Annals of Mathematical Statistics, 41(4),
1235-1248.
\end{document}
